\input{preamble}

% OK, start here.
%
\begin{document}

\title{Morphismes d'espaces algébriques}


\maketitle

\phantomsection
\label{section-phantom}

\tableofcontents

\section{Introduction}
\label{section-introduction}

\noindent
Dans ce chapitre, nous introduisons certains types de morphismes d'espaces
algébriques. On pourra consulter \cite{Kn}.

\medskip\noindent
Le but est d'étendre à la catégorie des espaces algébriques la définition
de chacun des types de morphismes de schémas définis dans les chapitres
Schémas et Morphismes de schémas. Chaque cas présente de légères différences
et il nous paraît préférable de les traiter séparément.





\section{Conventions}
\label{section-conventions}

\noindent
Nous supposons une fois pour toutes que tous les schémas appartiennent
à un grand site fppf $\Sch_{fppf}$. Tous les anneaux $A$ considérés
ont en outre la propriété que $\Spec(A)$ est (isomorphe à) un
objet de ce grand site.

\medskip\noindent
Soient $S$ un schéma et $X$ un espace algébrique au-dessus de $S$.
Dans ce chapitre et le suivant, nous noterons $X \times_S X$
le produit de $X$ par lui-même (dans la catégorie des espaces
algébriques au-dessus de $S$), au lieu de $X \times X$.





\section{Propriétés des morphismes représentables}
\label{section-representable}

\noindent
Soit $S$ un schéma.
Soit $f : X \to Y$ un morphisme représentable d'espaces algébriques. Dans
Espaces, section \ref{spaces-section-representable-properties},
nous avons défini ce que signifie pour $f$
posséder la propriété $\mathcal{P}$ lorsque $\mathcal{P}$ est une propriété
des morphismes de schémas qui
\begin{enumerate}
\item est préservée par tout changement de base,
voir Schémas, définition \ref{schemes-definition-preserved-by-base-change},
et
\item est locale sur la base pour la topologie fppf, voir
Descente, définition \ref{descent-definition-property-morphisms-local}.
\end{enumerate}
Plus précisément, dans ce cas, on dit que $f$ possède la propriété $\mathcal{P}$ si et seulement
si, pour tout schéma $U$ et tout morphisme $U \to Y$, le morphisme de schémas
$X \times_Y U \to U$ possède la propriété $\mathcal{P}$.

\medskip\noindent
D'après les listes de
Espaces, section \ref{spaces-section-lists},
ceci s'applique aux propriétés suivantes :
(1)(a) immersion fermée,
(1)(b) immersion ouverte,
(1)(c) immersion quasi-compacte,
(2) quasi-compact,
(3) universellement fermé,
(4) (quasi-)séparé,
(5) monomorphisme,
(6) surjectif,
(7) universellement injectif,
(8) affine,
(9) quasi-affine,
(10) (localement) de type fini,
(11) (localement) quasi-fini,
(12) (localement) de présentation finie,
(13) localement de type fini de dimension relative $d$,
(14) universellement ouvert,
(15) plat,
(16) syntomique,
(17) lisse,
(18) non ramifié (resp.\ G-non ramifié),
(19) \'etale,
(20) propre,
(21) fini ou entier,
(22) fini localement libre,
(23) universellement submersif,
(24) homéomorphisme universel, et
(25) immersion.

\medskip\noindent
Dans ce chapitre, nous redéfinirons ces notions pour les morphismes d'espaces
algébriques qui ne sont pas nécessairement représentables. Chaque fois, nous nous
assurerons que la nouvelle définition coïncide avec l'ancienne, afin d'éviter
toute ambiguïté.

\medskip\noindent
Remarquons que la définition ci-dessus s'applique lorsque $X$ est un schéma,
puisqu'un morphisme d'un schéma vers un espace algébrique est représentable.
Elle s'applique donc en particulier lorsque $X$ et $Y$ sont tous deux des schémas.
Dans
Espaces, lemme
\ref{spaces-lemma-morphism-schemes-gives-representable-transformation-property}
nous avons vu que, dans ce cas, les définitions
coïncident et qu'aucune ambiguïté ne se présente.

\medskip\noindent
En outre, dans
Espaces, lemme
\ref{spaces-lemma-base-change-representable-transformations-property}
nous avons vu que la propriété
ainsi définie pour les morphismes représentables d'espaces algébriques est stable par
tout changement de base par un morphisme d'espaces algébriques.
Enfin, dans
Espaces, lemmes
\ref{spaces-lemma-composition-representable-transformations-property} et
\ref{spaces-lemma-product-representable-transformations-property}
nous avons vu que, si $\mathcal{P}$ est stable par composition,
ce qui est le cas des propriétés
(1)(a), (1)(b), (1)(c), (2) -- (25), à l'exception de (13) ci-dessus, alors
les produits de morphismes représentables possèdent encore la propriété $\mathcal{P}$
et les composées de morphismes représentables possèdent encore la propriété $\mathcal{P}$.

\medskip\noindent
Nous utiliserons ces faits ci-dessous et, lorsque ce sera le cas, nous nous contenterons
de renvoyer à la présente section.















\section{Axiomes de séparation}
\label{section-separation-axioms}

\noindent
Il est naturel d'énumérer quelques propriétés a priori de la diagonale
d'un morphisme d'espaces algébriques.

\begin{lemma}
\label{lemma-properties-diagonal}
Soit $S$ un schéma appartenant à $\Sch_{fppf}$.
Soit $f : X \to Y$ un morphisme d'espaces algébriques au-dessus de $S$.
Soit $\Delta_{X/Y} : X \to X \times_Y X$ le morphisme diagonal.
Alors
\begin{enumerate}
\item $\Delta_{X/Y}$ est représentable,
\item $\Delta_{X/Y}$ est localement de type fini,
\item $\Delta_{X/Y}$ est un monomorphisme,
\item $\Delta_{X/Y}$ est séparé, et
\item $\Delta_{X/Y}$ est localement quasi-fini.
\end{enumerate}
\end{lemma}

\begin{proof}
Nous allons utiliser le fait que $\Delta_{X/S}$ est
représentable (par définition d'un espace algébrique) et qu'il
possède les propriétés (2) -- (5), voir
Espaces, lemme \ref{spaces-lemma-properties-diagonal}.
Remarquons que nous avons une factorisation
$$
X
\longrightarrow
X \times_Y X
\longrightarrow
X \times_S X
$$
de la diagonale $\Delta_{X/S} : X \to X \times_S X$. Puisque
$X \times_Y X \to X \times_S X$ est un monomorphisme et que
$\Delta_{X/S}$ est représentable, il s'ensuit formellement que
$\Delta_{X/Y}$ est représentable. En particulier, les autres
assertions ont désormais un sens, voir
la section \ref{section-representable}.

\medskip\noindent
Choisissons un morphisme \'etale surjectif $U \to X$, où $U$ est un schéma.
Considérons le diagramme
$$
\xymatrix{
R = U \times_X U \ar[r] \ar[d] &
U \times_Y U \ar[d] \ar[r] &
U \times_S U \ar[d] \\
X \ar[r] & X \times_Y X \ar[r] & X \times_S X
}
$$
Les deux carrés sont cartésiens; le rectangle extérieur l'est donc aussi.
La ligne supérieure est formée de schémas et les flèches verticales
sont des morphismes \'etales surjectifs. D'après
Espaces, lemme \ref{spaces-lemma-representable-morphisms-spaces-property},
les propriétés (2) -- (5) de $\Delta_{X/Y}$ sont équivalentes à celles de
$R \to U \times_Y U$. Dans la démonstration de
Espaces, lemme \ref{spaces-lemma-properties-diagonal},
nous avons vu que $R \to U \times_S U$ possède les propriétés (2) -- (5).
Le morphisme $U \times_Y U \to U \times_S U$ est un monomorphisme
de schémas. Il en résulte que $R \to U \times_Y U$ possède les
propriétés (2) -- (5).

\medskip\noindent
En effet, pour (3), remarquons que $R \to U \times_Y U$
est un monomorphisme, puisque la composée
$R \to U \times_S U$ en est un. Pour (2), remarquons que
$R \to U \times_Y U$ est localement de type fini, car la
composée $R \to U \times_S U$ est localement de type fini
(Morphismes, lemme \ref{morphisms-lemma-permanence-finite-type}).
Un monomorphisme localement de type fini est localement quasi-fini
parce qu'il a des fibres finies
(Morphismes, lemme \ref{morphisms-lemma-finite-fibre}); d'où (5).
Un monomorphisme est séparé
(Schémas, lemme \ref{schemes-lemma-monomorphism-separated}); d'où (4).
\end{proof}

\begin{definition}
\label{definition-separated}
Soit $S$ un schéma.
Soit $f : X \to Y$ un morphisme d'espaces algébriques au-dessus de $S$.
Soit $\Delta_{X/Y} : X \to X \times_Y X$ le morphisme diagonal.
\begin{enumerate}
\item On dit que $f$ est {\it séparé} si $\Delta_{X/Y}$ est une immersion fermée.
\item On dit que $f$ est {\it localement séparé}\footnote{Dans la littérature,
ce terme désigne souvent les morphismes quasi-séparés et localement séparés.}
si $\Delta_{X/Y}$ est une immersion.
\item On dit que $f$ est {\it quasi-séparé} si $\Delta_{X/Y}$ est quasi-compact.
\end{enumerate}
\end{definition}

\noindent
Cette définition a un sens puisque $\Delta_{X/Y}$ est représentable;
nous savons donc ce que signifie pour lui posséder l'une des propriétés
décrites dans la définition. Nous verrons plus bas
(lemme \ref{lemma-match-separated}) que cette définition coïncide avec celles
que nous avons déjà pour les morphismes de schémas et les morphismes représentables.

\begin{lemma}
\label{lemma-trivial-implications}
Soit $S$ un schéma. Soit $f : X \to Y$ un morphisme d'espaces algébriques
au-dessus de $S$. Si $f$ est séparé, alors $f$ est localement séparé et
$f$ est quasi-séparé.
\end{lemma}

\begin{proof}
Cela résulte du principe général de
Espaces,
lemme \ref{spaces-lemma-representable-transformations-property-implication},
puisqu'une immersion fermée de schémas est une immersion et est quasi-compacte.
\end{proof}

\begin{lemma}
\label{lemma-base-change-separated}
Tous les axiomes de séparation énumérés dans la définition \ref{definition-separated}
sont stables par changement de base.
\end{lemma}

\begin{proof}
Soient $f : X \to Y$ et $Y' \to Y$ des morphismes d'espaces algébriques.
Soit $f' : X' \to Y'$ le changement de base de $f$ par $Y' \to Y$. Alors
$\Delta_{X'/Y'}$ est le changement de base de $\Delta_{X/Y}$ par
le morphisme $X' \times_{Y'} X' \to X \times_Y X$. D'après les résultats de
la section \ref{section-representable},
chacune des propriétés de la diagonale utilisées dans la
définition \ref{definition-separated}
est stable par changement de base. D'où le lemme.
\end{proof}

\begin{lemma}
\label{lemma-fibre-product-after-map}
\begin{slogan}
La flèche supérieure d'un « diagramme magique » d'espaces algébriques possède
de bonnes propriétés du type immersion, qui se renforcent
sous des hypothèses de séparation.
\end{slogan}
Soit $S$ un schéma. Soient $f : X \to Z$, $g : Y \to Z$ et $Z \to T$
des morphismes d'espaces algébriques au-dessus de $S$. Considérons le morphisme induit
$i : X \times_Z Y \to X \times_T Y$. Alors
\begin{enumerate}
\item $i$ est représentable, localement de type fini, localement quasi-fini,
séparé et est un monomorphisme,
\item si $Z \to T$ est localement séparé, alors $i$ est une immersion,
\item si $Z \to T$ est séparé, alors $i$ est une immersion fermée, et
\item si $Z \to T$ est quasi-séparé, alors $i$ est quasi-compact.
\end{enumerate}
\end{lemma}

\begin{proof}
Par la théorie générale des catégories, le diagramme suivant
$$
\xymatrix{
X \times_Z Y \ar[r]_i \ar[d] & X \times_T Y \ar[d] \\
Z \ar[r]^-{\Delta_{Z/T}} \ar[r] & Z \times_T Z
}
$$
est un diagramme de produit fibré. Ainsi, $i$ est le changement de base du
morphisme diagonal $\Delta_{Z/T}$. Le lemme résulte donc
du lemme \ref{lemma-properties-diagonal} et des résultats de la
section \ref{section-representable}.
\end{proof}

\begin{lemma}
\label{lemma-semi-diagonal}
\begin{slogan}
Propriétés du graphe d'un morphisme d'espaces algébriques
déduites des propriétés de séparation du but.
\end{slogan}
Soit $S$ un schéma. Soit $T$ un espace algébrique au-dessus de $S$.
Soit $g : X \to Y$ un morphisme d'espaces algébriques au-dessus de $T$.
Considérons le morphisme graphe $i : X \to X \times_T Y$ de $g$. Alors
\begin{enumerate}
\item $i$ est représentable, localement de type fini, localement quasi-fini,
séparé et est un monomorphisme,
\item si $Y \to T$ est localement séparé, alors $i$ est une immersion,
\item si $Y \to T$ est séparé, alors $i$ est une immersion fermée, et
\item si $Y \to T$ est quasi-séparé, alors $i$ est quasi-compact.
\end{enumerate}
\end{lemma}

\begin{proof}
C'est le cas particulier du lemme \ref{lemma-fibre-product-after-map}
obtenu en l'appliquant au morphisme $X = X \times_Y Y \to X \times_T Y$.
\end{proof}

\begin{lemma}
\label{lemma-section-immersion}
Soit $S$ un schéma.
Soit $f : X \to T$ un morphisme d'espaces algébriques au-dessus de $S$.
Soit $s : T \to X$ une section de $f$ (autrement dit,
$f \circ s = \text{id}_T$). Alors
\begin{enumerate}
\item $s$ est représentable, localement de type fini, localement quasi-fini,
séparé et est un monomorphisme,
\item si $f$ est localement séparé, alors $s$ est une immersion,
\item si $f$ est séparé, alors $s$ est une immersion fermée, et
\item si $f$ est quasi-séparé, alors $s$ est quasi-compact.
\end{enumerate}
\end{lemma}

\begin{proof}
C'est le cas particulier du lemme \ref{lemma-semi-diagonal} appliqué à
$g = s$, le morphisme considéré étant $i = s : T \to T \times_T X$.
\end{proof}

\begin{lemma}
\label{lemma-composition-separated}
Tous les axiomes de séparation énumérés dans la définition \ref{definition-separated}
sont stables par composition des morphismes.
\end{lemma}

\begin{proof}
Soient $f : X \to Y$ et $g : Y \to Z$ des morphismes d'espaces algébriques
auxquels s'applique l'axiome considéré.
Le morphisme diagonal $\Delta_{X/Z}$ est la composée
$$
X \longrightarrow X \times_Y X \longrightarrow X \times_Z X.
$$
Notre axiome de séparation consiste à imposer à la diagonale
une certaine propriété $\mathcal{P}$. D'après le
lemme \ref{lemma-fibre-product-after-map} ci-dessus, nous voyons que
la seconde flèche possède aussi cette propriété. Le lemme en résulte,
puisque la composée de morphismes (représentables) possédant la propriété
$\mathcal{P}$ possède encore la propriété $\mathcal{P}$, voir
la section \ref{section-representable}.
\end{proof}

\begin{lemma}
\label{lemma-separated-over-separated}
Soit $S$ un schéma.
Soit $f : X \to Y$ un morphisme d'espaces algébriques au-dessus de $S$.
\begin{enumerate}
\item Si $Y$ est séparé et $f$ séparé, alors $X$ est séparé.
\item Si $Y$ est quasi-séparé et $f$ quasi-séparé, alors
$X$ est quasi-séparé.
\item Si $Y$ est localement séparé et $f$ localement séparé, alors
$X$ est localement séparé.
\item Si $Y$ est séparé au-dessus de $S$ et $f$ séparé, alors
$X$ est séparé au-dessus de $S$.
\item Si $Y$ est quasi-séparé au-dessus de $S$ et $f$ quasi-séparé, alors
$X$ est quasi-séparé au-dessus de $S$.
\item Si $Y$ est localement séparé au-dessus de $S$ et $f$ localement séparé, alors
$X$ est localement séparé au-dessus de $S$.
\end{enumerate}
\end{lemma}

\begin{proof}
Les assertions (4), (5) et (6) résultent immédiatement du
lemme \ref{lemma-composition-separated}
et de
Espaces, définition \ref{spaces-definition-separated}.
Les assertions (1), (2) et (3) se ramènent aux assertions (4), (5) et (6) en considérant
$X$ et $Y$ comme des espaces algébriques au-dessus de $\Spec(\mathbf{Z})$, voir
Propriétés des espaces, définition \ref{spaces-properties-definition-separated}.
\end{proof}

\begin{lemma}
\label{lemma-compose-after-separated}
Soit $S$ un schéma.
Soient $f : X \to Y$ et $g : Y \to Z$ des morphismes d'espaces algébriques au-dessus de $S$.
\begin{enumerate}
\item Si $g \circ f$ est séparé, alors $f$ l'est aussi.
\item Si $g \circ f$ est localement séparé, alors $f$ l'est aussi.
\item Si $g \circ f$ est quasi-séparé, alors $f$ l'est aussi.
\end{enumerate}
\end{lemma}

\begin{proof}
Considérons la factorisation
$$
X \to X \times_Y X \to X \times_Z X
$$
du morphisme diagonal de $g \circ f$. Dans tous les cas, le dernier morphisme
est un monomorphisme. Ainsi, pour tout schéma $T$ et tout morphisme
$T \to X \times_Y X$, on a l'égalité
$$
X \times_{(X \times_Y X)} T = X \times_{(X \times_Z X)} T.
$$
Le résultat est donc clair.
\end{proof}

\begin{lemma}
\label{lemma-separated-implies-morphism-separated}
Soit $S$ un schéma. Soit $X$ un espace algébrique au-dessus de $S$.
\begin{enumerate}
\item Si $X$ est séparé, alors $X$ est séparé au-dessus de $S$.
\item Si $X$ est localement séparé, alors $X$ est localement séparé au-dessus de $S$.
\item Si $X$ est quasi-séparé, alors $X$ est quasi-séparé au-dessus de $S$.
\end{enumerate}
Soit $f : X \to Y$ un morphisme d'espaces algébriques au-dessus de $S$.
\begin{enumerate}
\item[(4)] Si $X$ est séparé au-dessus de $S$, alors $f$ est séparé.
\item[(5)] Si $X$ est localement séparé au-dessus de $S$, alors $f$ est localement séparé.
\item[(6)] Si $X$ est quasi-séparé au-dessus de $S$, alors $f$ est quasi-séparé.
\end{enumerate}
\end{lemma}

\begin{proof}
Les assertions (4), (5) et (6) résultent immédiatement du
lemme \ref{lemma-compose-after-separated}
et de
Espaces, définition \ref{spaces-definition-separated}.
Les assertions (1), (2) et (3) résultent des assertions (4), (5) et (6) en
considérant $X$ et $Y$ comme des espaces algébriques au-dessus de
$\Spec(\mathbf{Z})$, voir
Propriétés des espaces, définition \ref{spaces-properties-definition-separated}.
\end{proof}

\begin{lemma}
\label{lemma-separated-local}
Soit $S$ un schéma.
Soit $f : X \to Y$ un morphisme d'espaces algébriques sur $S$.
Soit $\mathcal{P}$ l'un quelconque des axiomes de
séparation de la définition \ref{definition-separated}.
Les assertions suivantes sont équivalentes :
\begin{enumerate}
\item $f$ vérifie $\mathcal{P}$,
\item pour tout schéma $Z$ et tout morphisme $Z \to Y$, le
changement de base $Z \times_Y X \to Z$ de $f$ vérifie $\mathcal{P}$,
\item pour tout schéma affine $Z$ et tout morphisme $Z \to Y$, le
changement de base $Z \times_Y X \to Z$ de $f$ vérifie $\mathcal{P}$,
\item pour tout schéma affine $Z$ et tout morphisme $Z \to Y$,
l'espace algébrique $Z \times_Y X$ vérifie $\mathcal{P}$ (voir
Propriétés des espaces, définition \ref{spaces-properties-definition-separated}),
\item il existe un schéma $V$ et un morphisme \'etale surjectif
$V \to Y$ tels que le changement de base $V \times_Y X \to V$ vérifie
$\mathcal{P}$, et
\item il existe un recouvrement ouvert $Y = \bigcup Y_i$ tel que chacun
des morphismes $f^{-1}(Y_i) \to Y_i$ vérifie $\mathcal{P}$.
\end{enumerate}
\end{lemma}

\begin{proof}
Nous utiliserons à plusieurs reprises le
lemme \ref{lemma-base-change-separated}
sans le mentionner davantage. En particulier, il est clair que
(1) implique (2) et que (2) implique (3).

\medskip\noindent
Montrons que (3) et (4) sont équivalentes. Remarquons que si $Z$ est un schéma
affine, alors le morphisme $Z \to \Spec(\mathbf{Z})$ est séparé
en tant que morphisme d'espaces algébriques sur $\Spec(\mathbf{Z})$.
Si $Z \times_Y X \to Z$ vérifie $\mathcal{P}$, alors
$Z \times_Y X \to \Spec(\mathbf{Z})$ vérifie $\mathcal{P}$
comme composé (voir le
lemme \ref{lemma-composition-separated}). Par conséquent, l'espace
algébrique $Z \times_Y X$ vérifie $\mathcal{P}$. Réciproquement, si l'espace
algébrique $Z \times_Y X$ vérifie $\mathcal{P}$, alors
$Z \times_Y X \to \Spec(\mathbf{Z})$ vérifie $\mathcal{P}$, et
donc, d'après le
lemme \ref{lemma-compose-after-separated},
$Z \times_Y X \to Z$ vérifie $\mathcal{P}$.

\medskip\noindent
Montrons que (3) implique (5). Supposons (3). Soit $V$ un schéma,
et soit $V \to Y$ \'etale surjectif. Nous devons montrer que
$V \times_Y X \to V$ vérifie la propriété $\mathcal{P}$. Autrement dit,
nous devons montrer que le morphisme
$$
V \times_Y X \longrightarrow
(V \times_Y X) \times_V (V \times_Y X) = V \times_Y X \times_Y X
$$
possède la propriété correspondante (c'est-à-dire qu'il est une immersion fermée,
une immersion, ou qu'il est quasi-compact). Soit $V = \bigcup V_j$ un
recouvrement de $V$ par des ouverts affines. Par hypothèse, chacun des morphismes
$$
V_j \times_Y X \longrightarrow V_j \times_Y X \times_Y X
$$
possède bien la propriété correspondante. Puisque les propriétés d'être une immersion fermée,
une immersion, une immersion quasi-compacte ou un morphisme quasi-compact sont locales
sur le but pour la topologie de Zariski, et puisque les $V_j$ recouvrent $V$, nous obtenons la conclusion voulue.

\medskip\noindent
Montrons que (5) implique (1). Soit $V \to Y$ comme en (5).
Nous avons alors le diagramme de produit fibré
$$
\xymatrix{
V \times_Y X \ar[r] \ar[d] &
X \ar[d] \\
V \times_Y X \times_Y X \ar[r] &
X \times_Y X
}
$$
Par hypothèse, la flèche verticale de gauche est une immersion fermée,
une immersion, une immersion quasi-compacte, ou est quasi-compacte. Il résulte du
lemme \ref{spaces-lemma-descent-representable-transformations-property} du chapitre Espaces
que la flèche verticale de droite est également une immersion fermée,
une immersion, une immersion quasi-compacte, ou est quasi-compacte.

\medskip\noindent
Il est clair que (1) implique (6) en prenant le recouvrement $Y = Y$.
Supposons que $Y = \bigcup Y_i$ soit comme en (6). Choisissons des schémas $V_i$ et
des morphismes \'etales surjectifs $V_i \to Y_i$. Notons que les morphismes
$V_i \times_Y X \to V_i$ vérifient $\mathcal{P}$, puisqu'ils sont des changements de base
des morphismes $f^{-1}(Y_i) \to Y_i$. Posons $V = \coprod V_i$.
Alors $V \to Y$ est un morphisme comme en (5) (détails omis). Ainsi,
(6) implique (5), ce qui termine la démonstration.
\end{proof}

\begin{lemma}
\label{lemma-match-separated}
Soit $S$ un schéma.
Soit $f : X \to Y$ un morphisme représentable d'espaces algébriques sur $S$.
\begin{enumerate}
\item Le morphisme $f$ est localement séparé.
\item Le morphisme $f$ est (quasi-)séparé au sens de la
définition \ref{definition-separated}
ci-dessus si et seulement si $f$ est (quasi-)séparé au sens de la
section \ref{section-representable}.
\end{enumerate}
En particulier, si $f : X \to Y$ est un morphisme de schémas sur $S$, alors
$f$ est (quasi-)séparé au sens de la
définition \ref{definition-separated}
si et seulement si $f$ est (quasi-)séparé en tant que morphisme de schémas.
\end{lemma}

\begin{proof}
C'est l'équivalence de (1) et (2) du
lemme \ref{lemma-separated-local},
combinée au fait que tout morphisme de schémas est localement séparé ; voir le
lemme \ref{schemes-lemma-diagonal-immersion} du chapitre Schémas.
\end{proof}










\section{Morphismes surjectifs}
\label{section-surjective}

\noindent
Nous avons déjà défini à la section \ref{section-representable}
ce que signifie, pour un morphisme représentable d'espaces algébriques,
d'être surjectif.

\begin{lemma}
\label{lemma-surjective-representable}
Soit $S$ un schéma. Soit $f : X \to Y$ un morphisme représentable
d'espaces algébriques sur $S$. Alors
$f$ est surjectif (au sens de la section \ref{section-representable})
si et seulement si l'application $|f| : |X| \to |Y|$ est surjective.
\end{lemma}

\begin{proof}
En effet, si $f : X \to Y$ est représentable, il est surjectif si et seulement si,
pour tout schéma $T$ et tout morphisme $T \to Y$, le changement de base
$f_T : T \times_Y X \to T$ de $f$ est un morphisme surjectif de schémas,
autrement dit, si et seulement si $|f_T|$ est surjective. D'après le
lemme \ref{spaces-properties-lemma-points-cartesian} du chapitre Propriétés des espaces,
l'application $|T \times_Y X| \to |T| \times_{|Y|} |X|$ est toujours surjective.
Par conséquent, $|f_T| : |T \times_Y X| \to |T|$ est surjective si $|f| : |X| \to |Y|$
est surjective. Réciproquement, si $|f_T|$ est surjective
pour tout $T \to Y$ comme ci-dessus, alors, en prenant pour $T$ le spectre d'un
corps, nous concluons que $|X| \to |Y|$ est surjective.
\end{proof}

\noindent
Cela nous permet de poser la définition suivante.

\begin{definition}
\label{definition-surjective}
Soit $S$ un schéma. Soit $f : X \to Y$ un morphisme d'espaces
algébriques sur $S$. Nous disons que $f$ est {\it surjectif}
si l'application $|f| : |X| \to |Y|$ induite sur les espaces topologiques sous-jacents
est surjective.
\end{definition}

\begin{lemma}
\label{lemma-surjective-local}
Soit $S$ un schéma.
Soit $f : X \to Y$ un morphisme d'espaces algébriques sur $S$.
Les assertions suivantes sont équivalentes :
\begin{enumerate}
\item $f$ est surjectif,
\item pour tout schéma $Z$ et tout morphisme $Z \to Y$, le morphisme
$Z \times_Y X \to Z$ est surjectif,
\item pour tout schéma affine $Z$ et tout morphisme
$Z \to Y$, le morphisme $Z \times_Y X \to Z$ est surjectif,
\item il existe un schéma $V$ et un morphisme \'etale surjectif
$V \to Y$ tels que $V \times_Y X \to V$ soit un morphisme surjectif,
\item il existe un schéma $U$ et un morphisme \'etale surjectif
$\varphi : U \to X$ tels que le composé $f \circ \varphi$
soit surjectif,
\item il existe un diagramme commutatif
$$
\xymatrix{
U \ar[d] \ar[r] & V \ar[d] \\
X \ar[r] & Y
}
$$
où $U$, $V$ sont des schémas, les flèches verticales étant \'etales surjectives
et la flèche horizontale supérieure étant surjective, et
\item il existe un recouvrement ouvert $Y = \bigcup Y_i$ tel que
chacun des morphismes $f^{-1}(Y_i) \to Y_i$ est surjectif.
\end{enumerate}
\end{lemma}

\begin{proof}
Démonstration omise.
\end{proof}

\begin{lemma}
\label{lemma-composition-surjective}
Le composé de morphismes surjectifs est surjectif.
\end{lemma}

\begin{proof}
Cela résulte immédiatement de la définition.
\end{proof}

\begin{lemma}
\label{lemma-base-change-surjective}
Tout changement de base d'un morphisme surjectif est surjectif.
\end{lemma}

\begin{proof}
Cela résulte immédiatement du
lemme \ref{spaces-properties-lemma-points-cartesian} du chapitre Propriétés des espaces.
\end{proof}










\section{Morphismes ouverts}
\label{section-open}

\noindent
Pour un morphisme représentable d'espaces algébriques, nous avons déjà défini (à la
section \ref{section-representable})
ce que signifie être universellement ouvert. Avant de donner la définition
naturelle, vérifions donc qu'elle coïncide avec celle-ci dans le cas représentable.

\begin{lemma}
\label{lemma-characterize-representable-universally-open}
Soit $S$ un schéma. Soit $f : X \to Y$ un morphisme représentable
d'espaces algébriques sur $S$. Les assertions suivantes sont équivalentes :
\begin{enumerate}
\item $f$ est universellement ouvert
(au sens de la section \ref{section-representable}), et
\item pour tout morphisme d'espaces algébriques $Z \to Y$, l'application
$|Z \times_Y X| \to |Z|$ entre espaces topologiques est ouverte.
\end{enumerate}
\end{lemma}

\begin{proof}
Supposons (1), et soit $Z \to Y$ comme en (2). Choisissons un schéma $V$ et
un morphisme \'etale surjectif $V \to Y$. Par hypothèse, le morphisme de
schémas $V \times_Y X \to V$ est universellement ouvert. D'après la
section \ref{spaces-properties-section-points} du chapitre Propriétés des espaces,
dans le diagramme commutatif
$$
\xymatrix{
|V \times_Y X| \ar[r] \ar[d] & |Z \times_Y X| \ar[d] \\
|V| \ar[r] & |Z|
}
$$
les flèches horizontales sont ouvertes et surjectives ; de plus,
$$
|V \times_Y X| \longrightarrow |V| \times_{|Z|} |Z \times_Y X|
$$
est surjective. Puisque la flèche verticale de gauche
est ouverte, il s'ensuit que celle de droite est
ouverte. Cela démontre (2). L'implication (2) $\Rightarrow$ (1) résulte
immédiatement des définitions.
\end{proof}

\noindent
Nous pouvons donc adopter la définition naturelle suivante.

\begin{definition}
\label{definition-open}
Soit $S$ un schéma. Soit $f : X \to Y$ un morphisme d'espaces algébriques
sur $S$.
\begin{enumerate}
\item Nous disons que $f$ est {\it ouvert} si l'application entre espaces topologiques
$|f| : |X| \to |Y|$ est ouverte.
\item Nous disons que $f$ est {\it universellement ouvert} si, pour tout morphisme
d'espaces algébriques $Z \to Y$, l'application entre espaces topologiques
$$
|Z \times_Y X| \to |Z|
$$
est ouverte, c'est-à-dire si le changement de base $Z \times_Y X \to Z$ est ouvert.
\end{enumerate}
\end{definition}

\noindent
Notons qu'un morphisme \'etale d'espaces algébriques est universellement ouvert ;
voir la
définition \ref{spaces-properties-definition-etale} du chapitre Propriétés des espaces et les
lemmes \ref{spaces-properties-lemma-etale-open} et
\ref{spaces-properties-lemma-base-change-etale}.

\begin{lemma}
\label{lemma-base-change-universally-open}
Tout changement de base d'un morphisme universellement ouvert d'espaces algébriques
par un morphisme quelconque d'espaces algébriques est universellement ouvert.
\end{lemma}

\begin{proof}
Cela résulte immédiatement de la définition.
\end{proof}

\begin{lemma}
\label{lemma-composition-open}
Le composé de deux morphismes (universellement) ouverts d'espaces algébriques
est (universellement) ouvert.
\end{lemma}

\begin{proof}
Démonstration omise.
\end{proof}

\begin{lemma}
\label{lemma-universally-open-local}
Soit $S$ un schéma. Soit $f : X \to Y$ un morphisme d'espaces algébriques
sur $S$. Les assertions suivantes sont équivalentes :
\begin{enumerate}
\item $f$ est universellement ouvert,
\item pour tout schéma $Z$ et tout morphisme $Z \to Y$,
la projection $|Z \times_Y X| \to |Z|$ est ouverte,
\item pour tout schéma affine $Z$ et tout morphisme $Z \to Y$,
la projection $|Z \times_Y X| \to |Z|$ est ouverte, et
\item il existe un schéma $V$ et un morphisme \'etale surjectif
$V \to Y$ tels que $V \times_Y X \to V$ soit un morphisme universellement ouvert
d'espaces algébriques, et
\item il existe un recouvrement ouvert $Y = \bigcup Y_i$ tel que
chacun des morphismes $f^{-1}(Y_i) \to Y_i$ soit universellement ouvert.
\end{enumerate}
\end{lemma}

\begin{proof}
Nous omettons de démontrer que (1) implique (2) et que (2) implique (3).

\medskip\noindent
Supposons (3). Choisissons un morphisme \'etale surjectif $V \to Y$.
Nous allons montrer que $V \times_Y X \to V$ est un morphisme
universellement ouvert d'espaces algébriques. Soit $Z \to V$ un morphisme
d'un espace algébrique vers $V$. Soit $W \to Z$ un morphisme \'etale surjectif,
où $W = \coprod W_i$ est une somme disjointe de schémas affines ; voir le
lemme
\ref{spaces-properties-lemma-cover-by-union-affines} du chapitre Propriétés des espaces.
Nous avons alors le diagramme commutatif suivant :
$$
\xymatrix{
\coprod_i |W_i \times_Y X| \ar@{=}[r] \ar[d] &
|W \times_Y X| \ar[r] \ar[d] &
|Z \times_Y X| \ar[d] \ar@{=}[r] &
|Z \times_V (V \times_Y X)| \ar[ld] \\
\coprod |W_i| \ar@{=}[r] &
|W| \ar[r] &
|Z|
}
$$
Nous devons montrer que la flèche sud-est est ouverte. Les flèches horizontales
centrales sont surjectives et ouvertes
(lemme \ref{spaces-properties-lemma-etale-open} du chapitre Propriétés des espaces).
D'après l'hypothèse (3) et le fait que
$W_i$ est affine, les flèches verticales de gauche sont ouvertes. Par conséquent,
la flèche verticale de droite est ouverte.

\medskip\noindent
Supposons que $V \to Y$ soit comme en (4). Montrons que $f$ est universellement ouvert.
Soit $Z \to Y$ un morphisme d'espaces algébriques. Considérons le
diagramme
$$
\xymatrix{
|(V \times_Y Z) \times_V (V \times_Y X)| \ar@{=}[r] \ar[rd] &
|V \times_Y X| \ar[r] \ar[d] &
|Z \times_Y X| \ar[d] \\
 &
|V \times_Y Z| \ar[r] &
|Z|
}
$$
La flèche sud-ouest est ouverte par hypothèse. Les flèches horizontales sont
surjectives et ouvertes, car les morphismes correspondants
d'espaces algébriques sont \'etales (voir le
lemme \ref{spaces-properties-lemma-etale-open} du chapitre Propriétés des espaces).
Il s'ensuit que la flèche verticale de droite est ouverte.

\medskip\noindent
Bien entendu, (1) implique (5) en prenant le recouvrement $Y = Y$.
Supposons que $Y = \bigcup Y_i$ soit comme en (5). Alors, pour tout $Z \to Y$,
nous obtenons un recouvrement ouvert correspondant $Z = \bigcup Z_i$ tel que
le changement de base de $f$ à $Z_i$ soit ouvert. Un argument topologique élémentaire
montre alors que $Z \times_Y X \to Z$ est ouvert. Par conséquent, (1) est vérifiée.
\end{proof}

\begin{lemma}
\label{lemma-space-over-field-universally-open}
Soit $S$ un schéma. Soit $p : X \to \Spec(k)$ un morphisme
d'espaces algébriques sur $S$, où $k$ est un corps. Alors
$p : X \to \Spec(k)$ est universellement ouvert.
\end{lemma}

\begin{proof}
Choisissons un schéma $U$ et un morphisme \'etale surjectif $U \to X$.
Le composé $U \to \Spec(k)$ est universellement ouvert (en tant que morphisme
de schémas) d'après le
lemme \ref{morphisms-lemma-scheme-over-field-universally-open} du chapitre Morphismes.
Soit $Z \to \Spec(k)$ un morphisme de schémas. Alors
$U \times_{\Spec(k)} Z \to X \times_{\Spec(k)} Z$ est surjectif ;
voir le
lemme \ref{lemma-base-change-surjective}.
Par conséquent, la première des applications
$$
|U \times_{\Spec(k)} Z| \to |X \times_{\Spec(k)} Z| \to |Z|
$$
est surjective. Comme leur composé est ouvert d'après ce qui précède, nous en déduisons que
la seconde application est elle aussi ouverte. Ainsi, $p$ est universellement ouvert d'après le
lemme \ref{lemma-universally-open-local}.
\end{proof}








\section{Morphismes submersifs}
\label{section-submersive}

\noindent
Pour un morphisme représentable d'espaces algébriques, nous avons déjà défini (à la
section \ref{section-representable})
ce que signifie être universellement submersif. Avant de donner la définition
naturelle, vérifions donc qu'elle coïncide avec celle-ci dans le cas représentable.

\begin{lemma}
\label{lemma-characterize-representable-universally-submersive}
Soit $S$ un schéma. Soit $f : X \to Y$ un morphisme représentable
d'espaces algébriques sur $S$. Les assertions suivantes sont équivalentes :
\begin{enumerate}
\item $f$ est universellement submersif
(au sens de la section \ref{section-representable}), et
\item pour tout morphisme d'espaces algébriques $Z \to Y$, l'application
$|Z \times_Y X| \to |Z|$ entre espaces topologiques est submersive.
\end{enumerate}
\end{lemma}

\begin{proof}
Supposons (1), et soit $Z \to Y$ comme en (2). Choisissons un schéma $V$ et
un morphisme \'etale surjectif $V \to Y$. Par hypothèse, le morphisme de
schémas $V \times_Y X \to V$ est universellement submersif. D'après la
section \ref{spaces-properties-section-points} du chapitre Propriétés des espaces,
dans le diagramme commutatif
$$
\xymatrix{
|V \times_Y X| \ar[r] \ar[d] & |Z \times_Y X| \ar[d] \\
|V| \ar[r] & |Z|
}
$$
les flèches horizontales sont ouvertes et surjectives ; de plus,
$$
|V \times_Y X| \longrightarrow |V| \times_{|Z|} |Z \times_Y X|
$$
est surjective. Puisque la flèche verticale de gauche est submersive,
il s'ensuit que celle de droite est submersive.
Cela démontre (2). L'implication (2) $\Rightarrow$ (1) résulte
immédiatement des définitions.
\end{proof}

\noindent
Nous pouvons donc adopter la définition naturelle suivante.

\begin{definition}
\label{definition-submersive}
Soit $S$ un schéma.
Soit $f : X \to Y$ un morphisme d'espaces algébriques sur $S$.
\begin{enumerate}
\item Nous disons que $f$ est {\it submersif}\footnote{Ceci diffère profondément
de la notion de submersion de variétés différentielles.}
si l'application continue $|X| \to |Y|$ est submersive ; voir la
définition \ref{topology-definition-submersive} du chapitre Topologie.
\item Nous disons que $f$ est {\it universellement submersif} si, pour tout
morphisme d'espaces algébriques $Y' \to Y$, le changement de base
$Y' \times_Y X \to Y'$ est submersif.
\end{enumerate}
\end{definition}

\noindent
Remarquons qu'un morphisme submersif est en particulier surjectif.

\begin{lemma}
\label{lemma-base-change-universally-submersive}
Tout changement de base d'un morphisme universellement submersif d'espaces algébriques
par un morphisme quelconque d'espaces algébriques est universellement submersif.
\end{lemma}

\begin{proof}
Cela résulte immédiatement de la définition.
\end{proof}

\begin{lemma}
\label{lemma-composition-universally-submersive}
Le composé de deux morphismes (universellement) submersifs
d'espaces algébriques est (universellement) submersif.
\end{lemma}

\begin{proof}
Démonstration omise.
\end{proof}












\section{Morphismes quasi-compacts}
\label{section-quasi-compact}

\noindent
D'après la section \ref{section-representable}, nous savons ce que signifie
pour un morphisme représentable d'espaces algébriques d'être quasi-compact.
Afin de formuler la définition pour un morphisme quelconque
d'espaces algébriques, faisons l'observation suivante.

\begin{lemma}
\label{lemma-characterize-representable-quasi-compact}
Soit $S$ un schéma.
Soit $f : X \to Y$ un morphisme représentable d'espaces algébriques sur $S$.
Les conditions suivantes sont équivalentes :
\begin{enumerate}
\item $f$ est quasi-compact
(au sens de la section \ref{section-representable}), et
\item pour tout espace algébrique quasi-compact $Z$ et tout morphisme
$Z \to Y$, l'espace algébrique $Z \times_Y X$ est quasi-compact.
\end{enumerate}
\end{lemma}

\begin{proof}
Supposons (1), et soit $Z \to Y$ un morphisme d'espaces algébriques tel que
$Z$ soit quasi-compact. D'après
Propriétés des espaces,
définition \ref{spaces-properties-definition-quasi-compact},
il existe un schéma quasi-compact $U$ et un morphisme \'etale surjectif
$U \to Z$. Puisque $f$ est représentable et quasi-compact,
la définition montre que $U \times_Y X$ est un schéma et que
$U \times_Y X \to U$ est quasi-compact. Ainsi $U \times_Y X$ est
un schéma quasi-compact. Le morphisme $U \times_Y X \to Z \times_Y X$
est \'etale et surjectif (car c'est le changement de base du morphisme
représentable, \'etale et surjectif $U \to Z$ ; voir la
section \ref{section-representable}).
Par définition, $Z \times_Y X$ est donc quasi-compact.

\medskip\noindent
Supposons (2). Soit $Z \to Y$ un morphisme, où $Z$ est un schéma.
Il faut montrer que $p : Z \times_Y X \to Z$ est quasi-compact.
Soit $U \subset Z$ un ouvert affine. Alors $p^{-1}(U) = U \times_Y X$
et le schéma $U \times_Y X$ est quasi-compact par l'hypothèse (2).
Ainsi $p$ est quasi-compact ; voir
Schémas, section \ref{schemes-section-quasi-compact}.
\end{proof}

\noindent
Ceci motive la définition suivante.

\begin{definition}
\label{definition-quasi-compact}
Soit $S$ un schéma.
Soit $f : X \to Y$ un morphisme d'espaces algébriques sur $S$.
Nous disons que $f$ est {\it quasi-compact} si, pour tout espace
algébrique quasi-compact $Z$ et tout morphisme $Z \to Y$, le produit fibré
$Z \times_Y X$ est quasi-compact.
\end{definition}

\noindent
D'après le lemme \ref{lemma-characterize-representable-quasi-compact}
ci-dessus, cette définition coïncide avec la notion déjà existante
pour les morphismes représentables d'espaces algébriques.

\begin{lemma}
\label{lemma-quasi-compact-is-quasi-compact}
Soit $S$ un schéma. Si $f : X \to Y$ est un morphisme quasi-compact
d'espaces algébriques sur $S$, alors l'application sous-jacente
$|f| : |X| \to |Y|$ d'espaces topologiques est quasi-compacte.
\end{lemma}

\begin{proof}
Soit $V \subset |Y|$ un ouvert quasi-compact. D'après Propriétés des espaces,
lemme \ref{spaces-properties-lemma-open-subspaces},
il existe un sous-espace ouvert $Y' \subset Y$ tel que $V = |Y'|$.
Alors $Y'$ est un espace algébrique quasi-compact d'après
Propriétés des espaces, lemme \ref{spaces-properties-lemma-quasi-compact-space},
et donc $X' = Y' \times_Y X$ est un espace algébrique
quasi-compact d'après la définition \ref{definition-quasi-compact}.
D'autre part, $X' \subset X$ est un sous-espace ouvert
(Espaces, lemme \ref{spaces-lemma-base-change-immersions})
et $|X'| = |f|^{-1}(V)$ d'après
Propriétés des espaces, lemme \ref{spaces-properties-lemma-points-cartesian}.
En appliquant de nouveau
Propriétés des espaces, lemme \ref{spaces-properties-lemma-quasi-compact-space},
on conclut que $|X'|$ est un ouvert quasi-compact de $|X|$, comme souhaité.
\end{proof}

\begin{lemma}
\label{lemma-base-change-quasi-compact}
Tout changement de base d'un morphisme quasi-compact d'espaces algébriques
par un morphisme quelconque d'espaces algébriques est quasi-compact.
\end{lemma}

\begin{proof}
Démonstration omise. Indication : transitivité des produits fibrés.
\end{proof}

\begin{lemma}
\label{lemma-composition-quasi-compact}
Le composé de deux morphismes quasi-compacts d'espaces algébriques
est quasi-compact.
\end{lemma}

\begin{proof}
Démonstration omise. Indication : transitivité des produits fibrés.
\end{proof}

\begin{lemma}
\label{lemma-surjection-from-quasi-compact}
\begin{slogan}
L'image d'un espace algébrique quasi-compact par un morphisme surjectif
est quasi-compacte.
\end{slogan}
Soit $S$ un schéma.
\begin{enumerate}
\item Si $X \to Y$ est un morphisme surjectif d'espaces algébriques sur $S$,
et si $X$ est quasi-compact, alors $Y$ est quasi-compact.
\item Si
$$
\xymatrix{
X \ar[rr]_f \ar[rd]_p & &
Y \ar[dl]^q \\
& Z
}
$$
est un diagramme commutatif de morphismes d'espaces algébriques sur $S$
et si $f$ est surjectif et $p$ quasi-compact, alors $q$ est quasi-compact.
\end{enumerate}
\end{lemma}

\begin{proof}
Supposons $X$ quasi-compact et $X \to Y$ surjectif. D'après la
définition \ref{definition-surjective},
l'application $|X| \to |Y|$ est surjective ; ainsi $Y$ est quasi-compact d'après
Propriétés des espaces, lemme \ref{spaces-properties-lemma-quasi-compact-space},
et le fait topologique que l'image d'un espace quasi-compact par une
application continue est quasi-compacte ; voir
Topologie, lemme \ref{topology-lemma-image-quasi-compact}.
Soient $f, p, q$ comme en (2).
Soit $T \to Z$ un morphisme dont la source est un espace algébrique quasi-compact.
Par hypothèse, $T \times_Z X$ est quasi-compact. D'après le
lemme \ref{lemma-base-change-surjective},
le morphisme $T \times_Z X \to T \times_Z Y$ est surjectif.
La partie (1) montre donc que $T \times_Z Y$ est également quasi-compact.
Ainsi $q$ est quasi-compact.
\end{proof}

\begin{lemma}
\label{lemma-descent-quasi-compact}
Soit $S$ un schéma.
Soit $f : X \to Y$ un morphisme d'espaces algébriques sur $S$.
Soit $g : Y' \to Y$ un morphisme universellement ouvert et surjectif
d'espaces algébriques tel que le changement de base $f' : X' \to Y'$ soit quasi-compact.
Alors $f$ est quasi-compact.
\end{lemma}

\begin{proof}
Soit $Z \to Y$ un morphisme d'espaces algébriques avec $Z$ quasi-compact.
Comme $g$ est universellement ouvert et surjectif, on voit que
$Y' \times_Y Z \to Z$ est ouvert et surjectif. Tout point de
$|Y' \times_Y Z|$ admettant un système fondamental de voisinages ouverts
quasi-compacts (voir
Propriétés des espaces,
lemme \ref{spaces-properties-lemma-space-locally-quasi-compact}),
on peut trouver un ouvert quasi-compact $W \subset |Y' \times_Y Z|$
dont l'image est $Z$. Notons
$f'' : W \times_Y X \to W$ le changement de base de $f'$ par $W \to Y'$.
Par hypothèse, $W \times_Y X$ est quasi-compact. Comme $W \to Z$ est surjectif,
on voit que $W \times_Y X \to Z \times_Y X$ est surjectif.
Ainsi $Z \times_Y X$ est quasi-compact d'après le
lemme \ref{lemma-surjection-from-quasi-compact}.
Ainsi $f$ est quasi-compact.
\end{proof}

\begin{lemma}
\label{lemma-quasi-compact-local}
\begin{slogan}
Les morphismes quasi-compacts d'espaces algébriques sont stables par
changement de base et cette propriété est locale sur la base.
\end{slogan}
Soit $S$ un schéma.
Soit $f : X \to Y$ un morphisme d'espaces algébriques sur $S$.
Les conditions suivantes sont équivalentes :
\begin{enumerate}
\item $f$ est quasi-compact,
\item pour tout schéma $Z$ et tout morphisme $Z \to Y$, le morphisme
d'espaces algébriques $Z \times_Y X \to Z$ est quasi-compact,
\item pour tout schéma affine $Z$ et tout morphisme
$Z \to Y$, l'espace algébrique $Z \times_Y X$ est quasi-compact,
\item il existe un schéma $V$ et un morphisme \'etale surjectif
$V \to Y$ tels que $V \times_Y X \to V$ soit un morphisme quasi-compact
d'espaces algébriques, et
\item il existe un morphisme \'etale surjectif
$Y' \to Y$ d'espaces algébriques tel que $Y' \times_Y X \to Y'$
soit un morphisme quasi-compact d'espaces algébriques, et
\item il existe un recouvrement de Zariski $Y = \bigcup Y_i$ tel que
chacun des morphismes $f^{-1}(Y_i) \to Y_i$ soit quasi-compact.
\end{enumerate}
\end{lemma}

\begin{proof}
Nous utiliserons le lemme \ref{lemma-base-change-quasi-compact}
sans plus le mentionner.
Il est clair que (1) implique (2) et que (2) implique (3).
Supposons (3). Soit $Z$ un espace algébrique quasi-compact sur $S$,
et soit $Z \to Y$ un morphisme. D'après
Propriétés des espaces, lemme
\ref{spaces-properties-lemma-quasi-compact-affine-cover},
il existe un schéma affine $U$ et un morphisme \'etale surjectif
$U \to Z$. Alors $U \times_Y X \to Z \times_Y X$ est un morphisme
surjectif d'espaces algébriques ; voir le
lemme \ref{lemma-base-change-surjective}.
Par hypothèse, $|U \times_Y X|$ est quasi-compact et son image
est $|Z \times_Y X|$ ; on en déduit que $|Z \times_Y X|$
est quasi-compact ; voir
Topologie, lemme \ref{topology-lemma-image-quasi-compact}.
Ceci prouve que (3) implique (1).

\medskip\noindent
Les implications (1) $\Rightarrow$ (4) et (4) $\Rightarrow$ (5) sont claires.
L'implication (5) $\Rightarrow$ (1) découle du
lemme \ref{lemma-descent-quasi-compact}
et du fait qu'un morphisme \'etale d'espaces algébriques est universellement ouvert
(voir la discussion qui suit la
définition \ref{definition-open}).

\medskip\noindent
Bien entendu, (1) implique (6) en prenant le recouvrement $Y = Y$.
Supposons que $Y = \bigcup Y_i$ soit comme en (6). Soit $Z$ affine et soit
$Z \to Y$ un morphisme. Il existe alors un recouvrement affine standard fini
$Z = Z_1 \cup \ldots \cup Z_n$ tel que chaque $Z_j \to Y$
se factorise par $Y_{i_j}$ pour un certain $i_j$. Ainsi, l'espace algébrique
$$
Z_j \times_Y X = Z_j \times_{Y_{i_j}} f^{-1}(Y_{i_j})
$$
est quasi-compact. Puisque
$Z \times_Y X = \bigcup_{j = 1, \ldots, n} Z_j \times_Y X$
est un recouvrement de Zariski, on voit que
$|Z \times_Y X| = \bigcup_{j = 1, \ldots, n} |Z_j \times_Y X|$
(voir Propriétés des espaces, lemme \ref{spaces-properties-lemma-open-subspaces})
est une réunion finie d'espaces quasi-compacts, donc est quasi-compact.
Ainsi, (6) implique (3).
\end{proof}

\noindent
Le lemme ci-dessous (ainsi que le suivant) garantit en particulier qu'un morphisme
$X \to \Spec(A)$ est quasi-compact dès que
$X$ est un espace algébrique quasi-compact.

\begin{lemma}
\label{lemma-quasi-compact-permanence}
Soit $S$ un schéma.
Soient $f : X \to Y$ et $g : Y \to Z$ des morphismes d'espaces algébriques sur $S$.
Si $g \circ f$ est quasi-compact et si $g$ est quasi-séparé,
alors $f$ est quasi-compact.
\end{lemma}

\begin{proof}
Ceci est vrai parce que $f$ est le composé
$(1, f) : X \to X \times_Z Y \to Y$. Le premier morphisme
est quasi-compact d'après le lemme \ref{lemma-section-immersion},
car c'est une section du morphisme quasi-séparé $X \times_Z Y \to X$
(un changement de base de $g$ ; voir le lemme \ref{lemma-base-change-separated}).
Le second morphisme est quasi-compact puisqu'il
est le changement de base de $g \circ f$ ; voir le
lemme \ref{lemma-base-change-quasi-compact}.
Enfin, les composés de morphismes quasi-compacts
sont quasi-compacts ; voir le lemme \ref{lemma-composition-quasi-compact}.
\end{proof}

\begin{lemma}
\label{lemma-quasi-compact-quasi-separated-permanence}
Soit $f : X \to Y$ un morphisme d'espaces algébriques
sur un schéma $S$.
\begin{enumerate}
\item Si $X$ est quasi-compact et $Y$ quasi-séparé, alors $f$ est
quasi-compact.
\item Si $X$ est quasi-compact et quasi-séparé et si $Y$ est quasi-séparé,
alors $f$ est quasi-compact et quasi-séparé.
\item Un produit fibré d'espaces algébriques quasi-compacts et quasi-séparés
est quasi-compact et quasi-séparé.
\end{enumerate}
\end{lemma}

\begin{proof}
La partie (1) résulte du
lemme \ref{lemma-quasi-compact-permanence}
avec $Z = S = \Spec(\mathbf{Z})$.
La partie (2) résulte de (1) et du
lemme \ref{lemma-compose-after-separated}.
Pour (3), soient $X \to Y$ et $Z \to Y$ des morphismes
d'espaces algébriques quasi-compacts et quasi-séparés. Alors
$X \times_Y Z \to Z$ est quasi-compact et quasi-séparé en tant que
changement de base de $X \to Y$, en vertu de (2) et des
lemmes \ref{lemma-base-change-quasi-compact} et
\ref{lemma-base-change-separated}.
Ainsi $X \times_Y Z$ est quasi-compact et quasi-séparé comme
espace algébrique quasi-compact et quasi-séparé sur
$Z$ ; voir les
lemmes \ref{lemma-separated-over-separated} et
\ref{lemma-composition-quasi-compact}.
\end{proof}


\section{Morphismes universellement fermés}
\label{section-universally-closed}

\noindent
Pour un morphisme représentable d'espaces algébriques, nous avons déjà défini (à la
section \ref{section-representable})
ce que signifie être universellement fermé. Avant de donner la définition
naturelle, vérifions donc qu'elle coïncide avec celle-ci dans le cas représentable.

\begin{lemma}
\label{lemma-characterize-representable-universally-closed}
Soit $S$ un schéma. Soit $f : X \to Y$ un morphisme représentable
d'espaces algébriques sur $S$. Les conditions suivantes sont équivalentes :
\begin{enumerate}
\item $f$ est universellement fermé
(au sens de la section \ref{section-representable}), et
\item pour tout morphisme d'espaces algébriques $Z \to Y$, le morphisme
d'espaces topologiques $|Z \times_Y X| \to |Z|$ est fermé.
\end{enumerate}
\end{lemma}

\begin{proof}
Supposons (1), et soit $Z \to Y$ comme en (2). Choisissons un schéma $V$ et
un morphisme \'etale surjectif $V \to Z$. Par hypothèse, le morphisme
de schémas $V \times_Y X \to V$ est universellement fermé. D'après
Propriétés des espaces, section \ref{spaces-properties-section-points},
dans le diagramme commutatif
$$
\xymatrix{
|V \times_Y X| \ar[r] \ar[d] & |Z \times_Y X| \ar[d] \\
|V| \ar[r] & |Z|
}
$$
les flèches horizontales sont ouvertes et surjectives et, de plus,
$$
|V \times_Y X| \longrightarrow |V| \times_{|Z|} |Z \times_Y X|
$$
est surjectif. Puisque la flèche verticale
de gauche est fermée, il s'ensuit que celle de droite est
fermée. Ceci prouve (2). L'implication (2) $\Rightarrow$ (1) résulte
immédiatement des définitions.
\end{proof}

\noindent
Nous pouvons donc adopter la définition naturelle suivante.

\begin{definition}
\label{definition-closed}
Soit $S$ un schéma. Soit $f : X \to Y$ un morphisme d'espaces algébriques
sur $S$.
\begin{enumerate}
\item Nous disons que $f$ est {\it fermé} si l'application d'espaces
topologiques $|X| \to |Y|$ est fermée.
\item Nous disons que $f$ est {\it universellement fermé} si, pour tout morphisme
d'espaces algébriques $Z \to Y$, le morphisme d'espaces topologiques
$$
|Z \times_Y X| \to |Z|
$$
est fermé, c'est-à-dire si le changement de base $Z \times_Y X \to Z$ est fermé.
\end{enumerate}
\end{definition}

\begin{lemma}
\label{lemma-base-change-universally-closed}
Tout changement de base d'un morphisme universellement fermé d'espaces algébriques
par un morphisme quelconque d'espaces algébriques est universellement fermé.
\end{lemma}

\begin{proof}
Cela résulte immédiatement de la définition.
\end{proof}

\begin{lemma}
\label{lemma-composition-universally-closed}
Le composé de deux morphismes (universellement) fermés d'espaces algébriques
est (universellement) fermé.
\end{lemma}

\begin{proof}
Démonstration omise.
\end{proof}

\begin{lemma}
\label{lemma-universally-closed-local}
Soit $S$ un schéma. Soit $f : X \to Y$ un morphisme d'espaces algébriques
sur $S$. Les conditions suivantes sont équivalentes :
\begin{enumerate}
\item $f$ est universellement fermé,
\item pour tout schéma $Z$ et tout morphisme $Z \to Y$,
la projection $|Z \times_Y X| \to |Z|$ est fermée,
\item pour tout schéma affine $Z$ et tout morphisme $Z \to Y$,
la projection $|Z \times_Y X| \to |Z|$ est fermée,
\item il existe un schéma $V$ et un morphisme \'etale surjectif
$V \to Y$ tel que $V \times_Y X \to V$ soit un morphisme universellement fermé
d'espaces algébriques, et
\item il existe un recouvrement de Zariski $Y = \bigcup Y_i$ tel que
chacun des morphismes $f^{-1}(Y_i) \to Y_i$ soit universellement fermé.
\end{enumerate}
\end{lemma}

\begin{proof}
Nous omettons la preuve que (1) implique (2) et que (2) implique (3).

\medskip\noindent
Supposons (3). Choisissons un morphisme \'etale surjectif $V \to Y$.
Nous allons montrer que $V \times_Y X \to V$ est un morphisme
universellement fermé d'espaces algébriques. Soit $Z \to V$ un morphisme
d'un espace algébrique vers $V$. Soit $W \to Z$ un morphisme \'etale surjectif
où $W = \coprod W_i$ est une somme disjointe de schémas affines ; voir
Propriétés des espaces,
lemme \ref{spaces-properties-lemma-cover-by-union-affines}.
Nous avons alors le diagramme commutatif suivant :
$$
\xymatrix{
\coprod_i |W_i \times_Y X| \ar@{=}[r] \ar[d] &
|W \times_Y X| \ar[r] \ar[d] &
|Z \times_Y X| \ar[d] \ar@{=}[r] &
|Z \times_V (V \times_Y X)| \ar[ld] \\
\coprod |W_i| \ar@{=}[r] &
|W| \ar[r] &
|Z|
}
$$
Il faut montrer que la flèche sud-est est fermée. Les flèches horizontales
centrales sont surjectives et ouvertes
(Propriétés des espaces, lemme \ref{spaces-properties-lemma-etale-open}).
L'hypothèse (3), jointe au fait que
$W_i$ est affine, montre que les flèches verticales de gauche sont fermées. Il
s'ensuit que la flèche verticale de droite est fermée.

\medskip\noindent
Supposons (4). Montrons que $f$ est universellement fermé.
Soit $Z \to Y$ un morphisme d'espaces algébriques. Considérons le
diagramme
$$
\xymatrix{
|(V \times_Y Z) \times_V (V \times_Y X)| \ar@{=}[r] \ar[rd] &
|V \times_Y X| \ar[r] \ar[d] &
|Z \times_Y X| \ar[d] \\
 &
|V \times_Y Z| \ar[r] &
|Z|
}
$$
La flèche sud-ouest est fermée par hypothèse. Les flèches horizontales sont
surjectives et ouvertes, car les morphismes correspondants d'espaces
algébriques sont \'etales (voir
Propriétés des espaces, lemme \ref{spaces-properties-lemma-etale-open}).
Il s'ensuit que la flèche verticale de droite est fermée.

\medskip\noindent
Bien entendu, (1) implique (5) en prenant le recouvrement $Y = Y$.
Supposons que $Y = \bigcup Y_i$ soit comme en (5). Pour tout $Z \to Y$,
on obtient alors un recouvrement de Zariski correspondant $Z = \bigcup Z_i$ tel que
le changement de base de $f$ à $Z_i$ soit fermé. Un argument topologique
élémentaire montre que $Z \times_Y X \to Z$ est fermé. Ainsi (1) est satisfaite.
\end{proof}

\begin{example}
\label{example-strange-universally-closed}
Exemple étrange d'un morphisme universellement fermé.
Soit $\mathbf{Q} \subset k$ un corps de caractéristique zéro.
Soit $X = \mathbf{A}^1_k/\mathbf{Z}$ comme dans
Espaces, exemple \ref{spaces-example-affine-line-translation}.
Nous affirmons que le morphisme structural $p : X \to \Spec(k)$
est universellement fermé.
En effet, si $Z/k$ est un schéma et si $T \subset |X \times_k Z|$ est fermé,
alors $T$ correspond à une partie fermée $\mathbf{Z}$-invariante
$T' \subset |\mathbf{A}^1 \times Z|$. Il est facile de voir que
cela implique que $T'$ est l'image inverse d'une partie $T''$ de
$Z$. D'après
Morphismes, lemme \ref{morphisms-lemma-fpqc-quotient-topology},
la partie $T'' \subset Z$ est fermée.
Bien entendu, $T''$ est l'image de $T$. Le morphisme $p$ est donc universellement
fermé d'après le lemme \ref{lemma-universally-closed-local}.
\end{example}

\begin{lemma}
\label{lemma-universally-closed-quasi-compact}
Soit $S$ un schéma.
Tout morphisme universellement fermé d'espaces algébriques sur $S$ est quasi-compact.
\end{lemma}

\begin{proof}
Cette démonstration reprend celle du cas des schémas ; voir
Morphismes, lemme \ref{morphisms-lemma-universally-closed-quasi-compact}.
Soit $f : X \to Y$ un morphisme d'espaces algébriques sur $S$.
Supposons que $f$ ne soit pas quasi-compact.
Notre but est de montrer que $f$ n'est pas universellement fermé. D'après le
lemme \ref{lemma-quasi-compact-local},
il existe un schéma affine $Z$ et un morphisme $Z \to Y$
tels que $Z \times_Y X \to Z$ ne soit pas quasi-compact. Pour atteindre notre but,
il suffit de montrer que $Z \times_Y X \to Z$ n'est pas universellement fermé ;
nous pouvons donc supposer que $Y = \Spec(B)$ pour un certain anneau $B$.

\medskip\noindent
Écrivons $X = \bigcup_{i \in I} X_i$, où les $X_i$ sont des sous-espaces ouverts
quasi-compacts de $X$. Par exemple, choisissons un morphisme \'etale surjectif
$U \to X$, où $U$ est un schéma, puis un recouvrement ouvert affine
$U = \bigcup U_i$, et soit $X_i \subset X$ l'image de $U_i$.
Nous utiliserons plus loin que les morphismes $X_i \to Y$ sont quasi-compacts ; voir le
lemme \ref{lemma-quasi-compact-permanence}.
Soit $T = \Spec(B[a_i ; i \in I])$. Soit $T_i = D(a_i) \subset T$.
Soit $Z \subset T \times_Y X$ le sous-espace fermé réduit dont l'ensemble fermé
sous-jacent de points est
$|T \times_Y X| \setminus \bigcup_{i \in I} |T_i \times_Y X_i|$ ; voir
Propriétés des espaces,
lemme \ref{spaces-properties-lemma-reduced-closed-subspace}.
(Notons que $T_i \times_Y X_i$ est un sous-espace ouvert de $T \times_Y X$, car
$T_i \to T$ et $X_i \to X$ sont des immersions ouvertes ; voir
Espaces, lemmes \ref{spaces-lemma-base-change-immersions} et
\ref{spaces-lemma-composition-immersions}.) Voici un diagramme :
$$
\xymatrix{
Z \ar[r] \ar[rd] &
T \times_Y X \ar[d]^{f_T} \ar[r]_q &
X \ar[d]^f \\
& T \ar[r]^p & Y
}
$$
Il suffit de prouver que l'image $f_T(|Z|)$ n'est pas fermée dans $|T|$.

\medskip\noindent
Nous affirmons qu'il existe un point $y \in Y$ tel qu'il n'existe aucun
voisinage ouvert affine $V$ de $y$ dans $Y$ tel que $X_V$ soit quasi-compact.
Sinon, on pourrait recouvrir $Y$ par un nombre fini de tels $V$ et, pour
chaque $V$, le morphisme $X_V \to V$ serait quasi-compact d'après le
lemme \ref{lemma-quasi-compact-permanence} ;
alors le
lemme \ref{lemma-quasi-compact-local}
entraînerait que $f$ est quasi-compact, contradiction. Fixons un $y \in Y$ ainsi obtenu.

\medskip\noindent
Soit $t \in T$ le point au-dessus de $y$ tel que $\kappa(t) = \kappa(y)$
et tel que $a_i = 1$ dans $\kappa(t)$ pour tout $i$. Supposons $z \in |Z|$ avec
$f_T(z) = t$. Alors $q(z) \in X_i$ pour un certain $i$. Par conséquent $f_T(z) \not \in T_i$
par construction de $Z$, ce qui contredit le fait que $t \in T_i$ par
construction. Ainsi $t \in |T| \setminus f_T(|Z|)$.

\medskip\noindent
Supposons $f_T(|Z|)$ fermé dans $|T|$. Il existe alors un élément
$g \in B[a_i; i \in I]$ tel que $f_T(|Z|) \subset V(g)$ mais $t \not \in V(g)$.
L'image de $g$ dans $\kappa(t)$ est donc non nulle. En particulier, un
coefficient de $g$ a une image non nulle dans $\kappa(y)$. Ce coefficient est donc
inversible sur un voisinage ouvert affine $V$ de $y$. Soit $J$ l'ensemble fini
des $j \in I$ tels que la variable $a_j$ apparaisse dans $g$.
Puisque $X_V$ n'est pas quasi-compact et que chaque $X_{i, V}$ est quasi-compact,
on peut choisir un point $x \in |X_V| \setminus \bigcup_{j \in J} |X_{j, V}|$.
Autrement dit, $x \in |X| \setminus \bigcup_{j \in J} |X_j|$ et $x$ se trouve
au-dessus d'un certain $v \in V$. Puisque $g$ possède un coefficient inversible sur $V$,
on peut trouver un point $t' \in T$ au-dessus de $v$ tel que $t' \not \in V(g)$ et
$t' \in V(a_i)$ pour tout $i \notin J$. En effet,
$V(a_i; i \in I \setminus J) = \Spec(B[a_j; j\in J])$
et l'ensemble des points de ce schéma au-dessus de $v$ est en bijection
avec $\Spec(\kappa(v)[a_j; j \in J])$, et la restriction de $g$ est
un élément non nul de cet anneau de polynômes par construction.
Autrement dit, $t' \not \in T_i$ pour tout $i \not \in J$. D'après
Propriétés des espaces, lemme \ref{spaces-properties-lemma-points-cartesian},
on peut trouver un point $z$ de $X \times_Y T$ qui s'envoie sur $x \in X$ et sur
$t' \in T$. Puisque $x \not \in |X_j|$ pour $j \in J$ et $t' \not \in T_i$
pour $i \in I \setminus J$, on voit que $z \in |Z|$. D'autre part,
$f_T(z) = t' \not \in V(g)$, ce qui contredit $f_T(Z) \subset V(g)$.
L'hypothèse « $f_T(|Z|)$ est fermé » est donc fausse, et l'on conclut bien
que $f_T$ n'est pas fermé, comme souhaité.
\end{proof}

\noindent
L'espace d'arrivée d'un morphisme surjectif et universellement fermé issu d'un
espace algébrique séparé est séparé.

\begin{lemma}
\label{lemma-image-universally-closed-separated}
Soit $S$ un schéma. Soit $B$ un espace algébrique sur $S$.
Soit $f : X \to Y$ un morphisme surjectif et universellement fermé
d'espaces algébriques sur $B$.
\begin{enumerate}
\item Si $X$ est quasi-séparé, alors $Y$ est quasi-séparé.
\item Si $X$ est séparé, alors $Y$ est séparé.
\item Si $X$ est quasi-séparé sur $B$, alors $Y$ est quasi-séparé sur $B$.
\item Si $X$ est séparé sur $B$, alors $Y$ est séparé sur $B$.
\end{enumerate}
\end{lemma}

\begin{proof}
Les parties (1) et (2) résultent de (3) et (4) pour
$S = B = \Spec(\mathbf{Z})$ (voir
Propriétés des espaces,
définition \ref{spaces-properties-definition-separated}).
Considérons le diagramme commutatif
$$
\xymatrix{
X \ar[d] \ar[rr]_{\Delta_{X/B}} & & X \times_B X \ar[d] \\
Y \ar[rr]^{\Delta_{Y/B}} & & Y \times_B Y
}
$$
La flèche verticale de gauche est surjective (c'est-à-dire universellement surjective).
La flèche verticale de droite est universellement fermée comme composée
des morphismes universellement fermés
$X \times_B X \to X \times_B Y \to Y \times_B Y$. Elle est donc aussi
quasi-compacte ; voir le
lemme \ref{lemma-universally-closed-quasi-compact}.

\medskip\noindent
Supposons $X$ quasi-séparé sur $B$, c'est-à-dire  $\Delta_{X/B}$
quasi-compact. Si $Z$ est quasi-compact et si $Z \to Y \times_B Y$ est
un morphisme, alors $Z \times_{Y \times_B Y} X \to Z \times_{Y \times_B Y} Y$
est surjectif et $Z \times_{Y \times_B Y} X$ est quasi-compact d'après les remarques
ci-dessus. On en déduit que $\Delta_{Y/B}$ est quasi-compact, c'est-à-dire que $Y$
est quasi-séparé sur $B$.

\medskip\noindent
Supposons $X$ séparé sur $B$, c'est-à-dire que $\Delta_{X/B}$ est une immersion
fermée. Si $Z$ est affine et si $Z \to Y \times_B Y$ est
un morphisme, alors $Z \times_{Y \times_B Y} X \to Z \times_{Y \times_B Y} Y$
est surjectif et $Z \times_{Y \times_B Y} X \to Z$ est universellement fermé
d'après les remarques ci-dessus. On en déduit que $\Delta_{Y/B}$ est universellement fermé.
Il s'ensuit que $\Delta_{Y/B}$ est représentable, localement de type fini et est un
monomorphisme (voir le
lemme \ref{lemma-properties-diagonal}),
et qu'il est universellement fermé, donc une immersion fermée ; voir
Morphismes \'etales,
lemme \ref{etale-lemma-characterize-closed-immersion}
(ainsi que le principe abstrait
Espaces, lemme
\ref{spaces-lemma-representable-transformations-property-implication}).
Ainsi $Y$ est séparé sur $B$.
\end{proof}













\section{Monomorphismes}
\label{section-monomorphisms}

\noindent
Un morphisme représentable $X \to Y$ d'espaces algébriques est un monomorphisme
au sens de la section \ref{section-representable} si, pour tout schéma
$Z$ et tout morphisme $Z \to Y$, le morphisme $Z \times_Y X \to Z$
est représentable par un monomorphisme de schémas.
Cela signifie exactement que $Z \times_Y X \to Z$
est une application injective de faisceaux sur $(\Sch/S)_{fppf}$. Puisque ceci
doit valoir pour tout $Z$ et toute application $Z \to Y$, c'est encore
équivalent à dire que l'application $X \to Y$ est une application injective de faisceaux sur
$(\Sch/S)_{fppf}$. Nous pouvons donc définir les monomorphismes parmi les morphismes
éventuellement non représentables\footnote{Nous ignorons si tout monomorphisme
d'espaces algébriques est représentable. Pour une discussion, voir
Compléments sur les morphismes d'espaces, section
\ref{spaces-more-morphisms-section-monomorphisms}.})
d'espaces algébriques comme suit.

\begin{definition}
\label{definition-monomorphism}
Soit $S$ un schéma.
Un morphisme d'espaces algébriques sur $S$ est appelé {\it monomorphisme}
s'il est une application injective de faisceaux, c'est-à-dire un monomorphisme dans la catégorie
des faisceaux sur $(\Sch/S)_{fppf}$.
\end{definition}

\noindent
Le lemme suivant montre que cela signifie aussi qu'il est un monomorphisme
dans la catégorie des espaces algébriques sur $S$.

\begin{lemma}
\label{lemma-monomorphism}
Soit $S$ un schéma.
Soit $j : X \to Y$ un morphisme d'espaces algébriques sur $S$.
Les conditions suivantes sont équivalentes :
\begin{enumerate}
\item $j$ est un monomorphisme (au sens de la définition \ref{definition-monomorphism}),
\item $j$ est un monomorphisme dans la catégorie des espaces algébriques sur $S$, et
\item le morphisme diagonal $\Delta_{X/Y} : X \to X \times_Y X$ est
un isomorphisme.
\end{enumerate}
\end{lemma}

\begin{proof}
Notons que $X \times_Y X$ est à la fois le produit fibré dans la catégorie des
faisceaux sur $(\Sch/S)_{fppf}$ et le produit fibré dans la catégorie
des espaces algébriques sur $S$ ; voir
Espaces, lemme \ref{spaces-lemma-fibre-product-spaces}.
L'équivalence de (1) et (3) est une caractérisation générale
des applications injectives de faisceaux sur un site quelconque.
L'équivalence de (2) et (3) est une caractérisation des monomorphismes
dans toute catégorie possédant des produits fibrés.
\end{proof}

\begin{lemma}
\label{lemma-monomorphism-separated}
Un monomorphisme d'espaces algébriques est séparé.
\end{lemma}

\begin{proof}
Cela vient de ce qu'un isomorphisme est une immersion fermée, compte tenu du
lemme \ref{lemma-monomorphism} ci-dessus.
\end{proof}

\begin{lemma}
\label{lemma-composition-monomorphism}
Un composé de monomorphismes est un monomorphisme.
\end{lemma}

\begin{proof}
C'est vrai, car un composé d'applications injectives de faisceaux est injectif.
\end{proof}

\begin{lemma}
\label{lemma-base-change-monomorphism}
Tout changement de base d'un monomorphisme est un monomorphisme.
\end{lemma}

\begin{proof}
C'est un fait général sur les produits fibrés dans une catégorie de faisceaux.
\end{proof}

\begin{lemma}
\label{lemma-monomorphism-local}
Soit $S$ un schéma.
Soit $f : X \to Y$ un morphisme d'espaces algébriques sur $S$.
Les conditions suivantes sont équivalentes :
\begin{enumerate}
\item $f$ est un monomorphisme,
\item pour tout schéma $Z$ et tout morphisme $Z \to Y$, le
changement de base $Z \times_Y X \to Z$ de $f$ est un monomorphisme,
\item pour tout schéma affine $Z$ et tout morphisme $Z \to Y$, le
changement de base $Z \times_Y X \to Z$ de $f$ est un monomorphisme,
\item il existe un schéma $V$ et un morphisme \'etale surjectif
$V \to Y$ tel que le changement de base $V \times_Y X \to V$ soit un
monomorphisme, et
\item il existe un recouvrement de Zariski $Y = \bigcup Y_i$ tel que chacun
des morphismes $f^{-1}(Y_i) \to Y_i$ soit un monomorphisme.
\end{enumerate}
\end{lemma}

\begin{proof}
Nous utiliserons sans plus le mentionner que tout changement de base d'un monomorphisme
est un monomorphisme ; voir le
lemme \ref{lemma-base-change-monomorphism}.
En particulier, il est clair que
(1) $\Rightarrow$ (2) $\Rightarrow$ (3) $\Rightarrow$ (4)
(en prenant pour $V$ une somme disjointe de schémas affines \'etales sur $Y$ ; voir
Propriétés des espaces,
lemme \ref{spaces-properties-lemma-cover-by-union-affines}).
Soit $V$ un schéma et soit $V \to Y$ un morphisme \'etale surjectif.
Si $V \times_Y X \to V$ est un monomorphisme, il
s'ensuit que $X \to Y$ est un monomorphisme. En effet, étant donné un
diagramme cartésien quelconque de faisceaux
$$
\vcenter{
\xymatrix{
\mathcal{F} \ar[r]_a \ar[d]_b & \mathcal{G} \ar[d]^c \\
\mathcal{H} \ar[r]^d & \mathcal{I}
}
}
\quad
\quad
\mathcal{F} = \mathcal{H} \times_\mathcal{I} \mathcal{G}
$$
si $c$ est une surjection de faisceaux et si $a$ est injectif, alors
$d$ est également injectif. Ainsi (4) implique (1). La preuve de l'équivalence de
(5) et (1) est omise.
\end{proof}

\begin{lemma}
\label{lemma-immersions-monomorphisms}
Une immersion d'espaces algébriques est un monomorphisme.
En particulier, toute immersion est séparée.
\end{lemma}

\begin{proof}
Soit $f : X \to Y$ une immersion d'espaces algébriques.
Pour tout morphisme $Z \to Y$ avec $Z$ représentable, le changement de
base $Z \times_Y X \to Z$ est une immersion de schémas, donc
un monomorphisme ; voir
Schémas, lemme \ref{schemes-lemma-immersions-monomorphisms}.
Ainsi $f$ est représentable et est un monomorphisme.
\end{proof}

\noindent
Nous améliorerons le lemme suivant dans
Espaces décents, lemme
\ref{decent-spaces-lemma-monomorphism-toward-disjoint-union-dim-0-rings}.

\begin{lemma}
\label{lemma-monomorphism-toward-field}
Soit $S$ un schéma. Soit $k$ un corps et soit $Z \to \Spec(k)$
un monomorphisme d'espaces algébriques sur $S$. Alors, soit
$Z = \emptyset$, soit $Z = \Spec(k)$.
\end{lemma}

\begin{proof}
D'après les
lemmes \ref{lemma-monomorphism-separated} et
\ref{lemma-separated-over-separated},
$Z$ est un espace algébrique séparé. Il existe donc un
sous-espace ouvert dense $Z' \subset Z$ qui est un schéma ; voir
Propriétés des espaces, proposition
\ref{spaces-properties-proposition-locally-quasi-separated-open-dense-scheme}.
D'après
Schémas, lemme \ref{schemes-lemma-mono-towards-spec-field},
soit $Z' = \emptyset$, soit $Z' \cong \Spec(k)$.
Dans le premier cas, on conclut que $Z = \emptyset$ et, dans le
second, que $Z' = Z = \Spec(k)$,
car $Z \to \Spec(k)$ est un monomorphisme qui est un
isomorphisme au-dessus de $Z'$.
\end{proof}

\begin{lemma}
\label{lemma-monomorphism-injective-points}
Soit $S$ un schéma. Si $X \to Y$ est un monomorphisme d'espaces algébriques
sur $S$, alors $|X| \to |Y|$ est injectif.
\end{lemma}

\begin{proof}
Cela résulte immédiatement des définitions.
\end{proof}















\section{Image directe des faisceaux quasi-cohérents}
\label{section-pushforward}

\noindent
Nous démontrons d'abord un lemme simple qui relie l'image directe des faisceaux de modules
pour un morphisme d'espaces algébriques à l'image directe des faisceaux de modules pour
un morphisme de schémas.

\begin{lemma}
\label{lemma-compute-pushforward}
Soit $S$ un schéma.
Soit $f : X \to Y$ un morphisme d'espaces algébriques sur $S$.
Soit $U \to X$ un morphisme \'etale surjectif d'un schéma vers $X$.
Posons $R = U \times_X U$ et notons, comme d'habitude, $t, s : R \to U$ les morphismes
de projection. Notons $a : U \to Y$ et $b : R \to Y$ les morphismes
induits. Pour tout objet $\mathcal{F}$ de $\textit{Mod}(\mathcal{O}_X)$,
il existe une suite exacte
$$
0 \to f_*\mathcal{F} \to a_*(\mathcal{F}|_U) \to b_*(\mathcal{F}|_R)
$$
où la seconde flèche est la différence $t^* - s^*$.
\end{lemma}

\begin{proof}
Nous notons encore $\mathcal{F}$ son prolongement en un faisceau de modules sur
$X_{spaces, \etale}$ ; voir
Propriétés des espaces,
remarque \ref{spaces-properties-remark-explain-equivalence}.
Soit $V \to Y$ un objet de $Y_\etale$. Alors $V \times_Y X$ est un
objet de $X_{spaces, \etale}$ et, par définition,
$f_*\mathcal{F}(V) = \mathcal{F}(V \times_Y X)$. Puisque $U \to X$ est
\'etale surjectif, on voit que $\{V \times_Y U \to V \times_Y X\}$
est un recouvrement. De plus,
$(V \times_Y U) \times_X (V \times_Y U) = V \times_Y R$.
Par conséquent, la condition de faisceau pour $\mathcal{F}$ sur
$X_{spaces, \etale}$ donne une suite exacte courte
$$
0 \to \mathcal{F}(V \times_Y X)
\to \mathcal{F}(V \times_Y U) \to \mathcal{F}(V \times_Y R)
$$
où la seconde flèche est la différence des restrictions par $t$ et par $s$.
Cette suite exacte est fonctorielle en $V$, ce qui donne le lemme.
\end{proof}

\noindent
Soit $S$ un schéma. Soit $f : X \to Y$ un morphisme quasi-compact et
quasi-séparé entre les espaces algébriques représentables $X$
et $Y$ sur $S$. D'après
Descente,
proposition \ref{descent-proposition-equivalence-quasi-coherent-functorial},
le foncteur
$f_* : \QCoh(\mathcal{O}_X) \to \QCoh(\mathcal{O}_Y)$
coïncide avec le foncteur usuel lorsque l'on considère $X$ et $Y$ comme des schémas.

\medskip\noindent
Plus généralement, supposons que $f : X \to Y$ soit un morphisme représentable,
quasi-compact et quasi-séparé d'espaces algébriques sur $S$. Soit $V$ un schéma
et soit $V \to Y$ un morphisme \'etale surjectif. Soit $U = V \times_Y X$
et soit $f' : U \to V$ le changement de base de $f$. Pour tout
$\mathcal{O}_X$-module quasi-cohérent $\mathcal{F}$, on a
\begin{equation}
\label{equation-representable-pushforward}
f'_*(\mathcal{F}|_U) = (f_*\mathcal{F})|_V,
\end{equation}
voir
Propriétés des espaces,
lemme \ref{spaces-properties-lemma-pushforward-etale-base-change-modules}.
Comme $f' : U \to V$ est un morphisme quasi-compact et quasi-séparé
de schémas, la remarque du paragraphe précédent permet de
calculer $f'_*(\mathcal{F}|_U)$ en considérant $\mathcal{F}|_U$ comme un
faisceau quasi-cohérent sur le schéma $U$ et $f'$ comme un morphisme de schémas.
Nous l'utiliserons fréquemment sans plus le mentionner.

\medskip\noindent
Le degré de généralité suivant consiste à considérer un morphisme quelconque
quasi-compact et quasi-séparé d'espaces algébriques.

\begin{lemma}
\label{lemma-pushforward}
Soit $S$ un schéma.
Soit $f : X \to Y$ un morphisme d'espaces algébriques sur $S$.
Si $f$ est quasi-compact et quasi-séparé, alors $f_*$ transforme les
$\mathcal{O}_X$-modules quasi-cohérents en
$\mathcal{O}_Y$-modules quasi-cohérents.
\end{lemma}

\begin{proof}
Soit $\mathcal{F}$ un faisceau quasi-cohérent sur $X$. Il faut montrer que
$f_*\mathcal{F}$ est un faisceau quasi-cohérent sur $Y$. Pour cela, il suffit
de montrer que, pour tout schéma affine $V$ et tout morphisme \'etale $V \to Y$,
la restriction de $f_*\mathcal{F}$ à $V$ est quasi-cohérente ; voir
Propriétés des espaces,
lemme \ref{spaces-properties-lemma-characterize-quasi-coherent}.
Soit $f' : V \times_Y X \to V$
le changement de base de $f$ par $V \to Y$. Notons que $f'$ est encore
quasi-compact et quasi-séparé ; voir les
lemmes \ref{lemma-base-change-quasi-compact} et
\ref{lemma-base-change-separated}.
D'après (\ref{equation-representable-pushforward}),
la restriction de $f_*\mathcal{F}$ à $V$ est l'image directe par $f'_*$ de la
restriction de $\mathcal{F}$ à $V \times_Y X$. Ainsi,
on peut remplacer $f$ par $f'$ et supposer que $Y$ est un schéma affine.

\medskip\noindent
Supposons que $Y$ soit un schéma affine. Puisque $f$ est quasi-compact, $X$
est quasi-compact. On peut donc choisir un schéma affine $U$ et un morphisme
\'etale surjectif $U \to X$ ; voir
Propriétés des espaces,
lemme \ref{spaces-properties-lemma-quasi-compact-affine-cover}.
Le lemme \ref{lemma-compute-pushforward} donne une suite exacte
$$
0 \to f_*\mathcal{F} \to a_*(\mathcal{F}|_U) \to b_*(\mathcal{F}|_R).
$$
où $R = U \times_X U$.
Comme $X \to Y$ est quasi-séparé, $R \to U \times_Y U$
est un monomorphisme quasi-compact. Il en résulte que $R$ est un schéma
quasi-compact séparé (puisque $U$ et $Y$ sont affines à ce stade).
Ainsi $a : U \to Y$ et $b : R \to Y$ sont des morphismes quasi-compacts et
quasi-séparés de schémas. Par conséquent, d'après
Descente,
proposition \ref{descent-proposition-equivalence-quasi-coherent-functorial},
les faisceaux $a_*(\mathcal{F}|_U)$ et $b_*(\mathcal{F}|_R)$
sont quasi-cohérents (voir aussi la discussion qui précède ce lemme).
Il en résulte que $f_*\mathcal{F}$ est un noyau de
modules quasi-cohérents, et est donc lui-même quasi-cohérent ; voir
Propriétés des espaces,
lemme \ref{spaces-properties-lemma-properties-quasi-coherent}.
\end{proof}

\noindent
Les images directes supérieures sont étudiées dans
Cohomologie des espaces, section
\ref{spaces-cohomology-section-higher-direct-image}.









\section{Immersions}
\label{section-immersions}

\noindent
Les immersions ouvertes, fermées et localement fermées d'espaces algébriques ont été définies dans
Espaces, section \ref{spaces-section-Zariski}.
Plus précisément, un morphisme d'espaces algébriques est une
{\it immersion fermée} (resp. une {\it immersion ouverte}, resp.\ une {\it immersion})
s'il est représentable et est une immersion fermée (resp.\ une immersion ouverte,
resp.\ une immersion) au sens de la section \ref{section-representable}.

\medskip\noindent
En particulier, ces types de morphismes sont stables par changement de base
et par composition de morphismes dans la catégorie des espaces
algébriques sur $S$ ; voir
Espaces, lemmes \ref{spaces-lemma-composition-immersions} et
\ref{spaces-lemma-base-change-immersions}.

\begin{lemma}
\label{lemma-closed-immersion-local}
Soit $S$ un schéma. Soit $f : X \to Y$ un morphisme d'espaces algébriques
sur $S$. Les conditions suivantes sont équivalentes :
\begin{enumerate}
\item $f$ est une immersion fermée (resp.\ une immersion ouverte, resp.\ une immersion),
\item pour tout schéma $Z$ et tout morphisme $Z \to Y$, le morphisme
$Z \times_Y X \to Z$ est une immersion fermée (resp.\ une immersion ouverte,
resp.\ une immersion),
\item pour tout schéma affine $Z$ et tout morphisme
$Z \to Y$, le morphisme $Z \times_Y X \to Z$ est une immersion fermée
(resp.\ une immersion ouverte, resp.\ une immersion),
\item il existe un schéma $V$ et un morphisme \'etale surjectif
$V \to Y$ tel que $V \times_Y X \to V$ soit une immersion fermée
(resp.\ une immersion ouverte, resp.\ une immersion), et
\item il existe un recouvrement de Zariski $Y = \bigcup Y_i$ tel que
chacun des morphismes $f^{-1}(Y_i) \to Y_i$ soit une immersion fermée
(resp.\ une immersion ouverte, resp.\ une immersion).
\end{enumerate}
\end{lemma}

\begin{proof}
Puisque tout changement de base d'une
immersion fermée (resp.\ immersion ouverte, resp.\ immersion)
en est encore une, il est clair que (1) implique (2) et que (2) implique (3).
De plus, (3) implique (4), car on peut prendre pour $V$ une somme disjointe de
schémas affines ; voir
Propriétés des espaces,
lemme \ref{spaces-properties-lemma-cover-by-union-affines}.

\medskip\noindent
Supposons que $V \to Y$ soit comme en (4).
Soit $\mathcal{P}$ la propriété d'être une
immersion fermée (resp.\ une immersion ouverte, resp.\ une immersion)
pour les morphismes de schémas. Notons que la propriété $\mathcal{P}$
est stable par tout changement de base et locale pour la topologie fppf sur la
base (voir la section \ref{section-representable}).
De plus, les morphismes de type $\mathcal{P}$ sont séparés et
localement quasi-finis (dans chacun des trois cas ; voir
Schémas, lemme \ref{schemes-lemma-immersions-monomorphisms}, et
Morphismes, lemme \ref{morphisms-lemma-immersion-locally-quasi-finite}).
Ainsi, d'après
Compléments sur les morphismes, lemme
\ref{more-morphisms-lemma-separated-locally-quasi-finite-morphisms-fppf-descend}
les morphismes de type $\mathcal{P}$ satisfont à la descente pour les recouvrements fppf. Ainsi
Espaces, lemme \ref{spaces-lemma-morphism-sheaves-with-P-effective-descent-etale}
s'applique, et l'on voit que $X \to Y$ est représentable et possède la propriété
$\mathcal{P}$ ; autrement dit, (1) est satisfaite.

\medskip\noindent
L'équivalence de (1) et (5) résulte du fait que
$\mathcal{P}$ est locale sur le but pour la topologie de Zariski (puisque nous avons vu
ci-dessus que $\mathcal{P}$ est en fait locale sur le but pour la topologie fppf).
\end{proof}

\begin{lemma}
\label{lemma-immersion-permanence}
Soit $S$ un schéma.
Soient $Z \to Y \to X$ des morphismes d'espaces algébriques sur $S$.
\begin{enumerate}
\item Si $Z \to X$ est représentable, localement de type fini, localement
quasi-fini, séparé et est un monomorphisme, alors $Z \to Y$ est
représentable, localement de type fini, localement quasi-fini,
séparé et est un monomorphisme.
\item Si $Z \to X$ est une immersion et si $Y \to X$ est localement séparé,
alors $Z \to Y$ est une immersion.
\item Si $Z \to X$ est une immersion fermée et si $Y \to X$ est séparé,
alors $Z \to Y$ est une immersion fermée.
\end{enumerate}
\end{lemma}

\begin{proof}
Dans chaque cas, la preuve consiste à considérer le diagramme commutatif
$$
\xymatrix{
Z \ar[r] \ar[rd] & Y \times_X Z \ar[r] \ar[d] & Z \ar[d] \\
& Y \ar[r] & X
}
$$
où le composé des flèches horizontales supérieures est l'identité.
Démontrons (1). La première flèche horizontale est une section de
$Y \times_X Z \to Z$ ; elle est donc représentable, localement de type fini,
localement quasi-finie, séparée et est un monomorphisme d'après le
lemme \ref{lemma-section-immersion}.
La flèche $Y \times_X Z \to Y$ est un changement de base de $Z \to X$ ;
elle est donc représentable, localement de type fini,
localement quasi-finie, séparée et est un monomorphisme
(car chacune de ces propriétés des morphismes de schémas est stable
par changement de base ; voir
Espaces, remarque \ref{spaces-remark-list-properties-stable-base-change}).
Il en va donc de même du composé (car chacune de ces propriétés des
morphismes de schémas est stable par composition ; voir Espaces, remarque
\ref{spaces-remark-list-properties-stable-composition}).
Ceci prouve (1). Les autres assertions se démontrent exactement de la même façon.
\end{proof}

\begin{lemma}
\label{lemma-immersion-when-closed}
Soit $S$ un schéma. Soit $i : Z \to X$ une immersion d'espaces
algébriques sur $S$. Alors $|i| : |Z| \to |X|$ est un homéomorphisme sur une
partie localement fermée, et $i$ est une immersion fermée si et seulement si
l'image $|i|(|Z|) \subset |X|$ est une partie fermée.
\end{lemma}

\begin{proof}
La première assertion est le lemme
\ref{spaces-properties-lemma-subspace-induced-topology} de Propriétés des espaces.
Soit $U$ un schéma et soit $U \to X$ un morphisme \'etale surjectif.
Par hypothèse, $T = U \times_X Z$ est un schéma et le morphisme $j : T \to U$
est une immersion de schémas. D'après le lemme \ref{lemma-closed-immersion-local},
le morphisme $i$ est une immersion fermée si et seulement si $j$ est une immersion
fermée. D'après Schémas, lemme \ref{schemes-lemma-immersion-when-closed},
cela équivaut à ce que $j(T)$ soit fermé dans $U$.
Or la partie $j(T) \subset U$ est l'image inverse de
$|i|(|Z|) \subset |X|$ ; voir
Propriétés des espaces, lemme \ref{spaces-properties-lemma-points-cartesian}.
Ceci achève la démonstration.
\end{proof}

\begin{remark}
\label{remark-immersion}
Soit $S$ un schéma. Soit $i : Z \to X$ une immersion d'espaces
algébriques sur $S$. Puisque $i$ est un monomorphisme, on peut considérer $|Z|$ comme
une partie de $|X|$ ; c'est ce que nous ferons dans la suite de cette remarque.
Soit $\partial |Z|$ la frontière de $|Z|$ dans
l'espace topologique $|X|$. En formule,
$$
\partial |Z| = \overline{|Z|} \setminus |Z|.
$$
Soit $\partial Z$ le sous-espace fermé réduit de $X$ tel que
$|\partial Z| = \partial |Z|$,
obtenu en prenant la structure réduite induite de sous-espace fermé ; voir
Propriétés des espaces,
définition \ref{spaces-properties-definition-reduced-induced-space}.
Par construction, $|Z|$ est fermé dans
$|X| \setminus |\partial Z| = |X \setminus \partial Z|$.
Ainsi, toute immersion d'espaces algébriques peut bien se
factoriser en une immersion fermée suivie d'une immersion ouverte
(mais pas dans l'autre ordre en général ; voir
Morphismes, exemple \ref{morphisms-example-thibaut}).
\end{remark}

\begin{remark}
\label{remark-space-structure-locally-closed-subset}
Soit $S$ un schéma. Soit $X$ un espace algébrique sur $S$.
Soit $T \subset |X|$ une partie localement fermée.
Soit $\partial T$ la frontière de $T$ dans
l'espace topologique $|X|$. En formule,
$$
\partial T = \overline{T} \setminus T.
$$
Soit $U \subset X$ le sous-espace ouvert de $X$ tel que
$|U| = |X| \setminus \partial T$ ; voir
Propriétés des espaces, lemme \ref{spaces-properties-lemma-open-subspaces}.
Soit $Z$ le sous-espace fermé réduit de $U$ tel que
$|Z| = T$, obtenu en prenant la structure réduite induite
de sous-espace fermé ; voir
Propriétés des espaces,
définition \ref{spaces-properties-definition-reduced-induced-space}.
Par construction, $Z \to U$ est une immersion fermée d'espaces algébriques
et $U \to X$ est une immersion ouverte ; par conséquent,
$Z \to X$ est une immersion d'espaces algébriques sur $S$ (voir
Espaces, lemme \ref{spaces-lemma-composition-immersions}).
Notons que $Z$ est un espace algébrique réduit et que
$|Z| = T$ comme parties de $|X|$. On dit parfois que
$Z$ est la {\it structure réduite induite de sous-espace} sur $T$.
\end{remark}

\begin{lemma}
\label{lemma-factor-the-other-way}
Soit $S$ un schéma. Soit $Z \to X$ une immersion d'espaces algébriques sur
$S$. Supposons $Z \to X$ quasi-compact.
Il existe une factorisation $Z \to \overline{Z} \to X$ où
$Z \to \overline{Z}$ est une immersion ouverte et $\overline{Z} \to X$
est une immersion fermée.
\end{lemma}

\begin{proof}
Soit $U$ un schéma et soit $U \to X$ \'etale surjectif.
Notons comme d'habitude $R = U \times_X U$, avec pour projections
$s, t : R \to U$. Posons $T = Z \times_X U$. Soit $\overline{T} \subset U$
l'image schématique de $T \to U$. Notons que
$s^{-1}\overline{T} = t^{-1}\overline{T}$, car la formation des
images schématiques de morphismes quasi-compacts commute
au changement de base plat ; voir
Morphismes, lemme \ref{morphisms-lemma-flat-base-change-scheme-theoretic-image}.
On obtient donc un sous-espace fermé $\overline{Z} \subset X$ dont
le changement de base à $U$ est $\overline{T}$ ; voir
Propriétés des espaces, lemme
\ref{spaces-properties-lemma-subspaces-presentation}.
D'après Morphismes, lemme \ref{morphisms-lemma-quasi-compact-immersion},
le morphisme $T \to \overline{T}$
est une immersion ouverte. Il s'ensuit que $Z \to \overline{Z}$ est
une immersion ouverte, ce qui donne le résultat.
\end{proof}


















\section{Immersions fermées}
\label{section-closed-immersions}

\noindent
Dans cette section, nous précisons certains résultats obtenus précédemment sur les
immersions d'espaces algébriques. Voir
Espaces, section \ref{spaces-section-Zariski},
et la
section \ref{section-immersions} du présent chapitre.
Cette section est l'analogue de
Morphismes, section \ref{morphisms-section-closed-immersions},
pour les espaces algébriques.

\begin{lemma}
\label{lemma-closed-immersion-ideals}
Soit $S$ un schéma.
Soit $X$ un espace algébrique sur $S$.
Pour toute immersion fermée $i : Z \to X$, le faisceau
$i_*\mathcal{O}_Z$ est un $\mathcal{O}_X$-module quasi-cohérent ; l'application
$i^\sharp : \mathcal{O}_X \to i_*\mathcal{O}_Z$ est surjective et son
noyau est un faisceau quasi-cohérent d'idéaux. La règle
$Z \mapsto \Ker(\mathcal{O}_X \to i_*\mathcal{O}_Z)$
définit une bijection renversant les inclusions
$$
\begin{matrix}
\text{sous-espaces fermés}\\
Z \subset X
\end{matrix}
\longrightarrow
\begin{matrix}
\text{faisceaux quasi-cohérents}\\
\text{d'idéaux }\mathcal{I} \subset \mathcal{O}_X
\end{matrix}
$$
De plus, étant donné un sous-espace fermé $Z$ correspondant au faisceau
quasi-cohérent d'idéaux $\mathcal{I} \subset \mathcal{O}_X$, un morphisme d'espaces
algébriques $h : Y \to X$ se factorise par $Z$ si et seulement si l'application
$h^*\mathcal{I} \to h^*\mathcal{O}_X = \mathcal{O}_Y$ est nulle.
\end{lemma}

\begin{proof}
Soit $U \to X$ un morphisme \'etale surjectif dont la source est un schéma.
Considérons le diagramme
$$
\xymatrix{
U \times_X Z \ar[r] \ar[d]_{i'} & Z \ar[d]^i \\
U \ar[r] & X
}
$$
D'après le
lemme \ref{lemma-closed-immersion-local},
$i$ est une immersion fermée
si et seulement si $i'$ est une immersion fermée. D'après
Propriétés des espaces,
lemme \ref{spaces-properties-lemma-pushforward-etale-base-change-modules},
$i'_*\mathcal{O}_{U \times_X Z}$ est la restriction de
$i_*\mathcal{O}_Z$ à $U$. Ainsi, les assertions concernant
$\mathcal{O}_X \to i_*\mathcal{O}_Z$ sont équivalentes aux
assertions correspondantes concernant
$\mathcal{O}_U \to i'_*\mathcal{O}_{U \times_X Z}$.
Puisque $i'$ est une immersion fermée de schémas, ces résultats découlent de
Morphismes, lemme \ref{morphisms-lemma-closed-immersion}.

\medskip\noindent
Démontrons que, pour un faisceau quasi-cohérent
d'idéaux $\mathcal{I} \subset \mathcal{O}_X$, la formule
$$
Z(T) = \{h : T \to X \mid h^*\mathcal{I} \to \mathcal{O}_T
\text{ est nulle}\}
$$
définit un sous-espace fermé de $X$. C'est clairement un sous-foncteur de $X$.
Pour montrer que $Z \to X$ est représentable par des immersions fermées, soit
$\varphi : U \to X$ un morphisme d'un schéma vers $X$. Alors
$Z \times_X U$ est représenté par le sous-foncteur analogue de $U$ correspondant
au faisceau d'idéaux $\Im(\varphi^*\mathcal{I} \to \mathcal{O}_U)$. D'après
Propriétés des espaces,
lemme \ref{spaces-properties-lemma-pullback-quasi-coherent},
le $\mathcal{O}_U$-module $\varphi^*\mathcal{I}$ est quasi-cohérent
sur $U$, et donc $\Im(\varphi^*\mathcal{I} \to \mathcal{O}_U)$
est un faisceau quasi-cohérent d'idéaux sur $U$. D'après
Schémas, lemme \ref{schemes-lemma-characterize-closed-subspace},
on conclut que $Z \times_X U$ est représenté par le sous-schéma fermé
de $U$ associé à $\Im(\varphi^*\mathcal{I} \to \mathcal{O}_U)$.
Ainsi $Z$ est un sous-espace fermé de $X$.

\medskip\noindent
Dans la formule ci-dessus définissant $Z$, les arguments $T$ sont des schémas, puisque les espaces
algébriques sont des faisceaux sur $(\Sch/S)_{fppf}$. Nous omettons la vérification
que la même formule reste vraie lorsque $T$ est un espace algébrique.
\end{proof}

\begin{definition}
\label{definition-inverse-image-closed-subspace}
Soit $S$ un schéma. Soit $f : Y \to X$ un morphisme d'espaces algébriques
sur $S$. Soit $Z \subset X$ un sous-espace fermé. L'
{\it image inverse $f^{-1}(Z)$ du sous-espace fermé $Z$}
est le sous-espace fermé $Z \times_X Y$ de $Y$.
\end{definition}

\noindent
Cette définition a un sens d'après le lemme \ref{lemma-closed-immersion-local}.
Si $\mathcal{I} \subset \mathcal{O}_X$ est le faisceau quasi-cohérent
d'idéaux correspondant à $Z$ par le lemme \ref{lemma-closed-immersion-ideals}, alors
$f^{-1}\mathcal{I}\mathcal{O}_Y = \Im(f^*\mathcal{I} \to \mathcal{O}_Y)$
est le faisceau d'idéaux correspondant à $f^{-1}(Z)$.

\begin{lemma}
\label{lemma-closed-immersion-quasi-compact}
Une immersion fermée d'espaces algébriques est quasi-compacte.
\end{lemma}

\begin{proof}
Ceci résulte de
Schémas, lemme \ref{schemes-lemma-closed-immersion-quasi-compact},
par des principes généraux ; voir
Espaces, lemme
\ref{spaces-lemma-representable-transformations-property-implication}.
\end{proof}

\begin{lemma}
\label{lemma-closed-immersion-separated}
Une immersion fermée d'espaces algébriques est séparée.
\end{lemma}

\begin{proof}
Ceci résulte de
Schémas, lemme \ref{schemes-lemma-immersions-monomorphisms},
par des principes généraux ; voir
Espaces, lemme
\ref{spaces-lemma-representable-transformations-property-implication}.
\end{proof}

\begin{lemma}
\label{lemma-closed-immersion-push-pull}
Soit $S$ un schéma. Soit $i : Z \to X$ une immersion fermée d'espaces
algébriques sur $S$.
\begin{enumerate}
\item Le foncteur
$$
i_{small, *} :
\Sh(Z_\etale)
\longrightarrow
\Sh(X_\etale)
$$
est pleinement fidèle et son image essentielle est formée des faisceaux d'ensembles
$\mathcal{F}$ sur $X_\etale$ dont la restriction à $X \setminus Z$ est
isomorphe à $*$, et
\item le foncteur
$$
i_{small, *} :
\textit{Ab}(Z_\etale)
\longrightarrow
\textit{Ab}(X_\etale)
$$
est pleinement fidèle et son image essentielle est formée des faisceaux abéliens sur
$X_\etale$ dont le support est contenu dans $|Z|$.
\end{enumerate}
Dans les deux cas, $i_{small}^{-1}$ est un inverse à gauche du foncteur
$i_{small, *}$.
\end{lemma}

\begin{proof}
Soit $U$ un schéma et soit $U \to X$ \'etale surjectif.
Posons $V = Z \times_X U$. Alors $V$ est un schéma et $i' : V \to U$ est
une immersion fermée de schémas. D'après
Propriétés des espaces,
lemme \ref{spaces-properties-lemma-pushforward-etale-base-change},
pour tout faisceau $\mathcal{G}$ sur $Z$, on a
$$
(i_{small}^{-1}i_{small, *}\mathcal{G})|_V =
(i')_{small}^{-1}i'_{small, *}(\mathcal{G}|_V)
$$
D'après
Cohomologie \'etale, proposition
\ref{etale-cohomology-proposition-closed-immersion-pushforward},
l'application
$(i')_{small}^{-1}i'_{small, *}(\mathcal{G}|_V) \to \mathcal{G}|_V$
est un isomorphisme. Puisque $V \to Z$ est surjectif et \'etale, cela entraîne
que $i_{small}^{-1}i_{small, *}\mathcal{G} \to \mathcal{G}$ est un
isomorphisme. Ceci implique clairement que $i_{small, *}$ est pleinement fidèle ; voir
Sites, lemme \ref{sites-lemma-exactness-properties}.
Pour prouver l'assertion sur l'image essentielle, considérons un faisceau d'ensembles
$\mathcal{F}$ sur $X_\etale$ dont la restriction à $X \setminus Z$ est
isomorphe à $*$. Comme dans la preuve de
Cohomologie \'etale, proposition
\ref{etale-cohomology-proposition-closed-immersion-pushforward}
nous considérons le morphisme d'adjonction
$$
\mathcal{F} \longrightarrow i_{small, *}i_{small}^{-1}\mathcal{F}.
$$
Comme dans la première partie, le résultat correspondant pour le cas des
schémas montre que la restriction de ce morphisme à $U$ est un isomorphisme.
Puisque $U$ est un recouvrement \'etale de $X$, on
conclut qu'il s'agit d'un isomorphisme.
\end{proof}

\begin{lemma}
\label{lemma-stalk-push-closed}
Soit $S$ un schéma. Soit $i : Z \to X$ une immersion fermée d'espaces
algébriques sur $S$. Soit $\overline{z}$ un point géométrique de $Z$
d'image $\overline{x}$ dans $X$. Alors
$(i_{small, *}\mathcal{F})_{\overline{x}} = \mathcal{F}_{\overline{z}}$
pour tout faisceau $\mathcal{F}$ sur $Z_\etale$.
\end{lemma}

\begin{proof}
Choisissons un voisinage \'etale $(U, \overline{u})$ de $\overline{x}$.
Alors la fibre $(i_{small, *}\mathcal{F})_{\overline{x}}$
est la fibre de $i_{small, *}\mathcal{F}|_U$ en $\overline{u}$.
D'après Propriétés des espaces,
lemme \ref{spaces-properties-lemma-pushforward-etale-base-change},
on peut remplacer $X$ par $U$ et $Z$ par $Z \times_X U$.
Alors $Z \to X$ est une immersion fermée de schémas, et le résultat est
Cohomologie \'etale, lemme
\ref{etale-cohomology-lemma-stalk-pushforward-closed-immersion}.
\end{proof}

\noindent
Le lemme suivant vaut plus généralement dans le cadre d'une immersion fermée
de topos (insérer ici une référence future).

\begin{lemma}
\label{lemma-closed-immersion-rings}
Soit $S$ un schéma. Soit $i : Z \to X$ une immersion fermée d'espaces
algébriques sur $S$. Soit $\mathcal{A}$ un faisceau d'anneaux sur $X_\etale$.
Soit $\mathcal{B}$ un faisceau d'anneaux sur $Z_\etale$.
Soit $\varphi : \mathcal{A} \to i_{small, *}\mathcal{B}$
un homomorphisme de faisceaux d'anneaux, ce qui donne un
morphisme de topos annelés
$$
f : (\Sh(Z_\etale), \mathcal{B}) \longrightarrow (\Sh(X_\etale), \mathcal{A}).
$$
Pour un faisceau de $\mathcal{A}$-modules $\mathcal{F}$ et un
faisceau de $\mathcal{B}$-modules $\mathcal{G}$, l'application canonique
$$
\mathcal{F} \otimes_\mathcal{A} f_*\mathcal{G}
\longrightarrow
f_*(f^*\mathcal{F} \otimes_\mathcal{B} \mathcal{G})
$$
est un isomorphisme.
\end{lemma}

\begin{proof}
Cette application est l'adjointe de l'application
$$
f^*\mathcal{F} \otimes_\mathcal{B}
f^* f_*\mathcal{G} =
f^*(\mathcal{F} \otimes_\mathcal{A} f_*\mathcal{G})
\longrightarrow
f^*\mathcal{F} \otimes_\mathcal{B} \mathcal{G}
$$
provenant de $\text{id} : f^*\mathcal{F} \to f^*\mathcal{F}$
et du morphisme d'adjonction $f^* f_*\mathcal{G} \to \mathcal{G}$.
Pour voir que cette application est un isomorphisme, il suffit de le vérifier sur les fibres
(Propriétés des espaces, théorème
\ref{spaces-properties-theorem-exactness-stalks}).
Soit $\overline{z} : \Spec(k) \to Z$ un point géométrique
d'image $\overline{x} = i \circ \overline{z} : \Spec(k) \to X$.
En explicitant l'action de notre application sur les fibres, on voit qu'il
faut montrer
$$
\mathcal{F}_{\overline{x}}
\otimes_{\mathcal{A}_{\overline{x}}}
\mathcal{G}_{\overline{z}} =
(\mathcal{F}_{\overline{x}}
\otimes_{\mathcal{A}_{\overline{x}}}
\mathcal{B}_{\overline{z}}) \otimes_{\mathcal{B}_{\overline{z}}}
\mathcal{G}_{\overline{z}}
$$
ce qui est vrai. Nous avons utilisé ici que
la formation des produits tensoriels commute à celle des fibres, ainsi que le
comportement des fibres par image inverse (Propriétés des espaces,
lemme \ref{spaces-properties-lemma-stalk-pullback}) et par image directe
le long d'une immersion fermée (lemme
\ref{lemma-stalk-push-closed}).
\end{proof}






\section{Immersions fermées et faisceaux quasi-cohérents}
\label{section-closed-immersions-quasi-coherent}

\noindent
Cette section est l'analogue de
Morphismes, section \ref{morphisms-section-closed-immersions-quasi-coherent}.

\begin{lemma}
\label{lemma-i-star-equivalence}
Soit $S$ un schéma. Soit $i : Z \to X$ une immersion fermée d'espaces
algébriques sur $S$. Soit $\mathcal{I} \subset \mathcal{O}_X$ le faisceau
quasi-cohérent d'idéaux définissant $Z$.
\begin{enumerate}
\item Pour tout $\mathcal{O}_X$-module $\mathcal{F}$, le morphisme d'adjonction
$\mathcal{F} \to i_*i^*\mathcal{F}$ induit un isomorphisme
$\mathcal{F}/\mathcal{I}\mathcal{F} \cong i_*i^*\mathcal{F}$.
\item Le foncteur $i^*$ est un inverse à gauche de $i_*$, c'est-à-dire que, pour tout
$\mathcal{O}_Z$-module $\mathcal{G}$, le morphisme d'adjonction
$i^*i_*\mathcal{G} \to \mathcal{G}$ est un isomorphisme.
\item Le foncteur
$$
i_* :
\QCoh(\mathcal{O}_Z)
\longrightarrow
\QCoh(\mathcal{O}_X)
$$
est exact et pleinement fidèle, et son image essentielle est formée des
$\mathcal{O}_X$-modules quasi-cohérents $\mathcal{F}$ tels que $\mathcal{I}\mathcal{F} = 0$.
\end{enumerate}
\end{lemma}

\begin{proof}
Au cours de cette démonstration, nous travaillons exclusivement avec les faisceaux sur
les petits sites étales, et nous employons $i_*, i^{-1}, \ldots$
pour désigner l'image directe et l'image inverse des faisceaux de groupes abéliens
au lieu de $i_{small, *}, i_{small}^{-1}$.

\medskip\noindent
Soit $\mathcal{F}$ un $\mathcal{O}_X$-module. Le
lemme \ref{lemma-closed-immersion-rings}, appliqué avec
$\mathcal{A} = \mathcal{O}_X$ et
$\mathcal{G} = \mathcal{B} = \mathcal{O}_Z$, montre que
$i_*i^*\mathcal{F} = \mathcal{F} \otimes_{\mathcal{O}_X} i_*\mathcal{O}_Z$.
Le
lemme \ref{lemma-closed-immersion-ideals}
nous donne une suite exacte courte
$$
0 \to \mathcal{I} \to \mathcal{O}_X \to i_*\mathcal{O}_Z \to 0
$$
Il résulte des propriétés du produit tensoriel que
$\mathcal{F} \otimes_{\mathcal{O}_X} i_*\mathcal{O}_Z
= \mathcal{F}/\mathcal{I}\mathcal{F}$. Ceci démontre (1) (à ceci près
que nous omettons de vérifier que le morphisme est induit par le
morphisme d'adjonction).

\medskip\noindent
Soit $\mathcal{G}$ un $\mathcal{O}_Z$-module quelconque. Le
lemme \ref{lemma-closed-immersion-push-pull}
montre que $i^{-1}i_*\mathcal{G} = \mathcal{G}$.
Pour démontrer (2), il nous suffit donc de montrer que le morphisme canonique
$\mathcal{G} \otimes_{i^{-1}\mathcal{O}_X} \mathcal{O}_Z \to \mathcal{G}$
est un isomorphisme. Cela résulte des propriétés générales des produits tensoriels
dès que l'on montre que $i^{-1}\mathcal{O}_X \to \mathcal{O}_Z$ est surjectif. D'après le
lemme \ref{lemma-closed-immersion-push-pull},
il suffit de prouver que
$i_*i^{-1}\mathcal{O}_X \to i_*\mathcal{O}_Z$
est surjectif. Comme le morphisme surjectif $\mathcal{O}_X \to i_*\mathcal{O}_Z$
se factorise par ce morphisme, nous obtenons (2).

\medskip\noindent
Enfin, démontrons la partie la plus intéressante du lemme, à savoir (3).
Une immersion fermée est quasi-compacte et séparée ; voir les
lemmes \ref{lemma-closed-immersion-quasi-compact} et
\ref{lemma-closed-immersion-separated}. Le
lemme \ref{lemma-pushforward}
s'applique donc, et l'image directe d'un faisceau quasi-cohérent
sur $Z$ est bien un faisceau quasi-cohérent sur $X$.
Nous obtenons ainsi le foncteur
$i^{QCoh}_* : \QCoh(\mathcal{O}_Z)
\to \QCoh(\mathcal{O}_X)$.
Il résulte de (2) que $i^{QCoh}_*$ est pleinement fidèle, puisque
$i^*$ est son adjoint à gauche et que le morphisme d'adjonction de (2) est un isomorphisme.

\medskip\noindent
Décrivons maintenant l'image essentielle du
foncteur $i_*$. Il est clair que $\mathcal{I}(i_*\mathcal{G}) = 0$
pour tout $\mathcal{O}_Z$-module, puisque $\mathcal{I}$ est le noyau
du morphisme $\mathcal{O}_X \to i_*\mathcal{O}_Z$, qui est précisément celui
par lequel on munit d'une structure de $\mathcal{O}_X$-module le faisceau $i_*\mathcal{G}$.
Supposons ensuite que $\mathcal{F}$ soit un $\mathcal{O}_X$-module
quasi-cohérent tel que $\mathcal{I}\mathcal{F} = 0$.
Alors $\mathcal{F}$ est un $i_*\mathcal{O}_Z$-module,
car $i_*\mathcal{O}_Z = \mathcal{O}_X/\mathcal{I}$. En particulier,
son support est contenu dans $|Z|$. Le
lemme \ref{lemma-closed-immersion-push-pull}
montre que $\mathcal{F} \cong i_*\mathcal{G}$ pour un certain
$\mathcal{O}_Z$-module $\mathcal{G}$. Le seul détail restant consiste
à montrer que $\mathcal{G}$ est quasi-cohérent. Or,
$\mathcal{G} \cong i^*\mathcal{F}$ par (2), et
Propriétés des espaces,
lemme \ref{spaces-properties-lemma-pullback-quasi-coherent}, donne le résultat.
\end{proof}

\noindent
Soit $i : Z \to X$ une immersion fermée d'espaces algébriques.
D'après le lemme précédent, nous noterons souvent,
par abus de notation, $\mathcal{F}$ le faisceau $i_*\mathcal{F}$ sur $X$.

\begin{lemma}
\label{lemma-largest-quasi-coherent-subsheaf}
Soit $S$ un schéma. Soit $X$ un espace algébrique sur $S$.
Soit $\mathcal{F}$ un $\mathcal{O}_X$-module quasi-cohérent.
Soit $\mathcal{G} \subset \mathcal{F}$ un $\mathcal{O}_X$-sous-module.
Il existe un unique $\mathcal{O}_X$-sous-module quasi-cohérent
$\mathcal{G}' \subset \mathcal{G}$ ayant la propriété suivante :
pour tout $\mathcal{O}_X$-module quasi-cohérent $\mathcal{H}$, le morphisme
$$
\Hom_{\mathcal{O}_X}(\mathcal{H}, \mathcal{G}')
\longrightarrow
\Hom_{\mathcal{O}_X}(\mathcal{H}, \mathcal{G})
$$
est bijectif. En particulier, $\mathcal{G}'$ est le plus grand
$\mathcal{O}_X$-sous-module quasi-cohérent de $\mathcal{F}$ contenu dans $\mathcal{G}$.
\end{lemma}

\begin{proof}
Soit $\mathcal{G}_a$, $a \in A$, l'ensemble des
$\mathcal{O}_X$-sous-modules quasi-cohérents contenus dans $\mathcal{G}$.
Alors l'image $\mathcal{G}'$ de
$$
\bigoplus\nolimits_{a \in A} \mathcal{G}_a \longrightarrow \mathcal{F}
$$
est quasi-cohérente, car l'image d'un morphisme de faisceaux quasi-cohérents
sur $X$ est quasi-cohérente et toute somme directe de faisceaux quasi-cohérents
est quasi-cohérente ; voir le lemme de
Propriétés des espaces,
lemme \ref{spaces-properties-lemma-properties-quasi-coherent}.
Le module $\mathcal{G}'$ est contenu dans $\mathcal{G}$. C'est donc le
plus grand $\mathcal{O}_X$-sous-module quasi-cohérent contenu dans $\mathcal{G}$.

\medskip\noindent
Pour démontrer la formule, soit $\mathcal{H}$ un
$\mathcal{O}_X$-module quasi-cohérent et soit $\alpha : \mathcal{H} \to \mathcal{G}$
un morphisme de $\mathcal{O}_X$-modules. L'image du composé
$\mathcal{H} \to \mathcal{G} \to \mathcal{F}$ est quasi-cohérente,
comme image d'un morphisme de faisceaux quasi-cohérents. Elle est donc contenue
dans $\mathcal{G}'$. Par conséquent, $\alpha$ se factorise par $\mathcal{G}'$,
comme souhaité.
\end{proof}

\begin{lemma}
\label{lemma-i-upper-shriek}
Soit $S$ un schéma.
Soit $i : Z \to X$ une immersion fermée d'espaces algébriques sur $S$.
Il existe un foncteur\footnote{Cette notation n'est vraisemblablement pas standard.}
$i^! : \QCoh(\mathcal{O}_X) \to \QCoh(\mathcal{O}_Z)$
qui est un adjoint à droite de $i_*$. (Comparer avec
Modules, lemme \ref{modules-lemma-i-star-right-adjoint}.)
\end{lemma}

\begin{proof}
Étant donné un $\mathcal{O}_X$-module quasi-cohérent $\mathcal{G}$, considérons
le sous-faisceau $\mathcal{H}_Z(\mathcal{G})$ de $\mathcal{G}$ formé des sections locales
annulées par le faisceau d'idéaux $\mathcal{I}$ définissant $Z$. D'après le
lemme \ref{lemma-largest-quasi-coherent-subsheaf},
il existe canoniquement un plus grand $\mathcal{O}_X$-sous-module quasi-cohérent,
noté $\mathcal{H}_Z(\mathcal{G})'$. Par construction, nous avons
$$
\Hom_{\mathcal{O}_X}(i_*\mathcal{F}, \mathcal{H}_Z(\mathcal{G})')
=
\Hom_{\mathcal{O}_X}(i_*\mathcal{F}, \mathcal{G})
$$
pour tout $\mathcal{O}_Z$-module quasi-cohérent $\mathcal{F}$.
Nous pouvons donc poser $i^!\mathcal{G} = i^*(\mathcal{H}_Z(\mathcal{G})')$.
Nous omettons les détails.
\end{proof}

\noindent
Au moyen de la correspondance biunivoque ($1$-à-$1$) entre les faisceaux quasi-cohérents
d'idéaux et les sous-espaces fermés (voir le
lemme \ref{lemma-closed-immersion-ideals}),
nous pouvons définir les intersections et réunions schématiques
de sous-espaces fermés.

\begin{definition}
\label{definition-scheme-theoretic-intersection-union}
Soit $S$ un schéma. Soit $X$ un espace algébrique sur $S$.
Soient $Z, Y \subset X$ des sous-espaces fermés
correspondant aux faisceaux quasi-cohérents d'idéaux
$\mathcal{I}, \mathcal{J} \subset \mathcal{O}_X$.
L'{\it intersection schématique} de $Z$ et $Y$
est le sous-espace fermé de $X$ défini par $\mathcal{I} + \mathcal{J}$.
La {\it réunion schématique} de $Z$ et $Y$
est le sous-espace fermé de $X$ défini par
$\mathcal{I} \cap \mathcal{J}$.
\end{definition}

\noindent
Il est clair que la formation de l'intersection schématique
commute à la localisation étale, et qu'il en va de même pour la
réunion schématique.

\begin{lemma}
\label{lemma-scheme-theoretic-intersection}
Soit $S$ un schéma. Soit $X$ un espace algébrique sur $S$.
Soient $Z, Y \subset X$ des sous-espaces fermés.
Soit $Z \cap Y$ l'intersection schématique de $Z$ et $Y$.
Alors $Z \cap Y \to Z$ et $Z \cap Y \to Y$ sont des immersions fermées,
et
$$
\xymatrix{
Z \cap Y \ar[r] \ar[d] & Z \ar[d] \\
Y \ar[r] & X
}
$$
est un diagramme cartésien d'espaces algébriques sur $S$, c'est-à-dire que
$Z \cap Y = Z \times_X Y$.
\end{lemma}

\begin{proof}
Les morphismes $Z \cap Y \to Z$ et $Z \cap Y \to Y$ sont des immersions fermées
d'après le lemme \ref{lemma-closed-immersion-ideals}.
Comme la formation de l'intersection schématique commute
à la localisation étale, le cas des schémas montre que le diagramme est cartésien.
Voir le
lemme de Morphismes \ref{morphisms-lemma-scheme-theoretic-intersection}.
\end{proof}

\begin{lemma}
\label{lemma-scheme-theoretic-union}
Soit $S$ un schéma. Soit $X$ un espace algébrique sur $S$.
Soient $Y, Z \subset X$ des sous-espaces fermés.
Soit $Y \cup Z$ la réunion schématique de $Y$ et $Z$.
Soit $Y \cap Z$ l'intersection schématique de $Y$ et $Z$.
Alors $Y \to Y \cup Z$ et $Z \to Y \cup Z$ sont des immersions fermées,
il existe une suite exacte courte
$$
0 \to \mathcal{O}_{Y \cup Z} \to \mathcal{O}_Y \times \mathcal{O}_Z
\to \mathcal{O}_{Y \cap Z} \to 0
$$
de $\mathcal{O}_X$-modules, et le diagramme
$$
\xymatrix{
Y \cap Z \ar[r] \ar[d] & Y \ar[d] \\
Z \ar[r] & Y \cup Z
}
$$
est cocartésien dans la catégorie des espaces algébriques sur $S$, c'est-à-dire que
$Y \cup Z = Y \amalg_{Y \cap Z} Z$.
\end{lemma}

\begin{proof}
Les morphismes $Y \to Y \cup Z$ et $Z \to Y \cup Z$ sont des immersions fermées
d'après le lemme \ref{lemma-closed-immersion-ideals}. Dans la suite exacte courte,
nous employons l'équivalence du lemme \ref{lemma-i-star-equivalence} pour considérer les
modules quasi-cohérents sur les sous-espaces fermés de $X$ comme des modules quasi-cohérents
sur $X$. Le premier morphisme de la suite provient des morphismes canoniques
$\mathcal{O}_{Y \cup Z} \to \mathcal{O}_Y$ et
$\mathcal{O}_{Y \cup Z} \to \mathcal{O}_Z$,
et le second du morphisme canonique
$\mathcal{O}_Y \to \mathcal{O}_{Y \cap Z}$ et de
l'opposé du morphisme canonique
$\mathcal{O}_Z \to \mathcal{O}_{Y \cap Z}$. Pour vérifier
l'exactitude, nous pouvons travailler localement pour la topologie étale et la déduire
du cas des schémas
(Morphismes, lemme \ref{morphisms-lemma-scheme-theoretic-union}).

\medskip\noindent
Pour montrer que le diagramme est cocartésien, soient un espace algébrique
$T$ sur $S$ et des morphismes $f : Y \to T$, $g : Z \to T$ qui coïncident sur
$Y \cap Z \to T$. Il s'agit de montrer qu'il existe un unique morphisme
$h : Y \cup Z \to T$ dont les restrictions sont $f$ et $g$.
Pour construire $h$, nous pouvons travailler localement sur $Y \cup Z$ pour la topologie étale
(car $Y \cup Z$, étant un espace algébrique, est un faisceau étale).
Nous pouvons donc supposer que $X$ est un schéma.
Dans ce cas, nous savons que $Y \cup Z$ est la somme amalgamée
de $Y$ et $Z$ le long de $Y \cap Z$ dans la catégorie des schémas
d'après Morphismes, lemme \ref{morphisms-lemma-scheme-theoretic-union}.
Choisissons un schéma $T'$ et un morphisme étale surjectif $T' \to T$.
Posons $Y' = T' \times_{T, f} Y$ et $Z' = T' \times_{T, g} Z$.
Alors $Y'$ et $Z'$ sont des schémas, et nous disposons d'un isomorphisme canonique
$\varphi : Y' \times_Y (Y \cap Z) \to Z' \times_Z (Y \cap Z)$
de schémas. D'après Compléments sur les morphismes, lemme
\ref{more-morphisms-lemma-pushout-along-closed-immersions}
la somme amalgamée $W' = Y' \amalg_{Y' \times_Y (Y \cap Z), \varphi} Z'$
existe dans la catégorie des schémas.
Le morphisme $W' \to Y \cup Z$ est étale d'après
Compléments sur les morphismes, lemme
\ref{more-morphisms-lemma-pushout-along-closed-immersions-properties-above}.
Il est surjectif puisque $Y' \to Y$ et $Z' \to Z$ le sont.
Les morphismes $f' : Y' \to T'$ et $g' : Z' \to T'$
se recollent en un unique morphisme de schémas $h' : W' \to T'$.
Par unicité, le composé $W' \to T' \to T$
descend au morphisme souhaité $h : Y \cup Z \to T$.
Nous omettons certains détails.
\end{proof}






\section{Supports de modules}
\label{section-support}

\noindent
Dans cette section, nous rassemblons quelques résultats élémentaires sur les
supports des modules quasi-cohérents sur les espaces algébriques. Soit $X$ un
espace algébrique. Le support d'un faisceau abélien sur $X_\etale$
a été défini dans Propriétés des espaces, section
\ref{spaces-properties-section-support}.
Nous adoptons la même définition pour les supports de modules.
Le lemme suivant affirme qu'elle coïncide avec la notion
définie pour les modules quasi-cohérents sur les schémas.

\begin{lemma}
\label{lemma-support-covering}
Soit $S$ un schéma. Soit $X$ un espace algébrique sur $S$.
Soit $\mathcal{F}$ un $\mathcal{O}_X$-module quasi-cohérent.
Soit $U$ un schéma et soit $\varphi : U \to X$ un morphisme étale.
Alors
$$
\text{Supp}(\varphi^*\mathcal{F}) = |\varphi|^{-1}(\text{Supp}(\mathcal{F}))
$$
où le membre de gauche est le support de $\varphi^*\mathcal{F}$ en tant que
module quasi-cohérent sur le schéma $U$.
\end{lemma}

\begin{proof}
Soit $u\in U$ un point (usuel) et soit $\overline{x}$ un
point géométrique au-dessus de $u$. D'après
Propriétés des espaces, lemme \ref{spaces-properties-lemma-stalk-quasi-coherent},
nous avons
$(\varphi^*\mathcal{F})_u \otimes_{\mathcal{O}_{U, u}}
\mathcal{O}_{X, \overline{x}} = \mathcal{F}_{\overline{x}}$.
Comme $\mathcal{O}_{U, u} \to \mathcal{O}_{X, \overline{x}}$
est l'hensélisé strict d'après
Propriétés des espaces, lemme
\ref{spaces-properties-lemma-describe-etale-local-ring}
il est fidèlement plat (voir
Compléments d'algèbre, lemme
\ref{more-algebra-lemma-dumb-properties-henselization}).
Par conséquent, $(\varphi^*\mathcal{F})_u = 0$ si et seulement si
$\mathcal{F}_{\overline{x}} = 0$. Cela démontre le lemme.
\end{proof}

\noindent
Pour les modules quasi-cohérents de type fini, le support est fermé,
se vérifie sur les fibres et commute au changement de base.

\begin{lemma}
\label{lemma-support-finite-type}
Soit $S$ un schéma. Soit $X$ un espace algébrique sur $S$.
Soit $\mathcal{F}$ un $\mathcal{O}_X$-module quasi-cohérent de type fini.
Alors
\begin{enumerate}
\item Le support de $\mathcal{F}$ est fermé.
\item Pour tout point géométrique $\overline{x}$ au-dessus de $x \in |X|$, nous avons
$$
x \in \text{Supp}(\mathcal{F})
\Leftrightarrow
\mathcal{F}_{\overline{x}} \not = 0
\Leftrightarrow
\mathcal{F}_{\overline{x}} \otimes_{\mathcal{O}_{X, \overline{x}}}
\kappa(\overline{x}) \not = 0.
$$
\item Pour tout morphisme d'espaces algébriques $f : Y \to X$, l'image inverse
$f^*\mathcal{F}$ est également de type fini et nous avons
$\text{Supp}(f^*\mathcal{F}) = f^{-1}(\text{Supp}(\mathcal{F}))$.
\end{enumerate}
\end{lemma}

\begin{proof}
Choisissons un schéma $U$ et un morphisme étale surjectif $\varphi : U \to X$.
D'après le lemme \ref{lemma-support-covering}, l'image inverse du support de
$\mathcal{F}$ est le support de $\varphi^*\mathcal{F}$, qui est fermé d'après
Morphismes, lemme \ref{morphisms-lemma-support-finite-type}.
Ainsi, (1) découle de la définition de la topologie sur $|X|$.

\medskip\noindent
La première équivalence de (2) est la définition du support.
La seconde équivalence découle du lemme de Nakayama ; voir
Algèbre, lemme \ref{algebra-lemma-NAK}.

\medskip\noindent
Soit $f : Y \to X$ comme en (3). Notons que $f^*\mathcal{F}$ est de type fini
d'après Propriétés des espaces, section
\ref{spaces-properties-section-properties-modules}.
Pour la dernière assertion, soit $\overline{y}$ un point géométrique de $Y$
d'image le point géométrique $\overline{x}$ dans $X$. Rappelons que
$$
(f^*\mathcal{F})_{\overline{y}} =
\mathcal{F}_{\overline{x}} \otimes_{\mathcal{O}_{X, \overline{x}}}
\mathcal{O}_{Y, \overline{y}},
$$
voir Propriétés des espaces, lemme
\ref{spaces-properties-lemma-stalk-pullback-quasi-coherent}.
Par conséquent, $(f^*\mathcal{F})_{\overline{y}} \otimes \kappa(\overline{y})$
est non nul si et seulement si
$\mathcal{F}_{\overline{x}} \otimes \kappa(\overline{x})$ est non nul.
D'après (2), il s'ensuit que $x \in \text{Supp}(\mathcal{F})$ si et seulement
si $y \in \text{Supp}(f^*\mathcal{F})$, ce qui constitue
l'assertion (3).
\end{proof}

\noindent
Notre objectif suivant est de montrer que le support schématique
d'un module quasi-cohérent de type fini (voir
Morphismes, définition \ref{morphisms-definition-scheme-theoretic-support})
a également un sens pour les modules quasi-cohérents de type fini sur les
espaces algébriques.

\begin{lemma}
\label{lemma-scheme-theoretic-support}
Soit $S$ un schéma. Soit $X$ un espace algébrique sur $S$.
Soit $\mathcal{F}$ un $\mathcal{O}_X$-module quasi-cohérent de type fini.
Il existe un plus petit sous-espace fermé $i : Z \to X$ tel qu'il
existe un $\mathcal{O}_Z$-module quasi-cohérent $\mathcal{G}$ vérifiant
$i_*\mathcal{G} \cong \mathcal{F}$. De plus :
\begin{enumerate}
\item Si $U$ est un schéma et si $\varphi : U \to X$ est un morphisme étale,
alors $Z \times_X U$ est le support schématique de $\varphi^*\mathcal{F}$.
\item Le faisceau quasi-cohérent $\mathcal{G}$ est unique à isomorphisme
unique près.
\item Le faisceau quasi-cohérent $\mathcal{G}$ est de type fini.
\item Les supports de $\mathcal{G}$ et de $\mathcal{F}$ sont tous deux $|Z|$.
\end{enumerate}
\end{lemma}

\begin{proof}
Choisissons un schéma $U$ et un morphisme étale surjectif $\varphi : U \to X$.
Posons $R = U \times_X U$, de projections $s, t : R \to U$.
Soit $i' : Z' \to U$ le support schématique de $\varphi^*\mathcal{F}$,
et soit $\mathcal{G}'$ le $\mathcal{O}_{Z'}$-module quasi-cohérent de type fini
(unique à isomorphisme unique près)
tel que $i'_*\mathcal{G}' = \varphi^*\mathcal{F}$ ; voir
Morphismes, lemme \ref{morphisms-lemma-scheme-theoretic-support}.
Comme $s^*\varphi^*\mathcal{F} = t^*\varphi^*\mathcal{F}$, nous avons
$R' = s^{-1}Z' = t^{-1}Z'$ comme sous-schémas fermés de $R$, d'après
Morphismes, lemme \ref{morphisms-lemma-flat-pullback-support}.
Nous pouvons donc appliquer Propriétés des espaces, lemme
\ref{spaces-properties-lemma-subspaces-presentation}
pour obtenir un sous-espace fermé $i : Z \to X$ dont l'image inverse sur $U$ est $Z'$.
Notons $s', t' : R' \to Z'$ les projections et
$j' : R' \to R$ l'immersion fermée considérée ; nous avons
$$
j'_* (s')^*\mathcal{G}' = s^* i'_*\mathcal{G}' =
s^*\varphi^*\mathcal{F} = t^*\varphi^*\mathcal{F} =
t^*i'_*\mathcal{G}' = j'_*(t')^*\mathcal{G}'
$$
(la première et la dernière égalité résultant de Cohomologie des schémas,
lemme \ref{coherent-lemma-flat-base-change-cohomology}).
Le lemme
Morphismes, lemme \ref{morphisms-lemma-flat-pullback-support},
appliqué à $R' \to R$, fournit donc, par unicité, un isomorphisme
$\alpha : (t')^*\mathcal{G}' \to (s')^*\mathcal{G}'$
compatible avec l'isomorphisme canonique
$t^*\varphi^*\mathcal{F} = s^*\varphi^*\mathcal{F}$
par l'intermédiaire de $j'_*$. Il est clair que $\alpha$ satisfait la condition de cocycle ;
nous pouvons donc appliquer
Propriétés des espaces, Proposition
\ref{spaces-properties-proposition-quasi-coherent}
pour obtenir un module quasi-cohérent $\mathcal{G}$ sur $Z$ dont la restriction
à $Z'$ s'identifie à $\mathcal{G}'$, compatiblement à $\alpha$.
En utilisant à nouveau l'équivalence de la proposition précitée
(cette fois pour $X$), nous concluons que $i_*\mathcal{G} \cong \mathcal{F}$.

\medskip\noindent
Cela démontre l'existence. Les autres propriétés du lemme découlent
de la comparaison avec le résultat pour les schémas, au moyen du
lemme \ref{lemma-support-covering}.
Nous omettons les démonstrations détaillées.
\end{proof}

\begin{definition}
\label{definition-scheme-theoretic-support}
Soit $S$ un schéma. Soit $X$ un espace algébrique sur $S$.
Soit $\mathcal{F}$ un $\mathcal{O}_X$-module quasi-cohérent de type fini.
Le {\it support schématique de $\mathcal{F}$} est le sous-espace fermé
$Z \subset X$ construit dans le lemme \ref{lemma-scheme-theoretic-support}.
\end{definition}

\noindent
Dans cette situation, nous considérons souvent $\mathcal{F}$ comme un faisceau
quasi-cohérent de type fini sur $Z$ (au moyen de l'équivalence de catégories du
lemme \ref{lemma-i-star-equivalence}).






\section{Image schématique}
\label{section-scheme-theoretic-image}

\noindent
Attention : une partie des résultats de cette section est d'une généralité extrême
et se comporte autrement qu'on pourrait s'y attendre.

\begin{lemma}
\label{lemma-scheme-theoretic-image}
\begin{slogan}
L'image schématique d'un morphisme d'espaces algébriques existe.
\end{slogan}
Soit $S$ un schéma. Soit $f : X \to Y$ un morphisme d'espaces algébriques
sur $S$. Il existe un sous-espace fermé $Z \subset Y$ tel que $f$ se factorise
par $Z$ et tel que, pour tout autre sous-espace fermé $Z' \subset Y$ tel que
$f$ se factorise par $Z'$, nous ayons $Z \subset Z'$.
\end{lemma}

\begin{proof}
Posons $\mathcal{I} = \Ker(\mathcal{O}_Y \to f_*\mathcal{O}_X)$.
Si $\mathcal{I}$ est quasi-cohérent, il suffit de prendre pour $Z$ le
sous-espace fermé défini par $\mathcal{I}$ ; voir
le lemme \ref{lemma-closed-immersion-ideals}.
En général, il faut, pour démontrer le lemme, établir qu'il existe
un plus grand faisceau quasi-cohérent d'idéaux $\mathcal{I}'$ contenu dans
$\mathcal{I}$.
Cela découle du lemme \ref{lemma-largest-quasi-coherent-subsheaf}.
\end{proof}

\noindent
Supposons que, dans la situation du lemme \ref{lemma-scheme-theoretic-image}
ci-dessus, $X$ et $Y$ soient représentables. Alors le sous-espace fermé $Z \subset Y$
obtenu dans le lemme coïncide avec le sous-schéma fermé $Z \subset Y$ obtenu dans
Morphismes, lemme \ref{morphisms-lemma-scheme-theoretic-image}.
En effet, les sous-espaces fermés (ou sous-schémas fermés) sont en correspondance,
en renversant les inclusions, avec les faisceaux quasi-cohérents d'idéaux sur
$Y_\etale$ et $Y$. Comme les catégories des modules quasi-cohérents
sur $Y_\etale$ et $Y$ sont les mêmes
(Propriétés des espaces, section \ref{spaces-properties-section-quasi-coherent}),
nous pouvons conclure. Ainsi, la définition suivante coïncide avec la définition
antérieure pour les morphismes de schémas.

\begin{definition}
\label{definition-scheme-theoretic-image}
Soit $S$ un schéma. Soit $f : X \to Y$ un morphisme d'espaces algébriques
sur $S$. L'{\it image schématique} de $f$ est le plus petit sous-espace
fermé $Z \subset Y$ par lequel $f$
se factorise ; voir le lemme \ref{lemma-scheme-theoretic-image} ci-dessus.
\end{definition}

\noindent
Nous désignons souvent simplement par $f : X \to Z$ la factorisation de $f$.
Si le morphisme $f$ n'est pas quasi-compact, alors, en général, la
construction de l'image schématique ne commute pas à la
restriction aux sous-espaces ouverts de $Y$.

\begin{lemma}
\label{lemma-quasi-compact-scheme-theoretic-image}
Soit $S$ un schéma.
Soit $f : X \to Y$ un morphisme d'espaces algébriques sur $S$.
Soit $Z \subset Y$ l'image schématique de $f$.
Si $f$ est quasi-compact, alors
\begin{enumerate}
\item le faisceau d'idéaux
$\mathcal{I} = \Ker(\mathcal{O}_Y \to f_*\mathcal{O}_X)$
est quasi-cohérent,
\item l'image schématique $Z$ est le sous-espace fermé
correspondant à $\mathcal{I}$,
\item pour tout morphisme étale $V \to Y$, l'image schématique de
$X \times_Y V \to V$ est égale à $Z \times_Y V$, et
\item l'image $|f|(|X|) \subset |Z|$ est une partie dense de $|Z|$.
\end{enumerate}
\end{lemma}

\begin{proof}
Pour démontrer (3), il suffit de démontrer (1) et (2), puisque la
formation de $\mathcal{I}$ commute à la localisation étale.
Si (1) est vraie, nous avons établi (2) dans la preuve du lemme
\ref{lemma-scheme-theoretic-image}. Montrons que $\mathcal{I}$ est quasi-cohérent.
Comme la propriété d'être quasi-cohérent est locale pour la topologie étale, nous pouvons
supposer que $Y$ est un schéma affine. Comme $f$ est quasi-compact,
nous pouvons trouver un schéma affine $U$ et un morphisme étale surjectif
$U \to X$. Notons $f'$ le composé $U \to X \to Y$.
Alors $f_*\mathcal{O}_X$ est un sous-faisceau de $f'_*\mathcal{O}_U$,
et donc $\mathcal{I} = \Ker(\mathcal{O}_Y \to f'_*\mathcal{O}_U)$.
D'après le lemme \ref{lemma-pushforward},
le faisceau $f'_*\mathcal{O}_U$ est quasi-cohérent sur $Y$. Ainsi, $\mathcal{I}$
est quasi-cohérent comme noyau d'un morphisme de modules quasi-cohérents.
Enfin, la partie (4) découle des parties (1), (2) et (3), car l'idéal
$\mathcal{I}$ sera l'idéal unité en tout point de $|Y|$ qui
n'appartient pas à l'adhérence de $|f|(|X|)$.
\end{proof}

\begin{lemma}
\label{lemma-scheme-theoretic-image-reduced}
Soit $S$ un schéma. Soit $f : X \to Y$ un morphisme d'espaces algébriques
sur $S$. Supposons que $X$ soit réduit. Alors
\begin{enumerate}
\item l'image schématique $Z$ de $f$ est l'adhérence $\overline{|f|(|X|)}$, munie de la
structure réduite induite, et
\item pour tout morphisme étale $V \to Y$, l'image schématique de
$X \times_Y V \to V$ est égale à $Z \times_Y V$.
\end{enumerate}
\end{lemma}

\begin{proof}
La partie (1) est vraie parce que l'adhérence $\overline{|f|(|X|)}$, munie de la
structure réduite induite, est le plus petit sous-espace fermé
de $Y$ par lequel $f$ se factorise ; voir
Propriétés des espaces, lemme \ref{spaces-properties-lemma-map-into-reduction}.
La partie (2) découle de (1), du fait que l'application $|V| \to |Y|$ est ouverte et du
fait que la réduction est préservée par localisation étale.
\end{proof}

\begin{lemma}
\label{lemma-reach-points-scheme-theoretic-image}
Soit $S$ un schéma.
Soit $f : X \to Y$ un morphisme quasi-compact d'espaces algébriques sur $S$.
Soit $Z$ l'image schématique de $f$.
Soit $z \in |Z|$. Il existe un anneau de valuation $A$ de
corps des fractions $K$ et un diagramme commutatif
$$
\xymatrix{
\Spec(K) \ar[rr] \ar[d] & & X \ar[d] \ar[ld] \\
\Spec(A) \ar[r] & Z \ar[r] & Y
}
$$
tel que le point fermé de $\Spec(A)$ s'envoie sur $z$.
\end{lemma}

\begin{proof}
Choisissons un schéma affine $V$, un point $z' \in V$
et un morphisme étale $V \to Y$ envoyant $z'$ sur $z$.
Soit $Z' \subset V$ l'image schématique de $X \times_Y V \to V$.
D'après le lemme \ref{lemma-quasi-compact-scheme-theoretic-image}, nous avons
$Z' = Z \times_Y V$. Ainsi, $z' \in Z'$.
Comme $f$ est quasi-compact et $V$ affine, nous voyons que
$X \times_Y V$ est quasi-compact. Il existe donc
un schéma affine $W$ et un morphisme étale surjectif
$W \to X \times_Y V$. Alors $Z' \subset V$ est également
l'image schématique de $W \to V$.
D'après Morphismes, lemme \ref{morphisms-lemma-reach-points-scheme-theoretic-image},
nous pouvons choisir un diagramme
$$
\xymatrix{
\Spec(K) \ar[r] \ar[d] &
W \ar[r] \ar[d] &
X \times_Y V \ar[d] \ar[r] &
X \ar[d] \\
\Spec(A) \ar[r] &
Z' \ar[r] &
V \ar[r] &
Y
}
$$
tel que le point fermé de $\Spec(A)$ s'envoie sur $z'$.
En composant avec $Z' \to Z$ et $W \to X \times_Y V \to X$,
nous obtenons une solution.
\end{proof}

\begin{lemma}
\label{lemma-factor-factor}
Soit $S$ un schéma. Soit
$$
\xymatrix{
X_1 \ar[d] \ar[r]_{f_1} & Y_1 \ar[d] \\
X_2 \ar[r]^{f_2} & Y_2
}
$$
un diagramme commutatif d'espaces algébriques sur $S$.
Soient $Z_i \subset Y_i$, $i = 1, 2$, les
images schématiques de $f_i$. Alors le morphisme
$Y_1 \to Y_2$ induit un morphisme $Z_1 \to Z_2$ et un
diagramme commutatif
$$
\xymatrix{
X_1 \ar[r] \ar[d] & Z_1 \ar[d] \ar[r] & Y_1 \ar[d] \\
X_2 \ar[r] & Z_2 \ar[r] & Y_2
}
$$
\end{lemma}

\begin{proof}
L'image inverse schématique de $Z_2$ dans $Y_1$
est un sous-espace fermé de $Y_1$ par lequel
$f_1$ se factorise. Par conséquent, $Z_1$ y est contenu.
Cela démontre le lemme.
\end{proof}

\begin{lemma}
\label{lemma-scheme-theoretic-image-of-partial-section}
Soit $S$ un schéma.
Soit $f : X \to Y$ un morphisme séparé d'espaces algébriques sur $S$.
Soit $V \subset Y$ un sous-espace ouvert tel que $V \to Y$ soit quasi-compact.
Soit $s : V \to X$ un morphisme tel que $f \circ s = \text{id}_V$.
Soit $Y'$ l'image schématique de $s$.
Alors $Y' \to Y$ est un isomorphisme au-dessus de $V$.
\end{lemma}

\begin{proof}
D'après le lemme \ref{lemma-quasi-compact-permanence},
le morphisme $s : V \to X$ est quasi-compact.
La construction de l'image schématique $Y'$
de $s$ commute donc à la restriction aux ouverts, d'après le
lemme \ref{lemma-quasi-compact-scheme-theoretic-image}.
En particulier, nous voyons que $Y' \cap f^{-1}(V)$ est
l'image schématique d'une section du morphisme séparé
$f^{-1}(V) \to V$. Comme une section d'un morphisme séparé
est une immersion fermée
(lemme \ref{lemma-section-immersion}),
nous concluons que
$Y' \cap f^{-1}(V) \to V$ est un isomorphisme, comme souhaité.
\end{proof}








\section{Adhérence schématique et densité}
\label{section-scheme-theoretic-closure}

\noindent
Cette section est l'analogue de
Morphismes, section \ref{morphisms-section-scheme-theoretic-closure}.

\begin{lemma}
\label{lemma-scheme-theoretically-dense-representable}
Soit $S$ un schéma. Soit $W \subset S$ un sous-schéma ouvert
schématiquement dense
(Morphismes, définition \ref{morphisms-definition-scheme-theoretically-dense}).
Soit $f : X \to S$ un morphisme de schémas plat, localement de
présentation finie et localement quasi-fini.
Alors $f^{-1}(W)$ est schématiquement dense dans $X$.
\end{lemma}

\begin{proof}
Nous utiliserons la caractérisation de Morphismes, lemme
\ref{morphisms-lemma-characterize-scheme-theoretically-dense}.
Supposons que $V \subset X$ soit un ouvert et que $g \in \Gamma(V, \mathcal{O}_V)$
soit une fonction dont la restriction à $f^{-1}(W) \cap V$ est nulle.
Il faut montrer que $g = 0$. Supposons $g \not = 0$ afin d'obtenir une
contradiction. D'après
Compléments sur les morphismes, lemme \ref{more-morphisms-lemma-go-down-with-annihilators}
nous pouvons, quitte à rétrécir $V$, trouver un ouvert $U \subset S$ s'insérant dans un
diagramme commutatif
$$
\xymatrix{
V \ar[r] \ar[d]_\pi & X \ar[d]^f \\
U \ar[r] & S,
}
$$
un sous-faisceau quasi-cohérent $\mathcal{F} \subset \mathcal{O}_U$, un entier
$r > 0$ et un morphisme injectif de $\mathcal{O}_U$-modules
$\mathcal{F}^{\oplus r} \to \pi_*\mathcal{O}_V$
dont l'image contient $g|_V$. Soit
$(g_1, \ldots, g_r) \in \Gamma(U, \mathcal{F}^{\oplus r})$ un antécédent de $g$.
Alors $g_i|_{W \cap U} = 0$, car $g|_{f^{-1}W \cap V} = 0$.
Par conséquent, $g_i = 0$, puisque $\mathcal{F} \subset \mathcal{O}_U$ et que
$W$ est schématiquement dense dans $S$.
Il s'ensuit que $g = 0$, ce qui est la contradiction cherchée.
\end{proof}

\begin{lemma}
\label{lemma-scheme-theoretically-dense}
Soit $S$ un schéma.
Soit $X$ un espace algébrique sur $S$.
Soit $U \subset X$ un sous-espace ouvert.
Les conditions suivantes sont équivalentes :
\begin{enumerate}
\item pour tout morphisme étale $\varphi : V \to X$ (d'espaces algébriques),
l'adhérence schématique de $\varphi^{-1}(U)$ dans $V$ est égale à $V$,
\item il existe un schéma $V$ et un morphisme étale surjectif
$\varphi : V \to X$ tels que l'adhérence schématique de
$\varphi^{-1}(U)$ dans $V$ soit égale à $V$,
\end{enumerate}
\end{lemma}

\begin{proof}
Observons que, si $V \to V'$ est un morphisme entre espaces algébriques étales
sur $X$, et si $Z \subset V$, resp.\ $Z' \subset V'$, sont les adhérences schématiques
de $U \times_X V$, resp.\ $U \times_X V'$, dans $V$, resp.\ $V'$,
alors $Z$ s'envoie dans $Z'$. Ainsi, si $V \to V'$ est étale et surjectif,
$Z = V$ entraîne $Z' = V'$. Ensuite, notons qu'un morphisme étale est
plat, localement de présentation finie et localement quasi-fini
(voir Morphismes, section \ref{morphisms-section-etale}).
Le lemme \ref{lemma-scheme-theoretically-dense-representable}
implique donc que, si $V$ et $V'$ sont des schémas, $Z' = V'$ entraîne
$Z = V$. Un argument formel utilisant le fait que tout espace algébrique admet un
recouvrement étale par un schéma montre que (1) et (2) sont équivalentes.
\end{proof}

\noindent
Il résulte du
lemme \ref{lemma-scheme-theoretically-dense}
que la définition suivante est compatible avec celle
du cas des schémas.

\begin{definition}
\label{definition-scheme-theoretically-dense}
Soit $S$ un schéma.
Soit $X$ un espace algébrique sur $S$.
Soit $U \subset X$ un sous-espace ouvert.
\begin{enumerate}
\item L'image schématique du morphisme $U \to X$
est appelée l'{\it adhérence schématique de $U$ dans $X$}.
\item Nous disons que $U$ est {\it schématiquement dense dans $X$}
si les conditions équivalentes du
lemme \ref{lemma-scheme-theoretically-dense} sont satisfaites.
\end{enumerate}
\end{definition}

\noindent
Avec cette définition, il n'est {\bf pas} vrai que $U$ soit schématiquement
dense dans $X$ si et seulement si l'adhérence schématique
de $U$ est $X$. C'est quelque peu inélégant. Nous verrons cependant que,
sous des conditions de finitude convenables, cette équivalence est vraie.

\begin{lemma}
\label{lemma-scheme-theoretically-dense-quasi-compact}
Soit $S$ un schéma. Soit $X$ un espace algébrique sur $S$.
Soit $U \subset X$ un sous-espace ouvert.
Si $U \to X$ est quasi-compact, alors $U$
est schématiquement dense dans $X$ si et seulement si l'adhérence schématique
de $U$ dans $X$ est $X$.
\end{lemma}

\begin{proof}
Cela découle de la partie (3) du lemme \ref{lemma-quasi-compact-scheme-theoretic-image}.
\end{proof}

\begin{lemma}
\label{lemma-characterize-scheme-theoretically-dense}
Soit $S$ un schéma.
Soit $j : U \to X$ une immersion ouverte d'espaces algébriques sur $S$.
Alors $U$ est schématiquement dense dans $X$ si et seulement si
$\mathcal{O}_X \to j_*\mathcal{O}_U$ est injectif.
\end{lemma}

\begin{proof}
Si $\mathcal{O}_X \to j_*\mathcal{O}_U$ est injectif,
il en va de même après restriction à tout
espace algébrique $V$ étale sur $X$.
L'adhérence schématique de $U \times_X V$ dans $V$
est donc égale à $V$ ; voir la preuve du
lemme \ref{lemma-scheme-theoretic-image}.
Réciproquement, supposons que l'adhérence schématique
de $U \times_X V$ soit égale à $V$ pour tout $V$ étale sur $X$.
Supposons que $\mathcal{O}_X \to j_*\mathcal{O}_U$ ne soit pas injectif.
Nous pouvons alors trouver un schéma affine, disons $V = \Spec(A)$, étale sur $X$,
et un élément non nul $f \in A$ tel que $f$ s'envoie sur zéro dans
$\Gamma(V \times_X U, \mathcal{O})$. Dans ce cas, l'adhérence schématique
de $V \times_X U$ dans $V$ est clairement contenue dans $\Spec(A/(f))$,
ce qui est absurde.
\end{proof}

\begin{lemma}
\label{lemma-intersection-scheme-theoretically-dense}
Soit $S$ un schéma. Soit $X$ un espace algébrique sur $S$.
Si $U$ et $V$ sont des sous-espaces ouverts
schématiquement denses dans $X$, alors $U \cap V$ l'est aussi.
\end{lemma}

\begin{proof}
Soit $W \to X$ un morphisme étale quelconque. Considérons l'application
$\mathcal{O}(W) \to \mathcal{O}(W \times_X V)
\to \mathcal{O}(W \times_X (V \cap U))$.
D'après le lemme \ref{lemma-characterize-scheme-theoretically-dense},
les deux applications sont injectives. Leur composée est donc injective.
Par conséquent, d'après le lemme \ref{lemma-characterize-scheme-theoretically-dense},
$U \cap V$ est schématiquement dense dans $X$.
\end{proof}

\begin{lemma}
\label{lemma-quasi-compact-immersion}
Soit $S$ un schéma. Soit $h : Z \to X$ une immersion d'espaces algébriques
sur $S$. Supposons que $Z \to X$ soit quasi-compact ou que $Z$ soit réduit.
Soit $\overline{Z} \subset X$ l'image schématique de $h$.
Alors le morphisme $Z \to \overline{Z}$ est une immersion ouverte
qui identifie $Z$ à un sous-espace ouvert schématiquement dense
de $\overline{Z}$. De plus, $Z$ est topologiquement
dense dans $\overline{Z}$.
\end{lemma}

\begin{proof}
Dans les deux cas, la formation de $\overline{Z}$ commute à la
localisation étale ; voir les lemmes
\ref{lemma-quasi-compact-scheme-theoretic-image} et
\ref{lemma-scheme-theoretic-image-reduced}.
Le présent lemme découle donc du cas des schémas ; voir
Morphismes, lemme \ref{morphisms-lemma-quasi-compact-immersion}.
\end{proof}

\begin{lemma}
\label{lemma-equality-of-morphisms}
Soit $S$ un schéma. Soit $B$ un espace algébrique sur $S$.
Soient $f, g : X \to Y$ des morphismes d'espaces algébriques sur $B$.
Soit $U \subset X$ un sous-espace ouvert tel que
$f|_U = g|_U$. Si l'adhérence schématique de $U$
dans $X$ est $X$ et si $Y \to B$ est séparé, alors $f = g$.
\end{lemma}

\begin{proof}
Comme $Y \to B$ est séparé, le produit fibré
$Y \times_{\Delta, Y \times_B Y, (f, g)} X$ est un sous-espace fermé $Z \subset X$. 
Comme $f|_U = g|_U$, nous avons $U \subset Z$. Donc $Z = X$,
car l'hypothèse sur $U$ affirme que son adhérence schématique est $X$.
\end{proof}






\section{Morphismes dominants}
\label{section-dominant}

\noindent
Nous reprenons la définition d'un morphisme dominant de schémas afin d'obtenir
la notion de morphisme dominant d'espaces algébriques. Nous avertissons le
lecteur que cette définition ne se comporte pas bien à moins que le morphisme
ne soit quasi-compact et que les espaces algébriques ne satisfassent certains
axiomes de séparation.

\begin{definition}
\label{definition-dominant}
Soit $S$ un schéma. Un morphisme $f : X \to Y$ d'espaces algébriques sur $S$ est
dit {\it dominant} si l'image de $|f| : |X| \to |Y|$ est dense dans $|Y|$.
\end{definition}






\section{Morphismes universellement injectifs}
\label{section-universally-injective}

\noindent
Nous avons déjà défini à la section \ref{section-representable}
ce que signifie, pour un morphisme représentable d'espaces algébriques,
être universellement injectif. Pour un corps $K$ sur $S$ (rappelons que cela
signifie que l'on s'est donné un morphisme structural $\Spec(K) \to S$) et un
espace algébrique $X$ sur $S$, nous posons
$X(K) = \Mor_S(\Spec(K), X)$. Nous traduisons d'abord la
condition portant sur les morphismes représentables en une condition sur le foncteur
des points.

\begin{lemma}
\label{lemma-universally-injective-representable}
Soit $S$ un schéma. Soit $f : X \to Y$ un morphisme représentable
d'espaces algébriques sur $S$. Alors $f$ est universellement injectif
(au sens de la section \ref{section-representable})
si et seulement si, pour tout corps $K$, l'application $X(K) \to Y(K)$ est injective.
\end{lemma}

\begin{proof}
Nous utiliserons
Morphismes, lemme \ref{morphisms-lemma-universally-injective}
sans autre mention.
Supposons que $f$ soit universellement injectif. Alors, pour tout corps $K$ et tout
morphisme $\Spec(K) \to Y$, le morphisme de schémas
$\Spec(K) \times_Y X \to \Spec(K)$ est universellement injectif.
Il existe donc au plus une section du morphisme
$\Spec(K) \times_Y X \to \Spec(K)$. Ainsi, l'application
$X(K) \to Y(K)$ est injective. Réciproquement, supposons que, pour tout corps $K$,
l'application $X(K) \to Y(K)$ soit injective. Soit $T \to Y$ un morphisme d'un
schéma vers $Y$, et considérons le changement de base $f_T : T \times_Y X \to T$.
Pour tout corps $K$, on a
$$
(T \times_Y X)(K) = T(K) \times_{Y(K)} X(K)
$$
par définition du produit fibré ; ainsi, l'injectivité de
$X(K) \to Y(K)$ entraîne celle de
$(T \times_Y X)(K) \to T(K)$, ce qui signifie que $f_T$ est universellement injectif,
comme souhaité.
\end{proof}

\noindent
Traduisons maintenant l'injectivité de l'application sur les points à valeurs
dans des corps en une condition plus géométrique.

\begin{lemma}
\label{lemma-universally-injective}
Soit $S$ un schéma.
Soit $f : X \to Y$ un morphisme d'espaces algébriques sur $S$.
Les conditions suivantes sont équivalentes :
\begin{enumerate}
\item l'application $X(K) \to Y(K)$ est injective pour tout corps $K$ sur $S$ ;
\item pour tout morphisme $Y' \to Y$ d'espaces algébriques sur $S$,
l'application induite $|Y' \times_Y X| \to |Y'|$ est injective ; et
\item le morphisme diagonal $X \to X \times_Y X$ est surjectif.
\end{enumerate}
\end{lemma}

\begin{proof}
Supposons (1). Soit $g : Y' \to Y$ un morphisme d'espaces
algébriques, et notons $f' : Y' \times_Y X \to Y'$ le changement de base de $f$.
Soient $K_i$, $i = 1, 2$, des corps et soient
$\varphi_i : \Spec(K_i) \to Y' \times_Y X$ des morphismes
tels que $f' \circ \varphi_1$ et $f' \circ \varphi_2$ définissent le
même élément de $|Y'|$. Par définition, cela signifie qu'il existe un
corps $\Omega$ et des plongements $\alpha_i : K_i \subset \Omega$ tels que
les deux morphismes
$f' \circ \varphi_i \circ \alpha_i : \Spec(\Omega) \to Y'$ soient égaux.
Voici le diagramme commutatif correspondant :
$$
\xymatrix{
\Spec(\Omega) \ar@/_5ex/[ddrr] \ar[rd]^{\alpha_1} \ar[r]_{\alpha_2} &
\Spec(K_2) \ar[rd]^{\varphi_2} \\
& \Spec(K_1) \ar[r]^{\varphi_1} &
Y' \times_Y X  \ar[d]^{f'} \ar[r]^{g'} &
 X \ar[d]^f \\
& & Y' \ar[r]^g & Y.
}
$$
En particulier, les composés $g \circ f' \circ \varphi_i \circ \alpha_i$
sont égaux. D'après (1), cela implique que les morphismes
$g' \circ \varphi_i \circ \alpha_i$ sont égaux, où $g' : Y' \times_Y X \to X$
est la projection. Par la propriété universelle du produit fibré, nous en déduisons
que les morphismes
$\varphi_i \circ \alpha_i : \Spec(\Omega) \to Y' \times_Y X$ sont
égaux. Autrement dit, $\varphi_1$ et $\varphi_2$ définissent le même point
de $Y' \times_Y X$. Nous en concluons que (2) est vérifiée.

\medskip\noindent
Supposons (2). Soit $K$ un corps sur $S$, et soient $a, b : \Spec(K) \to X$
deux morphismes tels que $f \circ a = f \circ b$. Notons
$c : \Spec(K) \to Y$ leur valeur commune. Par hypothèse,
$|\Spec(K) \times_{c, Y} X| \to |\Spec(K)|$ est injective.
Cela signifie qu'il existe un corps $\Omega$ et des plongements
$\alpha_i : K \to \Omega$ tels que
$$
\xymatrix{
\Spec(\Omega) \ar[r]_{\alpha_1} \ar[d]_{\alpha_2} &
\Spec(K) \ar[d]^a \\
\Spec(K) \ar[r]^-b &
\Spec(K) \times_{c, Y} X
}
$$
soit commutatif. En composant avec la projection sur $\Spec(K)$,
on voit que $\alpha_1 = \alpha_2$. Notons leur valeur commune $\alpha$.
Alors la famille $\{\alpha : \Spec(\Omega) \to \Spec(K)\}$
est couvrante pour la topologie fpqc de $\Spec(K)$ et les deux morphismes
$a, b$ deviennent égaux après ce changement de base. D'après
Propriétés des espaces, Proposition
\ref{spaces-properties-proposition-sheaf-fpqc},
on a $a = b$. Nous en concluons que (1) est vérifiée.

\medskip\noindent
Supposons (3). Soient $x, x' \in |X|$ deux points tels que
$f(x) = f(x')$ dans $|Y|$. D'après
Propriétés des espaces, lemme \ref{spaces-properties-lemma-points-cartesian},
il existe $x'' \in |X \times_Y X|$ dont les projections
sont $x$ et $x'$. Par hypothèse et d'après
Propriétés des espaces,
lemme \ref{spaces-properties-lemma-characterize-surjective},
il existe $x''' \in |X|$ tel que $\Delta_{X/Y}(x''') = x''$.
Ainsi $x = x'$. Autrement dit, $f$ est injectif sur les points.
Puisque la condition (3) est stable par changement de base, on voit que $f$
vérifie (2).

\medskip\noindent
Supposons (2). Alors, en particulier, $|X \times_Y X| \to |X|$ est injective,
ce qui entraîne immédiatement que
$|\Delta_{X/Y}| : |X| \to |X \times_Y X|$ est surjective ; cela
entraîne que $\Delta_{X/Y}$ est surjectif d'après
Propriétés des espaces,
lemme \ref{spaces-properties-lemma-characterize-surjective}.
\end{proof}

\noindent
D'après les deux lemmes précédents, la définition suivante est compatible avec
la notion déjà définie de morphisme représentable universellement injectif
d'espaces algébriques.

\begin{definition}
\label{definition-universally-injective}
Soit $S$ un schéma. Soit $f : X \to Y$ un morphisme d'espaces
algébriques sur $S$. Nous disons que $f$ est {\it universellement injectif} si,
pour tout morphisme $Y' \to Y$, l'application induite
$|Y' \times_Y X| \to |Y'|$ est injective.
\end{definition}

\noindent
Précisément, cela signifie qu'une, et donc chacune, des conditions équivalentes du
lemme \ref{lemma-universally-injective}
est satisfaite.

\begin{remark}
\label{remark-universally-injective-not-separated}
Un morphisme universellement injectif de schémas est séparé ; voir
Morphismes, lemme
\ref{morphisms-lemma-universally-injective-separated}.
Ce n'est pas le cas pour les morphismes d'espaces algébriques.
En effet, l'espace algébrique
$X = \mathbf{A}^1_k/\{x \sim -x \mid x \not = 0\}$
construit dans
Espaces, exemple \ref{spaces-example-affine-line-involution}
est muni d'un morphisme $X \to \mathbf{A}^1_k$ qui envoie
le point de coordonnée $x$ sur le point de coordonnée $x^2$.
C'est un isomorphisme hors de $0$, et il existe un unique point
de $X$ au-dessus de $0$. Comme $X$ n'est pas séparé, il s'agit d'un morphisme
universellement injectif d'espaces algébriques qui n'est pas séparé.
\end{remark}

\begin{lemma}
\label{lemma-base-change-universally-injective}
Tout changement de base d'un morphisme universellement injectif est universellement injectif.
\end{lemma}

\begin{proof}
Omis. Indication : c'est formel.
\end{proof}

\begin{lemma}
\label{lemma-universally-injective-local}
Soit $S$ un schéma.
Soit $f : X \to Y$ un morphisme d'espaces algébriques sur $S$.
Les conditions suivantes sont équivalentes :
\begin{enumerate}
\item $f$ est universellement injectif ;
\item pour tout schéma $Z$ et tout morphisme $Z \to Y$, le morphisme
$Z \times_Y X \to Z$ est universellement injectif ;
\item pour tout schéma affine $Z$ et tout morphisme
$Z \to Y$, le morphisme $Z \times_Y X \to Z$ est universellement injectif ;
\item il existe un schéma $Z$ et un morphisme surjectif
$Z \to Y$ tels que $Z \times_Y X \to Z$ soit universellement injectif ; et
\item il existe un recouvrement de Zariski $Y = \bigcup Y_i$ tel que
chacun des morphismes $f^{-1}(Y_i) \to Y_i$ soit universellement injectif.
\end{enumerate}
\end{lemma}

\begin{proof}
Nous utiliserons sans autre mention dans cette démonstration le fait que la propriété
d'être universellement injectif est préservée par changement de base
(lemme \ref{lemma-base-change-universally-injective}).
Il est clair que (1) $\Rightarrow$ (2) $\Rightarrow$ (3)
$\Rightarrow$ (4).

\medskip\noindent
Supposons donné $g : Z \to Y$ comme dans (4). Soit $y : \Spec(K) \to Y$ un
morphisme du spectre d'un corps vers $Y$. Par hypothèse, on
peut trouver une extension de corps $\alpha : K \subset K'$ et un morphisme
$z : \Spec(K') \to Z$ tels que $y \circ \alpha = g \circ z$
(avec l'abus de notation évident). Par hypothèse, le
morphisme $Z \times_Y X \to Z$ est universellement injectif ; il existe donc
au plus un
relèvement de $g \circ z : \Spec(K') \to Y$ en un morphisme vers $X$.
Comme la famille $\{\alpha : \Spec(K') \to \Spec(K)\}$ est
couvrante pour la topologie fpqc, il existe au plus un relèvement de
$y : \Spec(K) \to Y$ en un morphisme vers $X$ ; voir
Propriétés des espaces, Proposition
\ref{spaces-properties-proposition-sheaf-fpqc}.
Ainsi, (1) est vérifiée.

\medskip\noindent
Nous omettons de vérifier que (5) est équivalente à (1).
\end{proof}

\begin{lemma}
\label{lemma-composition-universally-injective}
Tout composé de morphismes universellement injectifs est universellement injectif.
\end{lemma}

\begin{proof}
Omis.
\end{proof}











\section{Morphismes affines}
\label{section-affine}

\noindent
Nous avons déjà défini à la section \ref{section-representable}
ce que signifie, pour un morphisme représentable d'espaces algébriques,
être affine.

\begin{lemma}
\label{lemma-affine-representable}
Soit $S$ un schéma. Soit $f : X \to Y$ un morphisme représentable
d'espaces algébriques sur $S$. Alors
$f$ est affine (au sens de la section \ref{section-representable})
si et seulement si, pour tout schéma affine $Z$
et tout morphisme $Z \to Y$, le schéma $X \times_Y Z$ est affine.
\end{lemma}

\begin{proof}
Cela résulte directement de la définition d'un morphisme affine de schémas
(Morphismes, définition \ref{morphisms-definition-affine}).
\end{proof}

\noindent
Nous pouvons ainsi poser la définition suivante.

\begin{definition}
\label{definition-affine}
Soit $S$ un schéma.
Soit $f : X \to Y$ un morphisme d'espaces algébriques sur $S$.
Nous disons que $f$ est {\it affine} si, pour tout schéma affine $Z$ et tout
morphisme $Z \to Y$, l'espace algébrique $X \times_Y Z$ est représentable
par un schéma affine.
\end{definition}

\begin{lemma}
\label{lemma-affine-local}
Soit $S$ un schéma.
Soit $f : X \to Y$ un morphisme d'espaces algébriques sur $S$.
Les conditions suivantes sont équivalentes :
\begin{enumerate}
\item $f$ est représentable et affine ;
\item $f$ est affine ;
\item pour tout schéma affine $V$ et tout morphisme étale $V \to Y$,
le schéma $X \times_Y V$ est affine ;
\item il existe un schéma $V$ et un morphisme étale surjectif
$V \to Y$ tels que $V \times_Y X \to V$ soit affine ; et
\item il existe un recouvrement de Zariski $Y = \bigcup Y_i$ tel
que chacun des morphismes $f^{-1}(Y_i) \to Y_i$ soit affine.
\end{enumerate}
\end{lemma}

\begin{proof}
Il est clair que (1) implique (2), que (2) implique (3), et que
(3) implique (4) en prenant pour $V$ une somme disjointe de schémas affines
étales sur $Y$ ; voir Propriétés des espaces,
lemme \ref{spaces-properties-lemma-cover-by-union-affines}.
Supposons $V \to Y$ comme dans (4). Pour tout ouvert affine $W$ de $V$,
$W \times_Y X$ est alors un ouvert affine de $V \times_Y X$. Par
Propriétés des espaces, lemme \ref{spaces-properties-lemma-subscheme},
on en déduit que $V \times_Y X$ est un schéma. De plus, le morphisme
$V \times_Y X \to V$ est affine. Nous pouvons donc appliquer
Espaces,
lemme \ref{spaces-lemma-morphism-sheaves-with-P-effective-descent-etale},
car la classe des morphismes affines satisfait à toutes les propriétés
requises (voir
Morphismes, lemmes \ref{morphisms-lemma-base-change-affine} et
Descente, lemmes \ref{descent-lemma-descending-property-affine}
et \ref{descent-lemma-affine}). Ce lemme donne
que $f$ est représentable et affine, c'est-à-dire que (1) est vérifiée.

\medskip\noindent
L'équivalence de (1) et (5) résulte de ce que la propriété
d'être affine est locale sur le but pour la topologie de Zariski (la référence précédente montre
qu'elle est en fait locale sur le but pour la topologie fpqc).
\end{proof}

\begin{lemma}
\label{lemma-composition-affine}
Tout composé de morphismes affines est affine.
\end{lemma}

\begin{proof}
Omis. Indication : transitivité des produits fibrés.
\end{proof}

\begin{lemma}
\label{lemma-base-change-affine}
Tout changement de base d'un morphisme affine est affine.
\end{lemma}

\begin{proof}
Omis. Indication : transitivité des produits fibrés.
\end{proof}

\begin{lemma}
\label{lemma-closed-immersion-affine}
Toute immersion fermée est affine.
\end{lemma}

\begin{proof}
Cela résulte immédiatement de l'énoncé correspondant pour les morphismes de
schémas ; voir Morphismes, lemme \ref{morphisms-lemma-closed-immersion-affine}.
\end{proof}

\begin{lemma}
\label{lemma-affine-equivalence-algebras}
Soit $S$ un schéma. Soit $X$ un espace algébrique sur $S$.
Il existe une anti-équivalence de catégories
$$
\begin{matrix}
\text{espaces algébriques} \\
\text{affines sur }X
\end{matrix}
\longleftrightarrow
\begin{matrix}
\text{faisceaux quasi-cohérents} \\
\text{d'}\mathcal{O}_X\text{-algèbres}
\end{matrix}
$$
qui associe à $f : Y \to X$ le faisceau $f_*\mathcal{O}_Y$.
De plus, cette équivalence est compatible avec tout changement de base.
\end{lemma}

\begin{proof}
Ce lemme est l'analogue de Morphismes, lemme
\ref{morphisms-lemma-affine-equivalence-algebras}.
Soit $\mathcal{A}$ un faisceau quasi-cohérent de $\mathcal{O}_X$-algèbres.
Nous allons construire un morphisme affine d'espaces algébriques
$\pi : Y = \underline{\Spec}_X(\mathcal{A}) \to X$ tel que
$\pi_*\mathcal{O}_Y \cong \mathcal{A}$. Pour cela, choisissons un schéma
$U$ et un morphisme étale surjectif $\varphi : U \to X$. Posons comme d'habitude
$R = U \times_X U$, avec projections $s, t : R \to U$. Notons
$\psi : R \to X$ le composé $\psi = \varphi \circ s = \varphi \circ t$.
D'après le lemme précité, il existe des morphismes affines de schémas
$\pi_0 : V \to U$ et $\pi_1 : W \to R$ tels que
$\pi_{0, *}\mathcal{O}_V \cong \varphi^*\mathcal{A}$ et
$\pi_{1, *}\mathcal{O}_W \cong \psi^*\mathcal{A}$.
Puisque la construction est compatible avec le changement de base, il existe
des morphismes $s', t' : W \to V$ tels que les diagrammes
$$
\vcenter{
\xymatrix{
W \ar[r]_{s'} \ar[d] & V \ar[d] \\
R \ar[r]^s & U
}
}
\quad\text{et}\quad
\vcenter{
\xymatrix{
W \ar[r]_{t'} \ar[d] & V \ar[d] \\
R \ar[r]^t & U
}
}
$$
soient cartésiens. Il s'ensuit que $s', t'$ sont étales. Ce qui précède entraîne
formellement que $(t', s') : W \to V \times_S V$ est un
monomorphisme. Nous omettons de vérifier que $W \to V \times_S V$ est une
relation d'équivalence (indication : considérer l'image inverse de $\mathcal{A}$
sur $U \times_X U \times_X U = R \times_{s, U, t} R$).
Le faisceau quotient $Y = V/W$
est un espace algébrique ; voir
Espaces, théorème \ref{spaces-theorem-presentation}.
D'après Groupoïdes, lemme \ref{groupoids-lemma-criterion-fibre-product},
on a $Y \times_X U \cong V$. Ainsi, $Y \to X$ est affine par le
lemme \ref{lemma-affine-local}. Enfin, l'isomorphisme
$$
(Y \times_X U \to U)_*\mathcal{O}_{Y \times_X U} =
\pi_{0, *}\mathcal{O}_V \cong \varphi^*\mathcal{A}
$$
est compatible avec les isomorphismes de recollement, d'où
$(Y \to X)_*\mathcal{O}_Y \cong \mathcal{A}$ d'après
Propriétés des espaces, Proposition
\ref{spaces-properties-proposition-quasi-coherent}.
Nous omettons de vérifier que cette construction est compatible avec tout
changement de base.
\end{proof}

\begin{definition}
\label{definition-relative-spec}
Soit $S$ un schéma. Soit $X$ un espace algébrique sur $S$.
Soit $\mathcal{A}$ un faisceau quasi-cohérent de
$\mathcal{O}_X$-algèbres. Le {\it spectre relatif de $\mathcal{A}$ sur
$X$}, ou simplement le {\it spectre de $\mathcal{A}$ sur $X$}, est le
morphisme affine $\underline{\Spec}(\mathcal{A}) \to X$
correspondant à $\mathcal{A}$ par l'équivalence de catégories du
lemme \ref{lemma-affine-equivalence-algebras}.
\end{definition}

\noindent
La formation du spectre relatif commute à tout changement de base.

\begin{remark}
\label{remark-factorization-quasi-compact-quasi-separated}
Soit $S$ un schéma. Soit $f : Y \to X$ un morphisme quasi-compact et
quasi-séparé d'espaces algébriques sur $S$. Alors
$f$ admet une factorisation canonique
$$
Y \longrightarrow \underline{\Spec}_X(f_*\mathcal{O}_Y) \longrightarrow X
$$
Cette écriture a un sens car $f_*\mathcal{O}_Y$ est quasi-cohérent
par le lemme \ref{lemma-pushforward}. Le morphisme
$Y \to \underline{\Spec}_X(f_*\mathcal{O}_Y)$ provient du
morphisme canonique de $\mathcal{O}_Y$-algèbres
$f^*f_*\mathcal{O}_Y \to \mathcal{O}_Y$, auquel correspond
un morphisme canonique
$Y \to Y \times_X \underline{\Spec}_X(f_*\mathcal{O}_Y)$ sur $Y$ (voir
le lemme \ref{lemma-affine-equivalence-algebras}), d'où une factorisation
de $f$ comme ci-dessus.
\end{remark}

\begin{lemma}
\label{lemma-affine-equivalence-modules}
Soit $S$ un schéma.
Soit $f : Y \to X$ un morphisme affine d'espaces algébriques sur $S$.
Posons $\mathcal{A} = f_*\mathcal{O}_Y$.
Le foncteur $\mathcal{F} \mapsto f_*\mathcal{F}$ induit
une équivalence de catégories
$$
\left\{
\begin{matrix}
\text{catégorie des}\\
\mathcal{O}_Y\text{-modules quasi-cohérents}
\end{matrix}
\right\}
\longrightarrow
\left\{
\begin{matrix}
\text{catégorie des}\\
\mathcal{A}\text{-modules quasi-cohérents}
\end{matrix}
\right\}
$$
De plus, un $\mathcal{A}$-module est
quasi-cohérent comme $\mathcal{O}_X$-module si et seulement s'il
est quasi-cohérent comme $\mathcal{A}$-module.
\end{lemma}

\begin{proof}
Omis.
\end{proof}

\begin{lemma}
\label{lemma-affine-permanence}
Soit $S$ un schéma. Soit $B$ un espace algébrique sur $S$.
Supposons que $g : X \to Y$ soit un morphisme d'espaces algébriques sur $B$.
\begin{enumerate}
\item Si $X$ est affine sur $B$ et si $\Delta : Y \to Y \times_B Y$ est affine,
alors $g$ est affine.
\item Si $X$ est affine sur $B$ et si $Y$ est séparé sur $B$,
alors $g$ est affine.
\item Tout morphisme d'un schéma affine vers un espace algébrique dont la
diagonale sur $\mathbf{Z}$ est affine (au sens de Propriétés des espaces, définition
\ref{spaces-properties-definition-separated}) est affine.
\item Tout morphisme d'un schéma affine vers un espace algébrique séparé
est affine.
\end{enumerate}
\end{lemma}

\begin{proof}
Démonstration de (1). Le changement de base $X \times_B Y \to Y$ est affine par le
lemme \ref{lemma-base-change-affine}.
Le morphisme $(1, g) : X \to X \times_B Y$ est le changement de base de
$Y \to Y \times_B Y$ par le morphisme $X \times_B Y \to Y \times_B Y$.
Il est donc affine par le
lemme \ref{lemma-base-change-affine}.
Tout composé de morphismes affines est affine
(voir le lemme \ref{lemma-composition-affine}), d'où (1).
L'assertion (2) résulte de (1), car une immersion fermée est affine
(voir le lemme \ref{lemma-closed-immersion-affine}) et dire que $Y/B$ est séparé
signifie que $\Delta$ est une immersion fermée. Les assertions (3) et (4) sont des cas
particuliers de (1) et (2).
\end{proof}

\begin{lemma}
\label{lemma-Artinian-affine}
Soit $S$ un schéma. Soit $X$ un espace algébrique quasi-séparé sur $S$.
Soit $A$ un anneau artinien. Tout morphisme $\Spec(A) \to X$ est affine.
\end{lemma}

\begin{proof}
Soit $U \to X$ un morphisme étale, avec $U$ affine. Pour démontrer le
lemme, il suffit de montrer que $\Spec(A) \times_X U$ est affine ; voir le
lemme \ref{lemma-affine-local}. Comme $X$ est quasi-séparé, le
schéma $\Spec(A) \times_X U$ est quasi-compact. De plus, le
morphisme de projection $\Spec(A) \times_X U \to \Spec(A)$ est étale.
Ce morphisme a donc des fibres discrètes finies ; en outre,
la topologie de $\Spec(A)$ est discrète. Ainsi, $\Spec(A) \times_X U$
est un schéma dont l'espace topologique
sous-jacent est un ensemble discret fini. On conclut par
Schémas, lemme \ref{schemes-lemma-scheme-finite-discrete-affine}.
\end{proof}












\section{Morphismes quasi-affines}
\label{section-quasi-affine}

\noindent
Nous avons déjà défini à la section \ref{section-representable}
ce que signifie, pour un morphisme représentable d'espaces algébriques,
être quasi-affine.

\begin{lemma}
\label{lemma-quasi-affine-representable}
Soit $S$ un schéma. Soit $f : X \to Y$ un morphisme représentable
d'espaces algébriques sur $S$. Alors
$f$ est quasi-affine (au sens de la section \ref{section-representable})
si et seulement si, pour tout schéma affine $Z$
et tout morphisme $Z \to Y$, le schéma $X \times_Y Z$ est quasi-affine.
\end{lemma}

\begin{proof}
Cela résulte directement de la définition d'un morphisme quasi-affine
de schémas
(Morphismes, définition \ref{morphisms-definition-quasi-affine}).
\end{proof}

\noindent
Nous pouvons ainsi poser la définition suivante.

\begin{definition}
\label{definition-quasi-affine}
Soit $S$ un schéma.
Soit $f : X \to Y$ un morphisme d'espaces algébriques sur $S$.
Nous disons que $f$ est {\it quasi-affine} si, pour tout schéma affine $Z$ et tout
morphisme $Z \to Y$, l'espace algébrique $X \times_Y Z$ est représentable
par un schéma quasi-affine.
\end{definition}

\begin{lemma}
\label{lemma-quasi-affine-local}
Soit $S$ un schéma.
Soit $f : X \to Y$ un morphisme d'espaces algébriques sur $S$.
Les conditions suivantes sont équivalentes :
\begin{enumerate}
\item $f$ est représentable et quasi-affine ;
\item $f$ est quasi-affine ;
\item il existe un schéma $V$ et un morphisme étale surjectif
$V \to Y$ tels que $V \times_Y X \to V$ soit quasi-affine ; et
\item il existe un recouvrement de Zariski $Y = \bigcup Y_i$ tel
que chacun des morphismes $f^{-1}(Y_i) \to Y_i$ soit quasi-affine.
\end{enumerate}
\end{lemma}

\begin{proof}
Il est clair que (1) implique (2) et que (2) implique (3) en prenant
pour $V$ une somme disjointe de schémas affines étales sur $Y$ ; voir
Propriétés des espaces,
lemme \ref{spaces-properties-lemma-cover-by-union-affines}.
Supposons $V \to Y$ comme dans (3). Pour tout ouvert affine $W$ de $V$,
$W \times_Y X$ est alors un ouvert quasi-affine de $V \times_Y X$. Par
Propriétés des espaces, lemme \ref{spaces-properties-lemma-subscheme},
on en déduit que $V \times_Y X$ est un schéma. De plus, le morphisme
$V \times_Y X \to V$ est quasi-affine. Nous pouvons donc appliquer
Espaces,
lemme \ref{spaces-lemma-morphism-sheaves-with-P-effective-descent-etale},
car la classe des morphismes quasi-affines satisfait à toutes les
propriétés requises (voir
Morphismes, lemmes \ref{morphisms-lemma-base-change-quasi-affine} et
Descente, lemmes \ref{descent-lemma-descending-property-quasi-affine}
et \ref{descent-lemma-quasi-affine}). Ce lemme donne
que $f$ est représentable et quasi-affine, c'est-à-dire que (1) est vérifiée.

\medskip\noindent
L'équivalence de (1) et (4) résulte de ce que la propriété
d'être quasi-affine est locale sur le but pour la topologie de Zariski (la référence précédente
montre qu'elle est en fait locale sur le but pour la topologie fpqc).
\end{proof}

\begin{lemma}
\label{lemma-composition-quasi-affine}
Tout composé de morphismes quasi-affines est quasi-affine.
\end{lemma}

\begin{proof}
Omis.
\end{proof}

\begin{lemma}
\label{lemma-base-change-quasi-affine}
Tout changement de base d'un morphisme quasi-affine est quasi-affine.
\end{lemma}

\begin{proof}
Omis.
\end{proof}

\begin{lemma}
\label{lemma-characterize-quasi-affine}
Soit $S$ un schéma.
Un morphisme quasi-compact et quasi-séparé d'espaces algébriques
$f : Y \to X$ est quasi-affine si et seulement si la factorisation canonique
$Y \to \underline{\Spec}_X(f_*\mathcal{O}_Y)$
(remarque \ref{remark-factorization-quasi-compact-quasi-separated})
est une immersion ouverte.
\end{lemma}

\begin{proof}
Soit $U \to X$ un morphisme étale surjectif, où $U$ est un schéma.
Comme on peut vérifier que $f$ est quasi-affine après le changement de base à
$U$ (lemme \ref{lemma-quasi-affine-local}), que $f_*\mathcal{O}_Y|_U$
est égal à $(Y \times_X U \to U)_*\mathcal{O}_{Y \times_X U}$
(Propriétés des espaces, lemme
\ref{spaces-properties-lemma-pushforward-etale-base-change-modules}), et
que la formation du spectre relatif commute au changement de base
(lemme \ref{lemma-affine-equivalence-algebras}),
on voit que l'assertion se ramène au cas où $X$ est un schéma.
Si $X$ est un schéma et si $f$ est quasi-affine ou si
$Y \to \underline{\Spec}_X(f_*\mathcal{O}_Y)$ est une immersion ouverte,
alors $Y$ est également un schéma. On se ramène donc à
Morphismes, lemme \ref{morphisms-lemma-characterize-quasi-affine}.
\end{proof}







\section{Propriétés de morphismes locales sur la source et le but pour la topologie étale}
\label{section-local-source-target}

\noindent
Étant donnée une propriété de morphismes de schémas qui est {\it locale sur la
source et le but pour la topologie étale}, voir
Descente, définition \ref{descent-definition-local-source-target}
nous pouvons l'utiliser pour définir une propriété correspondante
de morphismes d'espaces algébriques, en imposant l'une ou l'autre des
conditions équivalentes du lemme ci-dessous.

\begin{lemma}
\label{lemma-local-source-target}
Soit $\mathcal{P}$ une propriété de morphismes de schémas
qui est locale sur la source et le but pour la topologie étale.
Soit $S$ un schéma.
Soit $f : X \to Y$ un morphisme d'espaces algébriques sur $S$.
Considérons des diagrammes commutatifs
$$
\xymatrix{
U \ar[d]_a \ar[r]_h & V \ar[d]^b \\
X \ar[r]^f & Y
}
$$
où $U$ et $V$ sont des schémas et où les flèches verticales sont étales.
Les conditions suivantes sont équivalentes :
\begin{enumerate}
\item pour tout diagramme comme ci-dessus, le morphisme $h$ possède la propriété
$\mathcal{P}$ ;
\item pour un certain diagramme comme ci-dessus dans lequel $a : U \to X$ est
surjectif, le morphisme $h$ possède la propriété $\mathcal{P}$.
\end{enumerate}
Si $X$ et $Y$ sont représentables, cela équivaut encore à ce que
$f$ (considéré comme morphisme de schémas) possède la propriété $\mathcal{P}$.
Si, de plus, $\mathcal{P}$ est préservée par tout changement de base et
locale sur la base pour la topologie fppf, cela équivaut encore, pour les
morphismes représentables $f$, à ce que $f$ possède la propriété $\mathcal{P}$
au sens de la section \ref{section-representable}.
\end{lemma}

\begin{proof}
Montrons l'équivalence de (1) et (2).
L'implication (1) $\Rightarrow$ (2) est immédiate (compte tenu de
Espaces, lemme \ref{spaces-lemma-lift-morphism-presentations}).
Supposons que
$$
\xymatrix{
U \ar[d] \ar[r]_h & V \ar[d] \\
X \ar[r]^f & Y
}
\quad\quad
\xymatrix{
U' \ar[d] \ar[r]_{h'} & V' \ar[d] \\
X \ar[r]^f & Y
}
$$
soient deux diagrammes comme dans le lemme. Supposons que $U \to X$ soit
surjectif et que $h$ possède la propriété $\mathcal{P}$. Pour montrer que (2)
implique (1), il faut prouver que $h'$ possède $\mathcal{P}$. Considérons
pour cela le diagramme
$$
\xymatrix{
U \ar[d]_h &
U \times_X U' \ar[l] \ar[d]^{(h, h')} \ar[r] &
U' \ar[d]^{h'} \\
V &
V \times_Y V' \ar[l] \ar[r] &
V'
}
$$
D'après
Descente, lemme \ref{descent-lemma-local-source-target-characterize}
le fait que $h$ possède $\mathcal{P}$ entraîne que $(h, h')$ possède $\mathcal{P}$ ;
comme $U \times_X U' \to U'$ est surjectif, cela entraîne (par le même
lemme) que $h'$ possède $\mathcal{P}$.

\medskip\noindent
Si $X$ et $Y$ sont représentables, alors
Descente, lemme \ref{descent-lemma-local-source-target-characterize}
s'applique et montre que (1) et (2) équivalent au fait que $f$ possède
$\mathcal{P}$.

\medskip\noindent
Enfin, supposons que $f$ soit représentable, que $U, V, a, b, h$ soient
comme dans la partie (2) du lemme et que $\mathcal{P}$ soit préservée par
tout changement de base. Il faut montrer que, pour tout schéma
$Z$ et tout morphisme $Z \to Y$, le changement de base $Z \times_Y X \to Z$
possède la propriété $\mathcal{P}$. Considérons le diagramme
$$
\xymatrix{
Z \times_Y U \ar[d] \ar[r] &
Z \times_Y V \ar[d] \\
Z \times_Y X \ar[r] &
Z
}
$$
Notons que la flèche horizontale supérieure est un changement de base de $h$
et possède donc la propriété $\mathcal{P}$. La flèche verticale gauche est étale
et surjective, tandis que la flèche verticale droite est étale. Ainsi,
Descente, lemme \ref{descent-lemma-local-source-target-characterize}
intervient une fois encore et montre que $Z \times_Y X \to Z$
possède la propriété $\mathcal{P}$.
\end{proof}

\begin{definition}
\label{definition-P}
Soit $S$ un schéma.
Soit $\mathcal{P}$ une propriété de morphismes de schémas
qui est locale sur la source et le but pour la topologie étale.
On dit qu'un morphisme $f : X \to Y$ d'espaces algébriques sur $S$
{\it possède la propriété $\mathcal{P}$} si les conditions équivalentes du
lemme \ref{lemma-local-source-target}
sont satisfaites.
\end{definition}

\noindent
Voici deux remarques immédiates.

\begin{remark}
\label{remark-composition-P}
Soit $S$ un schéma. Soit $\mathcal{P}$ une propriété de morphismes de schémas
qui est locale sur la source et le but pour la topologie étale. Supposons de plus
que $\mathcal{P}$ soit stable par composition. Alors la classe des morphismes
d'espaces algébriques possédant la propriété $\mathcal{P}$ est stable par composition.
\end{remark}

\begin{remark}
\label{remark-base-change-P}
Soit $S$ un schéma. Soit $\mathcal{P}$ une propriété de morphismes de schémas
qui est locale sur la source et le but pour la topologie étale. Supposons de plus
que $\mathcal{P}$ soit stable par changement de base. Alors la classe des morphismes
d'espaces algébriques possédant la propriété $\mathcal{P}$ est stable par changement de base.
\end{remark}

\noindent
Étant donnée une propriété de morphismes de germes de schémas qui est {\it locale
sur la source et le but pour la topologie étale}, voir
Descente, définition \ref{descent-definition-local-source-target-at-point}
nous pouvons l'utiliser pour définir une propriété correspondante de morphismes
d'espaces algébriques en un point, en imposant l'une ou l'autre des
conditions équivalentes du lemme ci-dessous.

\begin{lemma}
\label{lemma-local-source-target-at-point}
Soit $\mathcal{Q}$ une propriété de morphismes de germes qui est
locale sur la source et le but pour la topologie étale.
Soit $S$ un schéma.
Soit $f : X \to Y$ un morphisme d'espaces algébriques sur $S$.
Soit $x \in |X|$ un point de $X$.
Considérons les diagrammes
$$
\xymatrix{
U \ar[d]_a \ar[r]_h & V \ar[d]^b \\
X \ar[r]^f & Y
}
\quad\quad
\xymatrix{
u \ar[d] \ar[r] & v \ar[d] \\
x \ar[r] & y
}
$$
où $U$ et $V$ sont des schémas, où $a, b$ sont étales et où $u, v, x, y$
sont des points des espaces correspondants.
Les conditions suivantes sont équivalentes :
\begin{enumerate}
\item pour tout diagramme comme ci-dessus, on a $\mathcal{Q}((U, u) \to (V, v))$ ;
\item pour un certain diagramme comme ci-dessus, on a $\mathcal{Q}((U, u) \to (V, v))$.
\end{enumerate}
Si $X$ et $Y$ sont représentables, cela équivaut encore
à $\mathcal{Q}((X, x) \to (Y, y))$.
\end{lemma}

\begin{proof}
Omis. Indication : la démonstration est très semblable à celle du
lemme \ref{lemma-local-source-target}.
\end{proof}

\begin{definition}
\label{definition-P-at-point}
Soit $\mathcal{Q}$ une propriété de morphismes de germes
de schémas qui est locale sur la source et le but pour la topologie étale.
Soit $S$ un schéma.
Étant donnés un morphisme $f : X \to Y$ d'espaces algébriques sur $S$ et
un point $x \in |X|$, on dit que $f$
{\it possède la propriété $\mathcal{Q}$ en $x$} si les conditions équivalentes du
lemme \ref{lemma-local-source-target-at-point}
sont satisfaites.
\end{definition}

\noindent
Il ne faut pas utiliser aveuglément le lemme suivant pour passer d'une propriété
de morphismes à une propriété de morphismes en un point. Par exemple, si
$\mathcal{P}$ est la propriété d'être plat, alors la propriété
$\mathcal{Q}$ du lemme suivant signifie « $f$ est plat dans un voisinage ouvert
de $x$ », ce qui n'est pas équivalent à « $f$ est plat en $x$ ».

\begin{lemma}
\label{lemma-local-source-target-global-implies-local}
Soit $\mathcal{P}$ une propriété de morphismes de schémas
qui est locale sur la source et le but pour la topologie étale.
Considérons la propriété $\mathcal{Q}$ de morphismes
de germes associée à $\mathcal{P}$ dans
Descente, lemme \ref{descent-lemma-local-source-target-global-implies-local}.
Alors :
\begin{enumerate}
\item $\mathcal{Q}$ est locale sur la source et le but pour la topologie étale ;
\item étant donnés un morphisme d'espaces algébriques $f : X \to Y$ et $x \in |X|$,
les conditions suivantes sont équivalentes :
\begin{enumerate}
\item $f$ possède $\mathcal{Q}$ en $x$ ;
\item il existe un voisinage ouvert $X' \subset X$ de $x$
tel que $X' \to Y$ possède $\mathcal{P}$ ;
\end{enumerate}
\item étant donné un morphisme d'espaces algébriques $f : X \to Y$,
les conditions suivantes sont équivalentes :
\begin{enumerate}
\item $f$ possède $\mathcal{P}$ ;
\item pour tout $x \in |X|$, le morphisme $f$ possède $\mathcal{Q}$ en $x$.
\end{enumerate}
\end{enumerate}
\end{lemma}

\begin{proof}
Voir
Descente, lemme \ref{descent-lemma-local-source-target-global-implies-local}
pour (1). L'implication (2)(a) $\Rightarrow$ (2)(b) s'obtient en posant
$X' = a(U) \subset X$ dans un diagramme comme dans le
lemme \ref{lemma-local-source-target-at-point}.
L'implication (2)(b) $\Rightarrow$ (2)(a) est claire.
L'équivalence de (3)(a) et (3)(b) résulte de l'énoncé correspondant
pour les morphismes de schémas ; voir
Descente, lemme \ref{descent-lemma-local-source-target-local-implies-global}.
\end{proof}

\begin{remark}
\label{remark-when-apply}
Nous appliquerons le
lemme \ref{lemma-local-source-target-global-implies-local}
ci-dessus à tous les cas énumérés dans
Descente, remarque \ref{descent-remark-list-local-source-target}
sauf « plat ». Dans chaque cas, nous le ferons en disant que
$f$ possède la propriété $\mathcal{P}$ en $x$ si $f$ possède
$\mathcal{P}$ dans un voisinage de $x$.
\end{remark}




\section{Morphismes de type fini}
\label{section-finite-type}

\noindent
La propriété « localement de type fini » des morphismes de schémas
est locale sur la source et le but pour la topologie étale ; voir
Descente, remarque \ref{descent-remark-list-local-source-target}.
Elle est aussi stable par changement de base et locale sur le but pour
la topologie fpqc ; voir Morphismes, lemme \ref{morphisms-lemma-base-change-finite-type}, et
Descente, lemmes \ref{descent-lemma-descending-property-locally-finite-type}.
Ainsi, d'après le
lemme \ref{lemma-local-source-target}
ci-dessus, nous pouvons définir comme suit ce que signifie, pour un morphisme
d'espaces algébriques, être localement de type fini ;
cette définition coïncide avec la notion déjà définie à la
section \ref{section-representable}
lorsque le morphisme est représentable.

\begin{definition}
\label{definition-locally-finite-type}
Soit $S$ un schéma.
Soit $f : X \to Y$ un morphisme d'espaces algébriques sur $S$.
\begin{enumerate}
\item On dit que $f$ est
{\it localement de type fini} si les conditions équivalentes du
lemme \ref{lemma-local-source-target}
sont satisfaites pour
$\mathcal{P} = \text{localement de type fini}$.
\item Soit $x \in |X|$. On dit que $f$ est {\it de type fini en $x$}
s'il existe un voisinage ouvert $X' \subset X$ de $x$ tel
que $f|_{X'} : X' \to Y$ soit localement de type fini.
\item On dit que $f$ est
{\it de type fini} s'il est localement de type fini et quasi-compact.
\end{enumerate}
\end{definition}

\noindent
Considérons l'espace algébrique $\mathbf{A}^1_k/\mathbf{Z}$ de
Espaces, exemple \ref{spaces-example-affine-line-translation}.
Le morphisme $\mathbf{A}^1_k/\mathbf{Z} \to \Spec(k)$
est de type fini.

\begin{lemma}
\label{lemma-composition-finite-type}
Le composé de morphismes de type fini est de type fini.
Il en est de même des morphismes localement de type fini.
\end{lemma}

\begin{proof}
Voir la remarque \ref{remark-composition-P} et
Morphismes, lemme \ref{morphisms-lemma-composition-finite-type}.
\end{proof}

\begin{lemma}
\label{lemma-base-change-finite-type}
Tout changement de base d'un morphisme de type fini est de type fini.
Il en est de même des morphismes localement de type fini.
\end{lemma}

\begin{proof}
Voir la remarque \ref{remark-base-change-P} et
Morphismes, lemme \ref{morphisms-lemma-base-change-finite-type}.
\end{proof}

\begin{lemma}
\label{lemma-finite-type-local}
Soit $S$ un schéma.
Soit $f : X \to Y$ un morphisme d'espaces algébriques sur $S$.
Les conditions suivantes sont équivalentes :
\begin{enumerate}
\item $f$ est localement de type fini ;
\item pour tout $x \in |X|$, le morphisme $f$ est de type fini en $x$ ;
\item pour tout schéma $Z$ et tout morphisme $Z \to Y$, le morphisme
$Z \times_Y X \to Z$ est localement de type fini ;
\item pour tout schéma affine $Z$ et tout morphisme
$Z \to Y$, le morphisme $Z \times_Y X \to Z$ est localement de type fini ;
\item il existe un schéma $V$ et un morphisme étale surjectif
$V \to Y$ tels que $V \times_Y X \to V$ soit localement de type fini ;
\item il existe un schéma $U$ et un morphisme étale surjectif
$\varphi : U \to X$ tels que le composé $f \circ \varphi$
soit localement de type fini ;
\item pour tout diagramme commutatif
$$
\xymatrix{
U \ar[d] \ar[r] & V \ar[d] \\
X \ar[r] & Y
}
$$
où $U$, $V$ sont des schémas et où les flèches verticales sont étales,
la flèche horizontale supérieure est localement de type fini ;
\item il existe un diagramme commutatif
$$
\xymatrix{
U \ar[d] \ar[r] & V \ar[d] \\
X \ar[r] & Y
}
$$
où $U$, $V$ sont des schémas, où les flèches verticales sont étales, où $U \to X$
est surjectif et où la flèche horizontale supérieure est localement de type fini ;
\item il existe des recouvrements de Zariski $Y = \bigcup_{i \in I} Y_i$
et $f^{-1}(Y_i) = \bigcup X_{ij}$ tels que
chaque morphisme $X_{ij} \to Y_i$ soit localement de type fini.
\end{enumerate}
\end{lemma}

\begin{proof}
Chacune des conditions (2), (3), (4), (5), (6), (7) et (9)
implique directement la condition (8). Par exemple, si
(5) est satisfaite, nous pouvons choisir un schéma $V$ qui soit une somme
disjointe de schémas affines et un morphisme étale surjectif $V \to Y$
(voir Propriétés des espaces, lemme
\ref{spaces-properties-lemma-cover-by-union-affines}).
Alors $V \times_Y X \to V$ est localement de type fini d'après (5).
Choisissons un schéma $U$ et un morphisme étale surjectif
$U \to V \times_Y X$. Alors $U \to V$ est localement de type
fini d'après le lemme \ref{lemma-composition-finite-type}.
La condition (8) est donc satisfaite.

\medskip\noindent
Les conditions (1), (7) et (8) sont équivalentes par définition.

\medskip\noindent
Pour terminer la démonstration, montrons que (1) implique toutes les conditions
(2), (3), (4), (5), (6) et (9). Pour (2), c'est immédiat.
Pour (3), (4), (5) et (9), cela résulte du fait que la propriété d'être
localement de type fini est préservée par changement de base ; voir le
lemme \ref{lemma-base-change-finite-type}.
Pour (6), prenons une présentation étale surjective $U \to X$ par un schéma ; le morphisme composé est alors localement de type fini.
\end{proof}

\begin{lemma}
\label{lemma-locally-finite-type-locally-noetherian}
Soit $S$ un schéma.
Soit $f : X \to Y$ un morphisme d'espaces algébriques sur $S$.
Si $f$ est localement de type fini et si $Y$ est localement noethérien,
alors $X$ est localement noethérien.
\end{lemma}

\begin{proof}
Soit
$$
\xymatrix{
U \ar[d] \ar[r] & V \ar[d] \\
X \ar[r] & Y
}
$$
un diagramme commutatif où $U$, $V$ sont des schémas et où les flèches verticales
sont étales surjectives. Si $f$ est localement de type fini, alors
$U \to V$ est localement de type fini. Si $Y$ est localement noethérien, alors
$V$ est localement noethérien. D'après
Morphismes, lemme \ref{morphisms-lemma-finite-type-noetherian}
on voit que $U$ est localement noethérien, ce qui signifie que $X$ est localement
noethérien.
\end{proof}

\begin{lemma}
\label{lemma-permanence-finite-type}
Soit $S$ un schéma.
Soient $f : X \to Y$, $g : Y \to Z$ des morphismes d'espaces algébriques sur $S$.
Si $g \circ f : X \to Z$ est localement de type fini, alors $f : X \to Y$
est localement de type fini.
\end{lemma}

\begin{proof}
On peut trouver un diagramme
$$
\xymatrix{
U \ar[r] \ar[d] & V \ar[r] \ar[d] & W \ar[d] \\
X \ar[r] & Y \ar[r] & Z
}
$$
où $U$, $V$, $W$ sont des schémas et où les flèches verticales sont étales et surjectives ;
voir
Espaces, lemme \ref{spaces-lemma-lift-morphism-presentations}.
On peut alors utiliser le
lemme \ref{lemma-finite-type-local}
et
Morphismes, lemme \ref{morphisms-lemma-permanence-finite-type}
pour conclure.
\end{proof}

\begin{lemma}
\label{lemma-immersion-locally-finite-type}
Une immersion est localement de type fini.
\end{lemma}

\begin{proof}
Cela résulte du principe général
Espaces, lemme
\ref{spaces-lemma-representable-transformations-property-implication}
et
Morphismes, lemme \ref{morphisms-lemma-immersion-locally-finite-type}.
\end{proof}







\section{Points et points géométriques}
\label{section-points-fields}

\noindent
Dans cette section, nous faisons quelques remarques sur les points et les points géométriques (voir
Propriétés des espaces,
définition \ref{spaces-properties-definition-geometric-point}).
Une manière de concevoir un point géométrique de $X$ consiste à
considérer un point géométrique $\overline{s} : \Spec(k) \to S$ de $S$
et un relèvement de $\overline{s}$ en un morphisme $\overline{x}$ vers $X$.
On a le diagramme
$$
\xymatrix{
\Spec(k) \ar[r]_-{\overline{x}} \ar[rd]_{\overline{s}} & X \ar[d] \\
& S.
}
$$
Nous disons souvent « soit $k$ un corps algébriquement clos sur $S$ »
pour indiquer que $\Spec(k)$ est muni d'un morphisme
$\Spec(k) \to S$. Dans cette situation, nous écrivons
$$
X(k) = \Mor_S(\Spec(k), X) =
\{\overline{x} \in X\text{ situé au-dessus de }\overline{s}\}
$$
pour l'ensemble des points à valeurs dans $k$ de $X$. Dans ce cas, l'application
$X(k) \to |X|$ est à valeurs dans la partie $|X_s| \subset |X|$.
Ici, $X_s = \Spec(\kappa(s)) \times_S X$, où $s \in S$
est le point correspondant à $\overline{s}$. Comme
$\Spec(\kappa(s)) \to S$ est un monomorphisme, son changement de base
$X_s \to X$ est aussi un monomorphisme, et $|X_s|$ est bien une partie de $|X|$.

\begin{lemma}
\label{lemma-locally-finite-type-surjective-geometric-points}
Soit $S$ un schéma. Soit $f : X \to Y$ un morphisme d'espaces algébriques
sur $S$. Supposons que $f$ soit localement de type fini. Les conditions suivantes sont équivalentes :
\begin{enumerate}
\item $f$ est surjectif ;
\item pour tout corps algébriquement clos $k$ sur $S$, l'application induite
$X(k) \to Y(k)$ est surjective.
\end{enumerate}
\end{lemma}

\begin{proof}
Choisissons un diagramme
$$
\xymatrix{
U \ar[d] \ar[r] & V \ar[d] \\
X \ar[r] & Y
}
$$
où $U$, $V$ sont des schémas sur $S$ et où les flèches verticales sont surjectives
et étales ; voir
Espaces, lemme \ref{spaces-lemma-lift-morphism-presentations}.
Comme $f$ est localement de type fini, on voit que $U \to V$ est localement de
type fini.

\medskip\noindent
Supposons (1) et soit $\overline{y} \in Y(k)$. Alors $U \to Y$ est
surjectif et localement de type fini d'après les
lemmes \ref{lemma-composition-surjective} et
\ref{lemma-composition-finite-type}.
Posons $Z = U \times_{Y, \overline{y}} \Spec(k)$. C'est un schéma.
La projection $Z \to \Spec(k)$
est surjective et localement de type fini d'après les
lemmes \ref{lemma-base-change-surjective} et
\ref{lemma-base-change-finite-type}.
Il résulte de
Variétés, lemme \ref{varieties-lemma-locally-finite-type-Jacobson}
que $Z$ possède un point à valeurs dans $k$, noté $\overline{z}$. L'image
$\overline{x} \in X(k)$ de $\overline{z}$ s'envoie sur $\overline{y}$,
comme voulu.

\medskip\noindent
Supposons (2). D'après
Propriétés des espaces,
lemme \ref{spaces-properties-lemma-characterize-surjective}
il suffit de montrer que $|X| \to |Y|$ est surjective.
Soit $y \in |Y|$. Choisissons $u \in U$ s'envoyant sur $y$.
Soit $k \supset \kappa(u)$ une clôture algébrique. Notons
$\overline{u} \in U(k)$ le point correspondant et
$\overline{y} \in Y(k)$ son image. Par hypothèse, il existe
$\overline{x} \in X(k)$ s'envoyant sur $\overline{y}$.
Il est alors clair que l'image $x \in |X|$ de $\overline{x}$
s'envoie sur $y$.
\end{proof}

\noindent
Pour énoncer le lemme suivant, introduisons la
notation suivante. Étant donné un schéma $T$, posons
$$
\lambda(T) = \sup\{\aleph_0, |\kappa(t)| ; t \in T\}.
$$
Autrement dit, $\lambda(T)$ est le plus petit cardinal infini qui majore
toutes les cardinalités des corps résiduels de $T$. Notons que, si $R$
est un anneau, la cardinalité de tout corps résiduel $\kappa(\mathfrak p)$
de $R$ est majorée par celle de $R$ (détails omis).
Il s'ensuit que $\lambda(T) \leq \text{size}(T)$, où
$\text{size}(T)$ est la taille du schéma $T$ introduite dans
Ensembles, section \ref{sets-section-categories-schemes}.
Si $L/K$ est une extension de corps de type fini, alors
$|K| \leq |L| \leq \max\{\aleph_0, |K|\}$. Il s'ensuit que, si $T' \to T$
est un morphisme de schémas localement de type
fini, alors $\lambda(T') \leq \lambda(T)$ ; si $T' \to T$ est en outre
surjectif, il y a égalité. Supposons ensuite que $S$ soit un schéma
et que $X$ soit un espace algébrique sur $S$. Dans ce cas, posons
$$
\lambda(X) := \lambda(U)
$$
où $U$ est un schéma quelconque sur $S$ muni d'un morphisme étale surjectif
vers $X$. Cette définition est indépendante du choix de $U$ :
en effet, étant donnés deux tels schémas $U$ et $U'$, le produit fibré
$U \times_X U'$ est un schéma qui admet un morphisme étale surjectif
vers $U$ comme vers $U'$, d'où $\lambda(U) = \lambda(U \times_X U') =
\lambda(U')$ d'après ce qui précède.

\begin{lemma}
\label{lemma-large-enough}
Soit $S$ un schéma. Soient $X$, $Y$ des espaces algébriques sur $S$.
\begin{enumerate}
\item Lorsque $k$ parcourt les corps algébriquement clos sur $S$,
les points géométriques $\overline{y} \in Y(k)$ recouvrent tout $|Y|$.
\item Lorsque $k$ parcourt les corps algébriquement clos sur $S$ tels que
$|k| \geq \lambda(Y)$ et $|k| > \lambda(X)$, les points géométriques
$\overline{y} \in Y(k)$ recouvrent tout $|Y|$.
\item Pour tout point géométrique
$\overline{s} : \Spec(k) \to S$ tel que
$k$ soit de cardinalité $> \lambda(X)$, l'application
$$
X(k) \longrightarrow |X_s|
$$
est surjective.
\item Soit $X \to Y$ un morphisme d'espaces algébriques sur $S$.
Pour tout point géométrique $\overline{s} : \Spec(k) \to S$ tel que
$k$ soit de cardinalité $> \lambda(X)$, l'application
$$
X(k) \longrightarrow |X| \times_{|Y|} Y(k)
$$
est surjective.
\item Soit $X \to Y$ un morphisme d'espaces algébriques sur $S$.
Les conditions suivantes sont équivalentes :
\begin{enumerate}
\item l'application $X \to Y$ est surjective ;
\item pour tout corps algébriquement clos $k$ sur $S$ tel que
$|k| > \lambda(X)$ et $|k| \geq \lambda(Y)$, l'application $X(k) \to Y(k)$
est surjective.
\end{enumerate}
\end{enumerate}
\end{lemma}

\begin{proof}
Pour démontrer (1), choisissons un morphisme étale surjectif $V \to Y$, où
$V$ est un schéma. Pour chaque $v \in V$, choisissons une clôture algébrique
$\kappa(v) \subset k_v$. Considérons les morphismes
$\overline{y}_v : \Spec(k_v) \to V \to Y$.
Par construction de $|Y|$, ceux-ci recouvrent $|Y|$.

\medskip\noindent
Pour démontrer (2), nous utiliserons les deux faits suivants, dont nous omettons les démonstrations :
(\romannumeral1) si $K$ est un corps et $\overline{K}$ une clôture algébrique, alors
$|\overline{K}| \leq \max\{\aleph_0, |K|\}$ ;
(\romannumeral2) pour tout corps algébriquement clos $k$ et tout cardinal
$\aleph$ tel que $\aleph \geq |k|$, il existe une extension de corps
algébriquement clos $k'/k$ telle que $|k'| = \aleph$.
Posons maintenant $\aleph = \max\{\lambda(X), \lambda(Y)\}^+$.
Ici, $\lambda^+ > \lambda$ désigne le cardinal immédiatement supérieur ; voir
Ensembles, section \ref{sets-section-cardinals}.
D'après (\romannumeral1),
les corps $k_v$ construits au premier paragraphe
de la démonstration ont tous une cardinalité majorée par $\lambda(Y)$. Par suite,
d'après (\romannumeral2), il existe des extensions $k_v \subset k'_v$ telles que
$|k'_v| = \aleph$. Les morphismes
$\overline{y}'_v : \Spec(k'_v) \to Y$ recouvrent $|Y|$, comme voulu.
Pour achever réellement la démonstration de (2), il faut vérifier que les schémas
$\Spec(k'_v)$ sont (isomorphes à) des objets de $\Sch_{fppf}$,
car nos conventions imposent que tous les schémas soient des objets de
$\Sch_{fppf}$ ; le reste de ce paragraphe peut être omis par
qui ne s'intéresse pas aux considérations ensemblistes. Par
construction, il existe un objet $T$ de $\Sch_{fppf}$ tel que
$\lambda(X)$ et $\lambda(Y)$ soient majorés par $\text{size}(T)$.
D'après notre construction de la catégorie $\Sch_{fppf}$ dans
Topologies, définitions \ref{topologies-definition-big-fppf-site}
comme la catégorie $\Sch_\alpha$ construite dans
Ensembles, lemme \ref{sets-lemma-construct-category}
on voit que tout schéma dont la taille est $\leq \text{size}(T)^+$
est isomorphe à un objet de $\Sch_{fppf}$. Voir
l'expression de la fonction $Bound$ dans
Ensembles, équation (\ref{sets-equation-bound}).
Comme $\aleph \leq \text{size}(T)^+$, on conclut.

\medskip\noindent
Dans la partie (3), la notation $X_s$ désigne
le produit fibré $\Spec(\kappa(s)) \times_S X$, où $s \in S$
est le point correspondant à $\overline{s}$. La partie (3) résulte donc
de (4) en prenant $Y = \Spec(\kappa(s))$.

\medskip\noindent
Démontrons (4). Soit $X \to Y$ un morphisme d'espaces algébriques sur $S$.
Soit $k$ un corps algébriquement clos sur $S$, de cardinalité
$> \lambda(X)$. Soient $\overline{y} \in Y(k)$ et $x \in |X|$ ayant pour image le
même élément $y$ de $|Y|$.
Il faut trouver $\overline{x} \in X(k)$ s'envoyant sur $x$ et $\overline{y}$.
Choisissons un diagramme commutatif
$$
\xymatrix{
U \ar[d] \ar[r] & V \ar[d] \\
X \ar[r] & Y
}
$$
où $U$, $V$ sont des schémas sur $S$ et où les flèches verticales sont surjectives
et étales ; voir
Espaces, lemme \ref{spaces-lemma-lift-morphism-presentations}.
Choisissons $u \in |U|$ s'envoyant sur $x$, et notons $v \in |V|$ son image.
Nous considérerons $u = \Spec(\kappa(u))$ et
$v = \Spec(\kappa(v))$ comme des schémas.
Notons que $V \times_Y \Spec(k)$ est un schéma étale sur $k$.
C'est donc une somme disjointe de spectres d'extensions finies
séparables de $k$ ; voir
Morphismes, lemme \ref{morphisms-lemma-etale-over-field}.
Comme $v$ s'envoie sur $y$, on voit que $v \times_Y \Spec(k)$ est un
schéma non vide. Comme $v \to V$ est un monomorphisme, on voit que
$v \times_Y \Spec(k) \to V \times_Y \Spec(k)$
est un monomorphisme. Ainsi, $v \times_Y \Spec(k)$ est une somme
disjointe de spectres d'extensions finies séparables de $k$, d'après
Schémas, lemme \ref{schemes-lemma-mono-towards-spec-field}.
On en déduit que le morphisme
$v \times_Y \Spec(k) \to \Spec(k)$
possède une section, c'est-à-dire qu'il existe un morphisme
$\overline{v} : \Spec(k) \to V$ situé au-dessus de $v$ et de
$\overline{y}$. Considérons enfin le schéma
$$
u \times_{V, \overline{v}} \Spec(k)
=
\Spec(\kappa(u) \otimes_{\kappa(v)} k)
$$
où $\kappa(v) \to k$ est le morphisme de corps qui définit le morphisme
$\overline{v}$. Comme, par hypothèse, la cardinalité de $k$ est strictement supérieure à celle
de $\kappa(u)$, on peut appliquer
Algèbre, lemme \ref{algebra-lemma-base-change-Jacobson}
pour voir que tout idéal maximal
$\mathfrak m \subset \kappa(u) \otimes_{\kappa(v)} k$
a un corps résiduel algébrique sur $k$, donc égal à $k$.
Un tel idéal maximal fournit donc un morphisme
$\overline{u} : \Spec(k) \to U$ situé au-dessus de $u$ et s'envoyant sur
$\overline{v}$. Le composé $\Spec(k) \to U \to X$
est le point géométrique recherché $\overline{x} \in X(k)$.
Ceci achève la démonstration de (4).

\medskip\noindent
La partie (5) est une conséquence formelle des parties (2) et (4) et de
Propriétés des espaces,
lemme \ref{spaces-properties-lemma-characterize-surjective}.
\end{proof}








\section{Points de type fini}
\label{section-points-finite-type}

\noindent
Soit $S$ un schéma. Soit $X$ un espace algébrique sur $S$.
Un point de type fini $x \in |X|$ est un point qui peut être représenté par
un morphisme $\Spec(k) \to X$ localement de type fini.
Les points de type fini remplacent avantageusement les points fermés pour les espaces
algébriques et les champs algébriques. Par exemple, il y en a toujours « assez ».

\begin{lemma}
\label{lemma-point-finite-type}
Soit $S$ un schéma. Soit $X$ un espace algébrique sur $S$.
Soit $x \in |X|$. Les conditions suivantes sont équivalentes :
\begin{enumerate}
\item Il existe un morphisme $\Spec(k) \to X$
qui est localement de type fini et représente $x$.
\item Il existe un schéma $U$, un point fermé $u \in U$ et un morphisme étale
$\varphi : U \to X$ tels que $\varphi(u) = x$.
\end{enumerate}
\end{lemma}

\begin{proof}
Soient $u \in U$ et $U \to X$ comme en (2). Alors
$\Spec(\kappa(u)) \to U$ est de type fini, et $U \to X$ est
représentable et localement de type fini (d'après le principe général
Espaces, lemme
\ref{spaces-lemma-representable-transformations-property-implication}
et
Morphismes, lemmes \ref{morphisms-lemma-etale-locally-finite-presentation} et
\ref{morphisms-lemma-finite-presentation-finite-type}). On voit donc
que (1) est satisfaite d'après le
lemme \ref{lemma-composition-finite-type}.

\medskip\noindent
Réciproquement, supposons que $\Spec(k) \to X$ soit localement de type fini et
représente $x$. Soit $U \to X$ un morphisme étale surjectif, où $U$
est un schéma. Par hypothèse, $U \times_X \Spec(k) \to U$ est localement
de type fini. Choisissons un point de type fini $v$ de
$U \times_X \Spec(k)$ (il en existe au moins un ; voir
Morphismes, lemme \ref{morphisms-lemma-identify-finite-type-points}).
D'après
Morphismes, lemme \ref{morphisms-lemma-finite-type-points-morphism}
l'image $u \in U$ de $v$ est un point de type fini de $U$.
Par conséquent, d'après
Morphismes, lemme \ref{morphisms-lemma-identify-finite-type-points}
quitte à rétrécir $U$, on peut supposer que $u$ est un point fermé de $U$, c'est-à-dire que
(2) est satisfaite.
\end{proof}

\begin{definition}
\label{definition-finite-type-point}
Soit $S$ un schéma. Soit $X$ un espace algébrique sur $S$.
On dit qu'un point $x \in |X|$ est un {\it point de type fini}\footnote{C'est un
léger abus de langage : il serait peut-être plus correct de dire
« point localement de type fini ».} si les conditions équivalentes du
lemme \ref{lemma-point-finite-type}
sont satisfaites. On note $X_{\text{ft-pts}}$ l'ensemble des points de type fini
de $X$.
\end{definition}

\noindent
On peut décrire comme suit l'ensemble des points de type fini.

\begin{lemma}
\label{lemma-identify-finite-type-points}
Soit $S$ un schéma. Soit $X$ un espace algébrique sur $S$. On a
$$
X_{\text{ft-pts}} =
\bigcup\nolimits_{\varphi : U \to X\text{ étale }} |\varphi|(U_0)
$$
où $U_0$ est l'ensemble des points fermés de $U$.
Ici, on peut faire parcourir à $U$ tous les schémas étales sur $X$, ou tous les
schémas affines étales sur $X$.
\end{lemma}

\begin{proof}
Cela résulte immédiatement du
lemme \ref{lemma-point-finite-type}.
\end{proof}

\begin{lemma}
\label{lemma-finite-type-points-morphism}
Soit $S$ un schéma.
Soit $f : X \to Y$ un morphisme d'espaces algébriques sur $S$.
Si $f$ est localement de type fini, alors
$f(X_{\text{ft-pts}}) \subset Y_{\text{ft-pts}}$.
\end{lemma}

\begin{proof}
Prenons $x \in X_{\text{ft-pts}}$. Représentons $x$ par un morphisme localement de type fini
$x : \Spec(k) \to X$. Alors $f \circ x$ est localement de type fini d'après le
lemme \ref{lemma-composition-finite-type}.
Ainsi, $f(x) \in Y_{\text{ft-pts}}$.
\end{proof}

\begin{lemma}
\label{lemma-finite-type-points-surjective-morphism}
Soit $S$ un schéma.
Soit $f : X \to Y$ un morphisme d'espaces algébriques sur $S$.
Si $f$ est localement de type fini et surjectif, alors
$f(X_{\text{ft-pts}}) = Y_{\text{ft-pts}}$.
\end{lemma}

\begin{proof}
On a $f(X_{\text{ft-pts}}) \subset Y_{\text{ft-pts}}$ d'après le
lemme \ref{lemma-finite-type-points-morphism}.
Soit $y \in |Y|$ un point de type fini. Représentons $y$ par un morphisme
$\Spec(k) \to Y$ localement de type fini.
Comme $f$ est surjectif, l'espace algébrique
$X_k = \Spec(k) \times_Y X$ est non vide et possède donc
un point de type fini $x \in |X_k|$ d'après le
lemme \ref{lemma-identify-finite-type-points}.
Or $X_k \to X$ est un morphisme localement de type fini comme changement de base
de $\Spec(k) \to Y$
(lemme \ref{lemma-base-change-finite-type}).
Ainsi, l'image de $x$ dans $X$ est un point de type fini d'après le
lemme \ref{lemma-finite-type-points-morphism}
et s'envoie sur $y$ par construction.
\end{proof}

\begin{lemma}
\label{lemma-enough-finite-type-points}
Soit $S$ un schéma. Soit $X$ un espace algébrique sur $S$.
Pour toute partie localement fermée $T \subset |X|$, on a
$$
T \not = \emptyset
\Rightarrow
T \cap X_{\text{ft-pts}} \not = \emptyset.
$$
En particulier, pour toute partie fermée $T \subset |X|$,
$T \cap X_{\text{ft-pts}}$ est dense dans $T$.
\end{lemma}

\begin{proof}
Soit $i : Z \to X$ la structure réduite induite de sous-espace sur $T$ ; voir la
remarque \ref{remark-space-structure-locally-closed-subset}.
Toute immersion est localement de type fini ; voir le
lemme \ref{lemma-immersion-locally-finite-type}.
Ainsi, d'après le
lemme \ref{lemma-finite-type-points-morphism},
on a $Z_{\text{ft-pts}} \subset X_{\text{ft-pts}} \cap T$.
Enfin, tout schéma affine non vide $U$ muni d'un morphisme étale vers
$Z$ possède au moins un point fermé. Ainsi, $Z$ possède au moins un
point de type fini d'après le
lemme \ref{lemma-identify-finite-type-points}.
Le lemme en résulte.
\end{proof}

\noindent
Voici une autre caractérisation, plus technique, d'un point de type fini
d'un espace algébrique.

\begin{lemma}
\label{lemma-point-finite-type-monomorphism}
Soit $S$ un schéma. Soit $X$ un espace algébrique sur $S$.
Soit $x \in |X|$. Les conditions suivantes sont équivalentes :
\begin{enumerate}
\item $x$ est un point de type fini ;
\item il existe un espace algébrique $Z$ dont l'espace topologique sous-jacent
$|Z|$ est un singleton, et un morphisme $f : Z \to X$ localement de type
fini tel que $\{x\} = |f|(|Z|)$ ;
\item il existe un espace algébrique $Z$ et un morphisme $f : Z \to X$
possédant les propriétés suivantes :
\begin{enumerate}
\item il existe un morphisme étale surjectif
$z : \Spec(k) \to Z$, où $k$ est un corps ;
\item $f$ est localement de type fini ;
\item $f$ est un monomorphisme ;
\item $x = f(z)$.
\end{enumerate}
\end{enumerate}
\end{lemma}

\begin{proof}
Supposons que $x$ soit un point de type fini. Choisissons un schéma affine $U$,
un point fermé $u \in U$ et un morphisme étale $\varphi : U \to X$
avec $\varphi(u) = x$ ; voir le
lemme \ref{lemma-identify-finite-type-points}.
Posons comme d'habitude $u = \Spec(\kappa(u))$. Les morphismes de projection
$u \times_X u \to u$ sont les composés
$$
u \times_X u \to u \times_X U \to u \times_X X = u
$$
où la première flèche est une immersion fermée (un changement de base de
$u \to U$) et la seconde est étale (un changement de base du morphisme étale
$U \to X$). Ainsi, $u \times_X U$ est une somme disjointe de spectres
d'extensions finies séparables de $k$ (voir
Morphismes, lemme \ref{morphisms-lemma-etale-over-field})
et, par conséquent, le sous-schéma fermé $u \times_X u$ est une somme disjointe de
spectres d'extensions finies séparables de $k$ ; ainsi, $u \times_X u \to u$ est étale. D'après
Espaces, théorème \ref{spaces-theorem-presentation}
on voit que $Z = u/u \times_X u$ est un espace algébrique. Par construction,
le diagramme
$$
\xymatrix{
u \ar[d] \ar[r] & U \ar[d] \\
Z \ar[r] & X
}
$$
est commutatif et ses flèches verticales sont étales. Ainsi, $Z \to X$ est localement
de type fini (voir le
lemme \ref{lemma-finite-type-local}).
Par construction, le morphisme $Z \to X$ est un monomorphisme et
l'image de $u$ est $x$. La condition (3) est donc satisfaite.

\medskip\noindent
Il est clair que (3) implique (2).
Si (2) est satisfaite, alors $x$ est un point de type fini de $X$ d'après le
lemme \ref{lemma-finite-type-points-morphism}
(et le
lemme \ref{lemma-enough-finite-type-points}
montre que $Z_{\text{ft-pts}}$ est non vide, c'est-à-dire que l'unique point de
$Z$ est un point de type fini de $Z$).
\end{proof}





\section{Espaces de Nagata}
\label{section-nagata}

\noindent
Voir Propriétés des espaces, section
\ref{spaces-properties-section-types-properties} pour la définition
d'un espace algébrique de Nagata.

\begin{lemma}
\label{lemma-finite-type-nagata}
Soit $S$ un schéma.
Soit $f : X \to Y$ un morphisme d'espaces algébriques sur $S$.
Si $Y$ est de Nagata et si $f$ est localement de type fini, alors $X$ est de Nagata.
\end{lemma}

\begin{proof}
Soit $V$ un schéma et soit $V \to Y$ un morphisme étale surjectif.
Soit $U$ un schéma et soit $U \to X \times_Y V$ un morphisme étale
surjectif. Si $Y$ est de Nagata, alors $V$ est un schéma de Nagata.
Si $X \to Y$ est localement de type fini, alors $U \to V$ est localement
de type fini. Ainsi, $U$ est un schéma de Nagata d'après
Morphismes, lemme \ref{morphisms-lemma-finite-type-nagata}.
Alors $X$ est de Nagata par définition.
\end{proof}

\begin{lemma}
\label{lemma-ubiquity-nagata}
Les types suivants d'espaces algébriques sont de Nagata.
\begin{enumerate}
\item Tout espace algébrique localement de type fini sur un schéma de Nagata.
\item Tout espace algébrique localement de type fini sur un corps.
\item Tout espace algébrique localement de type fini sur un
anneau local noethérien complet.
\item Tout espace algébrique localement de type fini sur $\mathbf{Z}$.
\item Tout espace algébrique localement de type fini sur un anneau de Dedekind de
caractéristique nulle.
\item Et ainsi de suite.
\end{enumerate}
\end{lemma}

\begin{proof}
La première assertion résulte du lemme \ref{lemma-finite-type-nagata}.
Il en va donc de même des autres ; voir
Morphismes, lemme \ref{morphisms-lemma-ubiquity-nagata}.
\end{proof}








\section{Morphismes quasi-finis}
\label{section-quasi-finite}

\noindent
La propriété « localement quasi-fini » des morphismes de schémas est
locale sur la source et le but pour la topologie étale ; voir
Descente, remarque \ref{descent-remark-list-local-source-target}.
Elle est aussi stable par changement de base et locale sur le but pour la
topologie fpqc ; voir Morphismes, lemme \ref{morphisms-lemma-base-change-quasi-finite}, et
Descente, lemme \ref{descent-lemma-descending-property-quasi-finite}.
Ainsi, d'après le
lemme \ref{lemma-local-source-target}
ci-dessus, on peut définir comme suit ce que signifie, pour un morphisme
d'espaces algébriques, d'être localement quasi-fini ; cette définition
coïncide avec la notion déjà définie dans la
section \ref{section-representable}
lorsque le morphisme est représentable.

\begin{definition}
\label{definition-locally-quasi-finite}
Soit $S$ un schéma.
Soit $f : X \to Y$ un morphisme d'espaces algébriques sur $S$.
\begin{enumerate}
\item On dit que $f$ est
{\it localement quasi-fini} si les conditions équivalentes du
lemme \ref{lemma-local-source-target} sont satisfaites pour
$\mathcal{P} = \text{localement quasi-fini}$.
\item Soit $x \in |X|$. On dit que $f$ est {\it quasi-fini en $x$}
s'il existe un voisinage ouvert $X' \subset X$ de $x$ tel
que $f|_{X'} : X' \to Y$ soit localement quasi-fini.
\item Un morphisme d'espaces algébriques $f : X \to Y$ est
{\it quasi-fini} s'il est localement quasi-fini et quasi-compact.
\end{enumerate}
\end{definition}

\noindent
La dernière partie est compatible avec la notion de morphisme quasi-fini pour les morphismes
de schémas d'après
Morphismes, lemme \ref{morphisms-lemma-quasi-finite-locally-quasi-compact}.

\begin{lemma}
\label{lemma-base-change-quasi-finite-locus}
Soit $S$ un schéma. Soient $f : X \to Y$ et $g : Y' \to Y$ des morphismes
d'espaces algébriques sur $S$. Notons $f' : X' \to Y'$ le changement de base
de $f$ par $g$. Notons $g' : X' \to X$ la projection.
Supposons que $f$ soit localement de type fini.
Soient $W \subset |X|$, resp.\ $W' \subset |X'|$,
l'ensemble des points où $f$, resp.\ $f'$, est quasi-fini.
\begin{enumerate}
\item $W \subset |X|$ et $W' \subset |X'|$ sont ouverts ;
\item $W' = (g')^{-1}(W)$, c'est-à-dire que la formation du lieu où
$f$ est quasi-fini commute aux changements de base ;
\item le changement de base d'un morphisme localement quasi-fini est
localement quasi-fini ;
\item le changement de base d'un morphisme quasi-fini est quasi-fini.
\end{enumerate}
\end{lemma}

\begin{proof}
Choisissons un schéma $V$ et un morphisme étale surjectif $V \to Y$.
Choisissons un schéma $U$ et un morphisme étale surjectif $U \to V \times_Y X$.
Choisissons un schéma $V'$ et un morphisme étale surjectif $V' \to Y' \times_Y V$.
Posons $U' = V' \times_V U$, de sorte que $U' \to X'$ soit lui aussi un morphisme
étale surjectif. On a le diagramme
$$
\vcenter{
\xymatrix{
U' \ar[d] \ar[r] & U \ar[d] \\
V' \ar[r] & V
}
}
\quad\text{au-dessus de}\quad
\vcenter{
\xymatrix{
X' \ar[d] \ar[r] & X \ar[d] \\
Y' \ar[r] & Y
}
}
$$
Choisissons $u \in |U|$ d'image $x \in |X|$.
La propriété d'être « localement quasi-fini » est locale pour la topologie étale sur
la source et le but ; voir
Descente, remarque \ref{descent-remark-list-local-source-target}.
Les lemmes \ref{lemma-local-source-target-at-point} et
\ref{lemma-local-source-target-global-implies-local} s'appliquent donc, et
on voit que $f : X \to Y$ est quasi-fini en $x$
si et seulement si $U \to V$ est quasi-fini en $u$.
Il en va de même de $f' : X' \to Y'$ et du morphisme $U' \to V'$.
Ainsi, les assertions (1), (2) et (3) se ramènent à
Morphismes, lemmes \ref{morphisms-lemma-base-change-quasi-finite} et
\ref{morphisms-lemma-quasi-finite-points-open}.
L'assertion (4) résulte de (3) et du lemme \ref{lemma-base-change-quasi-compact}.
\end{proof}

\begin{lemma}
\label{lemma-composition-quasi-finite}
Le composé de morphismes quasi-finis est quasi-fini.
Il en va de même pour « localement quasi-fini ».
\end{lemma}

\begin{proof}
Voir la remarque \ref{remark-composition-P} et
Morphismes, lemme \ref{morphisms-lemma-composition-quasi-finite}.
\end{proof}

\begin{lemma}
\label{lemma-base-change-quasi-finite}
Un changement de base d'un morphisme quasi-fini est quasi-fini.
Il en va de même pour « localement quasi-fini ».
\end{lemma}

\begin{proof}
Conséquence immédiate du lemme \ref{lemma-base-change-quasi-finite-locus}.
\end{proof}

\noindent
Le lemme suivant caractérise les morphismes localement quasi-finis comme les
morphismes localement de type fini à « fibres discrètes ».
Cependant, cela ne revient pas à demander que $|X| \to |Y|$ ait
des fibres discrètes, comme le montre la discussion de
Exemples, section
\ref{examples-section-specializations-fibre-etale}.

\begin{lemma}
\label{lemma-locally-quasi-finite}
Soit $S$ un schéma. Soit $f : X \to Y$ un morphisme d'espaces algébriques.
Supposons que $f$ soit localement de type fini. Les conditions suivantes sont équivalentes :
\begin{enumerate}
\item $f$ est localement quasi-fini ;
\item pour tout morphisme $\Spec(k) \to Y$, où $k$ est un corps,
l'espace $|X_k|$ est discret. Ici, $X_k = \Spec(k) \times_Y X$.
\end{enumerate}
\end{lemma}

\begin{proof}
Supposons que $f$ soit localement quasi-fini. Soit $\Spec(k) \to Y$ comme
en (2). Choisissons un morphisme étale surjectif $U \to X$, où $U$ est un schéma.
Alors $U_k = \Spec(k) \times_Y U \to X_k$ est un
morphisme étale d'espaces algébriques d'après
Propriétés des espaces, lemme \ref{spaces-properties-lemma-base-change-etale}.
D'après le
lemme \ref{lemma-base-change-quasi-finite},
on voit que $X_k \to \Spec(k)$ est localement quasi-fini. Par
définition, cela signifie que $U_k \to \Spec(k)$ est localement
quasi-fini. Ainsi, $|U_k|$ est discret d'après
Morphismes, lemme \ref{morphisms-lemma-locally-quasi-finite-fibres}.
Comme $|U_k| \to |X_k|$ est surjective et ouverte, on en conclut que $|X_k|$
est discret.

\medskip\noindent
Réciproquement, supposons (2). Choisissons un morphisme étale surjectif $V \to Y$,
où $V$ est un schéma. Choisissons un morphisme étale surjectif
$U \to V \times_Y X$, où $U$ est un schéma.
Remarquons que $U \to V$ est localement de type fini puisque $f$ l'est.
On a le diagramme
$$
\xymatrix{
U \ar[r] \ar[rd] & X \times_Y V \ar[d] \ar[r] & V \ar[d] \\
& X \ar[r] & Y
}
$$
Si $f$ n'est pas localement quasi-fini, alors $U \to V$ ne l'est pas.
Il existe donc une spécialisation $u \leadsto u'$ pour certains $u, u' \in U$
au-dessus d'un même point $v \in V$ ; voir
Morphismes, lemme \ref{morphisms-lemma-quasi-finite-at-point-characterize}.
Nous affirmons que $u, u'$ n'ont pas la même image dans
$X_v = \Spec(\kappa(v)) \times_Y X$,
ce qui contredira, comme voulu, l'hypothèse que $|X_v|$ est discret.
Posons $d = \text{trdeg}_{\kappa(v)}(\kappa(u))$ et
$d' = \text{trdeg}_{\kappa(v)}(\kappa(u'))$.
On voit alors que $d > d'$ d'après
Morphismes, lemme \ref{morphisms-lemma-dimension-fibre-specialization}.
Remarquons que $U_v$ (la fibre de $U \to V$ au-dessus de $v$) est le produit fibré de
$U$ et de $X_v$ au-dessus de $X \times_Y V$ ; ainsi, $U_v \to X_v$ est étale (comme
changement de base du morphisme étale $U \to X \times_Y V$). Si
$u, u' \in U_v$ s'envoient sur le même élément de $|X_v|$, alors il existe un point
$r \in R_v = U_v \times_{X_v} U_v$ tel que $t(r) = u$ et $s(r) = u'$ ; voir
Propriétés des espaces, lemme \ref{spaces-properties-lemma-points-cartesian}.
Remarquons que $s, t : R_v \to U_v$ sont des morphismes étales de schémas
sur $\kappa(v)$ ; ainsi, les extensions $\kappa(u) \subset \kappa(r) \supset \kappa(u')$
sont finies séparables entre corps sur $\kappa(v)$ (voir
Morphismes, lemme \ref{morphisms-lemma-etale-over-field}).
On en conclut que les degrés de transcendance sont égaux.
Cette contradiction achève la démonstration.
\end{proof}

\begin{lemma}
\label{lemma-quasi-finite-local}
Soit $S$ un schéma.
Soit $f : X \to Y$ un morphisme d'espaces algébriques sur $S$.
Les conditions suivantes sont équivalentes :
\begin{enumerate}
\item $f$ est localement quasi-fini ;
\item pour tout $x \in |X|$, le morphisme $f$ est quasi-fini en $x$ ;
\item pour tout schéma $Z$ et tout morphisme $Z \to Y$, le morphisme
$Z \times_Y X \to Z$ est localement quasi-fini ;
\item pour tout schéma affine $Z$ et tout morphisme
$Z \to Y$, le morphisme $Z \times_Y X \to Z$ est localement quasi-fini ;
\item il existe un schéma $V$ et un morphisme étale surjectif
$V \to Y$ tels que $V \times_Y X \to V$ soit localement quasi-fini ;
\item il existe un schéma $U$ et un morphisme étale surjectif
$\varphi : U \to X$ tels que le composé $f \circ \varphi$
soit localement quasi-fini ;
\item pour tout diagramme commutatif
$$
\xymatrix{
U \ar[d] \ar[r] & V \ar[d] \\
X \ar[r] & Y
}
$$
où $U$ et $V$ sont des schémas et où les flèches verticales sont étales,
la flèche horizontale supérieure est localement quasi-finie ;
\item il existe un diagramme commutatif
$$
\xymatrix{
U \ar[d] \ar[r] & V \ar[d] \\
X \ar[r] & Y
}
$$
où $U$ et $V$ sont des schémas, où les flèches verticales sont étales et où
$U \to X$ est surjectif, tel que la flèche horizontale supérieure soit
localement quasi-finie ;
\item il existe des recouvrements de Zariski $Y = \bigcup_{i \in I} Y_i$
et $f^{-1}(Y_i) = \bigcup X_{ij}$ tels que
chaque morphisme $X_{ij} \to Y_i$ soit localement quasi-fini.
\end{enumerate}
\end{lemma}

\begin{proof}
La démonstration est omise.
\end{proof}

\begin{lemma}
\label{lemma-immersion-quasi-finite}
Une immersion est localement quasi-finie.
\end{lemma}

\begin{proof}
La démonstration est omise.
\end{proof}

\begin{lemma}
\label{lemma-permanence-quasi-finite}
Soit $S$ un schéma.
Soient $X \to Y \to Z$ des morphismes d'espaces algébriques sur $S$.
Si $X \to Z$ est localement quasi-fini, alors $X \to Y$
est localement quasi-fini.
\end{lemma}

\begin{proof}
Choisissons un diagramme commutatif
$$
\xymatrix{
U \ar[d] \ar[r] & V \ar[d] \ar[r] & W \ar[d] \\
X \ar[r] & Y \ar[r] & Z
}
$$
dont les flèches verticales sont étales et surjectives. (Voir
Espaces, lemme \ref{spaces-lemma-lift-morphism-presentations}.)
Appliquons
Morphismes, lemme \ref{morphisms-lemma-permanence-quasi-finite}
à la ligne supérieure.
\end{proof}

\begin{lemma}
\label{lemma-quasi-finite-at-a-finite-number-of-points}
Soit $S$ un schéma. Soit $f : X \to Y$ un morphisme de type fini
d'espaces algébriques sur $S$. Soit $y \in |Y|$.
Il n'y a qu'un nombre fini
de points de $|X|$ au-dessus de $y$ où $f$ est quasi-fini.
\end{lemma}

\begin{proof}
Choisissons un corps $k$ et un morphisme $\Spec(k) \to Y$ dans la classe
d'équivalence déterminée par $y$. La fibre $X_k = \Spec(k) \times_Y X$ est un
espace algébrique de type fini sur un corps, donc en particulier quasi-compact.
L'application $|X_k| \to |X|$ a pour image la fibre de $|X| \to |Y|$
au-dessus de $y$ (Propriétés des espaces, lemme
\ref{spaces-properties-lemma-points-cartesian}).
De plus, l'ensemble des points où $X_k \to \Spec(k)$ est
quasi-fini s'envoie sur l'ensemble des points au-dessus de $y$ où
$f$ est quasi-fini d'après le lemme \ref{lemma-base-change-quasi-finite-locus}.
Choisissons un schéma affine $U$ et un morphisme étale surjectif $U \to X_k$
(Propriétés des espaces, lemme
\ref{spaces-properties-lemma-quasi-compact-affine-cover}).
Alors $U \to \Spec(k)$ est un morphisme de type fini et il n'y a qu'un
nombre fini de points où ce morphisme est quasi-fini ;
voir Morphismes, lemme
\ref{morphisms-lemma-quasi-finite-at-a-finite-number-of-points}.
Comme $X_k \to \Spec(k)$ est quasi-fini en un point $x'$ si et seulement
si ce point est l'image d'un point de $U$ où $U \to \Spec(k)$ est
quasi-fini, on conclut.
\end{proof}

\begin{lemma}
\label{lemma-monomorphism-loc-finite-type-loc-quasi-finite}
Soit $S$ un schéma.
Soit $f : X \to Y$ un morphisme d'espaces algébriques sur $S$.
Si $f$ est localement de type fini et est un monomorphisme, alors $f$
est séparé et localement quasi-fini.
\end{lemma}

\begin{proof}
Un monomorphisme est séparé ; voir le
lemme \ref{lemma-monomorphism-separated}.
D'après le
lemme \ref{lemma-quasi-finite-local},
il suffit de démontrer le lemme après avoir effectué un changement de base
par $Z \to Y$, avec $Z$ affine. On peut donc supposer que $Y$ est un
schéma affine. Choisissons un schéma affine $U$ et un morphisme étale
$U \to X$. Comme $X \to Y$ est localement de type fini, le morphisme
de schémas affines $U \to Y$ est de type fini.
Comme $X \to Y$ est un monomorphisme, on a $U \times_X U = U \times_Y U$.
En particulier, les applications $U \times_Y U \to U$ sont étales.
Soit $y \in Y$. Alors soit $U_y$ est vide, soit
$\Spec(\kappa(u)) \times_{\Spec(\kappa(y))} U_y$
est isomorphe à la fibre de $U \times_Y U \to U$ au-dessus de $u$ pour
un certain $u \in U$ au-dessus de $y$. Il s'ensuit que les fibres de
$U \to Y$ sont des ensembles discrets finis (car $U \times_Y U \to U$
est un morphisme étale de schémas affines ; voir
Morphismes, lemme \ref{morphisms-lemma-etale-over-field}).
Ainsi, $U \to Y$ est quasi-fini ; voir
Morphismes, lemme \ref{morphisms-lemma-quasi-finite-at-point-characterize}.
Comme $U \to X$ était un morphisme étale quelconque, avec $U$ affine,
il en résulte que $X \to Y$ est localement quasi-fini.
\end{proof}













\section{Morphismes de présentation finie}
\label{section-finite-presentation}

\noindent
La propriété « localement de présentation finie » des morphismes de schémas est
locale sur la source et le but pour la topologie étale ; voir
Descente, remarque \ref{descent-remark-list-local-source-target}.
Elle est aussi stable par changement de base et locale sur le but pour la topologie fpqc ; voir
Morphismes, lemme \ref{morphisms-lemma-base-change-finite-presentation}, et
Descente,
lemme \ref{descent-lemma-descending-property-locally-finite-presentation}.
Ainsi, d'après le
lemme \ref{lemma-local-source-target}
ci-dessus, on peut définir comme suit ce que signifie, pour un morphisme
d'espaces algébriques, d'être localement de présentation
finie ; cette définition coïncide avec la notion déjà définie dans la
section \ref{section-representable}
lorsque le morphisme est représentable.

\begin{definition}
\label{definition-locally-finite-presentation}
Soit $S$ un schéma.
Soit $f : X \to Y$ un morphisme d'espaces algébriques sur $S$.
\begin{enumerate}
\item On dit que $f$ est {\it localement de présentation finie} si
les conditions équivalentes du
lemme \ref{lemma-local-source-target}
sont satisfaites pour $\mathcal{P} =$ « localement de présentation finie ».
\item Soit $x \in |X|$. On dit que $f$ est de {\it présentation finie en $x$}
s'il existe un voisinage ouvert $X' \subset X$ de $x$ tel
que $f|_{X'} : X' \to Y$ soit localement de
présentation finie\footnote{Il semblerait maladroit de dire « localement de présentation finie
en $x$ », mais la terminologie retenue peut prêter à confusion : en effet,
« de présentation finie en $x$ » {\bf ne signifie pas} qu'il existe
un voisinage ouvert $X' \subset X$ tel que $f|_{X'}$ soit de présentation
finie.}.
\item Un morphisme d'espaces algébriques $f : X \to Y$ est
{\it de présentation finie}
s'il est localement de présentation finie, quasi-compact et
quasi-séparé.
\end{enumerate}
\end{definition}

\noindent
Remarquons qu'un morphisme de présentation finie n'est {\bf pas} simplement un morphisme
quasi-compact et localement de présentation finie.

\begin{lemma}
\label{lemma-composition-finite-presentation}
Le composé de morphismes de présentation finie est de présentation finie.
Il en va de même pour « localement de présentation finie ».
\end{lemma}

\begin{proof}
Voir la remarque \ref{remark-composition-P} et
Morphismes, lemme \ref{morphisms-lemma-composition-finite-presentation}.
On utilise aussi le résultat pour les morphismes quasi-compacts et quasi-séparés
(lemmes \ref{lemma-composition-quasi-compact} et
\ref{lemma-composition-separated}).
\end{proof}

\begin{lemma}
\label{lemma-base-change-finite-presentation}
Un changement de base d'un morphisme de présentation finie est de présentation finie.
Il en va de même pour « localement de présentation finie ».
\end{lemma}

\begin{proof}
Voir la remarque \ref{remark-base-change-P} et
Morphismes, lemme \ref{morphisms-lemma-base-change-finite-presentation}.
On utilise aussi le résultat pour les morphismes quasi-compacts et quasi-séparés
(lemmes \ref{lemma-base-change-quasi-compact} et
\ref{lemma-base-change-separated}).
\end{proof}

\begin{lemma}
\label{lemma-finite-presentation-local}
Soit $S$ un schéma.
Soit $f : X \to Y$ un morphisme d'espaces algébriques sur $S$.
Les conditions suivantes sont équivalentes :
\begin{enumerate}
\item $f$ est localement de présentation finie ;
\item pour tout $x \in |X|$, le morphisme $f$ est de présentation finie en $x$ ;
\item pour tout schéma $Z$ et tout morphisme $Z \to Y$, le morphisme
$Z \times_Y X \to Z$ est localement de présentation finie ;
\item pour tout schéma affine $Z$ et tout morphisme
$Z \to Y$, le morphisme $Z \times_Y X \to Z$ est localement de présentation finie ;
\item il existe un schéma $V$ et un morphisme étale surjectif
$V \to Y$ tels que $V \times_Y X \to V$ soit
localement de présentation finie ;
\item il existe un schéma $U$ et un morphisme étale surjectif
$\varphi : U \to X$ tels que le composé $f \circ \varphi$
soit localement de présentation finie ;
\item pour tout diagramme commutatif
$$
\xymatrix{
U \ar[d] \ar[r] & V \ar[d] \\
X \ar[r] & Y
}
$$
où $U$ et $V$ sont des schémas et où les flèches verticales sont étales,
la flèche horizontale supérieure est localement de présentation finie ;
\item il existe un diagramme commutatif
$$
\xymatrix{
U \ar[d] \ar[r] & V \ar[d] \\
X \ar[r] & Y
}
$$
où $U$ et $V$ sont des schémas, où les flèches verticales sont étales et où
$U \to X$ est surjectif, tel que la flèche horizontale supérieure soit
localement de présentation finie ;
\item il existe des recouvrements de Zariski $Y = \bigcup_{i \in I} Y_i$
et $f^{-1}(Y_i) = \bigcup X_{ij}$ tels que
chaque morphisme $X_{ij} \to Y_i$ soit localement de présentation finie.
\end{enumerate}
\end{lemma}

\begin{proof}
La démonstration est omise.
\end{proof}

\begin{lemma}
\label{lemma-finite-presentation-finite-type}
Un morphisme localement de présentation finie est localement de type fini.
Un morphisme de présentation finie est de type fini.
\end{lemma}

\begin{proof}
Soit $f : X \to Y$ un morphisme d'espaces algébriques localement de
présentation finie. Cela signifie qu'il existe un diagramme comme dans le
lemme \ref{lemma-local-source-target},
où $h$ est localement de présentation finie et où $a$ est une flèche verticale surjective. D'après
Morphismes, lemme \ref{morphisms-lemma-finite-presentation-finite-type}
$h$ est localement de type fini.
Ainsi, $X \to Y$ est localement de type fini par définition.
Si $f$ est de présentation finie, il est quasi-compact et
il en résulte que $f$ est de type fini.
\end{proof}

\begin{lemma}
\label{lemma-finite-presentation-noetherian}
Soit $S$ un schéma.
Soit $f : X \to Y$ un morphisme d'espaces algébriques sur $S$.
Si $f$ est de présentation finie et si $Y$ est noethérien,
alors $X$ est noethérien.
\end{lemma}

\begin{proof}
Supposons que $f$ soit de présentation finie et que $Y$ soit noethérien. D'après les
lemmes \ref{lemma-finite-presentation-finite-type} et
\ref{lemma-locally-finite-type-locally-noetherian}
on voit que $X$ est localement noethérien. Comme $f$ est quasi-compact
et que $Y$ est quasi-compact, on voit que $X$ est quasi-compact.
Comme $f$ est de présentation finie, il est quasi-séparé (voir la
définition \ref{definition-locally-finite-presentation}),
et comme $Y$ est noethérien, il est quasi-séparé (voir
Propriétés des espaces,
définition \ref{spaces-properties-definition-noetherian}).
Ainsi, $X$ est quasi-séparé d'après le
lemme \ref{lemma-separated-over-separated}.
On a donc vérifié les trois conditions de
Propriétés des espaces,
définition \ref{spaces-properties-definition-noetherian}
et cela conclut.
\end{proof}

\begin{lemma}
\label{lemma-noetherian-finite-type-finite-presentation}
Soit $S$ un schéma.
Soit $f : X \to Y$ un morphisme d'espaces algébriques sur $S$.
\begin{enumerate}
\item Si $Y$ est localement noethérien et si $f$ est localement de type fini,
alors $f$ est localement de présentation finie.
\item Si $Y$ est localement noethérien et si $f$ est de type fini et quasi-séparé,
alors $f$ est de présentation finie.
\end{enumerate}
\end{lemma}

\begin{proof}
Supposons que $f : X \to Y$ soit localement de type fini et que $Y$ soit localement noethérien.
Cela signifie qu'il existe un diagramme comme dans le
lemme \ref{lemma-local-source-target},
où $h$ est localement de type fini et où $a$ est une flèche verticale surjective. D'après
Morphismes, lemme
\ref{morphisms-lemma-noetherian-finite-type-finite-presentation}
$h$ est localement de présentation finie.
Ainsi, $X \to Y$ est localement de présentation finie par définition.
Ceci démontre (1).
Si $f$ est de type fini et quasi-séparé, alors il est aussi
quasi-compact et quasi-séparé, et (2) en résulte immédiatement.
\end{proof}

\begin{lemma}
\label{lemma-finite-presentation-quasi-compact-quasi-separated}
Soit $S$ un schéma. Soit $Y$ un espace algébrique sur $S$ qui est
quasi-compact et quasi-séparé. Si $X$ est de présentation finie sur
$Y$, alors $X$ est quasi-compact et quasi-séparé.
\end{lemma}

\begin{proof}
La démonstration est omise.
\end{proof}

\begin{lemma}
\label{lemma-finite-presentation-permanence}
Soit $S$ un schéma.
Soient $f : X \to Y$ et $Y \to Z$ des morphismes d'espaces algébriques sur $S$.
Si $X$ est localement de présentation finie sur $Z$ et si
$Y$ est localement de type fini sur $Z$, alors $f$ est localement
de présentation finie.
\end{lemma}

\begin{proof}
Choisissons un schéma $W$ et un morphisme étale surjectif $W \to Z$.
Choisissons ensuite un schéma $V$ et un morphisme étale surjectif $V \to W \times_Z Y$.
Choisissons enfin un schéma $U$ et un morphisme étale surjectif
$U \to V \times_Y X$. Par définition, $U$ est localement de présentation finie
sur $W$ et $V$ est localement de type fini sur $W$. D'après
Morphismes, lemme \ref{morphisms-lemma-finite-presentation-permanence}
le morphisme $U \to V$ est localement de présentation finie.
Ainsi, $f$ est localement de présentation finie.
\end{proof}

\begin{lemma}
\label{lemma-diagonal-morphism-finite-type}
Soit $S$ un schéma. Soit $f : X \to Y$ un morphisme d'espaces algébriques
sur $S$, de diagonale $\Delta : X \to X \times_Y X$. Si $f$ est localement de
type fini, alors $\Delta$ est localement de présentation finie. Si $f$ est
quasi-séparé et localement de type fini, alors $\Delta$ est de
présentation finie.
\end{lemma}

\begin{proof}
Remarquons que $\Delta$ est un morphisme sur $X$ (par la seconde
projection $X \times_Y X \to X$). Supposons $f$ localement de type fini.
Remarquons que $X$ est de présentation finie sur $X$ et que $X \times_Y X$ est
de type fini sur $X$ (d'après le lemme \ref{lemma-base-change-finite-type}).
La première assertion résulte donc du
lemme \ref{lemma-finite-presentation-permanence}.
La seconde assertion résulte de la première, des définitions et
du fait qu'un morphisme diagonal est séparé
(lemme \ref{lemma-properties-diagonal}).
\end{proof}

\begin{lemma}
\label{lemma-open-immersion-locally-finite-presentation}
Une immersion ouverte d'espaces algébriques est localement de présentation finie.
\end{lemma}

\begin{proof}
Une immersion ouverte est représentable par définition ; on peut donc
utiliser le principe général
Espaces,
lemme \ref{spaces-lemma-representable-transformations-property-implication}
et
Morphismes,
lemme \ref{morphisms-lemma-open-immersion-locally-finite-presentation}.
\end{proof}

\begin{lemma}
\label{lemma-closed-immersion-finite-presentation}
Une immersion fermée $i : Z \to X$ est de présentation finie si et seulement si
le faisceau quasi-cohérent d'idéaux associé
$\mathcal{I} = \Ker(\mathcal{O}_X \to i_*\mathcal{O}_Z)$
est de type fini (comme $\mathcal{O}_X$-module).
\end{lemma}

\begin{proof}
Soit $U$ un schéma et soit $U \to X$ un morphisme étale surjectif.
D'après le lemme \ref{lemma-finite-presentation-local},
on voit que $i' : Z \times_X U \to U$ est de présentation finie si et
seulement si $i$ l'est. D'après Propriétés des espaces, section
\ref{spaces-properties-section-properties-modules}
on voit que $\mathcal{I}$ est de type fini si et seulement si
$\mathcal{I}|_U = \Ker(\mathcal{O}_U \to i'_*\mathcal{O}_{Z \times_X U})$
l'est. Le résultat découle donc du cas des schémas ; voir Morphismes,
lemme \ref{morphisms-lemma-closed-immersion-finite-presentation}.
\end{proof}







\section{Parties constructibles}
\label{section-constructible}

\noindent
Cette section fait suite à
Propriétés des espaces, section \ref{spaces-properties-section-constructible}.

\begin{lemma}
\label{lemma-inverse-image-constructible}
Soit $S$ un schéma.
Soit $f : X \to Y$ un morphisme d'espaces algébriques sur $S$.
Soit $E \subset |Y|$ une partie.
Si $E$ est localement constructible pour la topologie étale dans $Y$, alors
$f^{-1}(E)$ est localement constructible pour la topologie étale dans $X$.
\end{lemma}

\begin{proof}
Choisissons un schéma $V$ et un morphisme étale surjectif $\varphi : V \to Y$.
Choisissons un schéma $U$ et un morphisme étale surjectif
$U \to V \times_Y X$. Alors $U \to X$ est étale surjectif,
et l'image inverse de $f^{-1}(E)$ dans $U$ est l'image
inverse de $\varphi^{-1}(E)$ par $U \to V$. Le lemme découle donc
du cas des schémas pour $U \to V$
(Morphismes, lemme \ref{morphisms-lemma-inverse-image-constructible})
et de la définition (Propriétés des espaces, définition
\ref{spaces-properties-definition-locally-constructible}).
\end{proof}

\begin{theorem}[Théorème de Chevalley]
\label{theorem-chevalley}
Soit $S$ un schéma.
Soit $f : X \to Y$ un morphisme d'espaces algébriques sur $S$.
Supposons que $f$ soit quasi-compact et localement de présentation finie.
Alors l'image de toute partie de $|X|$ localement constructible pour la topologie étale est
une partie de $|Y|$ localement constructible pour la topologie étale.
\end{theorem}

\begin{proof}
Soit $E \subset |X|$ localement constructible pour la topologie étale.
Soit $V \to Y$ un morphisme étale, avec $V$ affine.
Il suffit de montrer que l'image inverse de $f(E)$ dans
$V$ est constructible ; voir Propriétés des espaces, définition
\ref{spaces-properties-definition-locally-constructible}.
Comme $f$ est quasi-compact, $V \times_Y X$ est un espace algébrique quasi-compact.
Choisissons un schéma affine $U$ et un morphisme étale surjectif
$U \to V \times_Y X$ (Propriétés des espaces, lemme
\ref{spaces-properties-lemma-quasi-compact-affine-cover}).
D'après Propriétés des espaces, lemme \ref{spaces-properties-lemma-points-cartesian},
l'image inverse de $f(E)$ dans $V$ est l'image par $U \to V$
de l'image inverse de $E$ dans $U$.
Le résultat découle donc du cas des schémas ; voir
Morphismes, lemme \ref{morphisms-lemma-chevalley}.
\end{proof}







\section{Morphismes plats}
\label{section-flat}

\noindent
La propriété « plat » des morphismes de schémas est
locale sur la source et le but pour la topologie étale ; voir
Descente, remarque \ref{descent-remark-list-local-source-target}.
Elle est aussi stable par changement de base et locale sur le but pour la topologie fpqc ; voir
Morphismes, lemme \ref{morphisms-lemma-base-change-flat} et
Descente, lemme \ref{descent-lemma-descending-property-flat}.
Ainsi, d'après le
lemme \ref{lemma-local-source-target}
ci-dessus, on peut définir comme suit la notion de morphisme plat d'espaces
algébriques ; cette définition coïncide avec la notion déjà définie dans la
section \ref{section-representable}
lorsque le morphisme est représentable.

\begin{definition}
\label{definition-flat}
Soit $S$ un schéma.
Soit $f : X \to Y$ un morphisme d'espaces algébriques sur $S$.
\begin{enumerate}
\item On dit que $f$ est {\it plat} si les conditions équivalentes du
lemme \ref{lemma-local-source-target} sont satisfaites pour
$\mathcal{P} =$ « plat ».
\item Soit $x \in |X|$. On dit que $f$ est {\it plat en $x$} si les
conditions équivalentes du
lemme \ref{lemma-local-source-target-at-point}
sont satisfaites pour $\mathcal{Q} =$ « l'application induite sur les anneaux locaux est plate ».
\end{enumerate}
Remarquons que la seconde partie a un sens d'après
Descente, lemme \ref{descent-lemma-flat-at-point}.
\end{definition}

\noindent
Vérifions rapidement la cohérence de cette définition.

\begin{lemma}
\label{lemma-flat-is-flat-at-all-points}
Soit $S$ un schéma. Soit $f : X \to Y$ un morphisme d'espaces algébriques
sur $S$. Alors $f$ est plat si et seulement si $f$ est plat en tout point de $|X|$.
\end{lemma}

\begin{proof}
Choisissons un diagramme commutatif
$$
\xymatrix{
U \ar[d]_a \ar[r]_h & V \ar[d]^b \\
X \ar[r]^f & Y
}
$$
où $U$ et $V$ sont des schémas, où les flèches verticales sont étales et où
$a$ est surjectif. Par définition, $f$ est plat si et seulement si $h$ est plat
(définition \ref{definition-P}).
Par définition, $f$ est plat en $x \in |X|$ si et seulement si $h$ est plat
en un certain (ce qui équivaut à tout) $u \in U$ s'envoyant sur $x$
(définition \ref{definition-P-at-point}).
Le lemme résulte donc du fait qu'un morphisme de schémas
est plat si et seulement s'il est plat en tout point de la source
(Morphismes, définition \ref{morphisms-definition-flat}).
\end{proof}

\begin{lemma}
\label{lemma-composition-flat}
Le composé de morphismes plats est plat.
\end{lemma}

\begin{proof}
Voir la remarque \ref{remark-composition-P} et
Morphismes, lemme \ref{morphisms-lemma-composition-flat}.
\end{proof}

\begin{lemma}
\label{lemma-base-change-flat}
Le changement de base d'un morphisme plat est plat.
\end{lemma}

\begin{proof}
Voir la remarque \ref{remark-base-change-P} et
Morphismes, lemme \ref{morphisms-lemma-base-change-flat}.
\end{proof}

\begin{lemma}
\label{lemma-flat-local}
Soit $S$ un schéma.
Soit $f : X \to Y$ un morphisme d'espaces algébriques sur $S$.
Les conditions suivantes sont équivalentes :
\begin{enumerate}
\item $f$ est plat ;
\item pour tout $x \in |X|$, le morphisme $f$ est plat en $x$ ;
\item pour tout schéma $Z$ et tout morphisme $Z \to Y$, le morphisme
$Z \times_Y X \to Z$ est plat ;
\item pour tout schéma affine $Z$ et tout morphisme
$Z \to Y$, le morphisme $Z \times_Y X \to Z$ est plat ;
\item il existe un schéma $V$ et un morphisme étale surjectif
$V \to Y$ tels que $V \times_Y X \to V$ soit plat ;
\item il existe un schéma $U$ et un morphisme étale surjectif
$\varphi : U \to X$ tels que le composé $f \circ \varphi$
soit plat ;
\item pour tout diagramme commutatif
$$
\xymatrix{
U \ar[d] \ar[r] & V \ar[d] \\
X \ar[r] & Y
}
$$
où $U$ et $V$ sont des schémas et où les flèches verticales sont étales,
la flèche horizontale supérieure est plate ;
\item il existe un diagramme commutatif
$$
\xymatrix{
U \ar[d] \ar[r] & V \ar[d] \\
X \ar[r] & Y
}
$$
où $U$ et $V$ sont des schémas, où les flèches verticales sont étales et où
$U \to X$ est surjectif, tel que la flèche horizontale supérieure soit plate ;
\item il existe des recouvrements de Zariski $Y = \bigcup Y_i$ et
$f^{-1}(Y_i) = \bigcup X_{ij}$ tels que
chaque morphisme $X_{ij} \to Y_i$ soit plat.
\end{enumerate}
\end{lemma}

\begin{proof}
La démonstration est omise.
\end{proof}

\begin{lemma}
\label{lemma-fppf-open}
Un morphisme plat localement de présentation finie est universellement ouvert.
\end{lemma}

\begin{proof}
Soit $f : X \to Y$ un morphisme plat et localement de présentation finie
d'espaces algébriques sur $S$. Choisissons un diagramme
$$
\xymatrix{
U \ar[r]_\alpha \ar[d] & V \ar[d] \\
X \ar[r] & Y
}
$$
où $U$ et $V$ sont des schémas et où les flèches verticales sont surjectives et
étales ; voir
Espaces, lemme \ref{spaces-lemma-lift-morphism-presentations}.
D'après les
lemmes \ref{lemma-flat-local} et \ref{lemma-finite-presentation-local},
le morphisme $\alpha$ est plat et localement de présentation finie.
Ainsi, d'après
Morphismes, lemme \ref{morphisms-lemma-fppf-open}
on voit que $\alpha$ est universellement ouvert.
Ainsi, $X \to Y$ est universellement ouvert d'après le
lemme \ref{lemma-universally-open-local}.
\end{proof}

\begin{lemma}
\label{lemma-fpqc-quotient-topology}
Soit $S$ un schéma.
Soit $f : X \to Y$ un morphisme plat, quasi-compact et surjectif
d'espaces algébriques sur $S$.
Une partie $T \subset |Y|$ est ouverte (resp.\ fermée) si et seulement si
$f^{-1}(T)$ est ouverte (resp.\ fermée) dans $|X|$.
Autrement dit, $f$ est submersif, et même universellement submersif.
\end{lemma}

\begin{proof}
Choisissons des schémas affines $V_i$ et des morphismes étales $V_i \to Y$ tels que
$g : V = \coprod V_i \to Y$ soit surjectif ; voir
Propriétés des espaces,
lemme \ref{spaces-properties-lemma-cover-by-union-affines}.
Pour tout $i$, l'espace algébrique $V_i \times_Y X$ est quasi-compact.
On peut donc trouver un schéma affine $U_i$ et un morphisme étale surjectif
$U_i \to V_i \times_Y X$ ; voir
Propriétés des espaces,
lemme \ref{spaces-properties-lemma-quasi-compact-affine-cover}.
Alors le composé $U_i \to V_i \times_Y X \to V_i$ est un morphisme surjectif
et plat entre schémas affines.
Bien sûr, $U = \coprod U_i \to X$ est alors étale surjectif,
et $U \to V \times_Y X$ est surjectif. De plus, le morphisme $U \to V$ est la
somme disjointe des morphismes $U_i \to V_i$. Ainsi, $U \to V$ est surjectif,
quasi-compact et plat. Considérons le diagramme
$$
\xymatrix{
U \ar[r] \ar[d] & X \ar[d] \\
V \ar[r] & Y
}
$$
Par définition de la topologie sur $|Y|$, la partie $T$ est fermée
(resp.\ ouverte) si et seulement si $g^{-1}(T) \subset |V|$ est fermée
(resp.\ ouverte). Il en va de même de
$f^{-1}(T)$ et de son image inverse dans $|U|$.
Comme $U \to V$ est quasi-compact, surjectif et plat, on conclut par
Morphismes, lemme \ref{morphisms-lemma-fpqc-quotient-topology}.
\end{proof}

\begin{lemma}
\label{lemma-flat-at-point-etale-local-rings}
Soit $S$ un schéma.
Soit $f : X \to Y$ un morphisme d'espaces algébriques sur $S$.
Soit $\overline{x}$ un point géométrique de $X$ au-dessus du point
$x \in |X|$. Posons $\overline{y} = f \circ \overline{x}$. Les conditions suivantes
sont équivalentes :
\begin{enumerate}
\item $f$ est plat en $x$ ;
\item l'application entre les anneaux locaux étales
$\mathcal{O}_{Y, \overline{y}} \to \mathcal{O}_{X, \overline{x}}$
est plate.
\end{enumerate}
\end{lemma}

\begin{proof}
Choisissons un diagramme commutatif
$$
\xymatrix{
U \ar[d]_a \ar[r]_h & V \ar[d]^b \\
X \ar[r]^f & Y
}
$$
où $U$ et $V$ sont des schémas, où $a, b$ sont étales, et où
$u \in U$ s'envoie sur $x$. On peut trouver un point géométrique
$\overline{u} : \Spec(k) \to U$ au-dessus de $u$ tel que
$\overline{x} = a \circ \overline{u}$ ; voir
Propriétés des espaces, lemme
\ref{spaces-properties-lemma-geometric-lift-to-usual}.
Posons $\overline{v} = h \circ \overline{u}$, d'image $v \in V$.
On sait que
$$
\mathcal{O}_{X, \overline{x}} = \mathcal{O}_{U, u}^{sh}
\quad\text{et}\quad
\mathcal{O}_{Y, \overline{y}} = \mathcal{O}_{V, v}^{sh}
$$
voir
Propriétés des espaces, lemme
\ref{spaces-properties-lemma-describe-etale-local-ring}.
On obtient un diagramme commutatif
$$
\xymatrix{
\mathcal{O}_{U, u} \ar[r] &
\mathcal{O}_{X, \overline{x}} \\
\mathcal{O}_{V, v} \ar[u] \ar[r] &
\mathcal{O}_{Y, \overline{y}} \ar[u]
}
$$
formé d'anneaux locaux et dont les flèches horizontales sont plates. Il faut montrer que la
flèche verticale gauche est plate si et seulement si la flèche verticale droite l'est.
Algèbre, lemme \ref{algebra-lemma-flatness-descends-more-general}
nous dit que $\mathcal{O}_{U, u}$ est plat sur $\mathcal{O}_{V, v}$
si et seulement si $\mathcal{O}_{X, \overline{x}}$ est plat sur
$\mathcal{O}_{V, v}$. Le résultat découle donc de
Compléments sur la platitude, lemme \ref{flat-lemma-flat-up-down-henselization}.
\end{proof}

\begin{lemma}
\label{lemma-flat-morphism-sites}
Soit $S$ un schéma.
Soit $f : X \to Y$ un morphisme d'espaces algébriques sur $S$.
Alors $f$ est plat si et seulement si le morphisme de sites
$
(f_{small}, f^\sharp) :
(X_\etale, \mathcal{O}_X)
\to
(Y_\etale, \mathcal{O}_Y)
$
associé à $f$ est plat.
\end{lemma}

\begin{proof}
La platitude de $(f_{small}, f^\sharp)$ est définie au moyen de la
platitude de $\mathcal{O}_X$ comme $f_{small}^{-1}\mathcal{O}_Y$-module.
Elle se vérifie sur les fibres ; voir
Modules sur les sites, lemme \ref{sites-modules-lemma-check-flat-stalks}
et
Propriétés des espaces, théorème \ref{spaces-properties-theorem-exactness-stalks}.
Or on a déjà vu que la platitude de $f$ se vérifie sur les fibres ; voir le
lemme \ref{lemma-flat-at-point-etale-local-rings}.
\end{proof}

\begin{lemma}
\label{lemma-flat-pullback-support}
Soit $S$ un schéma. Soit $f : Y \to X$ un morphisme d'espaces algébriques
sur $S$. Soit $\mathcal{F}$ un $\mathcal{O}_X$-module quasi-cohérent
de type fini de support schématique $Z \subset X$.
Si $f$ est plat, alors $f^{-1}(Z)$ est le support schématique de
$f^*\mathcal{F}$.
\end{lemma}

\begin{proof}
En utilisant la caractérisation du support schématique
donnée dans le lemme \ref{lemma-scheme-theoretic-support}
et la caractérisation des morphismes plats au moyen de
recouvrements étales donnée dans le lemme \ref{lemma-flat-local},
on se ramène au cas des schémas, qui est
Morphismes, lemme \ref{morphisms-lemma-flat-pullback-support}.
\end{proof}

\begin{lemma}
\label{lemma-flat-morphism-scheme-theoretically-dense-open}
Soit $S$ un schéma.
Soit $f : X \to Y$ un morphisme plat d'espaces algébriques sur $S$.
Soit $V \to Y$ une immersion ouverte quasi-compacte. Si $V$
est schématiquement dense dans $Y$, alors $f^{-1}V$
est schématiquement dense dans $X$.
\end{lemma}

\begin{proof}
En utilisant la caractérisation des ouverts schématiquement denses
du lemme \ref{lemma-scheme-theoretically-dense}
et la caractérisation des morphismes plats au moyen de
recouvrements étales donnée dans le lemme \ref{lemma-flat-local},
on se ramène au cas des schémas, qui est
Morphismes, lemme
\ref{morphisms-lemma-flat-morphism-scheme-theoretically-dense-open}.
\end{proof}

\begin{lemma}
\label{lemma-flat-base-change-scheme-theoretic-image}
Soit $S$ un schéma. Soit $f : X \to Y$ un morphisme plat d'espaces algébriques
sur $S$. Soit $g : V \to Y$ un morphisme quasi-compact d'espaces algébriques.
Soit $Z \subset Y$ l'image schématique de $g$, et soit $Z' \subset X$
l'image schématique du changement de base $V \times_Y X \to X$.
Alors $Z' = f^{-1}Z$.
\end{lemma}

\begin{proof}
Soit $Y' \to Y$ un morphisme étale surjectif tel que $Y'$ soit une
somme disjointe de schémas affines (Propriétés des espaces,
lemme \ref{spaces-properties-lemma-cover-by-union-affines}).
Soit $X' \to X \times_Y Y'$ un morphisme étale surjectif tel
que $X'$ soit une somme disjointe de schémas affines.
D'après le lemme \ref{lemma-flat-local}, le morphisme $X' \to Y'$ est plat.
Posons $V' = V \times_Y Y'$.
D'après le lemme \ref{lemma-quasi-compact-scheme-theoretic-image},
l'image inverse de $Z$ dans $Y'$ est l'image schématique
de $V' \to Y'$, et l'image inverse de $Z'$ dans $X'$
est l'image schématique de $V' \times_{Y'} X' \to X'$.
Comme $X' \to X$ est étale surjectif, il suffit de démontrer
le résultat pour les morphismes $X' \to Y'$ et $V' \to Y'$.
On peut donc supposer que $X$ et $Y$ sont des schémas affines.
Dans ce cas, $V$ est un espace algébrique quasi-compact.
Choisissons un schéma affine $W$ et un morphisme étale surjectif
$W \to V$ (Propriétés des espaces, lemme
\ref{spaces-properties-lemma-quasi-compact-affine-cover}).
Il est clair que l'image schématique de $V \to Y$
coïncide avec l'image schématique de $W \to Y$, et
de même pour $V \times_Y X \to X$ et $W \times_Y X \to X$.
On se ramène ainsi au cas des schémas, qui est
Morphismes, lemme
\ref{morphisms-lemma-flat-base-change-scheme-theoretic-image}.
\end{proof}





\section{Modules plats}
\label{section-flat-modules}

\noindent
Dans cette section, nous définissons ce que signifie, pour un module, être plat
en un point. Nous utiliserons pour cela la notion de fibre d'un faisceau sur
le petit site étale $X_\etale$ d'un espace algébrique ; voir
Propriétés des espaces, définition \ref{spaces-properties-definition-stalk}.

\begin{lemma}
\label{lemma-flat-at-point}
Soit $S$ un schéma. Soit $f : X \to Y$ un morphisme d'espaces
algébriques sur $S$. Soit $\mathcal{F}$ un faisceau quasi-cohérent sur $X$.
Soit $x \in |X|$. Les conditions suivantes sont équivalentes :
\begin{enumerate}
\item il existe un diagramme commutatif
$$
\xymatrix{
U \ar[d]_a \ar[r]_h & V \ar[d]^b \\
X \ar[r]^f & Y
}
$$
où $U$ et $V$ sont des schémas, où $a, b$ sont étales et où
$u \in U$ s'envoie sur $x$, tel que le module $a^*\mathcal{F}$ soit plat en $u$ sur $V$ ;
\item la fibre $\mathcal{F}_{\overline{x}}$ est plate sur
l'anneau local étale $\mathcal{O}_{Y, \overline{y}}$,
où $\overline{x}$ est un point géométrique quelconque au-dessus de
$x$ et $\overline{y} = f \circ \overline{x}$.
\end{enumerate}
\end{lemma}

\begin{proof}
Au cours de cette démonstration, fixons un point géométrique
$\overline{x} : \Spec(k) \to X$ au-dessus de $x$ et
notons $\overline{y} = f \circ \overline{x}$ son image dans $Y$.
Étant donné un diagramme comme en (1), on peut trouver un point géométrique
$\overline{u} : \Spec(k) \to U$ au-dessus de $u$ tel que
$\overline{x} = a \circ \overline{u}$ ; voir
Propriétés des espaces, lemme
\ref{spaces-properties-lemma-geometric-lift-to-usual}.
Posons $\overline{v} = h \circ \overline{u}$, d'image $v \in V$.
On sait que
$$
\mathcal{O}_{X, \overline{x}} = \mathcal{O}_{U, u}^{sh}
\quad\text{et}\quad
\mathcal{O}_{Y, \overline{y}} = \mathcal{O}_{V, v}^{sh}
$$
voir
Propriétés des espaces, lemme
\ref{spaces-properties-lemma-describe-etale-local-ring}.
On obtient un diagramme commutatif
$$
\xymatrix{
\mathcal{O}_{U, u} \ar[r] &
\mathcal{O}_{X, \overline{x}} \\
\mathcal{O}_{V, v} \ar[u] \ar[r] &
\mathcal{O}_{Y, \overline{y}} \ar[u]
}
$$
d'anneaux locaux. Enfin, on a
$$
\mathcal{F}_{\overline{x}} =
(a^*\mathcal{F})_u \otimes_{\mathcal{O}_{U, u}}
\mathcal{O}_{X, \overline{x}}
$$
d'après
Propriétés des espaces, lemme \ref{spaces-properties-lemma-stalk-quasi-coherent}.
Ainsi,
Algèbre, lemme \ref{algebra-lemma-flatness-descends-more-general}
nous dit que $(a^*\mathcal{F})_u$ est plat sur $\mathcal{O}_{V, v}$
si et seulement si $\mathcal{F}_{\overline{x}}$ est plate sur $\mathcal{O}_{V, v}$.
Le résultat découle donc de
Compléments sur la platitude, lemme \ref{flat-lemma-flat-up-down-henselization}.
\end{proof}

\begin{definition}
\label{definition-flat-module}
Soit $S$ un schéma.
Soit $f : X \to Y$ un morphisme d'espaces algébriques sur $S$.
Soit $\mathcal{F}$ un faisceau quasi-cohérent sur $X$.
\begin{enumerate}
\item Soit $x \in |X|$. On dit que $\mathcal{F}$ est {\it plate en $x$ sur $Y$}
si les conditions équivalentes du
lemme \ref{lemma-flat-at-point}
sont satisfaites.
\item On dit que $\mathcal{F}$ est {\it plate sur $Y$} si $\mathcal{F}$ est
plate sur $Y$ en tout $x \in |X|$.
\end{enumerate}
\end{definition}

\noindent
Cette définition posée, on dispose de l'inévitable lemme de changement de base.
Ce lemme entraîne que la formation du lieu de platitude d'un faisceau
quasi-cohérent commute au changement de base plat.

\begin{lemma}
\label{lemma-base-change-module-flat}
Soit $S$ un schéma. Soit
$$
\xymatrix{
X' \ar[d]_{f'} \ar[r]_{g'} & X \ar[d]^f \\
Y' \ar[r]^g & Y
}
$$
un diagramme cartésien d'espaces algébriques sur $S$. Soit $x' \in |X'|$
d'image $x \in |X|$. Soit $\mathcal{F}$ un faisceau quasi-cohérent
sur $X$, et posons $\mathcal{F}' = (g')^*\mathcal{F}$.
\begin{enumerate}
\item Si $\mathcal{F}$ est plate en $x$ sur $Y$,
alors $\mathcal{F}'$ est plate en $x'$ sur $Y'$.
\item Si $g$ est plat en $f'(x')$ et si
$\mathcal{F}'$ est plate en $x'$ sur $Y'$, alors
$\mathcal{F}$ est plate en $x$ sur $Y$.
\end{enumerate}
En particulier, si $\mathcal{F}$ est plate sur $Y$, alors
$\mathcal{F}'$ est plate sur $Y'$.
\end{lemma}

\begin{proof}
Choisissons un schéma $V$ et un morphisme étale surjectif $V \to Y$.
Choisissons un schéma $U$ et un morphisme étale surjectif $U \to V \times_Y X$.
Choisissons un schéma $V'$ et un morphisme étale surjectif $V' \to V \times_Y Y'$.
Alors $U' = V' \times_V U$ est un schéma muni d'un morphisme étale
surjectif $U' = V' \times_V U \to Y' \times_Y X = X'$. Choisissons $u' \in U'$
s'envoyant sur $x' \in |X'|$. On peut alors vérifier la platitude de
$\mathcal{F}'$ en $x'$ sur $Y'$ au moyen de la platitude de
$\mathcal{F}'|_{U'}$ en $u'$ sur $V'$. Le lemme découle donc de
Compléments sur les morphismes, lemme \ref{more-morphisms-lemma-flat-locus-base-change}.
\end{proof}

\noindent
Le lemme suivant traite de la « composition » des morphismes plats en
termes de modules. Il montre aussi que la platitude satisfait une sorte de
descente du haut vers le bas.

\begin{lemma}
\label{lemma-composition-module-flat}
Soit $S$ un schéma. Soient $X \to Y \to Z$ des morphismes d'espaces algébriques
sur $S$. Soit $\mathcal{F}$ un faisceau quasi-cohérent sur $X$.
Soit $x \in |X|$, d'image $y \in |Y|$.
\begin{enumerate}
\item Si $\mathcal{F}$ est plate en $x$ sur $Y$ et si
$Y$ est plat en $y$ sur $Z$, alors $\mathcal{F}$ est plate en
$x$ sur $Z$.
\item Soit $x : \Spec(K) \to X$ un représentant de $x$. Si
\begin{enumerate}
\item $\mathcal{F}$ est plate en $x$ sur $Y$ ;
\item $x^*\mathcal{F} \not = 0$ ;
\item $\mathcal{F}$ est plate en $x$ sur $Z$,
\end{enumerate}
alors $Y$ est plat en $y$ sur $Z$.
\item Soit $\overline{x}$ un point géométrique de $X$ au-dessus de $x$,
d'image $\overline{y}$ dans $Y$. Si $\mathcal{F}_{\overline{x}}$ est un
$\mathcal{O}_{Y, \overline{y}}$-module fidèlement plat et si
$\mathcal{F}$ est plate en $x$ sur $Z$, alors
$Y$ est plat en $y$ sur $Z$.
\end{enumerate}
\end{lemma}

\begin{proof}
Choisissons $\overline{x}$ et $\overline{y}$ comme dans (3), et notons
$\overline{z}$ le point géométrique de $Z$ ainsi induit. Au moyen de la
caractérisation de la platitude donnée dans les
lemmes \ref{lemma-flat-at-point} et
\ref{lemma-flat-at-point-etale-local-rings}
le lemme se ramène à une question purement algébrique sur l'homomorphisme
d'anneaux locaux $\mathcal{O}_{Z, \overline{z}} \to \mathcal{O}_{Y, \overline{y}}$
et sur le module $\mathcal{F}_{\overline{x}}$.
La partie (1) découle de
Algèbre, lemme \ref{algebra-lemma-composition-flat}.
Remarquons que la condition (2)(b) garantit que
$\mathcal{F}_{\overline{x}}/
\mathfrak m_{\overline{y}} \mathcal{F}_{\overline{x}}$
est non nul. Ainsi, (2)(a) $+$ (2)(b) entraînent que $\mathcal{F}_{\overline{x}}$
est un $\mathcal{O}_{Y, \overline{y}}$-module fidèlement plat ; voir
Algèbre, lemme \ref{algebra-lemma-ff}.
Ainsi, (2) est un cas particulier de (3).
Enfin, (3) découle de
Algèbre, lemme \ref{algebra-lemma-flat-permanence}.
\end{proof}

\noindent
Il arrive que le changement de base se fasse « en haut ». Voici un énoncé précis.

\begin{lemma}
\label{lemma-flat-permanence}
Soit $S$ un schéma. Soient $f : X \to Y$, $g : Y \to Z$ des morphismes
d'espaces algébriques sur $S$. Soit $\mathcal{G}$ un faisceau quasi-cohérent sur $Y$.
Soit $x \in |X|$, d'image $y \in |Y|$.
Si $f$ est plat en $x$, alors
$$
\mathcal{G}\text{ est plate sur }Z\text{ en }y
\Leftrightarrow
f^*\mathcal{G}\text{ est plate sur }Z\text{ en }x.
$$
En particulier, si $f$ est surjectif et plat, alors
$\mathcal{G}$ est plate sur $Z$ si et seulement si
$f^*\mathcal{G}$ est plate sur $Z$.
\end{lemma}

\begin{proof}
Choisissons un point géométrique $\overline{x}$ de $X$, et notons
$\overline{y}$ son image dans $Y$ et $\overline{z}$ son image dans $Z$.
Au moyen de la caractérisation de la platitude donnée dans les
lemmes \ref{lemma-flat-at-point} et
\ref{lemma-flat-at-point-etale-local-rings}
et de la description de la fibre de $f^*\mathcal{G}$ en $\overline{x}$ dans
Propriétés des espaces,
lemme \ref{spaces-properties-lemma-stalk-pullback-quasi-coherent}
le lemme se ramène à une question purement algébrique sur les homomorphismes
d'anneaux locaux
$\mathcal{O}_{Z, \overline{z}} \to \mathcal{O}_{Y, \overline{y}}
\to \mathcal{O}_{X, \overline{x}}$
et sur le module $\mathcal{G}_{\overline{y}}$.
Cet énoncé algébrique est
Algèbre, lemme \ref{algebra-lemma-flatness-descends-more-general}.
\end{proof}

\begin{lemma}
\label{lemma-pf-flat-module-open}
Soit $S$ un schéma.
Soit $f : X \to Y$ un morphisme d'espaces algébriques sur $S$.
Soit $\mathcal{F}$ un $\mathcal{O}_X$-module quasi-cohérent.
Supposons que $f$ soit localement de présentation finie, que $\mathcal{F}$ soit de
type fini, que $X = \text{Supp}(\mathcal{F})$ et que
$\mathcal{F}$ soit plate sur $Y$. Alors $f$ est universellement ouvert.
\end{lemma}

\begin{proof}
Choisissons un morphisme étale surjectif $\varphi : V \to Y$, où $V$ est un schéma.
Choisissons un morphisme étale surjectif $U \to V \times_Y X$, où $U$ est un schéma.
Il suffit alors de démontrer le lemme pour $U \to V$ et le
$\mathcal{O}_U$-module quasi-cohérent $\mathcal{F}|_U$.
Le lemme découle donc du cas des schémas ; voir
Morphismes, lemme \ref{morphisms-lemma-pf-flat-module-open}.
\end{proof}






\section{Platitude générique}
\label{section-generic-flatness}

\noindent
Cette section est l'analogue de
Morphismes, section \ref{morphisms-section-generic-flatness}.

\begin{proposition}
\label{proposition-generic-flatness-reduced}
Soit $S$ un schéma.
Soit $f : X \to Y$ un morphisme d'espaces algébriques sur $S$.
Soit $\mathcal{F}$ un faisceau quasi-cohérent de $\mathcal{O}_X$-modules.
Supposons que
\begin{enumerate}
\item $Y$ soit réduit ;
\item $f$ soit de type fini ;
\item $\mathcal{F}$ soit un $\mathcal{O}_X$-module de type fini.
\end{enumerate}
Alors il existe un sous-espace ouvert dense $W \subset Y$ tel que
le changement de base $X_W \to W$ de $f$ soit plat, localement de présentation finie et
quasi-compact, et tel que $\mathcal{F}|_{X_W}$ soit plate sur $W$ et de
présentation finie sur $\mathcal{O}_{X_W}$.
\end{proposition}

\begin{proof}
Soit $V$ un schéma et soit $V \to Y$ un morphisme étale surjectif.
Posons $X_V = V \times_Y X$ et notons $\mathcal{F}_V$ la restriction de
$\mathcal{F}$ à $X_V$. Supposons le résultat vrai pour le morphisme
$X_V \to V$ et le faisceau $\mathcal{F}_V$. Il existe alors un sous-schéma ouvert
$V' \subset V$ tel que $X_{V'} \to V'$ soit plat et de présentation finie,
et que $\mathcal{F}_{V'}$ soit un $\mathcal{O}_{X_{V'}}$-module de
présentation finie plat sur $V'$. Soit $W \subset Y$ l'image du
morphisme étale $V' \to Y$ ; voir
Propriétés des espaces, lemme \ref{spaces-properties-lemma-etale-image-open}.
Alors $V' \to W$ est un morphisme étale surjectif ; on voit donc que
$X_W \to W$ est plat, localement de présentation finie et quasi-compact d'après les
lemmes \ref{lemma-finite-presentation-local},
\ref{lemma-flat-local} et
\ref{lemma-quasi-compact-local}.
D'après la discussion de
Propriétés des espaces, section
\ref{spaces-properties-section-properties-modules}
on voit que $\mathcal{F}_W$ est de présentation finie comme
$\mathcal{O}_{X_W}$-module et, d'après le
lemme \ref{lemma-base-change-module-flat},
que $\mathcal{F}_W$ est plate sur $W$. Cet argument
ramène la proposition au cas où $Y$ est un schéma.

\medskip\noindent
Supposons que l'on sache démontrer la proposition lorsque $Y$ est un schéma affine.
Soit $f : X \to Y$ un morphisme de type fini d'espaces algébriques
sur $S$, où $Y$ est un schéma, et soit $\mathcal{F}$ un $\mathcal{O}_X$-module
quasi-cohérent de type fini. Choisissons un recouvrement ouvert affine
$Y = \bigcup V_j$. Par hypothèse, on peut trouver des ouverts denses $W_j \subset V_j$
tels que $X_{W_j} \to W_j$ soit plat, localement de présentation finie et
quasi-compact, et tels que $\mathcal{F}|_{X_{W_j}}$ soit plate sur $W_j$
et de présentation finie comme $\mathcal{O}_{X_{W_j}}$-module. Dans cette
situation, il suffit de prendre $W = \bigcup W_j$. On ramène ainsi
la proposition au cas où $Y$ est un schéma affine.

\medskip\noindent
Soit $Y$ un schéma affine sur $S$, soit
$f : X \to Y$ un morphisme de type fini d'espaces algébriques
sur $S$, et soit $\mathcal{F}$ un $\mathcal{O}_X$-module quasi-cohérent
de type fini. Comme $f$ est de type fini,
il est quasi-compact ; par suite, $X$ est quasi-compact. On peut donc trouver
un schéma affine $U$ et un morphisme étale surjectif $U \to X$ ; voir
Propriétés des espaces, lemme
\ref{spaces-properties-lemma-quasi-compact-affine-cover}.
Notons que $U \to Y$ est de type fini (c'est précisément ce que signifie ici
que $f$ est de type fini). On peut donc appliquer
Morphismes, Proposition \ref{morphisms-proposition-generic-flatness-reduced}
pour voir qu'il existe un ouvert dense $W \subset Y$ tel que
$U_W \to W$ soit plat et de présentation finie, et tel que
$\mathcal{F}|_{U_W}$ soit plate sur $W$ et de présentation finie
comme $\mathcal{O}_{U_W}$-module. D'après nos définitions, cela signifie que
le changement de base $X_W \to W$ de $f$ est plat, localement de présentation finie
et quasi-compact, et que $\mathcal{F}|_{X_W}$ est plate sur $W$ et de
présentation finie sur $\mathcal{O}_{X_W}$.
\end{proof}

\noindent
On ne peut renforcer le résultat de la proposition précédente en exigeant que
$X_W \to W$ soit de présentation finie, car
$\mathbf{A}^1_{\mathbf{Q}}/\mathbf{Z} \to \Spec(\mathbf{Q})$
fournit un contre-exemple. Le problème est que le morphisme diagonal
$\Delta_{X/Y}$ peut ne pas être quasi-compact, autrement dit que $f$ peut ne pas être
quasi-séparé. C'est manifestement le seul problème.

\begin{proposition}
\label{proposition-generic-flatness-reduced-quasi-separated}
Soit $S$ un schéma.
Soit $f : X \to Y$ un morphisme d'espaces algébriques sur $S$.
Soit $\mathcal{F}$ un faisceau quasi-cohérent de $\mathcal{O}_X$-modules.
Supposons que
\begin{enumerate}
\item $Y$ soit réduit ;
\item $f$ soit quasi-séparé ;
\item $f$ soit de type fini ;
\item $\mathcal{F}$ soit un $\mathcal{O}_X$-module de type fini.
\end{enumerate}
Alors il existe un sous-espace ouvert dense $W \subset Y$ tel que
le changement de base $X_W \to W$ de $f$ soit plat et de présentation finie,
et tel que $\mathcal{F}|_{X_W}$ soit plate sur $W$ et de
présentation finie sur $\mathcal{O}_{X_W}$.
\end{proposition}

\begin{proof}
Cela résulte immédiatement de la
Proposition \ref{proposition-generic-flatness-reduced}
et du fait que « de présentation finie » $=$
« localement de présentation finie » $+$ « quasi-compact » $+$
« quasi-séparé ».
\end{proof}















\section{Dimension relative}
\label{section-relative-dimension}

\noindent
Dans cette section, nous définissons la dimension relative
d'un morphisme d'espaces algébriques en un point, ainsi que quelques
propriétés étroitement liées.

\begin{definition}
\label{definition-dimension-fibre}
Soit $S$ un schéma.
Soit $f : X \to Y$ un morphisme d'espaces algébriques sur $S$.
Soit $x \in |X|$.
Soient $d, r \in \{0, 1, 2, \ldots, \infty\}$.
\begin{enumerate}
\item On dit que la
{\it dimension de l'anneau local de la fibre de $f$ en $x$} vaut $d$
si les conditions équivalentes du
lemme \ref{lemma-local-source-target-at-point}
sont satisfaites pour la propriété
$\mathcal{P}_d$ décrite dans
Descente, lemme \ref{descent-lemma-dimension-local-ring-fibre}.
\item On dit que le
{\it degré de transcendance de $x/f(x)$} vaut $r$
si les conditions équivalentes du
lemme \ref{lemma-local-source-target-at-point}
sont satisfaites pour la propriété
$\mathcal{P}_r$ décrite dans
Descente, lemme \ref{descent-lemma-transcendence-degree-at-point}.
\item On dit que
{\it $f$ est de dimension relative $d$ en $x$}
si les conditions équivalentes du
lemme \ref{lemma-local-source-target-at-point}
sont satisfaites pour la propriété
$\mathcal{P}_d$ décrite dans
Descente, lemme \ref{descent-lemma-dimension-at-point}.
\end{enumerate}
\end{definition}

\noindent
Précisons ce que cela signifie. Choisissons des
diagrammes
$$
\xymatrix{
U \ar[d]_a \ar[r]_h & V \ar[d]^b \\
X \ar[r]^f & Y
}
\quad\quad
\xymatrix{
u \ar[d] \ar[r] & v \ar[d] \\
x \ar[r] & y
}
$$
comme dans le
lemme \ref{lemma-local-source-target-at-point}.
Alors
$$
\begin{matrix}
\text{dimension relative de }f\text{ en }x & = &
\dim_u (U_v) \\
\text{dimension de l'anneau local de la fibre de }f\text{ en }x & = &
\dim(\mathcal{O}_{U_v, u})\\
\text{degré de transcendance de }x/f(x) & = &
\text{trdeg}_{\kappa(v)}(\kappa(u))
\end{matrix}
$$
Notons que si $Y = \Spec(k)$ est le spectre d'un corps, alors
la dimension relative de $X/Y$ en $x$ est égale à $\dim_x(X)$,
le degré de transcendance de $x/f(x)$ est le degré de transcendance
sur $k$, et la dimension de l'anneau local de la fibre de $f$
en $x$ est simplement la dimension de l'anneau local en $x$ ; autrement dit, les
notions relatives deviennent alors des notions absolues.

\begin{definition}
\label{definition-relative-dimension}
Soit $S$ un schéma.
Soit $f : X \to Y$ un morphisme d'espaces algébriques sur $S$.
Soit $d \in \{0, 1, 2, \ldots\}$.
\begin{enumerate}
\item On dit que $f$ est de {\it dimension relative $\leq d$} si
$f$ est de dimension relative $\leq d$ en tout $x \in |X|$.
\item On dit que $f$ est de {\it dimension relative $d$} si
$f$ est de dimension relative $d$ en tout $x \in |X|$.
\end{enumerate}
\end{definition}

\noindent
Être de dimension relative {\it égale} à $d$ signifie, en gros, que toutes les
fibres non vides sont équidimensionnelles de dimension $d$.

\begin{lemma}
\label{lemma-compare-tr-deg}
Soit $S$ un schéma. Soient $X \to Y \to Z$ des morphismes d'espaces
algébriques sur $S$. Soit $x \in |X|$, et soient $y \in |Y|$, $z \in |Z|$
ses images. Supposons que $X \to Y$ soit localement quasi-fini et que $Y \to Z$ soit
localement de type fini. Alors le degré de transcendance de $x/z$
est égal au degré de transcendance de $y/z$.
\end{lemma}

\begin{proof}
On peut choisir des diagrammes commutatifs
$$
\xymatrix{
U \ar[d] \ar[r] & V \ar[d] \ar[r] & W \ar[d] \\
X \ar[r] & Y \ar[r] & Z
}
\quad\quad
\xymatrix{
u \ar[d] \ar[r] & v \ar[d] \ar[r] & w \ar[d] \\
x \ar[r] & y \ar[r] & z
}
$$
où $U, V, W$ sont des schémas et où les flèches verticales sont étales.
Par définition, le morphisme $U \to V$ est localement quasi-fini,
ce qui entraîne que $\kappa(v) \subset \kappa(u)$ est finie ; voir
Morphismes, lemme \ref{morphisms-lemma-residue-field-quasi-finite}.
Le résultat est donc clair.
\end{proof}

\begin{lemma}
\label{lemma-jacobson-finite-type-points}
Soit $S$ un schéma. Soit $f : X \to Y$ un morphisme d'espaces algébriques
sur $S$. Si $f$ est localement de type fini, si $Y$ est de Jacobson
(Propriétés des espaces, remarque
\ref{spaces-properties-remark-list-properties-local-etale-topology}),
et si $x \in |X|$ est un point de type fini de $X$,
alors le degré de transcendance de $x/f(x)$ vaut $0$.
\end{lemma}

\begin{proof}
Choisissons un schéma $V$ et un morphisme étale surjectif $V \to Y$.
Choisissons un schéma $U$ et un morphisme étale surjectif $U \to X \times_Y V$.
D'après le lemme \ref{lemma-finite-type-points-surjective-morphism},
on peut trouver un point de type fini $u \in U$ s'envoyant sur $x$.
Après avoir rétréci $U$, on peut supposer que $u \in U$ est fermé
(Morphismes, lemme \ref{morphisms-lemma-identify-finite-type-points}).
Soit $v \in V$ l'image de $u$. D'après
Morphismes, lemme \ref{morphisms-lemma-jacobson-finite-type-points}
l'extension $\kappa(u)/\kappa(v)$ est finie.
Cela achève la démonstration.
\end{proof}

\begin{lemma}
\label{lemma-rel-dimension-dimension}
Soit $S$ un schéma. Soit $f : X \to Y$ un morphisme d'espaces algébriques
localement noethériens sur $S$, plat, localement de type fini et de
dimension relative $d$. Pour tout point $x$ de $|X|$, d'image
$y$ dans $|Y|$, on a $\dim_x(X) = \dim_y(Y) + d$.
\end{lemma}

\begin{proof}
Par définition de la dimension d'un espace algébrique
en un point (Propriétés des espaces, définition
\ref{spaces-properties-definition-dimension-at-point})
et par définition de la dimension relative $d$,
on se ramène à l'énoncé correspondant pour les schémas
(Morphismes, lemme \ref{morphisms-lemma-rel-dimension-dimension}).
\end{proof}







\section{Morphismes et dimensions des fibres}
\label{section-dimension-fibres}

\noindent
Cette section est l'analogue de
Morphismes, section \ref{morphisms-section-dimension-fibres}.
Les formulations de cette section sont un peu malcommodes, puisque
nous ne disposons pas d'anneaux locaux des espaces algébriques en leurs points.

\begin{lemma}
\label{lemma-dimension-fibre-at-a-point}
Soit $S$ un schéma.
Soit $f : X \to Y$ un morphisme d'espaces algébriques sur $S$.
Soit $x \in |X|$.
Supposons que $f$ soit localement de type fini.
Alors
$$
\begin{matrix}
\text{dimension relative de }f\text{ en }x \\
= \\
\text{dimension de l'anneau local de la fibre de }f\text{ en }x \\
+ \\
\text{degré de transcendance de }x/f(x)
\end{matrix}
$$
où les notations sont celles de la
définition \ref{definition-dimension-fibre}.
\end{lemma}

\begin{proof}
Cela résulte immédiatement de
Morphismes, lemme \ref{morphisms-lemma-dimension-fibre-at-a-point}
appliqué à $h : U \to V$ et $u \in U$
comme dans le
lemme \ref{lemma-local-source-target-at-point}.
\end{proof}

\begin{lemma}
\label{lemma-dimension-fibre-at-a-point-additive}
Soit $S$ un schéma.
Soient $f : X \to Y$ et $g : Y \to Z$ des morphismes d'espaces algébriques sur $S$.
Soit $x \in |X|$ et posons $y = f(x)$.
Supposons que $f$ et $g$ soient localement de type fini.
Alors :
\begin{enumerate}
\item
$$
\begin{matrix}
\text{dimension relative de }g \circ f\text{ en }x \\
\leq \\
\text{dimension relative de }f\text{ en }x \\
+ \\
\text{dimension relative de }g\text{ en }y
\end{matrix}
$$
\item il y a égalité dans (1) si, pour un certain morphisme $\Spec(k) \to Z$
depuis le spectre d'un corps représentant la classe de $g(f(x)) = g(y)$,
le morphisme $X_k \to Y_k$ est plat en $x$, par exemple si $f$ est plat en $x$ ;
\item
$$
\begin{matrix}
\text{degré de transcendance de }x/g(f(x)) \\
= \\
\text{degré de transcendance de }x/f(x) \\
+ \\
\text{degré de transcendance de }f(x)/g(f(x))
\end{matrix}
$$
\end{enumerate}
\end{lemma}

\begin{proof}
Choisissons un diagramme
$$
\xymatrix{
U \ar[d] \ar[r] & V \ar[d] \ar[r] & W \ar[d] \\
X \ar[r] & Y \ar[r] & Z
}
$$
où $U, V, W$ sont des schémas et où les flèches verticales sont étales et surjectives. (Voir
Espaces, lemme \ref{spaces-lemma-lift-morphism-presentations}.)
Choisissons $u \in U$ s'envoyant sur $x$. Notons $v, w$ les images
de $u$ dans $V, W$.
Appliquons
Morphismes, lemme \ref{morphisms-lemma-dimension-fibre-at-a-point-additive}
à la ligne supérieure et aux points $u, v, w$. Nous omettons les détails.
\end{proof}

\begin{lemma}
\label{lemma-dimension-fibre-after-base-change}
Soit $S$ un schéma. Soit
$$
\xymatrix{
X' \ar[r]_{g'} \ar[d]_{f'} & X \ar[d]^f \\
Y' \ar[r]^g & Y
}
$$
un diagramme de produit fibré d'espaces algébriques sur $S$.
Soit $x' \in |X'|$. Posons $x = g'(x')$. Supposons que $f$ soit localement de type fini.
Alors :
\begin{enumerate}
\item
$$
\begin{matrix}
\text{dimension relative de }f\text{ en }x \\
= \\
\text{dimension relative de }f'\text{ en }x'
\end{matrix}
$$
\item on a
$$
\begin{matrix}
\text{dimension de l'anneau local de la fibre de }f'\text{ en }x' \\
- \\
\text{dimension de l'anneau local de la fibre de }f\text{ en }x \\
= \\
\text{degré de transcendance de }x/f(x) \\
- \\
\text{degré de transcendance de }x'/f'(x')
\end{matrix}
$$
et cette valeur commune est $\geq 0$ ;
\item étant donnés $x$ et $y' \in |Y'|$ s'envoyant sur un même $y \in |Y|$,
on peut choisir $x'$ de sorte que l'entier de (2) soit $0$.
\end{enumerate}
\end{lemma}

\begin{proof}
Choisissons un morphisme étale surjectif $V \to Y$, où $V$ est un schéma.
Choisissons un morphisme étale surjectif $U \to V \times_Y X$, où $U$ est un schéma.
Choisissons un morphisme étale surjectif $V' \to V \times_Y Y'$, où $V'$ est un schéma.
Posons $U' = V' \times_V U$.
Alors le morphisme induit $U' \to X'$ est lui aussi étale et
surjectif (nous omettons l'argument). Choisissons $u' \in U'$
s'envoyant sur $x'$. Les parties (1) et (2) résultent alors de l'application de
Morphismes, lemme \ref{morphisms-lemma-dimension-fibre-after-base-change}
au diagramme de schémas faisant intervenir $U', U, V', V$ et le point $u'$.
Pour démontrer (3), choisissons d'abord $v \in V$ s'envoyant sur $y$.
Puis, à l'aide de Propriétés des espaces, lemme
\ref{spaces-properties-lemma-points-cartesian}
nous pouvons choisir $v' \in V'$ s'envoyant sur $y'$ et $v$, ainsi que
$u \in U$ s'envoyant sur $x$ et $v$. Enfin, d'après
Morphismes, lemme \ref{morphisms-lemma-dimension-fibre-after-base-change}
nous pouvons choisir $u' \in U'$ s'envoyant sur $v'$ et $u$ de sorte que
l'entier soit nul. Il suffit alors de prendre pour $x' \in |X'|$ l'image de $u'$.
\end{proof}

\begin{lemma}
\label{lemma-openness-bounded-dimension-fibres}
Soit $S$ un schéma.
Soit $f : X \to Y$ un morphisme d'espaces algébriques sur $S$.
Soit $n \geq 0$. Supposons que $f$ soit localement de type fini.
La partie
$$
W_n = \{x \in |X|
\text{ tels que la dimension relative de }f\text{ en } x \leq n\}
$$
est ouverte dans $|X|$.
\end{lemma}

\begin{proof}
Choisissons un diagramme
$$
\xymatrix{
U \ar[r]_h \ar[d]_a & V \ar[d] \\
X \ar[r] & Y
}
$$
où $U$ et $V$ sont des schémas et où les flèches verticales sont surjectives et
étales ; voir
Espaces, lemme \ref{spaces-lemma-lift-morphism-presentations}.
D'après
Morphismes, lemme \ref{morphisms-lemma-openness-bounded-dimension-fibres}
la partie $U_n$ des points où $h$ est de dimension relative
$\leq n$ est ouverte dans $U$. D'après notre définition de la dimension relative
des morphismes d'espaces algébriques en leurs points, on voit que
$U_n = a^{-1}(W_n)$.
Le lemme résulte de la définition de la topologie sur $|X|$.
\end{proof}

\begin{lemma}
\label{lemma-openness-bounded-dimension-fibres-finite-presentation}
Soit $S$ un schéma.
Soit $f : X \to Y$ un morphisme d'espaces algébriques sur $S$.
Soit $n \geq 0$. Supposons que $f$ soit localement de présentation finie.
L'ouvert
$$
W_n = \{x \in |X|
\text{ tels que la dimension relative de }f\text{ en } x \leq n\}
$$
du lemme \ref{lemma-openness-bounded-dimension-fibres}
est rétrocompact dans $|X|$. (Voir
Topologie, définition \ref{topology-definition-quasi-compact}.)
\end{lemma}

\begin{proof}
Choisissons un diagramme
$$
\xymatrix{
U \ar[r]_h \ar[d]_a & V \ar[d] \\
X \ar[r] & Y
}
$$
où $U$ et $V$ sont des schémas et où les flèches verticales sont surjectives et
étales ; voir
Espaces, lemme \ref{spaces-lemma-lift-morphism-presentations}.
Dans la démonstration du
lemme \ref{lemma-openness-bounded-dimension-fibres},
nous avons vu que $a^{-1}(W_n) = U_n$ est la partie correspondante
pour le morphisme $h$. D'après
Morphismes, lemme
\ref{morphisms-lemma-openness-bounded-dimension-fibres-finite-presentation}
on voit que $U_n$ est rétrocompact dans $U$.
Le lemme résulte de la définition de la topologie sur $|X|$ ; comparer à
Propriétés des espaces,
lemme \ref{spaces-properties-lemma-space-locally-quasi-compact}
et à sa démonstration.
\end{proof}

\begin{lemma}
\label{lemma-locally-quasi-finite-rel-dimension-0}
Soit $S$ un schéma.
Soit $f : X \to Y$ un morphisme d'espaces algébriques sur $S$.
Supposons que $f$ soit localement de type fini.
Alors $f$ est localement quasi-fini si et seulement si $f$ est de
dimension relative $0$ en tout $x \in |X|$.
\end{lemma}

\begin{proof}
Choisissons un diagramme
$$
\xymatrix{
U \ar[r]_h \ar[d]_a & V \ar[d] \\
X \ar[r] & Y
}
$$
où $U$ et $V$ sont des schémas et où les flèches verticales sont surjectives et
étales ; voir
Espaces, lemme \ref{spaces-lemma-lift-morphism-presentations}.
Les définitions entraînent que
$h$ est localement quasi-fini si et seulement si $f$ l'est,
et que $f$ est de dimension relative $0$ en tout $x \in |X|$ si et
seulement si $h$ est de dimension relative $0$ en tout $u \in U$.
Le résultat découle donc du résultat pour $h$, à savoir
Morphismes, lemme \ref{morphisms-lemma-locally-quasi-finite-rel-dimension-0}.
\end{proof}

\begin{lemma}
\label{lemma-locally-finite-type-quasi-finite-part}
Soit $S$ un schéma.
Soit $f : X \to Y$ un morphisme d'espaces algébriques sur $S$.
Supposons que $f$ soit localement de type fini.
Alors il existe un sous-espace ouvert canonique $X' \subset X$ tel que
$f|_{X'} : X' \to Y$ soit localement quasi-fini, et tel que la
dimension relative de $f$ en tout $x \in |X|$, $x \not \in |X'|$, soit
$\geq 1$. La formation de $X'$ commute à tout changement de base.
\end{lemma}

\begin{proof}
Combiner les
lemmes \ref{lemma-openness-bounded-dimension-fibres},
\ref{lemma-locally-quasi-finite-rel-dimension-0} et
\ref{lemma-dimension-fibre-after-base-change}.
\end{proof}

\begin{lemma}
\label{lemma-quasi-finite-at-point}
Soit $S$ un schéma. Considérons un diagramme cartésien
$$
\xymatrix{
X \ar[d] & F \ar[l]^p \ar[d] \\
Y & \Spec(k) \ar[l]
}
$$
où $X \to Y$ est un morphisme d'espaces algébriques sur $S$
localement de type fini et où $k$ est un corps sur $S$.
Soit $z \in |F|$ tel que $\dim_z(F) = 0$. Alors, après avoir remplacé $X$
par un sous-espace ouvert contenant $p(z)$, le morphisme
$$
X \longrightarrow Y
$$
est localement quasi-fini.
\end{lemma}

\begin{proof}
Soit $X' \subset X$ le sous-espace ouvert sur lequel $f$ est localement quasi-fini,
construit dans le
lemme \ref{lemma-locally-finite-type-quasi-finite-part}.
Comme la formation de $X'$ commute à tout changement de base, on voit
que $z \in X' \times_Y \Spec(k)$. Le lemme est donc clair.
\end{proof}





\section{La formule de dimension}
\label{section-dimension-formula}

\noindent
L'analogue de la formule de dimension
(Morphismes, lemme \ref{morphisms-lemma-dimension-formula})
est un peu délicat à formuler, car il faudrait définir
les espaces algébriques intègres (ce que nous ferons plus loin), ainsi que
les espaces algébriques universellement caténaires. La
version suivante est toutefois immédiate.

\begin{lemma}
\label{lemma-dimension-formula-general}
Soit $S$ un schéma. Soit $f : X \to Y$ un morphisme d'espaces
algébriques sur $S$. Supposons que $Y$ soit localement noethérien et que $f$ soit localement
de type fini. Soit $x \in |X|$, d'image $y \in |Y|$.
Alors
\begin{align*}
& \text{la dimension de l'anneau local de }X\text{ en }x \leq \\
& \text{la dimension de l'anneau local de }Y\text{ en }y + E - \\
& \text{ le degré de transcendance de }x/y
\end{align*}
où $E$ est le maximum des degrés de transcendance de $\xi/f(\xi)$,
lorsque $\xi \in |X|$ parcourt les points se spécialisant en $x$ en lesquels
l'anneau local de $X$ est de dimension $0$.
\end{lemma}

\begin{proof}
Choisissons un schéma affine $V$, un morphisme étale $V \to Y$ et un point
$v \in V$ s'envoyant sur $y$. Choisissons un schéma affine $U$, un morphisme étale
$U \to X \times_Y V$ et un point $u \in U$ s'envoyant sur $v$ dans $V$ et sur $x$
dans $X$. En explicitant la définition \ref{definition-dimension-fibre} et
Propriétés des espaces, définition
\ref{spaces-properties-definition-dimension-local-ring}
il faut montrer que
$$
\dim(\mathcal{O}_{U, u}) \leq
\dim(\mathcal{O}_{V, v}) + E - \text{trdeg}_{\kappa(v)}(\kappa(u))
$$
Soit $\xi_U \in U$ un point générique d'une composante irréductible de
$U$ contenant $u$. Alors $\xi_U$ s'envoie sur un point $\xi \in |X|$
figurant dans la liste utilisée pour définir la quantité $E$ et, en fait,
chaque $\xi$ intervenant dans la définition de $E$ s'obtient ainsi
(nous omettons ce petit détail). En particulier, il n'y a qu'un nombre fini
de tels $\xi$, et l'on peut prendre le maximum (c'est donc bien
un maximum et non une borne supérieure).
Le degré de transcendance de $\xi$ sur $f(\xi)$ est
$\text{trdeg}_{\kappa(\xi_V)}(\kappa(\xi_U))$, où $\xi_V \in V$
est l'image de $\xi_U$. Le lemme découle donc de
Morphismes, lemme \ref{morphisms-lemma-dimension-formula-general}.
\end{proof}

\begin{lemma}
\label{lemma-alteration-dimension-general}
Soit $S$ un schéma.
Soit $f : X \to Y$ un morphisme d'espaces algébriques sur $S$. Supposons
que $Y$ soit localement noethérien et que $f$ soit localement de type fini.
Alors
$$
\dim(X) \leq \dim(Y) + E
$$
où $E$ est la borne supérieure des degrés de transcendance de
$\xi/f(\xi)$, lorsque $\xi$ parcourt les points en
lesquels l'anneau local de $X$ est de dimension $0$.
\end{lemma}

\begin{proof}
Conséquence immédiate du lemme \ref{lemma-dimension-formula-general}
et de Propriétés des espaces, lemme \ref{spaces-properties-lemma-dimension}.
\end{proof}







\section{Morphismes syntomiques}
\label{section-syntomic}

\noindent
La propriété « syntomique » des morphismes de schémas est
locale sur la source et le but pour la topologie étale ; voir
Descente, remarque \ref{descent-remark-list-local-source-target}.
Elle est aussi stable par changement de base et locale sur le but pour la topologie fpqc ; voir
Morphismes, lemme \ref{morphisms-lemma-base-change-syntomic} et
Descente, lemme \ref{descent-lemma-descending-property-syntomic}.
Ainsi, d'après le
lemme \ref{lemma-local-source-target}
ci-dessus, on peut définir comme suit la notion de morphisme syntomique d'espaces
algébriques ; cette définition coïncide avec la notion déjà définie dans la
section \ref{section-representable}
lorsque le morphisme est représentable.

\begin{definition}
\label{definition-syntomic}
Soit $S$ un schéma.
Soit $f : X \to Y$ un morphisme d'espaces algébriques sur $S$.
\begin{enumerate}
\item On dit que $f$ est {\it syntomique} si les conditions équivalentes du
lemme \ref{lemma-local-source-target}
sont satisfaites pour $\mathcal{P} =$ « syntomique ».
\item Soit $x \in |X|$. On dit que $f$ est {\it syntomique en $x$} s'il
existe un voisinage ouvert $X' \subset X$ de $x$ tel
que $f|_{X'} : X' \to Y$ soit syntomique.
\end{enumerate}
\end{definition}

\begin{lemma}
\label{lemma-composition-syntomic}
Le composé de morphismes syntomiques est syntomique.
\end{lemma}

\begin{proof}
Voir la remarque \ref{remark-composition-P} et
Morphismes, lemme \ref{morphisms-lemma-composition-syntomic}.
\end{proof}

\begin{lemma}
\label{lemma-base-change-syntomic}
Le changement de base d'un morphisme syntomique est syntomique.
\end{lemma}

\begin{proof}
Voir la remarque \ref{remark-base-change-P} et
Morphismes, lemme \ref{morphisms-lemma-base-change-syntomic}.
\end{proof}

\begin{lemma}
\label{lemma-syntomic-local}
Soit $S$ un schéma.
Soit $f : X \to Y$ un morphisme d'espaces algébriques sur $S$.
Les conditions suivantes sont équivalentes :
\begin{enumerate}
\item $f$ est syntomique ;
\item pour tout $x \in |X|$, le morphisme $f$ est syntomique en $x$ ;
\item pour tout schéma $Z$ et tout morphisme $Z \to Y$, le morphisme
$Z \times_Y X \to Z$ est syntomique ;
\item pour tout schéma affine $Z$ et tout morphisme
$Z \to Y$, le morphisme $Z \times_Y X \to Z$ est syntomique ;
\item il existe un schéma $V$ et un morphisme étale surjectif
$V \to Y$ tels que $V \times_Y X \to V$ soit syntomique ;
\item il existe un schéma $U$ et un morphisme étale surjectif
$\varphi : U \to X$ tels que le composé $f \circ \varphi$
soit syntomique ;
\item pour tout diagramme commutatif
$$
\xymatrix{
U \ar[d] \ar[r] & V \ar[d] \\
X \ar[r] & Y
}
$$
où $U$ et $V$ sont des schémas et où les flèches verticales sont étales,
la flèche horizontale supérieure est syntomique ;
\item il existe un diagramme commutatif
$$
\xymatrix{
U \ar[d] \ar[r] & V \ar[d] \\
X \ar[r] & Y
}
$$
où $U$ et $V$ sont des schémas, où les flèches verticales sont étales et où
$U \to X$ est surjectif, tel que la flèche horizontale supérieure soit syntomique ;
\item il existe des recouvrements de Zariski $Y = \bigcup_{i \in I} Y_i$
et $f^{-1}(Y_i) = \bigcup X_{ij}$ tels que
chaque morphisme $X_{ij} \to Y_i$ soit syntomique.
\end{enumerate}
\end{lemma}

\begin{proof}
La démonstration est omise.
\end{proof}

\begin{lemma}
\label{lemma-syntomic-locally-finite-presentation}
Un morphisme syntomique est localement de présentation finie.
\end{lemma}

\begin{proof}
Cela résulte immédiatement du cas des schémas
(Morphismes, lemme \ref{morphisms-lemma-syntomic-locally-finite-presentation}).
\end{proof}

\begin{lemma}
\label{lemma-syntomic-flat}
Un morphisme syntomique est plat.
\end{lemma}

\begin{proof}
Cela résulte immédiatement du cas des schémas
(Morphismes, lemme \ref{morphisms-lemma-syntomic-flat}).
\end{proof}

\begin{lemma}
\label{lemma-syntomic-open}
Un morphisme syntomique est universellement ouvert.
\end{lemma}

\begin{proof}
Combiner les
lemmes \ref{lemma-syntomic-locally-finite-presentation},
\ref{lemma-syntomic-flat} et
\ref{lemma-fppf-open}.
\end{proof}





\section{Morphismes lisses}
\label{section-smooth}

\noindent
La propriété « lisse » des morphismes de schémas est
locale sur la source et le but pour la topologie étale ; voir
Descente, remarque \ref{descent-remark-list-local-source-target}.
Elle est aussi stable par changement de base et locale sur le but pour la topologie fpqc ; voir
Morphismes, lemme \ref{morphisms-lemma-base-change-smooth}
et
Descente, lemme \ref{descent-lemma-descending-property-smooth}.
Ainsi, d'après le
lemme \ref{lemma-local-source-target}
ci-dessus, on peut définir comme suit la notion de morphisme lisse d'espaces
algébriques ; cette définition coïncide avec la notion déjà définie dans la
section \ref{section-representable}
lorsque le morphisme est représentable.

\begin{definition}
\label{definition-smooth}
Soit $S$ un schéma.
Soit $f : X \to Y$ un morphisme d'espaces algébriques sur $S$.
\begin{enumerate}
\item On dit que $f$ est {\it lisse} si les conditions équivalentes du
lemme \ref{lemma-local-source-target} sont satisfaites pour
$\mathcal{P} =$ « lisse ».
\item Soit $x \in |X|$. On dit que $f$ est {\it lisse en $x$} s'il existe
un voisinage ouvert $X' \subset X$ de $x$ tel que $f|_{X'} : X' \to Y$
soit lisse.
\end{enumerate}
\end{definition}

\begin{lemma}
\label{lemma-composition-smooth}
Le composé de morphismes lisses est lisse.
\end{lemma}

\begin{proof}
Voir la remarque \ref{remark-composition-P} et
Morphismes, lemme \ref{morphisms-lemma-composition-smooth}.
\end{proof}

\begin{lemma}
\label{lemma-base-change-smooth}
Le changement de base d'un morphisme lisse est lisse.
\end{lemma}

\begin{proof}
Voir la remarque \ref{remark-base-change-P} et
Morphismes, lemme \ref{morphisms-lemma-base-change-smooth}.
\end{proof}

\begin{lemma}
\label{lemma-smooth-local}
Soit $S$ un schéma.
Soit $f : X \to Y$ un morphisme d'espaces algébriques sur $S$.
Les conditions suivantes sont équivalentes :
\begin{enumerate}
\item $f$ est lisse ;
\item pour tout $x \in |X|$, le morphisme $f$ est lisse en $x$ ;
\item pour tout schéma $Z$ et tout morphisme $Z \to Y$, le morphisme
$Z \times_Y X \to Z$ est lisse ;
\item pour tout schéma affine $Z$ et tout morphisme
$Z \to Y$, le morphisme $Z \times_Y X \to Z$ est lisse ;
\item il existe un schéma $V$ et un morphisme étale surjectif
$V \to Y$ tels que $V \times_Y X \to V$ soit un morphisme lisse ;
\item il existe un schéma $U$ et un morphisme étale surjectif
$\varphi : U \to X$ tels que le composé $f \circ \varphi$
soit lisse ;
\item pour tout diagramme commutatif
$$
\xymatrix{
U \ar[d] \ar[r] & V \ar[d] \\
X \ar[r] & Y
}
$$
où $U$ et $V$ sont des schémas et où les flèches verticales sont étales,
la flèche horizontale supérieure est lisse ;
\item il existe un diagramme commutatif
$$
\xymatrix{
U \ar[d] \ar[r] & V \ar[d] \\
X \ar[r] & Y
}
$$
où $U$ et $V$ sont des schémas, où les flèches verticales sont étales et où
$U \to X$ est surjectif, tel que la flèche horizontale supérieure soit lisse ;
\item il existe des recouvrements de Zariski $Y = \bigcup_{i \in I} Y_i$
et $f^{-1}(Y_i) = \bigcup X_{ij}$ tels que
chaque morphisme $X_{ij} \to Y_i$ soit lisse.
\end{enumerate}
\end{lemma}

\begin{proof}
La démonstration est omise.
\end{proof}

\begin{lemma}
\label{lemma-smooth-locally-finite-presentation}
Un morphisme lisse d'espaces algébriques est localement de présentation finie.
\end{lemma}

\begin{proof}
Soit $X \to Y$ un morphisme lisse d'espaces algébriques. Par
définition, cela signifie qu'il existe un diagramme comme dans le
lemme \ref{lemma-local-source-target},
où $h$ est lisse et où $a$ est une flèche verticale surjective. D'après
Morphismes, lemme \ref{morphisms-lemma-smooth-locally-finite-presentation}
$h$ est localement de présentation finie. Ainsi, $X \to Y$ est
localement de présentation finie par définition.
\end{proof}

\begin{lemma}
\label{lemma-smooth-locally-finite-type}
Un morphisme lisse d'espaces algébriques est localement de type fini.
\end{lemma}

\begin{proof}
Combiner les
lemmes \ref{lemma-smooth-locally-finite-presentation} et
\ref{lemma-finite-presentation-finite-type}.
\end{proof}

\begin{lemma}
\label{lemma-smooth-flat}
Un morphisme lisse d'espaces algébriques est plat.
\end{lemma}

\begin{proof}
Soit $X \to Y$ un morphisme lisse d'espaces algébriques. Par
définition, cela signifie qu'il existe un diagramme comme dans le
lemme \ref{lemma-local-source-target},
où $h$ est lisse et où $a$ est une flèche verticale surjective. D'après
Morphismes, lemme \ref{morphisms-lemma-smooth-flat},
$h$ est plat. Ainsi, $X \to Y$ est plat par définition.
\end{proof}

\begin{lemma}
\label{lemma-smooth-syntomic}
Un morphisme lisse d'espaces algébriques est syntomique.
\end{lemma}

\begin{proof}
Soit $X \to Y$ un morphisme lisse d'espaces algébriques. Par
définition, cela signifie qu'il existe un diagramme comme dans le
lemme \ref{lemma-local-source-target},
où $h$ est lisse et où $a$ est une flèche verticale surjective. D'après
Morphismes, lemme \ref{morphisms-lemma-smooth-syntomic}
$h$ est syntomique. Ainsi, $X \to Y$ est syntomique par définition.
\end{proof}

\begin{lemma}
\label{lemma-where-smooth}
Soit $S$ un schéma. Soit $f : X \to Y$ un morphisme d'espaces
algébriques sur $S$. Il existe un plus grand sous-espace ouvert $U \subset X$
tel que $f|_U : U \to Y$ soit lisse. De plus, la formation de
cet ouvert commute au changement de base par
\begin{enumerate}
\item les morphismes plats et
localement de présentation finie ;
\item les morphismes plats, pourvu que $f$ soit localement de présentation finie.
\end{enumerate}
\end{lemma}

\begin{proof}
L'existence de $U$ résulte du fait que la propriété
d'être lisse est locale sur la source pour la topologie de Zariski (et même étale) ; voir le
lemme \ref{lemma-smooth-local}. De plus, ce lemme permet
de ramener les propriétés (1) et (2) au cas
des morphismes de schémas. Le cas des schémas est traité dans
Morphismes, lemme \ref{morphisms-lemma-set-points-where-fibres-smooth}.
Certains détails sont omis.
\end{proof}

\begin{lemma}
\label{lemma-smoothness-dimension-spaces}
Soient $X$ et $Y$ des espaces algébriques localement noethériens sur un schéma
$S$, et soit $f : X \to Y$ un morphisme lisse.
Pour tout point $x \in |X|$ d'image $y \in |Y|$, on a
$$
\dim_x(X) = \dim_y(Y) + \dim_x(X_y)
$$
où $\dim_x(X_y)$ est la dimension relative de $f$ en $x$ au sens de la
définition \ref{definition-dimension-fibre}.
\end{lemma}

\begin{proof}
Par définition de la dimension d'un espace algébrique
en un point (Propriétés des espaces, définition
\ref{spaces-properties-definition-dimension-at-point}),
on se ramène à l'énoncé correspondant pour les schémas
(Morphismes, lemme \ref{morphisms-lemma-smoothness-dimension}).
\end{proof}



\section{Morphismes non ramifiés}
\label{section-unramified}

\noindent
La propriété « non ramifié » (resp.\ « G-non ramifié »)
des morphismes de schémas est locale sur la source et le but pour la topologie étale ; voir
Descente, remarque \ref{descent-remark-list-local-source-target}.
Elle est aussi stable par changement de base et locale sur le but pour la topologie fpqc ; voir
Morphismes, lemme \ref{morphisms-lemma-base-change-unramified} et
Descente, lemme \ref{descent-lemma-descending-property-unramified}.
Ainsi, d'après le
lemme \ref{lemma-local-source-target}
ci-dessus, on peut définir comme suit la notion de morphisme non ramifié
(resp.\ de morphisme G-non ramifié) d'espaces algébriques ;
cette définition coïncide avec la notion déjà définie dans la
section \ref{section-representable}
lorsque le morphisme est représentable.

\begin{definition}
\label{definition-unramified}
Soit $S$ un schéma.
Soit $f : X \to Y$ un morphisme d'espaces algébriques sur $S$.
\begin{enumerate}
\item On dit que $f$ est {\it non ramifié} si les conditions équivalentes du
lemme \ref{lemma-local-source-target}
sont satisfaites pour $\mathcal{P} = \text{non ramifié}$.
\item Soit $x \in |X|$. On dit que $f$ est {\it non ramifié en $x$} s'il
existe un voisinage ouvert $X' \subset X$ de $x$ tel que
$f|_{X'} : X' \to Y$ soit non ramifié.
\item On dit que $f$ est {\it G-non ramifié} si les conditions équivalentes du
lemme \ref{lemma-local-source-target}
sont satisfaites pour $\mathcal{P} = \text{G-non ramifié}$.
\item Soit $x \in |X|$. On dit que $f$ est {\it G-non ramifié en $x$} s'il
existe un voisinage ouvert $X' \subset X$ de $x$ tel que
$f|_{X'} : X' \to Y$ soit G-non ramifié.
\end{enumerate}
\end{definition}

\noindent
En raison du lemme suivant, nous ne développerons désormais que la théorie
des morphismes non ramifiés et, chaque fois que nous voudrons utiliser un morphisme
G-non ramifié, nous dirons simplement « un morphisme non ramifié localement de
présentation finie ».

\begin{lemma}
\label{lemma-unramified-G-unramified}
Soit $S$ un schéma.
Soit $f : X \to Y$ un morphisme d'espaces algébriques sur $S$.
Alors $f$ est G-non ramifié si et seulement si $f$ est non ramifié et
localement de présentation finie.
\end{lemma}

\begin{proof}
Considérons un diagramme quelconque comme dans le
lemme \ref{lemma-local-source-target}.
L'assertion revient alors à dire que le morphisme $h$ est
G-non ramifié si et seulement s'il est non ramifié et localement de
présentation finie. Cela résulte immédiatement de
Morphismes, définition \ref{morphisms-definition-unramified}.
\end{proof}

\begin{lemma}
\label{lemma-composition-unramified}
Le composé de morphismes non ramifiés est non ramifié.
\end{lemma}

\begin{proof}
Voir la remarque \ref{remark-composition-P} et
Morphismes, lemme \ref{morphisms-lemma-composition-unramified}.
\end{proof}

\begin{lemma}
\label{lemma-base-change-unramified}
Le changement de base d'un morphisme non ramifié est non ramifié.
\end{lemma}

\begin{proof}
Voir la remarque \ref{remark-base-change-P} et
Morphismes, lemme \ref{morphisms-lemma-base-change-unramified}.
\end{proof}

\begin{lemma}
\label{lemma-unramified-local}
Soit $S$ un schéma.
Soit $f : X \to Y$ un morphisme d'espaces algébriques sur $S$.
Les conditions suivantes sont équivalentes :
\begin{enumerate}
\item $f$ est non ramifié ;
\item pour tout $x \in |X|$, le morphisme $f$ est non ramifié en $x$ ;
\item pour tout schéma $Z$ et tout morphisme $Z \to Y$, le morphisme
$Z \times_Y X \to Z$ est non ramifié ;
\item pour tout schéma affine $Z$ et tout morphisme
$Z \to Y$, le morphisme $Z \times_Y X \to Z$ est non ramifié ;
\item il existe un schéma $V$ et un morphisme étale surjectif
$V \to Y$ tels que $V \times_Y X \to V$ soit un morphisme non ramifié ;
\item il existe un schéma $U$ et un morphisme étale surjectif
$\varphi : U \to X$ tels que le composé $f \circ \varphi$
soit non ramifié ;
\item pour tout diagramme commutatif
$$
\xymatrix{
U \ar[d] \ar[r] & V \ar[d] \\
X \ar[r] & Y
}
$$
où $U$ et $V$ sont des schémas et où les flèches verticales sont étales,
la flèche horizontale supérieure est non ramifiée ;
\item il existe un diagramme commutatif
$$
\xymatrix{
U \ar[d] \ar[r] & V \ar[d] \\
X \ar[r] & Y
}
$$
où $U$ et $V$ sont des schémas, où les flèches verticales sont étales et où
$U \to X$ est surjectif, tel que la flèche horizontale supérieure soit non ramifiée ;
\item il existe des recouvrements de Zariski $Y = \bigcup_{i \in I} Y_i$
et $f^{-1}(Y_i) = \bigcup X_{ij}$ tels que
chaque morphisme $X_{ij} \to Y_i$ soit non ramifié.
\end{enumerate}
\end{lemma}

\begin{proof}
La démonstration est omise.
\end{proof}


\begin{lemma}
\label{lemma-unramified-locally-finite-type}
Un morphisme non ramifié d'espaces algébriques est localement de type fini.
\end{lemma}

\begin{proof}
Au moyen d'un diagramme comme dans le
lemme \ref{lemma-local-source-target},
l'assertion se ramène à
Morphismes, lemme \ref{morphisms-lemma-unramified-locally-finite-type}.
\end{proof}

\begin{lemma}
\label{lemma-unramified-quasi-finite}
Si $f$ est non ramifié en $x$, alors $f$ est quasi-fini en $x$.
En particulier, un morphisme non ramifié est localement quasi-fini.
\end{lemma}

\begin{proof}
Au moyen d'un diagramme comme dans le
lemme \ref{lemma-local-source-target},
l'assertion se ramène à
Morphismes, lemme \ref{morphisms-lemma-unramified-quasi-finite}.
\end{proof}

\begin{lemma}
\label{lemma-immersion-unramified}
Une immersion d'espaces algébriques est non ramifiée.
\end{lemma}

\begin{proof}
Soit $i : X \to Y$ une immersion d'espaces algébriques. Choisissons un schéma
$V$ et un morphisme étale surjectif $V \to Y$. Alors $V \times_Y X \to V$
est une immersion de schémas, donc est non ramifié (voir
Morphismes, lemmes \ref{morphisms-lemma-open-immersion-unramified} et
\ref{morphisms-lemma-closed-immersion-unramified}).
Ainsi, $i$ est non ramifié par définition.
\end{proof}

\begin{lemma}
\label{lemma-diagonal-unramified-morphism}
Soit $S$ un schéma.
Soit $f : X \to Y$ un morphisme d'espaces algébriques sur $S$.
\begin{enumerate}
\item Si $f$ est non ramifié, alors le morphisme diagonal
$\Delta_{X/Y} : X \to X \times_Y X$ est une immersion ouverte.
\item Si $f$ est localement de type fini
et si $\Delta_{X/Y}$ est une immersion ouverte, alors $f$ est non ramifié.
\end{enumerate}
\end{lemma}

\begin{proof}
On sait dans tous les cas que $\Delta_{X/Y}$ est un monomorphisme représentable ; voir le
lemme \ref{lemma-properties-diagonal}.
Choisissons un schéma $V$ et un morphisme étale surjectif $V \to Y$.
Choisissons un schéma $U$ et un morphisme étale surjectif $U \to X \times_Y V$.
Considérons le diagramme commutatif
$$
\xymatrix{
U \ar[d] \ar[rr]_-{\Delta_{U/V}} & &
U \times_V U \ar[d] \ar[r] &
V \ar[d]^{\Delta_{V/Y}} \\
X \ar[rr]^-{\Delta_{X/Y}} & &
X \times_Y X \ar[r] &
V \times_Y V
}
$$
dont le carré de droite est cartésien. La flèche verticale de gauche est étale surjective.
La flèche verticale de droite est étale en tant que morphisme entre des schémas
étales sur $Y$ ; voir
Propriétés des espaces,
lemme \ref{spaces-properties-lemma-etale-permanence}.
Par conséquent, la flèche verticale du milieu est également étale (mais elle n'est pas
nécessairement surjective).

\medskip\noindent
Supposons que $f$ soit non ramifié. Alors $U \to V$ est non ramifié, donc
$\Delta_{U/V}$ est une immersion ouverte d'après
Morphismes, lemme \ref{morphisms-lemma-diagonal-unramified-morphism}.
En considérant le carré de gauche du diagramme ci-dessus, on conclut que
$\Delta_{X/Y}$ est un morphisme étale ; voir
Propriétés des espaces,
lemme \ref{spaces-properties-lemma-etale-local}.
Ainsi, $\Delta_{X/Y}$ est un monomorphisme étale représentable, ce qui
implique que c'est une immersion ouverte d'après
\'Morphismes étales, théorème \ref{etale-theorem-etale-radicial-open}.
(Voir aussi
Espaces, lemme
\ref{spaces-lemma-representable-transformations-property-implication}
pour le passage du langage des schémas à celui des foncteurs.)

\medskip\noindent
Supposons que $f$ soit localement de type fini et que $\Delta_{X/Y}$
soit une immersion ouverte. Cela implique que $U \to V$ est aussi localement de type
fini (par définition d'un morphisme d'espaces algébriques qui est
localement de type fini). En considérant le diagramme ci-dessus,
on conclut que $\Delta_{U/V}$ est étale en tant que morphisme entre des
schémas étales sur $X \times_Y X$ ; voir
Propriétés des espaces,
lemme \ref{spaces-properties-lemma-etale-permanence}.
Mais, puisque $\Delta_{U/V}$ est la diagonale d'un morphisme entre schémas,
on voit que c'est en tout cas une immersion ; voir
Schémas, lemme \ref{schemes-lemma-diagonal-immersion}.
Il s'agit donc d'une immersion ouverte, et l'on conclut
que $U \to V$ est non ramifié d'après
Morphismes, lemme \ref{morphisms-lemma-diagonal-unramified-morphism}.
Cela signifie à son tour que $f$ est non ramifié par définition.
\end{proof}

\begin{lemma}
\label{lemma-where-unramified}
Soit $S$ un schéma. Considérons un diagramme commutatif
$$
\xymatrix{
X \ar[rr]_f \ar[rd]_p & & Y \ar[ld]^q \\
& Z
}
$$
d'espaces algébriques sur $S$. Supposons que $X \to Z$ soit localement de
type fini. Alors il existe un sous-espace ouvert $U(f) \subset X$
tel que $|U(f)| \subset |X|$ soit l'ensemble des points où $f$ est non ramifié.
De plus, pour tout morphisme d'espaces algébriques $Z' \to Z$, si
$f' : X' \to Y'$ est le changement de base de $f$ par $Z' \to Z$, alors
$U(f')$ est l'image inverse de $U(f)$ par la projection $X' \to X$.
\end{lemma}

\begin{proof}
Ce lemme est l'analogue de
Morphismes, lemme \ref{morphisms-lemma-set-points-where-fibres-unramified}
et nous allons en fait l'en déduire. D'après la
définition \ref{definition-unramified},
l'ensemble $\{x \in |X| : f \text{ est non ramifié en }x\}$ est
ouvert dans $X$. Il suffit donc de démontrer la dernière assertion. D'après le
lemme \ref{lemma-permanence-finite-type},
le morphisme $X \to Y$ est localement de type fini. D'après le
lemme \ref{lemma-base-change-finite-type},
le morphisme $X' \to Y'$ est localement de type fini.

\medskip\noindent
Choisissons un schéma $W$ et un morphisme étale surjectif $W \to Z$.
Choisissons un schéma $V$ et un morphisme étale surjectif $V \to W \times_Z Y$.
Choisissons un schéma $U$ et un morphisme étale surjectif $U \to V \times_Y X$.
Enfin, choisissons un schéma $W'$ et un morphisme étale surjectif
$W' \to W \times_Z Z'$.
Posons $V' = W' \times_W V$ et $U' = W' \times_W U$ ; on obtient ainsi des
morphismes étales surjectifs $V' \to Y'$ et $U' \to X'$.
Nous utiliserons sans autre commentaire le fait qu'un morphisme étale d'espaces algébriques
induit une application ouverte entre les espaces topologiques associés (voir
Propriétés des espaces, lemme
\ref{spaces-properties-lemma-etale-open}).
Combiné avec le
lemme \ref{lemma-unramified-local}, cela
implique que $U(f)$ est l'image dans $|X|$ de l'ensemble $T$ des points de $U$
où le morphisme $U \to V$ est non ramifié. De même, $U(f')$ est l'image
dans $|X'|$ de l'ensemble $T'$ des points de $U'$ où le morphisme $U' \to V'$
est non ramifié. Or, par construction, le diagramme
$$
\xymatrix{
U' \ar[r] \ar[d] & U \ar[d] \\
V' \ar[r] & V
}
$$
est cartésien (dans la catégorie des schémas). Le lemme précité
Morphismes, lemme \ref{morphisms-lemma-set-points-where-fibres-unramified}
montre donc que $T'$ est l'image inverse de $T$. Puisque
$|U'| \to |X'|$ est surjective, cela démontre le lemme.
\end{proof}

\begin{lemma}
\label{lemma-permanence-unramified}
Soit $S$ un schéma.
Soient $X \to Y \to Z$ des morphismes d'espaces algébriques sur $S$.
Si $X \to Z$ est non ramifié, alors $X \to Y$ est non ramifié.
\end{lemma}

\begin{proof}
Choisissons un diagramme commutatif
$$
\xymatrix{
U \ar[d] \ar[r] & V \ar[d] \ar[r] & W \ar[d] \\
X \ar[r] & Y \ar[r] & Z
}
$$
dont les flèches verticales sont étales et surjectives. (Voir
Espaces, lemme \ref{spaces-lemma-lift-morphism-presentations}.)
Appliquer
Morphismes, lemme \ref{morphisms-lemma-unramified-permanence}
à la ligne supérieure.
\end{proof}














\section{Morphismes étales}
\label{section-etale}

\noindent
La notion de morphisme étale d'espaces algébriques a été définie dans
Propriétés des espaces, définition \ref{spaces-properties-definition-etale}.
Voici ce que signifie, pour un morphisme, être étale en un point.

\begin{definition}
\label{definition-etale}
Soit $S$ un schéma.
Soit $f : X \to Y$ un morphisme d'espaces algébriques sur $S$.
Soit $x \in |X|$. On dit que $f$ est {\it étale en $x$} s'il
existe un voisinage ouvert $X' \subset X$ de $x$ tel que
$f|_{X'} : X' \to Y$ soit étale.
\end{definition}

\begin{lemma}
\label{lemma-etale-local}
Soit $S$ un schéma.
Soit $f : X \to Y$ un morphisme d'espaces algébriques sur $S$.
Les conditions suivantes sont équivalentes :
\begin{enumerate}
\item $f$ est étale ;
\item pour tout $x \in |X|$, le morphisme $f$ est étale en $x$ ;
\item pour tout schéma $Z$ et tout morphisme $Z \to Y$, le morphisme
$Z \times_Y X \to Z$ est étale ;
\item pour tout schéma affine $Z$ et tout morphisme
$Z \to Y$, le morphisme $Z \times_Y X \to Z$ est étale ;
\item il existe un schéma $V$ et un morphisme étale surjectif
$V \to Y$ tels que $V \times_Y X \to V$ soit un morphisme étale ;
\item il existe un schéma $U$ et un morphisme étale surjectif
$\varphi : U \to X$ tels que le composé $f \circ \varphi$
soit étale ;
\item pour tout diagramme commutatif
$$
\xymatrix{
U \ar[d] \ar[r] & V \ar[d] \\
X \ar[r] & Y
}
$$
où $U$ et $V$ sont des schémas et où les flèches verticales sont étales,
la flèche horizontale supérieure est étale ;
\item il existe un diagramme commutatif
$$
\xymatrix{
U \ar[d] \ar[r] & V \ar[d] \\
X \ar[r] & Y
}
$$
où $U$ et $V$ sont des schémas, où les flèches verticales sont étales et où
$U \to X$ est surjectif, tel que la flèche horizontale supérieure soit étale ;
\item il existe des recouvrements de Zariski $Y = \bigcup Y_i$ et
$f^{-1}(Y_i) = \bigcup X_{ij}$ tels que chaque morphisme
$X_{ij} \to Y_i$ soit étale.
\end{enumerate}
\end{lemma}

\begin{proof}
Combiner
Propriétés des espaces, lemmes
\ref{spaces-properties-lemma-etale-local},
\ref{spaces-properties-lemma-base-change-etale} et
\ref{spaces-properties-lemma-composition-etale}.
Certains détails sont omis.
\end{proof}

\begin{lemma}
\label{lemma-composition-etale}
Le composé de deux morphismes étales d'espaces algébriques
est étale.
\end{lemma}

\begin{proof}
Il s'agit d'une reprise de
Propriétés des espaces, lemme \ref{spaces-properties-lemma-composition-etale}.
\end{proof}

\begin{lemma}
\label{lemma-base-change-etale}
Le changement de base d'un morphisme étale d'espaces algébriques
par un morphisme quelconque d'espaces algébriques est étale.
\end{lemma}

\begin{proof}
Il s'agit d'une reprise de
Propriétés des espaces, lemme \ref{spaces-properties-lemma-base-change-etale}.
\end{proof}

\begin{lemma}
\label{lemma-etale-locally-quasi-finite}
Un morphisme étale d'espaces algébriques est localement quasi-fini.
\end{lemma}

\begin{proof}
Soit $X \to Y$ un morphisme étale d'espaces algébriques ; voir
Propriétés des espaces, définition \ref{spaces-properties-definition-etale}.
D'après
Propriétés des espaces, lemme \ref{spaces-properties-lemma-etale-local}
on voit que cela signifie qu'il existe un diagramme comme dans le
lemme \ref{lemma-local-source-target},
où $h$ est étale et où $a$ est une flèche verticale surjective. D'après
Morphismes, lemme \ref{morphisms-lemma-etale-locally-quasi-finite}
$h$ est localement quasi-fini. Ainsi, $X \to Y$ est localement quasi-fini
par définition.
\end{proof}

\begin{lemma}
\label{lemma-etale-smooth}
Un morphisme étale d'espaces algébriques est lisse.
\end{lemma}

\begin{proof}
La démonstration est identique à celle du
lemme \ref{lemma-etale-locally-quasi-finite}.
Elle utilise le fait qu'un morphisme étale de schémas est lisse
(par définition d'un morphisme étale de schémas).
\end{proof}

\begin{lemma}
\label{lemma-etale-flat}
Un morphisme étale d'espaces algébriques est plat.
\end{lemma}

\begin{proof}
La démonstration est identique à celle du
lemme \ref{lemma-etale-locally-quasi-finite}.
Elle utilise
Morphismes, lemme \ref{morphisms-lemma-etale-flat}.
\end{proof}

\begin{lemma}
\label{lemma-etale-locally-finite-presentation}
\begin{slogan}
Étale implique localement de présentation finie.
\end{slogan}
Un morphisme étale d'espaces algébriques est localement de présentation finie.
\end{lemma}

\begin{proof}
La démonstration est identique à celle du
lemme \ref{lemma-etale-locally-quasi-finite}.
Elle utilise
Morphismes, lemme \ref{morphisms-lemma-etale-locally-finite-presentation}.
\end{proof}

\begin{lemma}
\label{lemma-etale-locally-finite-type}
Un morphisme étale d'espaces algébriques est localement de type fini.
\end{lemma}

\begin{proof}
Un morphisme étale est localement de présentation finie,
et un morphisme localement de présentation finie est localement de type fini ;
voir les
lemmes \ref{lemma-etale-locally-finite-presentation} et
\ref{lemma-finite-presentation-finite-type}.
\end{proof}

\begin{lemma}
\label{lemma-etale-unramified}
Un morphisme étale d'espaces algébriques est non ramifié.
\end{lemma}

\begin{proof}
La démonstration est identique à celle du
lemme \ref{lemma-etale-locally-quasi-finite}.
Elle utilise
Morphismes, lemme \ref{morphisms-lemma-etale-smooth-unramified}.
\end{proof}

\begin{lemma}
\label{lemma-etale-permanence}
Soit $S$ un schéma. Soient $X, Y$ des espaces algébriques étales sur
un espace algébrique $Z$. Tout morphisme $X \to Y$ sur $Z$ est étale.
\end{lemma}

\begin{proof}
Il s'agit d'une reprise de
Propriétés des espaces, lemme \ref{spaces-properties-lemma-etale-permanence}.
\end{proof}

\begin{lemma}
\label{lemma-unramified-flat-lfp-etale}
Un morphisme d'espaces algébriques localement de présentation finie,
plat et non ramifié est étale.
\end{lemma}

\begin{proof}
Soit $X \to Y$ un morphisme d'espaces algébriques localement de présentation
finie, plat et non ramifié. D'après
Propriétés des espaces, lemme \ref{spaces-properties-lemma-etale-local}
on voit que cela signifie qu'il existe un diagramme comme dans le
lemme \ref{lemma-local-source-target},
où $h$ est localement de présentation finie, plat et non ramifié,
et où $a$ est une flèche verticale surjective. D'après
Morphismes, lemme \ref{morphisms-lemma-flat-unramified-etale}
$h$ est étale. Ainsi, $X \to Y$ est étale par définition.
\end{proof}






\section{Morphismes propres}
\label{section-proper}

\noindent
La notion de morphisme propre joue un rôle important en géométrie algébrique.
Voici la définition d'un morphisme propre d'espaces algébriques.

\begin{definition}
\label{definition-proper}
Soit $S$ un schéma.
Soit $f : X \to Y$ un morphisme d'espaces algébriques sur $S$.
On dit que $f$ est {\it propre} si $f$ est séparé, de type fini et
universellement fermé.
\end{definition}

\begin{lemma}
\label{lemma-proper-local}
Soit $S$ un schéma. Soit $f : X \to Y$ un morphisme d'espaces algébriques
sur $S$. Les conditions suivantes sont équivalentes :
\begin{enumerate}
\item $f$ est propre ;
\item pour tout schéma $Z$ et tout morphisme $Z \to Y$,
la projection $Z \times_Y X \to Z$ est propre ;
\item pour tout schéma affine $Z$ et tout morphisme $Z \to Y$,
la projection $Z \times_Y X \to Z$ est propre ;
\item il existe un schéma $V$ et un morphisme étale surjectif
$V \to Y$ tel que $V \times_Y X \to V$ soit propre ;
\item il existe un recouvrement de Zariski $Y = \bigcup Y_i$ tel que
chacun des morphismes $f^{-1}(Y_i) \to Y_i$ soit propre.
\end{enumerate}
\end{lemma}

\begin{proof}
Combiner les lemmes \ref{lemma-separated-local},
\ref{lemma-finite-type-local},
\ref{lemma-quasi-compact-local} et
\ref{lemma-universally-closed-local}.
\end{proof}

\begin{lemma}
\label{lemma-base-change-proper}
Un changement de base d'un morphisme propre est propre.
\end{lemma}

\begin{proof}
Voir les
lemmes \ref{lemma-base-change-separated},
\ref{lemma-base-change-finite-type} et
\ref{lemma-base-change-universally-closed}.
\end{proof}

\begin{lemma}
\label{lemma-composition-proper}
Un composé de morphismes propres est propre.
\end{lemma}

\begin{proof}
Voir les
lemmes \ref{lemma-composition-separated},
\ref{lemma-composition-finite-type} et
\ref{lemma-composition-universally-closed}.
\end{proof}

\begin{lemma}
\label{lemma-closed-immersion-proper}
Une immersion fermée d'espaces algébriques est un morphisme propre
d'espaces algébriques.
\end{lemma}

\begin{proof}
Comme une immersion fermée est représentable par définition, cela résulte de
Espaces,
lemme \ref{spaces-lemma-representable-transformations-property-implication}
et du résultat correspondant pour les morphismes de schémas ; voir
Morphismes, lemme \ref{morphisms-lemma-closed-immersion-proper}.
\end{proof}

\begin{lemma}
\label{lemma-universally-closed-permanence}
Soit $S$ un schéma.
Considérons un diagramme commutatif d'espaces algébriques
$$
\xymatrix{
X \ar[rr] \ar[rd] & &
Y \ar[ld] \\
& B &
}
$$
sur $S$.
\begin{enumerate}
\item Si $X \to B$ est universellement fermé et si $Y \to B$ est
séparé, alors le morphisme $X \to Y$ est universellement fermé.
En particulier, l'image de $|X|$ dans $|Y|$ est fermée.
\item Si $X \to B$ est propre et si $Y \to B$ est séparé, alors
le morphisme $X \to Y$ est propre.
\end{enumerate}
\end{lemma}

\begin{proof}
Supposons que $X \to B$ soit universellement fermé et que $Y \to B$ soit séparé.
Factorisons le morphisme en $X \to X \times_B Y \to Y$.
Le premier morphisme est une immersion fermée ; voir le
lemme \ref{lemma-semi-diagonal} ;
il est donc universellement fermé.
La projection $X \times_B Y \to Y$ est le changement de base
d'un morphisme universellement fermé et est donc
universellement fermée ; voir le
lemme \ref{lemma-base-change-universally-closed}.
Ainsi, $X \to Y$ est universellement fermé en tant que composé
de morphismes universellement fermés ; voir le
lemme \ref{lemma-composition-universally-closed}.
Cela démontre (1). Pour en déduire (2), combiner (1) avec les
lemmes \ref{lemma-compose-after-separated},
\ref{lemma-quasi-compact-permanence} et
\ref{lemma-permanence-finite-type}.
\end{proof}

\begin{lemma}
\label{lemma-image-proper-is-proper}
Soit $S$ un schéma. Soit $B$ un espace algébrique sur $S$.
Soit $f : X \to Y$ un morphisme d'espaces algébriques sur $B$.
Si $X$ est universellement fermé sur $B$ et si $f$ est surjectif, alors
$Y$ est universellement fermé sur $B$. En particulier, si $Y$ est aussi
séparé et de type fini sur $B$, alors $Y$ est propre sur $B$.
\end{lemma}

\begin{proof}
Supposons que $X$ soit universellement fermé et que $f$ soit surjectif.
Notons $p : X \to B$ et $q : Y \to B$ les morphismes structuraux.
Soit $B' \to B$ un morphisme d'espaces algébriques sur $S$.
Le changement de base $f' : X_{B'} \to Y_{B'}$ est surjectif
(lemme \ref{lemma-base-change-surjective}), et le changement de
base $p' : X_{B'} \to B'$ est fermé.
Si $T \subset Y_{B'}$ est fermé, alors $(f')^{-1}(T) \subset X_{B'}$
est fermé ; par conséquent, $p'((f')^{-1}(T)) = q'(T)$ est fermé.
Ainsi, $q'$ est fermé.
\end{proof}

\begin{lemma}
\label{lemma-scheme-theoretic-image-is-proper}
Soit $S$ un schéma. Soit
$$
\xymatrix{
X \ar[rr]_h \ar[rd]_f & & Y \ar[ld]^g \\
& B
}
$$
un diagramme commutatif de morphismes d'espaces algébriques sur $S$.
Supposons que
\begin{enumerate}
\item $X \to B$ soit un morphisme propre ;
\item $Y \to B$ soit séparé et localement de type fini.
\end{enumerate}
Alors l'image schématique $Z \subset Y$ de $h$
est propre sur $B$ et $X \to Z$ est surjectif.
\end{lemma}

\begin{proof}
L'image schématique de $h$ est construite à la section
\ref{section-scheme-theoretic-image}.
Observons que $h$ est quasi-compact
(lemme \ref{lemma-quasi-compact-quasi-separated-permanence})
et que, par conséquent, $|h|(|X|) \subset |Z|$
est dense (lemme \ref{lemma-quasi-compact-scheme-theoretic-image}).
D'autre part, $|h|(|X|)$ est fermé dans $|Y|$
(lemme \ref{lemma-universally-closed-permanence})
et donc $X \to Z$ est surjectif.
Ainsi, $Z \to B$ est propre (lemme \ref{lemma-image-proper-is-proper}).
\end{proof}

\begin{lemma}
\label{lemma-separated-diagonal-proper}
Soit $S$ un schéma.
Soit $f : X \to Y$ un morphisme d'espaces algébriques sur $S$.
Les conditions suivantes sont équivalentes :
\begin{enumerate}
\item $f$ est séparé ;
\item $\Delta_{X/Y} : X \to X \times_Y X$ est universellement fermé ;
\item $\Delta_{X/Y} : X \to X \times_Y X$ est propre.
\end{enumerate}
\end{lemma}

\begin{proof}
L'implication (1) $\Rightarrow$ (3) résulte du
lemme \ref{lemma-closed-immersion-proper}.
Nous utiliserons
Espaces, lemme
\ref{spaces-lemma-representable-transformations-property-implication}
sans autre mention dans la suite de la démonstration.
Rappelons que $\Delta_{X/Y}$ est un monomorphisme
représentable qui est localement de type fini ; voir le
lemme \ref{lemma-properties-diagonal}.
Puisque propre $\Rightarrow$ universellement fermé pour les morphismes de schémas,
on conclut que (3) implique (2).
Si $\Delta_{X/Y}$ est universellement fermé, alors
\'Morphismes étales,
lemme \ref{etale-lemma-characterize-closed-immersion}
implique que c'est une immersion fermée. Ainsi, (2) $\Rightarrow$ (1),
ce qui achève la démonstration.
\end{proof}





\section{Critères valuatifs}
\label{section-valuative}

\noindent
Cette section introduit les notions de base sur les critères valuatifs pour les morphismes
d'espaces algébriques. Voici une liste de références vers des résultats complémentaires :
\begin{enumerate}
\item le critère valuatif de fermeture universelle se trouve à la
section \ref{section-valuative-criterion-universally-closed} ;
\item le critère valuatif de séparation se trouve à la
section \ref{section-valuative-separatedness} ;
\item le critère valuatif de propreté se trouve à la
section \ref{section-valuative-criterion-properness} ;
\item des réciproques supplémentaires se trouvent dans
Espaces décents, section
\ref{decent-spaces-section-valuative-criterion-universally-closed}
et Espaces décents, lemme
\ref{decent-spaces-lemma-re-characterize-universally-closed} ;
\item dans le cas noethérien, il suffit de vérifier le critère pour les
anneaux de valuation discrète, comme le montrent Cohomologie des espaces, section
\ref{spaces-cohomology-section-Noetherian-valuative-criterion} et
Limites d'espaces, section
\ref{spaces-limits-section-Noetherian-valuative-criterion} ;
\item des versions affinées des critères valuatifs dans le cas noethérien
se trouvent dans Limites d'espaces, section
\ref{spaces-limits-section-refined-valuative-criteria}.
\end{enumerate}
Nous énonçons d'abord formellement la définition, puis nous expliquons en quoi elle
diffère du cas des morphismes de schémas.

\begin{definition}
\label{definition-valuative-criterion}
Soit $S$ un schéma.
Soit $f : X \to Y$ un morphisme d'espaces algébriques sur $S$.
On dit que $f$ {\it satisfait la partie d'unicité du critère valuatif}
si, pour tout diagramme commutatif en traits pleins
$$
\xymatrix{
\Spec(K) \ar[r] \ar[d] & X \ar[d] \\
\Spec(A) \ar[r] \ar@{-->}[ru] & Y
}
$$
où $A$ est un anneau de valuation de corps des fractions $K$, il existe
au plus une flèche pointillée (sans en exiger l'existence).
On dit que $f$ {\it satisfait la partie d'existence du critère valuatif}
si, pour tout diagramme en traits pleins comme ci-dessus, il existe une extension
de corps $K'/K$, un anneau de valuation $A' \subset K'$ qui domine
$A$ et un morphisme $\Spec(A') \to X$ tels que le diagramme
suivant soit commutatif :
$$
\xymatrix{
\Spec(K') \ar[r] \ar[d] & \Spec(K) \ar[r] & X \ar[d] \\
\Spec(A') \ar[r] \ar[rru] & \Spec(A) \ar[r] & Y
}
$$
On dit que $f$ {\it satisfait le critère valuatif}
si $f$ satisfait à la fois les parties d'existence et d'unicité.
\end{definition}

\noindent
La formulation de la partie d'existence du critère valuatif est
légèrement différente pour les morphismes d'espaces algébriques, puisqu'il peut être
nécessaire d'étendre le corps des fractions de l'anneau de valuation.
En pratique, cette différence ne joue presque jamais de rôle.
\begin{enumerate}
\item La vérification de la partie d'unicité du critère valuatif ne
fait jamais intervenir d'extension du corps des fractions ; elle est donc exactement la même
que dans le cas des schémas.
\item Il est nécessaire d'autoriser en général des extensions de corps ; voir
l'exemple \ref{example-finite-separable-needed}.
\item Pour les morphismes d'espaces algébriques, il suffit toujours de
prendre une extension finie séparable $K'/K$ dans la partie d'existence
du critère valuatif ; voir le lemme \ref{lemma-finite-separable-enough}.
\item Si $f : X \to Y$ est un morphisme séparé d'espaces algébriques, on
peut toujours prendre $K = K'$ lors de la vérification de la partie d'existence du
critère valuatif ; voir le lemme \ref{lemma-usual-enough}.
\item Pour un morphisme quasi-compact et quasi-séparé
$f : X \to Y$, on obtient l'équivalence entre
« $f$ est séparé et universellement fermé » et « $f$ satisfait
le critère valuatif usuel » ; voir le
lemme \ref{lemma-characterize-separated-and-universally-closed}.
Le critère valuatif de propreté est le critère usuel ; voir le
lemme \ref{lemma-characterize-proper}.
\end{enumerate}
Comme première étape de la théorie, montrons que le critère est identique
à celui formulé pour les morphismes de schémas
lorsque le morphisme d'espaces algébriques est représentable.

\begin{lemma}
\label{lemma-valuative-criterion-representable}
Soit $S$ un schéma.
Soit $f : X \to Y$ un morphisme d'espaces algébriques sur $S$.
Supposons que $f$ soit représentable. Les conditions suivantes sont équivalentes :
\begin{enumerate}
\item $f$ satisfait la partie d'existence du critère valuatif
au sens de la définition \ref{definition-valuative-criterion} ;
\item pour tout diagramme commutatif en traits pleins
$$
\xymatrix{
\Spec(K) \ar[r] \ar[d] & X \ar[d] \\
\Spec(A) \ar[r] \ar@{-->}[ru] & Y
}
$$
où $A$ est un anneau de valuation de corps des fractions $K$, il existe
une flèche pointillée, c'est-à-dire que $f$ satisfait la partie d'existence du critère
valuatif au sens de
Schémas, définition \ref{schemes-definition-valuative-criterion}.
\end{enumerate}
\end{lemma}

\begin{proof}
Il suffit de montrer que, pour tout diagramme commutatif de la forme
$$
\xymatrix{
\Spec(K') \ar[r] \ar[d] & \Spec(K) \ar[r] & X \ar[d] \\
\Spec(A') \ar[r] \ar[rru]^\varphi & \Spec(A) \ar[r] & Y
}
$$
comme dans la définition \ref{definition-valuative-criterion}, on peut aussi
trouver un morphisme $\Spec(A) \to X$ qui s'insère dans le diagramme.
Posons $X_A = \Spec(A) \times_Y X$. Comme $f$ est représentable, on voit
que $X_A$ est un schéma. Le morphisme $\varphi$ donne un morphisme
$\varphi' : \Spec(A') \to X_A$. Soit $x \in X_A$ l'image du
point fermé par $\varphi' : \Spec(A') \to X_A$. On dispose alors du
diagramme commutatif d'anneaux suivant :
$$
\xymatrix{
K' & K \ar[l] & \mathcal{O}_{X_A, x} \ar[l] \ar[lld] \\
A' \ar[u] & A \ar[l] & A \ar[l] \ar[u]
}
$$
Puisque $A$ est un anneau de valuation et que $A'$ domine $A$, on voit
que $K \cap A' = A$. L'image de l'homomorphisme d'anneaux $\mathcal{O}_{X_A, x} \to K$
est donc contenue dans $A$. On obtient ainsi un morphisme $\Spec(A) \to X_A$ (voir
Schémas, section \ref{schemes-section-points})
comme souhaité.
\end{proof}

\begin{lemma}
\label{lemma-finite-separable-enough}
Soit $S$ un schéma.
Soit $f : X \to Y$ un morphisme d'espaces algébriques sur $S$.
Les conditions suivantes sont équivalentes :
\begin{enumerate}
\item $f$ satisfait la partie d'existence du critère valuatif
au sens de la définition \ref{definition-valuative-criterion} ;
\item $f$ satisfait la partie d'existence du critère valuatif
au sens de la définition \ref{definition-valuative-criterion}, modifiée en
exigeant que l'extension $K'/K$ soit finie séparable.
\end{enumerate}
\end{lemma}

\begin{proof}
Il faut montrer que (1) implique (2). Supposons donné un diagramme
$$
\xymatrix{
\Spec(K') \ar[r] \ar[d] & \Spec(K) \ar[r] & X \ar[d] \\
\Spec(A') \ar[r] \ar[rru] & \Spec(A) \ar[r] & Y
}
$$
comme dans la définition \ref{definition-valuative-criterion}, avec $K \subset K'$
arbitraire. Choisissons un schéma $U$ et un morphisme étale surjectif $U \to X$.
Alors
$$
\Spec(A') \times_X U \longrightarrow \Spec(A')
$$
est étale surjectif. Soit $p$ un point de $\Spec(A') \times_X U$
qui s'envoie sur le point fermé de $\Spec(A')$. Soit $p' \leadsto p$
une générisation de $p$ qui s'envoie sur le point générique de $\Spec(A')$.
Une telle générisation existe parce que les générisations se relèvent le long des morphismes
plats de schémas ; voir
Morphismes, lemme \ref{morphisms-lemma-generalizations-lift-flat}.
Alors $p'$ correspond à un point du schéma $\Spec(K') \times_X U$.
Notons que
$$
\Spec(K') \times_X U
=
\Spec(K') \times_{\Spec(K)} (\Spec(K) \times_X U)
$$
Par conséquent, $p'$ s'envoie sur un point $q' \in \Spec(K) \times_X U$ dont le
corps résiduel est une extension finie séparable de $K$. Enfin,
$p' \leadsto p$ s'envoie sur une spécialisation $u' \leadsto u$ dans le
schéma $U$. Avec toutes ces notations, on obtient le diagramme d'anneaux
suivant :
$$
\xymatrix{
\kappa(p') & & \kappa(q') \ar[ll] & \kappa(u') \ar[l] \\
& \mathcal{O}_{\Spec(A') \times_X U, p} \ar[lu] & &
\mathcal{O}_{U, u} \ar[ll] \ar[u] \\
K' \ar[uu] & A' \ar[l] \ar[u] & A \ar[l] \ar'[u][uu]
}
$$
Cela signifie que l'anneau $B \subset \kappa(q')$ engendré par
les images de $A$ et de $\mathcal{O}_{U, u}$ s'envoie dans un sous-anneau
de $\kappa(p')$ contenu dans l'image $B'$ de
$\mathcal{O}_{\Spec(A') \times_X U, p} \to \kappa(p')$.
Notons que $B'$ est un anneau local. Soit $\mathfrak m \subset B'$
son idéal maximal. Par construction, $A \cap \mathfrak m$
(resp.\ $\mathcal{O}_{U, u} \cap \mathfrak m$, resp.\ $A' \cap \mathfrak m$)
est l'idéal maximal de $A$ (resp.\ de $\mathcal{O}_{U, u}$, resp.\ de $A'$).
Posons $\mathfrak q = B \cap \mathfrak m$. C'est un
idéal premier tel que $A \cap \mathfrak q$ soit l'idéal maximal de $A$.
Ainsi, $B_{\mathfrak q} \subset \kappa(q')$ est un anneau local qui domine
$A$. D'après
Algèbre, lemme \ref{algebra-lemma-dominate}
on peut trouver un anneau de valuation $A_1 \subset \kappa(q')$
de corps des fractions $\kappa(q')$
qui domine $B_{\mathfrak q}$. L'homomorphisme (local) d'anneaux
$\mathcal{O}_{U, u} \to A_1$ donne un morphisme
$\Spec(A_1) \to U \to X$
tel que le diagramme
$$
\xymatrix{
\Spec(\kappa(q')) \ar[r] \ar[d] & \Spec(K) \ar[r] & X \ar[d] \\
\Spec(A_1) \ar[r] \ar[rru] & \Spec(A) \ar[r] & Y
}
$$
soit commutatif. Puisque le corps des fractions de $A_1$ est $\kappa(q')$
et que $\kappa(q')/K$ est une extension finie
séparable par construction, le lemme est démontré.
\end{proof}

\begin{lemma}
\label{lemma-push-down-solution}
Soit $S$ un schéma. Soit $f : X \to Y$ un morphisme séparé d'espaces
algébriques sur $S$. Supposons donné un diagramme
$$
\xymatrix{
\Spec(K') \ar[r] \ar[d] & \Spec(K) \ar[r] & X \ar[d] \\
\Spec(A') \ar[r] \ar[rru] & \Spec(A) \ar[r] \ar@{-->}[ru] & Y
}
$$
comme dans la définition \ref{definition-valuative-criterion}, avec $K \subset K'$
arbitraire. Alors il existe une flèche pointillée qui rend le diagramme commutatif.
\end{lemma}

\begin{proof}
Il faut montrer que l'on peut trouver un morphisme $\Spec(A) \to X$ qui s'insère
dans le diagramme.

\medskip\noindent
Considérons le changement de base $X_A = \Spec(A) \times_Y X$ de $X$.
Alors $X_A \to \Spec(A)$ est un morphisme séparé d'espaces
algébriques (lemme \ref{lemma-base-change-separated}). En changeant la base de
tous les morphismes du diagramme ci-dessus, on obtient
$$
\xymatrix{
\Spec(K') \ar[r] \ar[d] & \Spec(K) \ar[r] & X_A \ar[d] \\
\Spec(A') \ar[r] \ar[rru] & \Spec(A) \ar@{=}[r] & \Spec(A)
}
$$
On peut donc remplacer $X$ par $X_A$, supposer que $Y = \Spec(A)$ et
que l'on dispose d'un diagramme comme ci-dessus. On remplace alors $X$ par
un sous-espace ouvert quasi-compact contenant l'image de $|\Spec(A')| \to |X|$.

\medskip\noindent
Le morphisme $\Spec(A') \to X$ est quasi-compact d'après le
lemme \ref{lemma-quasi-compact-permanence}.
Soit $Z \subset X$ l'image schématique de
$\Spec(A') \to X$. Alors $Z$ est un espace algébrique réduit
(lemme \ref{lemma-scheme-theoretic-image-reduced}),
quasi-compact (en tant que sous-espace fermé de $X$) et séparé
(en tant que sous-espace fermé de $X$) sur $A$.
Considérons le changement de base
$$
\Spec(K') = \Spec(A') \times_{\Spec(A)} \Spec(K) \to
X \times_{\Spec(A)} \Spec(K) = X_K
$$
du morphisme $\Spec(A') \to X$ par le morphisme plat de schémas
$\Spec(K) \to \Spec(A)$. D'après le
lemme \ref{lemma-flat-base-change-scheme-theoretic-image},
on voit que l'image schématique de ce morphisme
est le changement de base $Z_K$ de $Z$. D'autre part, par
hypothèse (c'est-à-dire d'après le diagramme commutatif ci-dessus),
ce morphisme se factorise par un morphisme
$\Spec(K) \to Z_K$ qui est une section du morphisme structural
$Z_K \to \Spec(K)$. Comme $Z_K$ est séparé, cette section est une
immersion fermée (lemme \ref{lemma-section-immersion}).
On conclut que $Z_K = \Spec(K)$.

\medskip\noindent
Soit $V \to Z$ un morphisme étale surjectif avec $V$ un
schéma affine (Propriétés des espaces, lemme
\ref{spaces-properties-lemma-quasi-compact-affine-cover}).
Écrivons $V = \Spec(B)$. Alors $V \times_Z \Spec(A') = \Spec(C)$
est affine puisque $Z$ est séparé. Notons que $B \to C$ est injectif,
car $V$ est l'image schématique de $V \times_Z \Spec(A') \to V$
d'après le lemme \ref{lemma-quasi-compact-scheme-theoretic-image}.
D'autre part, $A' \to C$ est étale, car il correspond
au changement de base de $V \to Z$.
Puisque $A'$ est un $A$-module sans torsion,
la platitude de $A' \to C$ implique que $C$ est un
$A$-module sans torsion ; ainsi, $B$ est un $A$-module sans torsion.
Notons qu'être sans torsion comme $A$-module équivaut
à être plat (Compléments d'algèbre, lemme
\ref{more-algebra-lemma-valuation-ring-torsion-free-flat}).
Ensuite, écrivons
$$
V \times_Z V = \Spec(B')
$$
Notons que les deux homomorphismes d'anneaux $B \to B'$ sont étales puisque $V \to Z$
est étale. L'homomorphisme canonique surjectif $B \otimes_A B \to B'$
devient un isomorphisme après tensorisation par $K$ sur $A$,
car $Z_K = \Spec(K)$. Cependant, $B \otimes_A B$ est sans torsion
comme $A$-module d'après les remarques précédentes. Ainsi, $B' = B \otimes_A B$.
Il en résulte que le changement de base de l'homomorphisme d'anneaux
$A \to B$ par l'homomorphisme fidèlement plat $A \to B$
est étale (notons que $\Spec(B) \to \Spec(A)$ est surjectif,
puisque $Z \to \Spec(A)$ est surjectif). Ainsi, $A \to B$ est
étale (Descente, lemme \ref{descent-lemma-descending-property-etale}),
autrement dit, $V \to \Spec(A)$ est étale.
Comme $V \times_Z V = V \times_{\Spec(A)} V$,
on conclut que $Z = \Spec(A)$ en tant qu'espaces algébriques
(par exemple d'après Espaces, lemme \ref{spaces-lemma-space-presentation}),
ce qui achève la démonstration.
\end{proof}

\begin{lemma}
\label{lemma-usual-enough}
Soit $S$ un schéma. Soit $f : X \to Y$ un morphisme séparé d'espaces
algébriques sur $S$. Les conditions suivantes sont équivalentes :
\begin{enumerate}
\item $f$ satisfait la partie d'existence du critère valuatif
au sens de la définition \ref{definition-valuative-criterion} ;
\item pour tout diagramme commutatif en traits pleins
$$
\xymatrix{
\Spec(K) \ar[r] \ar[d] & X \ar[d] \\
\Spec(A) \ar[r] \ar@{-->}[ru] & Y
}
$$
où $A$ est un anneau de valuation de corps des fractions $K$, il existe
une flèche pointillée, c'est-à-dire que $f$ satisfait la partie d'existence du critère
valuatif au sens de
Schémas, définition \ref{schemes-definition-valuative-criterion}.
\end{enumerate}
\end{lemma}

\begin{proof}
Il faut montrer que (1) implique (2). Supposons donné un diagramme commutatif
$$
\xymatrix{
\Spec(K) \ar[r] \ar[d] & X \ar[d] \\
\Spec(A) \ar[r] & Y
}
$$
comme dans la partie (2). D'après (1), il existe un diagramme commutatif
$$
\xymatrix{
\Spec(K') \ar[r] \ar[d] & \Spec(K) \ar[r] & X \ar[d] \\
\Spec(A') \ar[r] \ar[rru] & \Spec(A) \ar[r] & Y
}
$$
comme dans la définition \ref{definition-valuative-criterion}, avec $K \subset K'$
arbitraire. D'après le lemme \ref{lemma-push-down-solution}, on peut trouver un morphisme
$\Spec(A) \to X$ qui s'insère dans le diagramme, c'est-à-dire que (2) est satisfaite.
\end{proof}

\begin{example}
\label{example-finite-separable-needed}
Considérons l'espace algébrique $X$ construit dans
Espaces, exemple \ref{spaces-example-non-representable-descent}.
Rappelons qu'il s'agit d'une forme tordue galoisienne de la droite affine à origine dédoublée.
La torsion galoisienne est relative à une extension galoisienne de degré deux
$k'/k$ de corps. Elle est ainsi munie d'un morphisme
$$
\pi : X \longrightarrow S = \mathbf{A}^1_k
$$
qui est quasi-compact. Affirmons que $\pi$ est universellement fermé.
En effet, après le changement de base par $\Spec(k') \to \Spec(k)$,
le morphisme $\pi$ s'identifie au morphisme
$$
\text{droite affine à origine dédoublée}
\longrightarrow
\text{droite affine}
$$
qui est universellement fermé (certains détails sont omis). Puisque le morphisme
$\Spec(k') \to \Spec(k)$ est universellement fermé et
surjectif, une chasse au diagramme montre que $\pi$ est universellement fermé.
D'autre part, considérons le diagramme
$$
\xymatrix{
\Spec(k((x))) \ar[r] \ar[d] & X \ar[d]^\pi \\
\Spec(k[[x]]) \ar[r] \ar@{..>}[ru] & \mathbf{A}^1_k
}
$$
Puisque l'unique point de $X$ au-dessus de $0 \in \mathbf{A}^1_k$
correspond à un monomorphisme $\Spec(k') \to X$,
il est clair qu'il ne peut exister de flèche pointillée ! Cela montre qu'une
extension finie séparable de corps est nécessaire en général.
\end{example}

\begin{lemma}
\label{lemma-base-change-valuative-criteria}
Le changement de base d'un morphisme d'espaces algébriques qui satisfait la
partie d'existence (resp.\ la partie d'unicité) du critère valuatif
par un morphisme quelconque d'espaces algébriques satisfait la
partie d'existence (resp.\ la partie d'unicité) du critère valuatif.
\end{lemma}

\begin{proof}
Soit $f : X \to Y$ un morphisme d'espaces algébriques sur le schéma $S$.
Soit $Z \to Y$ un morphisme quelconque d'espaces algébriques sur $S$.
Considérons un diagramme commutatif en traits pleins de la forme suivante :
$$
\xymatrix{
\Spec(K) \ar[r] \ar[d] & Z \times_Y X \ar[r] \ar[d] & X \ar[d] \\
\Spec(A) \ar[r] \ar@{-->}[ru] \ar@{-->}[rru] & Z \ar[r] & Y
}
$$
L'ensemble des flèches pointillées nord-ouest qui rendent le diagramme commutatif
est en correspondance biunivoque avec l'ensemble des flèches pointillées ouest-nord-ouest
qui le rendent commutatif. Cela démontre le lemme dans le cas de
« l'unicité ». Pour la partie d'existence, supposons que $f$ satisfasse la partie d'existence
du critère valuatif. Si l'on se donne un diagramme commutatif en traits pleins
comme ci-dessus, il existe par hypothèse une extension de corps $K'/K$,
un anneau de valuation $A' \subset K'$ qui domine $A$ et
un morphisme $\Spec(A') \to X$ qui s'insère dans le diagramme commutatif
suivant :
$$
\xymatrix{
\Spec(K') \ar[r] \ar[d] &
\Spec(K) \ar[r] & Z \times_Y X \ar[r] & X \ar[d] \\
\Spec(A') \ar[r] \ar[rrru] & \Spec(A) \ar[r] & Z \ar[r] & Y
}
$$
Et, d'après les remarques précédentes, la flèche oblique correspond à une flèche
$\Spec(A') \to Z \times_Y X$, comme souhaité.
\end{proof}

\begin{lemma}
\label{lemma-composition-valuative-criteria}
Le composé de deux morphismes d'espaces algébriques qui satisfont la
partie d'existence (resp.\ la partie d'unicité) du critère valuatif
satisfait la partie d'existence (resp.\ la partie d'unicité) du critère
valuatif.
\end{lemma}

\begin{proof}
Soient $f : X \to Y$ et $g : Y \to Z$ des morphismes d'espaces algébriques sur le
schéma $S$. Considérons un diagramme commutatif en traits pleins de la forme suivante :
$$
\xymatrix{
\Spec(K) \ar[dd] \ar[r] & X \ar[d]^f \\
& Y \ar[d]^g \\
\Spec(A) \ar[r] \ar@{-->}[ru] \ar@{-->}[ruu] & Z
}
$$
Si l'on dispose de la partie d'unicité pour $g$, il existe au
plus une flèche pointillée nord-ouest qui rende le diagramme commutatif.
Si l'on dispose aussi de la partie d'unicité pour $f$, il existe
au plus une flèche pointillée nord-nord-ouest qui rende le diagramme
commutatif. Dans le cas de l'existence, la démonstration résulte de l'examen
du diagramme suivant :
$$
\xymatrix{
\Spec(K'') \ar[r] \ar[dd] &
\Spec(K') \ar[r] &
\Spec(K) \ar[r] &
X \ar[d]^f \\
& & & Y \ar[d]^g \\
\Spec(A'') \ar[r] \ar[rrruu] &
\Spec(A') \ar[r] \ar[rru] &
\Spec(A) \ar[r] &
Z
}
$$
En effet, la partie d'existence pour $g$ fournit l'extension $K'$, l'anneau
de valuation $A'$ et la flèche $\Spec(A') \to Y$, après quoi
la partie d'existence pour $f$ fournit l'extension $K''$, l'anneau
de valuation $A''$ et la flèche $\Spec(A'') \to X$.
\end{proof}






\section{Critère valuatif de fermeture universelle}
\label{section-valuative-criterion-universally-closed}

\noindent
La partie d'existence du critère valuatif entraîne qu'un morphisme
quasi-compact est universellement fermé ; voir le
lemme \ref{lemma-quasi-compact-existence-universally-closed}.
Dans le cas des schémas, il s'agit d'une équivalence,
mais cet énoncé est faux pour les morphismes d'espaces algébriques.
L'exemple \ref{example-strange-universally-closed}
montre que $\mathbf{A}^1_k/\mathbf{Z} \to \Spec(k)$ est universellement
fermé, alors qu'il est facile de voir que la partie d'existence
du critère valuatif n'est pas satisfaite. Nous revenons sur ce sujet dans
Espaces décents, section
\ref{decent-spaces-section-valuative-criterion-universally-closed}
et y montrons que la réciproque vaut lorsque la source du morphisme est un espace décent
(voir aussi
Espaces décents,
lemme \ref{decent-spaces-lemma-re-characterize-universally-closed},
pour une version relative).

\begin{lemma}
\label{lemma-quasi-compact-existence-universally-closed}
Soit $S$ un schéma.
Soit $f : X \to Y$ un morphisme d'espaces algébriques sur $S$.
Supposons que
\begin{enumerate}
\item $f$ soit quasi-compact, et que
\item $f$ satisfasse la partie d'existence du critère valuatif.
\end{enumerate}
Alors $f$ est universellement fermé.
\end{lemma}

\begin{proof}
D'après les lemmes \ref{lemma-base-change-quasi-compact}
et \ref{lemma-base-change-valuative-criteria},
les propriétés (1) et (2) sont préservées par
tout changement de base. D'après le lemme \ref{lemma-universally-closed-local},
il suffit de montrer que $|T \times_Y X| \to |T|$ est fermé
chaque fois que $T$ est un schéma affine sur $S$ muni d'un morphisme vers $Y$. Il
suffit donc de prouver ceci : si $Y$ est un schéma affine, si $f : X \to Y$ est quasi-compact
et satisfait la partie d'existence du critère valuatif, alors
$f : |X| \to |Y|$ est fermé. Dans cette situation, $X$ est un espace
algébrique quasi-compact. D'après
Propriétés des espaces,
lemme \ref{spaces-properties-lemma-quasi-compact-affine-cover},
il existe un schéma affine $U$ et un morphisme étale surjectif
$\varphi : U \to X$. Soit $T \subset |X|$ fermé. L'image inverse
$\varphi^{-1}(T) \subset U$ est fermée ; c'est donc l'ensemble des points
d'un sous-schéma fermé affine $Z \subset U$. Ainsi, d'après
Algèbre, lemme \ref{algebra-lemma-image-stable-specialization-closed}
la partie $f(T) = f(\varphi(|Z|)) \subset |Y|$ est fermée dès qu'elle est
stable par spécialisation.

\medskip\noindent
Soit $y' \leadsto y$ une spécialisation dans $Y$ telle que $y' \in f(T)$.
Choisissons un point $x' \in T \subset |X|$ qui s'envoie sur $y'$ par $f$.
On peut représenter $x'$ par un morphisme $\Spec(K) \to X$
pour un certain corps $K$. On obtient donc le diagramme suivant
$$
\xymatrix{
\Spec(K) \ar[r]_-{x'} \ar[d] & X \ar[d]^f \\
\Spec(\mathcal{O}_{Y, y}) \ar[r] & Y,
}
$$
voir
Schémas, section \ref{schemes-section-points}
pour l'existence de la flèche verticale de gauche.
Choisissons un anneau de valuation $A \subset K$ qui domine l'image de
l'homomorphisme d'anneaux $\mathcal{O}_{Y, y} \to K$ (c'est possible puisque
cette image est un anneau local qui n'est pas un corps, car $y' \not = y$ ; voir
Algèbre, lemme \ref{algebra-lemma-dominate}).
Par hypothèse, il existe une extension de corps $K'/K$, un
anneau de valuation $A' \subset K'$ qui domine $A$, et un morphisme
$\Spec(A') \to X$ qui complète le diagramme commutatif.
Comme $A'$ domine $A$ et que $A$ domine $\mathcal{O}_{Y, y}$,
le point fermé de $\Spec(A')$ s'envoie sur
un point $x \in X$ tel que $f(x) = y$ et qui est une spécialisation de $x'$.
Ainsi $x \in T$, puisque $T$ est fermé, et donc $y \in f(T)$, comme voulu.
\end{proof}

\noindent
Le lemme suivant sera généralisé dans Espaces décents, lemme
\ref{decent-spaces-lemma-re-characterize-universally-closed}.

\begin{lemma}
\label{lemma-characterize-universally-closed-quasi-separated}
Soit $S$ un schéma. Soit $f : X \to Y$ un
morphisme d'espaces algébriques sur $S$.
\begin{enumerate}
\item Si $f$ est quasi-séparé et universellement fermé, alors
$f$ satisfait la partie d'existence du critère valuatif.
\item Si $f$ est quasi-compact et quasi-séparé, alors
$f$ est universellement fermé si et seulement si la partie d'existence du
critère valuatif est satisfaite.
\end{enumerate}
\end{lemma}

\begin{proof}
Si (1) est vrai, sa combinaison avec le
lemme \ref{lemma-quasi-compact-existence-universally-closed}
donne (2). Supposons que $f$ soit quasi-séparé et universellement fermé.
Considérons un diagramme
$$
\xymatrix{
\Spec(K) \ar[r] \ar[d] & X \ar[d] \\
\Spec(A) \ar[r] & Y
}
$$
comme dans la définition \ref{definition-valuative-criterion}.
Un argument formel montre que l'existence du diagramme souhaité
$$
\xymatrix{
\Spec(K') \ar[r] \ar[d] & \Spec(K) \ar[r] & X \ar[d] \\
\Spec(A') \ar[r] \ar[rru] & \Spec(A) \ar[r] & Y
}
$$
résulte de l'existence d'un tel diagramme pour le morphisme $X_A \to \Spec(A)$.
Puisque les propriétés d'être quasi-séparé et universellement fermé sont préservées
par changement de base, le lemme résulte de ce qui est démontré au paragraphe suivant.

\medskip\noindent
Considérons un diagramme en traits pleins
$$
\xymatrix{
\Spec(K) \ar[r]_-x \ar[d] & X \ar[d]^f \\
\Spec(A) \ar@{=}[r] \ar@{..>}[ru] & \Spec(A)
}
$$
où $A$ est un anneau de valuation de corps des fractions $K$.
D'après le lemme \ref{lemma-quasi-compact-permanence} et puisque
$f$ est quasi-séparé, le
morphisme $x$ est quasi-compact.
Comme $f$ est universellement fermé,
$|f|(\overline{\{x\}})$ est en particulier fermé dans $\Spec(A)$.
Puisque cette image contient le point générique de $\Spec(A)$,
il existe un point $x' \in |X|$, dans l'adhérence
de $x$, qui s'envoie sur le point fermé de $\Spec(A)$.
D'après le lemme \ref{lemma-reach-points-scheme-theoretic-image},
on peut trouver un diagramme commutatif
$$
\xymatrix{
\Spec(K') \ar[r] \ar[d] & \Spec(K) \ar[d] \\
\Spec(A') \ar[r] & X
}
$$
tel que le point fermé de $\Spec(A')$ s'envoie sur $x' \in |X|$.
Il s'ensuit que $\Spec(A') \to \Spec(A)$ envoie le point fermé
sur le point fermé ; autrement dit, $A'$ domine $A$, ce qui achève la preuve.
\end{proof}

\begin{lemma}
\label{lemma-characterize-universally-closed-separated}
Soit $S$ un schéma. Soit $f : X \to Y$ un morphisme d'espaces algébriques
sur $S$. Supposons que $f$ soit quasi-compact et séparé. Alors les assertions suivantes
sont équivalentes :
\begin{enumerate}
\item $f$ est universellement fermé ;
\item la partie d'existence du critère valuatif est satisfaite
au sens de la définition \ref{definition-valuative-criterion} ;
\item pour tout diagramme commutatif en traits pleins
$$
\xymatrix{
\Spec(K) \ar[r] \ar[d] & X \ar[d] \\
\Spec(A) \ar[r] \ar@{-->}[ru] & Y
}
$$
où $A$ est un anneau de valuation de corps des fractions $K$, il existe
une flèche pointillée ; autrement dit, $f$ satisfait la partie d'existence du critère
valuatif au sens de
Schémas, définition \ref{schemes-definition-valuative-criterion}.
\end{enumerate}
\end{lemma}

\begin{proof}
Puisque $f$ est séparé, les assertions (2) et (3) sont équivalentes d'après le
lemme \ref{lemma-usual-enough}.
L'équivalence de (3) et (1) résulte du
lemme \ref{lemma-characterize-universally-closed-quasi-separated}.
\end{proof}

\begin{lemma}
\label{lemma-lift-valuation-ring-through-flat-morphism}
Soit $S$ un schéma. Soit $f : X \to Y$ un morphisme plat
d'espaces algébriques sur $S$.
Soit $\Spec(A) \to Y$ un morphisme, où $A$ est un
anneau de valuation. Si le point fermé de $\Spec(A)$ s'envoie sur un
point de $|Y|$ appartenant à l'image de $|X| \to |Y|$, alors il existe
un diagramme commutatif
$$
\xymatrix{
\Spec(A') \ar[r] \ar[d] & X \ar[d] \\
\Spec(A) \ar[r] & Y
}
$$
où $A \to A'$ est une extension d'anneaux de valuation
(Compléments d'algèbre, définition
\ref{more-algebra-definition-extension-valuation-rings}).
\end{lemma}

\begin{proof}
Le changement de base $X_A \to \Spec(A)$ est plat
(lemme \ref{lemma-base-change-flat}) et le point fermé de
$\Spec(A)$ appartient à l'image de $|X_A| \to |\Spec(A)|$
(Propriétés des espaces, lemme \ref{spaces-properties-lemma-points-cartesian}).
On peut donc supposer que $Y = \Spec(A)$. Soit $U \to X$ un morphisme
étale surjectif, où $U$ est un schéma. Soit $u \in U$ au-dessus
du point fermé de $\Spec(A)$. Considérons l'homomorphisme local plat
$A \to B = \mathcal{O}_{U, u}$. D'après
Algèbre, lemme \ref{algebra-lemma-ff-rings},
il existe un idéal premier $\mathfrak q \subset B$ tel que
$\mathfrak q$ soit au-dessus de $(0) \subset A$. D'après
Algèbre, lemme \ref{algebra-lemma-dominate},
on peut trouver un anneau de valuation $A' \subset \kappa(\mathfrak q)$
qui domine $B/\mathfrak q$. Le morphisme induit
$\Spec(A') \to U \to X$ résout le problème
posé par le lemme.
\end{proof}

\begin{lemma}
\label{lemma-refined-valuative-criterion-universally-closed}
Soit $S$ un schéma. Soient $f : X \to Y$ et $h : U \to X$ des morphismes
d'espaces algébriques sur $S$. Si
\begin{enumerate}
\item $f$ et $h$ sont quasi-compacts,
\item $|h|(|U|)$ est dense dans $|X|$,
\end{enumerate}
et si, pour tout diagramme commutatif en traits pleins
$$
\xymatrix{
\Spec(K) \ar[r] \ar[d] & U \ar[r] & X \ar[d] \\
\Spec(A) \ar[rr] \ar@{-->}[rru] & & Y
}
$$
où $A$ est un anneau de valuation de corps des fractions $K$,
\begin{enumerate}
\item[(3)] il existe au plus une flèche pointillée rendant le diagramme
commutatif, et
\item[(4)] il existe une extension de corps $K'/K$, un
anneau de valuation $A' \subset K'$ qui domine $A$ et un morphisme
$\Spec(A') \to X$ tels que le diagramme suivant soit commutatif
$$
\xymatrix{
\Spec(K') \ar[r] \ar[d] & \Spec(K) \ar[r] & U \ar[r] & X \ar[d] \\
\Spec(A') \ar[r] \ar[rrru] & \Spec(A) \ar[rr] & & Y
}
$$
\end{enumerate}
alors $f$ est universellement fermé. Si, de plus,
\begin{enumerate}
\item[(5)] $f$ est quasi-séparé,
\end{enumerate}
alors $f$ est séparé et universellement fermé.
\end{lemma}

\begin{proof}
Supposons (1), (2), (3) et (4).
Nous allons vérifier la partie d'existence du critère valuatif pour $f$,
ce qui entraînera que $f$ est universellement fermé d'après le
lemme \ref{lemma-quasi-compact-existence-universally-closed}.
Pour cela, considérons un diagramme commutatif
\begin{equation}
\label{equation-start-with}
\vcenter{
\xymatrix{
\Spec(K) \ar[r] \ar[d] & X \ar[d] \\
\Spec(A) \ar[r] & Y
}
}
\end{equation}
où $A$ est un anneau de valuation et $K$ le corps des fractions de $A$.
Notons que, les anneaux de valuation et les corps étant réduits, on peut
remplacer $U$, $X$ et $Y$ par leurs réductions respectives d'après
Propriétés des espaces, lemme \ref{spaces-properties-lemma-map-into-reduction}.
Dans ce cas, l'hypothèse que $h(U)$ est dense signifie que
l'image schématique de $h : U \to X$ est $X$ ; voir le
lemme \ref{lemma-scheme-theoretic-image-reduced}.

\medskip\noindent
Réduction au cas où $Y$ est affine. Choisissons un morphisme étale
$\Spec(R) \to Y$ tel que l'image du point fermé de $\Spec(A)$
appartienne à $\Im(|\Spec(R)| \to |Y|)$. D'après le
lemme \ref{lemma-lift-valuation-ring-through-flat-morphism},
on peut trouver un homomorphisme local $A \to A'$ d'anneaux de valuation
et un morphisme $\Spec(A') \to \Spec(R)$ s'insérant dans un
diagramme commutatif
$$
\xymatrix{
\Spec(A') \ar[r] \ar[d] & \Spec(R) \ar[d] \\
\Spec(A) \ar[r] & Y
}
$$
Puisque la définition \ref{definition-valuative-criterion}
autorise les extensions d'anneaux de valuation,
on peut clairement remplacer $A$ par $A'$, $Y$ par $\Spec(R)$,
$X$ par $X \times_Y \Spec(R)$ et $U$ par $U \times_Y \Spec(R)$.

\medskip\noindent
Nous supposons désormais que $Y = \Spec(R)$ est un schéma affine.
Soit $\Spec(B) \to X$ un morphisme étale dont la source est un schéma affine
tel que le morphisme $\Spec(K) \to X$ appartienne à l'image de
$|\Spec(B)| \to |X|$. Puisque l'on peut remplacer $K$ par une extension
$K' \supset K$ et $A$ par un anneau de valuation $A' \subset K'$
qui domine $A$ (un tel anneau existe d'après
Algèbre, lemme \ref{algebra-lemma-dominate}),
on peut supposer que le morphisme $\Spec(K) \to X$ se factorise par $\Spec(B)$
(par définition de $|X|$). Autrement dit, on peut considérer $K$ comme une $B$-algèbre.
Choisissons une algèbre polynomiale $P$ sur $B$ et une surjection de $B$-algèbres
$P \to K$. Alors $\Spec(P) \to X$ est plat, comme composée
$\Spec(P) \to \Spec(B) \to X$. Par conséquent, l'image schématique
du morphisme $U \times_X \Spec(P) \to \Spec(P)$ est $\Spec(P)$ d'après le
lemme \ref{lemma-flat-base-change-scheme-theoretic-image}.
D'après le lemme \ref{lemma-reach-points-scheme-theoretic-image},
on peut trouver un diagramme commutatif
$$
\xymatrix{
\Spec(K') \ar[r] \ar[d] & U \times_X \Spec(P) \ar[d] \\
\Spec(A') \ar[r] & \Spec(P)
}
$$
où $A'$ est un anneau de valuation et $K'$ le corps des fractions de $A'$ ;
le point fermé de $\Spec(A')$ s'envoie sur
$\Spec(K) \subset \Spec(P)$. Autrement dit, il existe un homomorphisme de $B$-algèbres
$\varphi : K \to A'/\mathfrak m_{A'}$. Choisissons un anneau de valuation
$A'' \subset A'/\mathfrak m_{A'}$ qui domine $\varphi(A)$ et dont le
corps des fractions est $K'' = A'/\mathfrak m_{A'}$
(Algèbre, lemme \ref{algebra-lemma-dominate}). Posons
$$
C = \{\lambda \in A' \mid \lambda \bmod \mathfrak m_{A'} \in A''\}
$$
qui est un anneau de valuation d'après
Algèbre, lemme \ref{algebra-lemma-stack-valuation-rings}.
Comme $C$ est une $R$-algèbre de corps des fractions $K'$, on obtient un diagramme
commutatif en traits pleins
$$
\xymatrix{
\Spec(K'_1) \ar@{-->}[r] \ar@{-->}[d] &
\Spec(K') \ar[r] \ar[d] & U \ar[r] & X \ar[d] \\
\Spec(C_1) \ar@{-->}[r] \ar@{-->}[rrru] & \Spec(C) \ar[rr] & & Y
}
$$
comme dans l'énoncé du lemme. L'hypothèse (4) fournit donc
$C \to C_1$ et les flèches pointillées qui rendent le diagramme commutatif.
Soit $A_1' = (C_1)_\mathfrak p$ le localisé de $C_1$
en un idéal premier $\mathfrak p \subset C_1$ situé au-dessus de
$\mathfrak m_{A'} \subset C$. Comme $C \to C_1$ est plat d'après
Compléments d'algèbre, lemme
\ref{more-algebra-lemma-valuation-ring-torsion-free-flat}
un tel idéal premier $\mathfrak p$ existe d'après
Algèbre, lemmes \ref{algebra-lemma-local-flat-ff} et
\ref{algebra-lemma-ff-rings}.
Notons que $A'$ est le localisé de $C$ en $\mathfrak m_{A'}$
et que $A'_1$ est un anneau de valuation
(Algèbre, lemme \ref{algebra-lemma-make-valuation-rings}).
Autrement dit, $A' \to A'_1$ est un homomorphisme local d'anneaux de
valuation. L'hypothèse (3) implique que le diagramme
$$
\xymatrix{
\Spec(A'_1) \ar[r] \ar[d] & \Spec(C_1) \ar[r] & X \\
\Spec(A') \ar[r] & \Spec(P) \ar[r] & \Spec(B) \ar[u]
}
$$
est commutatif. Ainsi, le morphisme induit par $\Spec(C_1) \to X$
sur $\Spec(C_1/\mathfrak p)$ a pour restriction la composée
$$
\Spec(\kappa(\mathfrak p)) \to
\Spec(A'/\mathfrak m_{A'}) = \Spec(K'') \to
\Spec(K) \to X
$$
au point générique de $\Spec(C_1/\mathfrak p)$. De plus,
$C_1/\mathfrak p$ est un anneau de valuation
(Algèbre, lemme \ref{algebra-lemma-make-valuation-rings})
qui domine $A''$, lequel domine $A$.
Ainsi, le morphisme $\Spec(C_1/\mathfrak p) \to X$ fournit la solution exigée par la
partie d'existence du critère valuatif pour le diagramme
(\ref{equation-start-with}), comme voulu.

\medskip\noindent
Supposons enfin que (5) soit également satisfaite ; autrement dit, que le morphisme
$\Delta : X \to X \times_Y X$ soit quasi-compact. Dans ce cas,
les hypothèses (1) -- (4) sont satisfaites pour $h$ et $\Delta$. La première
partie de la preuve montre donc que $\Delta$ est universellement fermé.
D'après le lemme \ref{lemma-separated-diagonal-proper}, on en conclut que
$f$ est séparé.
\end{proof}

















\section{Critère valuatif de séparation}
\label{section-valuative-separatedness}

\noindent
Nous démontrons d'abord une réciproque, puis nous énonçons le critère.

\begin{lemma}
\label{lemma-separated-implies-valuative}
Soit $S$ un schéma.
Soit $f : X \to Y$ un morphisme d'espaces algébriques sur $S$.
Si $f$ est séparé, alors $f$ satisfait la partie
d'unicité du critère valuatif.
\end{lemma}

\begin{proof}
Considérons un diagramme comme dans la définition \ref{definition-valuative-criterion}.
Supposons qu'il existe deux morphismes distincts
$a, b : \Spec(A) \to X$ qui complètent ce diagramme.
Soit $Z \subset \Spec(A)$ l'égalisateur de $a$ et $b$.
Alors $Z = \Spec(A) \times_{(a, b), X \times_Y X, \Delta} X$.
Si $f$ est séparé, alors $\Delta$ est une immersion fermée, et
C'est un sous-schéma fermé de $\Spec(A)$. Par hypothèse, il contient
le point générique de $\Spec(A)$. Comme $A$ est intègre,
cela implique $Z = \Spec(A)$. Ainsi $a = b$, comme voulu.
\end{proof}

\begin{lemma}[Critère valuatif de séparation]
\label{lemma-valuative-criterion-separatedness}
Soit $S$ un schéma.
Soit $f : X \to Y$ un morphisme d'espaces algébriques sur $S$.
Supposons que
\begin{enumerate}
\item le morphisme $f$ soit quasi-séparé, et que
\item le morphisme $f$ satisfasse la partie d'unicité
du critère valuatif.
\end{enumerate}
Alors $f$ est séparé.
\end{lemma}

\begin{proof}
L'hypothèse (1) signifie que $\Delta_{X/Y}$ est quasi-compact.
Nous affirmons que le morphisme
$\Delta_{X/Y} : X \to X \times_Y X$ satisfait la partie d'existence
du critère valuatif.
Considérons un diagramme commutatif en traits pleins
$$
\xymatrix{
\Spec(K) \ar[r] \ar[d] & X \ar[d] \\
\Spec(A) \ar[r] \ar@{-->}[ru] & X \times_Y X
}
$$
La flèche horizontale inférieure correspond à une
paire de morphismes $a, b : \Spec(A) \to X$ sur $Y$.
D'après l'hypothèse (2), on a $a = b$. Prendre $a$ pour flèche
pointillée convient. Par conséquent, le
lemme \ref{lemma-quasi-compact-existence-universally-closed}
s'applique, et $\Delta_{X/Y}$ est universellement fermé.
Puisque $\Delta_{X/Y}$ est toujours localement de type fini et
séparé, il résulte de
Compléments sur les morphismes, lemme \ref{more-morphisms-lemma-characterize-finite}
que $\Delta_{X/Y}$ est un morphisme fini (on peut aussi utiliser le
principe général de
Espaces, lemme
\ref{spaces-lemma-representable-transformations-property-implication}).
À ce stade, $\Delta_{X/Y}$ est un monomorphisme représentable et fini,
donc une immersion fermée d'après
Morphismes, lemme \ref{morphisms-lemma-finite-monomorphism-closed}.
\end{proof}

\begin{lemma}
\label{lemma-characterize-separated-and-universally-closed}
Soit $S$ un schéma. Soit $f : X \to Y$ un morphisme d'espaces algébriques
sur $S$. Supposons que $f$ soit quasi-compact et quasi-séparé. Alors les
assertions suivantes sont équivalentes :
\begin{enumerate}
\item $f$ est séparé et universellement fermé ;
\item le critère valuatif est satisfait au sens de la définition
\ref{definition-valuative-criterion},
\item pour tout diagramme commutatif en traits pleins
$$
\xymatrix{
\Spec(K) \ar[r] \ar[d] & X \ar[d] \\
\Spec(A) \ar[r] \ar@{-->}[ru] & Y
}
$$
où $A$ est un anneau de valuation de corps des fractions $K$, il existe
une unique flèche pointillée ; autrement dit, $f$ satisfait le critère
valuatif au sens de
Schémas, définition \ref{schemes-definition-valuative-criterion}.
\end{enumerate}
\end{lemma}

\begin{proof}
Puisque $f$ est quasi-séparé, la partie d'unicité
du critère valuatif implique que $f$ est séparé
(lemme \ref{lemma-valuative-criterion-separatedness}).
Réciproquement, si $f$ est séparé, il satisfait la
partie d'unicité du critère valuatif
(lemme \ref{lemma-separated-implies-valuative}).
Cela étant, dans chacun des trois cas, le
morphisme $f$ est séparé et satisfait la partie d'unicité
du critère valuatif. Dans ce cas, le lemme est une conséquence
formelle du
lemme \ref{lemma-characterize-universally-closed-separated}.
\end{proof}



\section{Critère valuatif de propreté}
\label{section-valuative-criterion-properness}

\noindent
Voici l'énoncé.

\begin{lemma}[Critère valuatif de propreté]
\label{lemma-characterize-proper}
Soit $S$ un schéma. Soit $f : X \to Y$ un morphisme d'espaces algébriques
sur $S$. Supposons que $f$ soit de type fini et quasi-séparé. Alors les
assertions suivantes sont équivalentes :
\begin{enumerate}
\item $f$ est propre ;
\item le critère valuatif est satisfait au sens de la définition
\ref{definition-valuative-criterion},
\item pour tout diagramme commutatif en traits pleins
$$
\xymatrix{
\Spec(K) \ar[r] \ar[d] & X \ar[d] \\
\Spec(A) \ar[r] \ar@{-->}[ru] & Y
}
$$
où $A$ est un anneau de valuation de corps des fractions $K$, il existe
une unique flèche pointillée ; autrement dit, $f$ satisfait le critère
valuatif au sens de
Schémas, définition \ref{schemes-definition-valuative-criterion}.
\end{enumerate}
\end{lemma}

\begin{proof}
Cela résulte formellement du
lemme \ref{lemma-characterize-separated-and-universally-closed}
et des définitions.
\end{proof}





\section{Morphismes entiers et finis}
\label{section-integral}

\noindent
Nous avons déjà défini dans la section \ref{section-representable}
ce que signifie, pour un morphisme représentable d'espaces algébriques,
être entier (resp.\ fini).

\begin{lemma}
\label{lemma-integral-representable}
Soit $S$ un schéma. Soit $f : X \to Y$ un morphisme représentable
d'espaces algébriques sur $S$. Alors
$f$ est entier, resp.\ fini
(au sens de la section \ref{section-representable}),
si et seulement si, pour tout schéma affine $Z$
et tout morphisme $Z \to Y$, le schéma $X \times_Y Z$ est affine et
entier, resp.\ fini, sur $Z$.
\end{lemma}

\begin{proof}
Cela résulte directement de la définition d'un morphisme entier (resp.\ fini)
de schémas
(Morphismes, définition \ref{morphisms-definition-integral}).
\end{proof}

\noindent
Cela permet de poser la définition suivante.

\begin{definition}
\label{definition-integral}
Soit $S$ un schéma.
Soit $f : X \to Y$ un morphisme d'espaces algébriques sur $S$.
\begin{enumerate}
\item Nous disons que $f$ est {\it entier} si, pour tout schéma affine $Z$
et tout morphisme $Z \to Y$, l'espace algébrique $X \times_Y Z$ est
représentable par un schéma affine entier sur $Z$.
\item Nous disons que $f$ est {\it fini} si, pour tout schéma affine $Z$
et tout morphisme $Z \to Y$, l'espace algébrique $X \times_Y Z$ est
représentable par un schéma affine fini sur $Z$.
\end{enumerate}
\end{definition}

\begin{lemma}
\label{lemma-integral-local}
Soit $S$ un schéma.
Soit $f : X \to Y$ un morphisme d'espaces algébriques sur $S$.
Les assertions suivantes sont équivalentes :
\begin{enumerate}
\item $f$ est représentable et entier (resp.\ fini) ;
\item $f$ est entier (resp.\ fini) ;
\item il existe un schéma $V$ et un morphisme étale surjectif
$V \to Y$ tel que $V \times_Y X \to V$ soit entier (resp. fini) ;
\item il existe un recouvrement de Zariski $Y = \bigcup Y_i$ tel que
chacun des morphismes $f^{-1}(Y_i) \to Y_i$ soit entier (resp.\ fini).
\end{enumerate}
\end{lemma}

\begin{proof}
Il est clair que (1) implique (2) et que (2) implique (3) en prenant pour
$V$ une somme disjointe de schémas affines étales sur $Y$ ; voir
Propriétés des espaces,
lemme \ref{spaces-properties-lemma-cover-by-union-affines}.
Supposons que $V \to Y$ soit comme en (3). Pour tout ouvert affine $W$ de $V$,
$W \times_Y X$ est alors un ouvert affine de $V \times_Y X$. D'après
Propriétés des espaces, lemme \ref{spaces-properties-lemma-subscheme}
on en conclut que $V \times_Y X$ est un schéma. De plus, le morphisme
$V \times_Y X \to V$ est affine. On peut donc appliquer
Espaces,
lemme \ref{spaces-lemma-morphism-sheaves-with-P-effective-descent-etale}
car la classe des morphismes entiers (resp.\ finis)
satisfait toutes les propriétés requises (voir
Morphismes, lemmes \ref{morphisms-lemma-base-change-finite} et
Descente, lemmes
\ref{descent-lemma-descending-property-integral},
\ref{descent-lemma-descending-property-finite},
et \ref{descent-lemma-affine}).
La conclusion de ce lemme est que $f$ est représentable
et entier (resp.\ fini), c'est-à-dire que (1) est satisfaite.

\medskip\noindent
L'équivalence de (1) et (4) résulte du fait que la propriété d'être
entier (resp.\ fini) est locale pour la topologie de Zariski sur le but (la
référence ci-dessus montre qu'en fait, la propriété d'être entier ou fini est
locale pour la topologie fpqc sur le but).
\end{proof}

\begin{lemma}
\label{lemma-composition-integral}
La composée de morphismes entiers (resp.\ finis) est entière
(resp.\ finie).
\end{lemma}

\begin{proof}
Omis.
\end{proof}

\begin{lemma}
\label{lemma-base-change-integral}
Tout changement de base d'un morphisme entier (resp.\ fini) est entier
(resp.\ fini).
\end{lemma}

\begin{proof}
Omis.
\end{proof}

\begin{lemma}
\label{lemma-finite-integral}
Un morphisme fini d'espaces algébriques est entier.
Un morphisme entier d'espaces algébriques
qui est localement de type fini est fini.
\end{lemma}

\begin{proof}
Dans les deux cas, le morphisme est représentable, et la condition se vérifie
après un changement de base par un schéma affine muni d'un morphisme vers $Y$ ; voir le
lemme \ref{lemma-integral-local}. Le résultat découle donc du
lemme analogue pour les schémas ; voir
Morphismes, lemme \ref{morphisms-lemma-finite-integral}.
\end{proof}

\begin{lemma}
\label{lemma-integral-universally-closed}
Soit $S$ un schéma.
Soit $f : X \to Y$ un morphisme d'espaces algébriques sur $S$.
Les assertions suivantes sont équivalentes :
\begin{enumerate}
\item $f$ est entier ;
\item $f$ est affine et universellement fermé.
\end{enumerate}
\end{lemma}

\begin{proof}
Dans les deux cas, le morphisme est représentable, et la condition se vérifie
après un changement de base par un schéma affine muni d'un morphisme vers $Y$ ; voir les
lemmes \ref{lemma-integral-local},
\ref{lemma-affine-local} et
\ref{lemma-universally-closed-local}.
Le résultat découle donc de
Morphismes, lemme \ref{morphisms-lemma-integral-universally-closed}.
\end{proof}

\begin{lemma}
\label{lemma-finite-quasi-finite}
Un morphisme fini d'espaces algébriques est quasi-fini.
\end{lemma}

\begin{proof}
Soit $f : X \to Y$ un morphisme d'espaces algébriques.
D'après la
définition \ref{definition-integral}
et les
lemmes \ref{lemma-quasi-compact-local} et
\ref{lemma-quasi-finite-local}
les deux propriétés se vérifient après changement de base à un schéma affine sur $Y$ ;
autrement dit, on peut supposer $Y$ affine.
Si $f$ est fini, alors $X$ est un schéma.
Le résultat découle donc du résultat correspondant pour les schémas ; voir
Morphismes, lemme \ref{morphisms-lemma-finite-quasi-finite}.
\end{proof}

\begin{lemma}
\label{lemma-finite-proper}
Soit $S$ un schéma. Soit $f : X \to Y$ un morphisme d'espaces algébriques sur
$S$. Les assertions suivantes sont équivalentes :
\begin{enumerate}
\item $f$ est fini ;
\item $f$ est affine et propre.
\end{enumerate}
\end{lemma}

\begin{proof}
Dans les deux cas, le morphisme est représentable, et la condition se vérifie
après changement de base à un schéma affine muni d'un morphisme vers $Y$ ; voir les
lemmes \ref{lemma-integral-local}, \ref{lemma-affine-local} et
\ref{lemma-proper-local}. Le résultat découle donc de
Morphismes, lemme \ref{morphisms-lemma-finite-proper}.
\end{proof}

\begin{lemma}
\label{lemma-closed-immersion-finite}
Une immersion fermée est finie (et a fortiori entière).
\end{lemma}

\begin{proof}
Omis.
\end{proof}

\begin{lemma}
\label{lemma-finite-union-finite}
Soit $S$ un schéma.
Soient $X_i \to Y$, $i = 1, \ldots, n$, des morphismes finis
d'espaces algébriques sur $S$.
Alors $X_1 \amalg \ldots \amalg X_n \to Y$ est également fini.
\end{lemma}

\begin{proof}
Cela résulte du cas des schémas
(Morphismes, lemme \ref{morphisms-lemma-finite-union-finite})
par localisation étale.
\end{proof}

\begin{lemma}
\label{lemma-finite-permanence}
Soit $S$ un schéma.
Soient $f : X \to Y$ et $g : Y \to Z$ des morphismes d'espaces algébriques sur $S$.
\begin{enumerate}
\item Si $g \circ f$ est fini et si $g$ est séparé, alors $f$ est fini.
\item Si $g \circ f$ est entier et si $g$ est séparé, alors $f$ est entier.
\end{enumerate}
\end{lemma}

\begin{proof}
Supposons que $g \circ f$ soit fini (resp.\ entier) et que $g$ soit séparé.
Le changement de base $X \times_Z Y \to Y$ est fini (resp.\ entier) d'après le
lemme \ref{lemma-base-change-integral}.
Le morphisme $X \to X \times_Z Y$ est
une immersion fermée puisque $Y \to Z$ est séparé ; voir le
lemme \ref{lemma-section-immersion}.
Une immersion fermée est finie (resp.\ entière) ;
voir le lemme \ref{lemma-closed-immersion-finite}.
La composée de morphismes finis (resp.\ entiers) est finie
(resp.\ entière) ;
voir le lemme \ref{lemma-composition-integral}. D'où le résultat.
\end{proof}





\section{Morphismes finis localement libres}
\label{section-finite-locally-free}

\noindent
Nous avons déjà défini dans la section \ref{section-representable}
ce que signifie, pour un morphisme représentable d'espaces algébriques,
être fini localement libre.

\begin{lemma}
\label{lemma-finite-locally-free-representable}
Soit $S$ un schéma. Soit $f : X \to Y$ un morphisme représentable
d'espaces algébriques sur $S$. Alors $f$ est fini localement libre
(au sens de la section \ref{section-representable})
si et seulement si $f$ est affine et si le faisceau $f_*\mathcal{O}_X$ est
un $\mathcal{O}_Y$-module fini localement libre.
\end{lemma}

\begin{proof}
Supposons que $f$ soit fini localement libre (au sens de la
section \ref{section-representable}). Cela signifie que,
pour tout morphisme $V \to Y$ dont la source est un schéma, le
changement de base $f' : V \times_Y X \to V$ est un morphisme fini localement libre
de schémas. Cela signifie à son tour (par définition d'un morphisme fini
localement libre de schémas) que
$f'_*\mathcal{O}_{V \times_Y X}$
est un $\mathcal{O}_V$-module fini localement libre. On peut choisir $V \to Y$
étale et surjectif. D'après
Propriétés des espaces,
lemme \ref{spaces-properties-lemma-pushforward-etale-base-change-modules}
la restriction de $f_*\mathcal{O}_X$ à $V$ est
finie localement libre. Par conséquent, d'après
Modules sur les sites, lemme \ref{sites-modules-lemma-local-final-object}
appliqué au faisceau $f_*\mathcal{O}_X$ sur $Y_{spaces, \etale}$,
on en conclut que $f_*\mathcal{O}_X$ est fini localement libre.

\medskip\noindent
Réciproquement, supposons que $f$ soit affine et que $f_*\mathcal{O}_X$ soit un
$\mathcal{O}_Y$-module fini localement libre. Soit $V$ un schéma et soit
$V \to Y$ un morphisme étale surjectif. De nouveau, d'après
Propriétés des espaces,
lemme \ref{spaces-properties-lemma-pushforward-etale-base-change-modules}
$f'_*\mathcal{O}_{V \times_Y X}$ est fini localement libre.
Ainsi, $f' : V \times_Y X \to V$ est fini localement libre (puisqu'il est aussi affine).
D'après
Espaces,
lemme \ref{spaces-lemma-morphism-sheaves-with-P-effective-descent-etale}
on en conclut que $f$ est fini localement libre (utiliser
Morphismes, lemme \ref{morphisms-lemma-base-change-finite-locally-free}
et Descente, lemmes \ref{descent-lemma-descending-property-finite-locally-free}
et \ref{descent-lemma-affine}). D'où le résultat.
\end{proof}

\noindent
Cela permet de poser la définition suivante.

\begin{definition}
\label{definition-finite-locally-free}
Soit $S$ un schéma.
Soit $f : X \to Y$ un morphisme d'espaces algébriques sur $S$.
Nous disons que $f$ est {\it fini localement libre} si $f$ est affine
et si $f_*\mathcal{O}_X$ est un $\mathcal{O}_Y$-module fini localement libre.
Dans ce cas, nous disons que $f$ est
de {\it rang} ou de {\it degré} $d$
si le faisceau $f_*\mathcal{O}_X$ est fini localement libre de rang $d$.
\end{definition}

\begin{lemma}
\label{lemma-finite-locally-free-local}
Soit $S$ un schéma.
Soit $f : X \to Y$ un morphisme d'espaces algébriques sur $S$.
Les assertions suivantes sont équivalentes :
\begin{enumerate}
\item $f$ est représentable et fini localement libre ;
\item $f$ est fini localement libre ;
\item il existe un schéma $V$ et un morphisme étale surjectif
$V \to Y$ tel que $V \times_Y X \to V$ soit fini localement libre ;
\item il existe un recouvrement de Zariski $Y = \bigcup Y_i$ tel que
chaque morphisme $f^{-1}(Y_i) \to Y_i$ soit fini localement libre.
\end{enumerate}
\end{lemma}

\begin{proof}
Il est clair que (1) implique (2) et que (2) implique (3) en prenant pour
$V$ une somme disjointe de schémas affines étales sur $Y$ ; voir
Propriétés des espaces,
lemme \ref{spaces-properties-lemma-cover-by-union-affines}.
Supposons que $V \to Y$ soit comme en (3). Pour tout ouvert affine $W$ de $V$,
$W \times_Y X$ est alors un ouvert affine de $V \times_Y X$. D'après
Propriétés des espaces, lemme \ref{spaces-properties-lemma-subscheme}
on en conclut que $V \times_Y X$ est un schéma. De plus, le morphisme
$V \times_Y X \to V$ est affine. On peut donc appliquer
Espaces,
lemme \ref{spaces-lemma-morphism-sheaves-with-P-effective-descent-etale}
car la classe des morphismes finis localement libres
satisfait toutes les propriétés requises (voir
Morphismes, lemme \ref{morphisms-lemma-base-change-finite-locally-free}
et Descente, lemmes \ref{descent-lemma-descending-property-finite-locally-free}
et \ref{descent-lemma-affine}).
La conclusion de ce lemme est que $f$ est représentable
et fini localement libre, c'est-à-dire que (1) est satisfaite.

\medskip\noindent
L'équivalence de (1) et (4) résulte du fait que la propriété d'être
fini localement libre est locale pour la topologie de Zariski sur le but (la référence ci-dessus montre
qu'en fait, la propriété d'être fini localement libre est locale pour la topologie fpqc sur le but).
\end{proof}

\begin{lemma}
\label{lemma-composition-finite-locally-free}
La composée de morphismes finis localement libres est finie localement libre.
\end{lemma}

\begin{proof}
Omis.
\end{proof}

\begin{lemma}
\label{lemma-base-change-finite-locally-free}
Tout changement de base d'un morphisme fini localement libre est fini localement libre.
\end{lemma}

\begin{proof}
Omis.
\end{proof}

\begin{lemma}
\label{lemma-finite-flat}
Soit $S$ un schéma.
Soit $f : X \to Y$ un morphisme d'espaces algébriques sur $S$.
Les assertions suivantes sont équivalentes :
\begin{enumerate}
\item $f$ est fini localement libre ;
\item $f$ est fini, plat et localement de présentation finie.
\end{enumerate}
Si $Y$ est localement noethérien, ces assertions sont aussi équivalentes à
\begin{enumerate}
\item[(3)] $f$ est fini et plat.
\end{enumerate}
\end{lemma}

\begin{proof}
Dans chacun des trois cas, le morphisme est représentable et la propriété
se vérifie après changement de base par un morphisme étale surjectif
$V \to Y$ ; voir les
lemmes \ref{lemma-integral-local},
\ref{lemma-finite-locally-free-local},
\ref{lemma-flat-local} et
\ref{lemma-finite-presentation-local}.
Si $Y$ est localement noethérien, alors $V$ est localement noethérien.
Le résultat découle donc du résultat correspondant
pour les schémas ; voir
Morphismes, lemme \ref{morphisms-lemma-finite-flat}.
\end{proof}





\section{Applications rationnelles}
\label{section-rational-maps}

\noindent
Cette section est l'analogue de
Morphismes, section \ref{morphisms-section-rational-maps}.
Nous utiliserons sans autre mention le fait que l'intersection d'ouverts denses
d'un espace topologique est un ouvert dense.

\begin{definition}
\label{definition-rational-map}
Soit $S$ un schéma. Soient $X$, $Y$ des espaces algébriques sur $S$.
\begin{enumerate}
\item Soient $f : U \to Y$, $g : V \to Y$ des morphismes d'espaces algébriques
sur $S$, définis sur des sous-espaces ouverts denses $U$, $V$ de $X$. On dit que $f$ est
{\it équivalent} à $g$ si $f|_W = g|_W$ pour un sous-espace ouvert dense
$W \subset U \cap V$.
\item Une {\it application rationnelle de $X$ dans $Y$}
est une classe d'équivalence pour la relation d'équivalence définie en (1).
\item Étant donnés des morphismes $X \to B$ et $Y \to B$ d'espaces algébriques sur $S$,
on dit qu'une application rationnelle de $X$ dans $Y$ est une
{\it $B$-application rationnelle de $X$ dans $Y$}
s'il existe un représentant $f : U \to Y$ de la classe d'équivalence
qui soit un morphisme au-dessus de $B$.
\end{enumerate}
\end{definition}

\noindent
On dira que deux morphismes $f$, $g$ comme en (1) de la définition
définissent la même application rationnelle, plutôt que de les dire équivalents.
Dans de nombreux cas considérés par la suite, les espaces algébriques
$X$ et $Y$ contiendront des sous-espaces ouverts denses $X'$ et $Y'$ qui
sont des schémas. Alors une application rationnelle de $X$ dans $Y$ est la même chose
qu'une $S$-application rationnelle de $X'$ dans $Y'$ au sens de
Morphismes, définition \ref{definition-rational-map}.
Toute la théorie développée pour les schémas s'applique alors.

\begin{definition}
\label{definition-rational-function}
Soit $S$ un schéma. Soit $X$ un espace algébrique sur $S$. Une
{\it fonction rationnelle sur $X$} est une application rationnelle de $X$ dans $\mathbf{A}^1_S$.
\end{definition}

\noindent
En considérant la discussion qui suit
Morphismes, définition \ref{morphisms-definition-rational-function},
on voit qu'il s'agit de la même notion lorsque
$X$ est un schéma.

\medskip\noindent
Rappelons l'identification canonique
$$
\Mor_S(T, \mathbf{A}^1_S) = \Mor(T, \mathbf{A}^1_\mathbf{Z}) =
\Gamma(T, \mathcal{O}_T)
$$
pour tout schéma $T$ sur $S$ ; voir
Schémas, exemple \ref{schemes-example-global-sections}.
Ainsi, les lois d'addition et de multiplication font de $\mathbf{A}^1_S$ un anneau dans la
catégorie des schémas sur $S$. Autrement dit, l'addition
et la multiplication définissent des morphismes
$$
+ : \mathbf{A}^1_S \times_S \mathbf{A}^1_S \to \mathbf{A}^1_S
\quad\text{et}\quad
* : \mathbf{A}^1_S \times_S \mathbf{A}^1_S \to \mathbf{A}^1_S
$$
qui satisfont aux axiomes de l'addition et de la multiplication dans un anneau
(commutatif, d'élément unité $1$, comme toujours). L'ensemble des applications rationnelles
à valeurs dans $\mathbf{A}^1_S$ est donc lui aussi muni d'une structure naturelle d'anneau.

\begin{definition}
\label{definition-ring-of-rational-functions}
Soit $S$ un schéma.
Soit $X$ un espace algébrique sur $S$.
L'{\it anneau des fonctions rationnelles sur $X$}
est l'anneau $R(X)$ dont les éléments sont les fonctions rationnelles, muni de
l'addition et de la multiplication qui viennent d'être décrites.
\end{definition}

\noindent
Nous définirons plus loin le corps des fonctions d'un espace algébrique intègre ; voir
Espaces sur les corps, section \ref{spaces-over-fields-section-integral-spaces}.

\begin{definition}
\label{definition-domain-of-definition}
Soit $S$ un schéma. Soit $\varphi$ une application rationnelle entre deux
espaces algébriques $X$ et $Y$ sur $S$. On dit que
$\varphi$ est {\it définie au point $x \in |X|$} s'il existe un
représentant $(U, f)$ de $\varphi$ tel que $x \in |U|$. Le
{\it domaine de définition} de $\varphi$ est l'ensemble des points où
$\varphi$ est définie.
\end{definition}

\noindent
Le domaine de définition est considéré comme un sous-espace ouvert de $X$ grâce à
Propriétés des espaces, lemme \ref{spaces-properties-lemma-open-subspaces}.
Avec cette définition, il n'est pas vrai en général que $\varphi$ possède un
représentant défini sur tout son domaine de définition.

\begin{lemma}
\label{lemma-rational-map-from-reduced-to-separated}
Soit $S$ un schéma. Soient $X$ et $Y$ des espaces algébriques sur $S$.
Supposons $X$ réduit et $Y$ séparé sur $S$. Soit
$\varphi$ une application rationnelle de $X$ dans $Y$, de domaine de définition
$U \subset X$. Il existe alors un unique morphisme $f : U \to Y$
d'espaces algébriques qui représente $\varphi$.
\end{lemma}

\begin{proof}
Soient $(V, g)$ et $(V', g')$ des représentants de $\varphi$. Alors
Les morphismes $g, g'$ coïncident sur un sous-espace ouvert dense $W \subset V \cap V'$.
D'autre part, l'égalisateur $E$ de $g|_{V \cap V'}$ et $g'|_{V \cap V'}$
est un sous-espace fermé de $V \cap V'$, car il s'obtient par changement de base
de $\Delta : Y \to Y \times_S Y$ par le morphisme
$V \cap V' \to Y \times_S Y$ défini par $g|_{V \cap V'}$ et $g'|_{V \cap V'}$.
Or $W \subset E$ implique que $|E| = |V \cap V'|$. Comme $V \cap V'$
est réduit, on en déduit $E = V \cap V'$ schématiquement, c'est-à-dire
$g|_{V \cap V'} = g'|_{V \cap V'}$ ; voir
Propriétés des espaces, lemme \ref{spaces-properties-lemma-map-into-reduction}.
On peut donc recoller les représentants $g : V \to Y$ de $\varphi$
en un morphisme $f : U \to Y$, puisque $\coprod V \to U$ est un morphisme surjectif de
faisceaux fppf et que
$\coprod_{V, V'} V \cap V' = (\coprod V) \times_U (\coprod V)$.
\end{proof}

\noindent
En général, la composition d'applications rationnelles n'a pas de sens. En effet,
l'image d'un représentant de la première application rationnelle peut
ne pas rencontrer le domaine de définition de la seconde.
Toutefois, si les espaces sont irréductibles et si l'on considère
des applications rationnelles dominantes, on peut les composer.

\begin{definition}
\label{definition-dominant-rational}
Soit $S$ un schéma. Soient $X$ et $Y$ des espaces algébriques sur $S$.
Supposons $|X|$ et $|Y|$ irréductibles. Une application rationnelle de $X$ dans $Y$
est dite {\it dominante} si tout représentant $f : U \to Y$ est un morphisme dominant
au sens de la définition \ref{definition-dominant}.
\end{definition}

\noindent
On peut composer une application rationnelle dominante
$\varphi$ entre des espaces algébriques irréductibles
$X$ et $Y$ avec une application rationnelle arbitraire
$\psi$ de $Y$ dans $Z$. Choisissons en effet des
représentants $f : U \to Y$, avec $|U| \subset |X|$ ouvert dense,
et $g : V \to Z$, avec $|V| \subset |Y|$ ouvert dense. Alors
$W = |f|^{-1}(V) \subset |X|$ est ouvert non vide (car
l'image de $|f|$ est dense et rencontre donc l'ouvert non vide $V$),
donc dense puisque $|X|$ est irréductible. On définit $\psi \circ \varphi$
comme la classe d'équivalence de $g \circ f|_W : W \to Z$. Nous omettons la vérification
que cette définition ne dépend pas des choix.

\medskip\noindent
On obtient ainsi une catégorie dont les objets sont les espaces algébriques
irréductibles sur $S$ et dont les morphismes sont les applications rationnelles dominantes.

\begin{definition}
\label{definition-birational-spaces}
Soit $S$ un schéma. Soient $X$ et $Y$ des espaces algébriques
sur $S$, avec $|X|$ et $|Y|$ irréductibles.
On dit que $X$ et $Y$ sont {\it birationnels} si $X$ et $Y$ sont isomorphes
dans la catégorie des espaces algébriques irréductibles sur $S$
et des applications rationnelles dominantes.
\end{definition}

\noindent
Si $X$ et $Y$ sont des espaces algébriques irréductibles birationnels,
l'ensemble des applications rationnelles de $X$ dans $Z$ est en bijection
avec l'ensemble des applications rationnelles de $Y$ dans $Z$, pour tout espace
algébrique $Z$ (fonctoriellement en $Z$). Pour des espaces algébriques irréductibles « généraux »,
ce n'est là qu'une définition possible. On pourrait aussi
exiger que $X$ et $Y$ aient des anneaux de fonctions
rationnelles isomorphes ; ces deux notions sont parfois équivalentes
(insérer ici un renvoi futur).

\begin{lemma}
\label{lemma-birational}
Soit $S$ un schéma. Soient $X$ et $Y$ des espaces
algébriques sur $S$, avec $|X|$ et $|Y|$ irréductibles.
Alors $X$ et $Y$ sont birationnels si et seulement s'il existe
des sous-espaces ouverts non vides $U \subset X$ et $V \subset Y$
qui sont isomorphes comme espaces algébriques sur $S$.
\end{lemma}

\begin{proof}
Supposons $X$ et $Y$ birationnels. Soient $f : U \to Y$ et $g : V \to X$
des représentants d'applications rationnelles dominantes inverses de $X$ dans $Y$ et de $Y$ dans $X$.
Après avoir rétréci $U$, on peut supposer que $f : U \to Y$ se factorise par $V$.
Comme $g \circ f$ est l'identité en tant qu'application rationnelle dominante,
la composée $U \to V \to X$ est l'identité sur un ouvert dense de $U$.
En remplaçant $U$ par un ouvert plus petit, on peut donc supposer que
$U \to V \to X$ est l'inclusion de $U$ dans $X$.
Par symétrie, il existe un sous-espace ouvert $V' \subset V$
tel que $g|_{V'} : V' \to X$ se factorise par $U \subset X$
et que $V' \to U \to Y$ soit l'identité.
L'image réciproque de $|V'|$ par $|U| \to |V|$ est un ouvert de $|U|$,
donc elle est égale à $|U'|$ pour un sous-espace ouvert $U' \subset U$ ; voir
Propriétés des espaces, lemme \ref{spaces-properties-lemma-open-subspaces}.
Alors $U' \subset U \to V$ se factorise en $U' \to V'$.
De même, $V' \to U$ se factorise en $V' \to U'$.
On vérifie que $U' \to V'$
et $V' \to U'$ sont des morphismes
d'espaces algébriques sur $S$ inverses l'un de l'autre,
ce qui achève la démonstration.
\end{proof}









\section{Normalisation relative des espaces algébriques}
\label{section-normalization-X-in-Y}

\noindent
Cette section est l'analogue de
Morphismes, section \ref{morphisms-section-normalization-X-in-Y}.

\begin{lemma}
\label{lemma-integral-closure}
Soit $S$ un schéma. Soit $X$ un espace algébrique sur $S$.
Soit $\mathcal{A}$ un faisceau quasi-cohérent de $\mathcal{O}_X$-algèbres.
Il existe un sous-faisceau quasi-cohérent de $\mathcal{O}_X$-algèbres
$\mathcal{A}' \subset \mathcal{A}$ tel que,
pour tout objet affine $U$ de $X_\etale$, l'anneau
$\mathcal{A}'(U) \subset \mathcal{A}(U)$ soit
la clôture intégrale de $\mathcal{O}_X(U)$ dans $\mathcal{A}(U)$.
\end{lemma}

\begin{proof}
Soit $U$ un objet de $X_\etale$. Alors $U$ est un schéma. Notons
$\mathcal{A}|_U$ la restriction au site de Zariski. Le faisceau
$\mathcal{A}|_U$ est un faisceau quasi-cohérent de $\mathcal{O}_U$-algèbres ;
on peut donc appliquer Morphismes, lemme \ref{morphisms-lemma-integral-closure},
pour obtenir une sous-algèbre quasi-cohérente $\mathcal{A}'_U \subset \mathcal{A}|_U$
telle que la valeur de $\mathcal{A}'_U$ sur tout ouvert affine $W \subset U$
est celle indiquée dans l'énoncé du lemme. Si $f : U' \to U$ est un morphisme
dans $X_\etale$, alors $\mathcal{A}|_{U'} = f^*(\mathcal{A}|_U)$,
où $f^*$ désigne l'image inverse par le morphisme $f$ pour la topologie de Zariski ;
cela résulte de la quasi-cohérence de $\mathcal{A}$ (voir l'introduction de
Propriétés des espaces, section \ref{spaces-properties-section-quasi-coherent},
ainsi que les renvois à la discussion du chapitre sur la
descente des schémas). Comme $f$ est étale, le lemme
Compléments sur les morphismes,
\ref{more-morphisms-lemma-integral-closure-smooth-pullback},
donne un isomorphisme canonique
$f^*(\mathcal{A}'_U) = \mathcal{A}'_{U'}$.
Cela montre aussitôt que l'on obtient un sous-préfaisceau
$\mathcal{A}' \subset \mathcal{A}$ de $\mathcal{O}_X$-algèbres sur $X_\etale$
qui est un faisceau pour la topologie de Zariski
et prend les valeurs voulues sur les objets affines.
Le fait que chaque $\mathcal{A}'_U$ soit quasi-cohérent
sur le schéma $U$ et que, pour $f : U' \to U$ étale, on ait
$\mathcal{A}'_{U'} = f^*(\mathcal{A}'_U)$ implique en outre que $\mathcal{A}'$
est quasi-cohérent sur $X_\etale$ (il s'agit en effet d'une propriété locale,
et les références données ci-dessus décrivent précisément de cette manière les modules quasi-cohérents
sur $U_\etale$).
\end{proof}

\begin{definition}
\label{definition-integral-closure}
Soit $S$ un schéma. Soit $X$ un espace algébrique sur $S$.
Soit $\mathcal{A}$ un faisceau quasi-cohérent de $\mathcal{O}_X$-algèbres.
La {\it clôture intégrale de $\mathcal{O}_X$ dans $\mathcal{A}$} est la
$\mathcal{O}_X$-sous-algèbre quasi-cohérente $\mathcal{A}' \subset \mathcal{A}$
construite dans le lemme \ref{lemma-integral-closure} ci-dessus.
\end{definition}

\noindent
Nous appliquerons notamment ceci lorsque $\mathcal{A} = f_*\mathcal{O}_Y$
pour un morphisme quasi-compact et quasi-séparé d'espaces algébriques
$f : Y \to X$ (voir le lemme \ref{lemma-pushforward}). On peut alors prendre
le spectre relatif de la $\mathcal{O}_X$-algèbre quasi-cohérente
(lemme \ref{lemma-affine-equivalence-algebras}) pour obtenir la
normalisation de $X$ dans $Y$.

\begin{definition}
\label{definition-normalization-X-in-Y}
Soit $S$ un schéma. Soit $f : Y \to X$ un morphisme quasi-compact et quasi-séparé
d'espaces algébriques sur $S$. Soit $\mathcal{O}'$ la fermeture
intégrale de $\mathcal{O}_X$ dans $f_*\mathcal{O}_Y$. La {\it normalisation de
$X$ dans $Y$} est le morphisme d'espaces algébriques
$$
\nu : X' = \underline{\Spec}_X(\mathcal{O}') \to X
$$
sur $S$. Elle est munie d'une factorisation naturelle
$$
Y \xrightarrow{f'} X' \xrightarrow{\nu} X
$$
du morphisme initial $f$.
\end{definition}

\noindent
Pour obtenir cette factorisation, on utilise la remarque
\ref{remark-factorization-quasi-compact-quasi-separated}
et la fonctorialité de la construction $\underline{\Spec}$.

\begin{lemma}
\label{lemma-properties-normalization}
Soit $S$ un schéma. Soit $f : Y \to X$ un morphisme quasi-compact et quasi-séparé
d'espaces algébriques sur $S$. Soit $Y \to X' \to X$ la
normalisation de $X$ dans $Y$.
\begin{enumerate}
\item Si $W \to X$ est un morphisme étale d'espaces algébriques sur $S$,
alors $W \times_X X'$ est la normalisation de $W$ dans $W \times_X Y$.
\item Si $Y$ et $X$ sont représentables, alors $X'$ est représentable
et canoniquement isomorphe à la normalisation du schéma $X$
dans le schéma $Y$ construite dans
Morphismes, section \ref{morphisms-section-normalization}.
\end{enumerate}
\end{lemma}

\begin{proof}
Il résulte immédiatement de la construction que la formation
de la normalisation de $X$ dans $Y$ commute au changement de base
étale, ce qui prouve (1). D'autre part, si
$X$ et $Y$ sont des schémas, alors, pour tout ouvert affine $U \subset X$,
$f_*\mathcal{O}_Y(U) = \mathcal{O}_Y(f^{-1}(U))$ ; par conséquent,
$\nu^{-1}(U)$ est le spectre du même anneau que celui
obtenu dans la construction correspondante pour les schémas.
\end{proof}

\noindent
Voici une caractérisation de cette construction.

\begin{lemma}
\label{lemma-characterize-normalization}
Soit $S$ un schéma. Soit $f : Y \to X$ un morphisme quasi-compact et quasi-séparé
d'espaces algébriques sur $S$. La factorisation $f = \nu \circ f'$,
où $\nu : X' \to X$ est la normalisation de $X$ dans $Y$, est caractérisée
par les deux propriétés suivantes :
\begin{enumerate}
\item le morphisme $\nu$ est entier ;
\item pour toute factorisation $f = \pi \circ g$, où $\pi : Z \to X$
est entier, il existe un diagramme commutatif
$$
\xymatrix{
Y \ar[d]_{f'} \ar[r]_g & Z \ar[d]^\pi \\
X' \ar[ru]^h \ar[r]^\nu & X
}
$$
pour un unique morphisme $h : X' \to Z$.
\end{enumerate}
De plus, dans (2), le morphisme $h : X' \to Z$ est la normalisation de
$Z$ dans $Y$.
\end{lemma}

\begin{proof}
Soit $\mathcal{O}' \subset f_*\mathcal{O}_Y$ la clôture intégrale de
$\mathcal{O}_X$ définie comme dans la définition \ref{definition-normalization-X-in-Y}.
Le morphisme $\nu$ est entier par construction, ce qui prouve (1).
Supposons donnée une factorisation $f = \pi \circ g$ avec $\pi : Z \to X$
entier comme en (2). D'après la définition \ref{definition-integral},
$\pi$ est affine ; ainsi, $Z$ est le spectre relatif
d'un faisceau quasi-cohérent de $\mathcal{O}_X$-algèbres $\mathcal{B}$.
Le morphisme $g : Y \to Z$ correspond à un morphisme de $\mathcal{O}_X$-algèbres
$\chi : \mathcal{B} \to f_*\mathcal{O}_Y$. Comme $\mathcal{B}(U)$ est
entière sur $\mathcal{O}_X(U)$ pour tout objet affine $U$ étale sur $X$
(d'après la définition \ref{definition-integral}),
le lemme \ref{lemma-integral-closure} montre
que $\chi(\mathcal{B}) \subset \mathcal{O}'$.
Par fonctorialité du spectre relatif,
le lemme \ref{lemma-affine-equivalence-algebras}
fournit alors un unique morphisme
$h : X' \to Z$. Nous omettons la vérification de la commutativité du diagramme.

\medskip\noindent
Il est clair que (1) et (2) caractérisent la
factorisation $f = \nu \circ f'$, puisqu'ils la caractérisent
comme objet initial d'une catégorie. Le morphisme $h$ de (2)
est entier d'après le lemme \ref{lemma-finite-permanence}.
Étant donnée une factorisation $g = \pi' \circ g'$ avec $\pi' : Z' \to Z$
entier, on obtient une factorisation $f = (\pi \circ \pi') \circ g'$ et
un morphisme $h' : X' \to Z'$. L'unicité implique que
$\pi' \circ h' = h$. La caractérisation (1), (2) s'applique donc
au morphisme $h : X' \to Z$, ce qui donne la dernière assertion du lemme.
\end{proof}

\begin{lemma}
\label{lemma-normalization-in-reduced}
Soit $S$ un schéma. Soit $f : Y \to X$ un morphisme quasi-compact et
quasi-séparé d'espaces algébriques sur $S$.
Soit $X' \to X$ la normalisation de $X$ dans $Y$.
Si $Y$ est réduit, alors $X'$ l'est aussi.
\end{lemma}

\begin{proof}
Cela résulte du fait qu'un sous-anneau d'un anneau réduit est réduit.
Nous omettons quelques détails.
\end{proof}

\begin{lemma}
\label{lemma-normalization-generic}
Soit $S$ un schéma. Soit $f : Y \to X$ un morphisme quasi-compact et quasi-séparé
de schémas. Soit $X' \to X$ la normalisation de $X$ dans $Y$.
Si $x' \in |X'|$ est un point de codimension $0$
(Propriétés des espaces, définition
\ref{spaces-properties-definition-dimension-local-ring}), alors
$x'$ est l'image d'un point $y \in |Y|$ de codimension $0$.
\end{lemma}

\begin{proof}
D'après le lemme \ref{lemma-properties-normalization} et les définitions, on peut
supposer que $X = \Spec(A)$ est affine. Alors $X' = \Spec(A')$, où $A'$ est
la clôture intégrale de $A$ dans $\Gamma(Y, \mathcal{O}_Y)$, et $x'$ correspond
à un idéal premier minimal de $A'$. Choisissons un morphisme étale surjectif
$V \to Y$, où $V = \Spec(B)$ est affine. Alors
$A' \to B$ est injectif ; tout idéal premier minimal de $A'$ est donc
l'image d'un idéal premier minimal de $B$ ; voir
Algèbre, lemme \ref{algebra-lemma-injective-minimal-primes-in-image}.
Le résultat en découle.
\end{proof}

\begin{lemma}
\label{lemma-normalization-in-disjoint-union}
Soit $S$ un schéma.
Soit $f : Y \to X$ un morphisme quasi-compact et quasi-séparé
d'espaces algébriques sur $S$.
Supposons que $Y = Y_1 \amalg Y_2$ soit la somme disjointe de deux
espaces algébriques.
Écrivons $f_i = f|_{Y_i}$. Soit $X_i'$ la normalisation de $X$ dans $Y_i$.
Alors $X_1' \amalg X_2'$ est la normalisation de $X$ dans $Y$.
\end{lemma}

\begin{proof}
Démonstration omise.
\end{proof}

\begin{lemma}
\label{lemma-normalization-in-universally-closed}
Soit $S$ un schéma.
Soit $f : X \to Y$ un morphisme quasi-compact, quasi-séparé et
universellement fermé d'espaces algébriques sur $S$.
Alors $f_*\mathcal{O}_X$ est entière sur $\mathcal{O}_Y$. Autrement
dit, la normalisation de $Y$ dans $X$ est égale à la factorisation
$$
X \longrightarrow \underline{\Spec}_Y(f_*\mathcal{O}_X)
\longrightarrow Y
$$
de la remarque \ref{remark-factorization-quasi-compact-quasi-separated}.
\end{lemma}

\begin{proof}
La question est locale sur $Y$ pour la topologie étale ; on peut donc se ramener au
cas où $Y = \Spec(R)$ est affine. Soit $h \in \Gamma(X, \mathcal{O}_X)$.
Il faut montrer que $h$ satisfait une équation unitaire sur $R$. Considérons $h$
comme un morphisme dans le diagramme commutatif suivant :
$$
\xymatrix{
X \ar[rr]_h \ar[rd]_f & & \mathbf{A}^1_Y \ar[ld] \\
& Y &
}
$$
Soit $Z \subset \mathbf{A}^1_Y$ l'image schématique de $h$ ;
voir la définition \ref{definition-scheme-theoretic-image}.
Le morphisme $h$ est quasi-compact, car $f$ est quasi-compact et
$\mathbf{A}^1_Y \to Y$ est séparé ; voir
le lemme \ref{lemma-quasi-compact-permanence}.
D'après le lemme \ref{lemma-quasi-compact-scheme-theoretic-image}, le
morphisme $X \to Z$ a une image dense sur les espaces topologiques sous-jacents. D'après
le lemme \ref{lemma-universally-closed-permanence}, le morphisme
$X \to Z$ est fermé. Ainsi $h(X) = Z$ ensemblistement.
On peut donc appliquer
le lemme \ref{lemma-image-proper-is-proper}
pour conclure que $Z \to Y$ est universellement fermé (et même propre).
Comme $Z \subset \mathbf{A}^1_Y$, le morphisme $Z \to Y$ est affine
et propre, donc entier d'après le lemme \ref{lemma-integral-universally-closed}.
En écrivant $\mathbf{A}^1_Y = \Spec(R[T])$, on en déduit que
l'idéal $I \subset R[T]$ de $Z$ contient un polynôme unitaire
$P(T) \in R[T]$. Ainsi $P(h) = 0$, ce qui prouve le résultat.
\end{proof}

\begin{lemma}
\label{lemma-normalization-in-integral}
Soit $S$ un schéma. Soit $f : Y \to X$ un morphisme entier
d'espaces algébriques sur $S$.
Alors la normalisation de $X$ dans $Y$ est égale à $Y$.
\end{lemma}

\begin{proof}
D'après le lemme \ref{lemma-integral-universally-closed}, c'est un cas particulier du
lemme \ref{lemma-normalization-in-universally-closed}.
\end{proof}

\begin{lemma}
\label{lemma-nagata-normalization-finite}
Soit $S$ un schéma. Soit $f : X \to Y$ un morphisme d'espaces algébriques
sur $S$. Supposons que
\begin{enumerate}
\item $Y$ soit de Nagata ;
\item $f$ soit quasi-séparé et de type fini ;
\item $X$ soit réduit.
\end{enumerate}
Alors la normalisation $\nu : Y' \to Y$ de $Y$ dans $X$ est finie.
\end{lemma}

\begin{proof}
La question est locale sur $Y$ pour la topologie étale ; voir
le lemme \ref{lemma-properties-normalization}.
On peut donc supposer que $Y = \Spec(R)$ est affine.
Alors $R$ est un anneau de Nagata noethérien,
et il faut montrer que la clôture intégrale de
$R$ dans $\Gamma(X, \mathcal{O}_X)$ est finie sur $R$.
Comme $f$ est quasi-compact, l'espace $X$ est quasi-compact.
Choisissons un schéma affine $U$ et un morphisme étale surjectif
$U \to X$ (Propriétés des espaces, lemme
\ref{spaces-properties-lemma-quasi-compact-affine-cover}).
Alors $\Gamma(X, \mathcal{O}_X) \subset \Gamma(U, \mathcal{O}_U)$.
Comme $R$ est noethérien, il suffit de montrer que
la clôture intégrale de $R$ dans $\Gamma(U, \mathcal{O}_U)$
est finie sur $R$. Puisque $U \to Y$ est de type fini,
cela résulte de
Morphismes, lemme \ref{morphisms-lemma-nagata-normalization-finite}.
\end{proof}







\section{Normalisation}
\label{section-normalization}

\noindent
Cette section est l'analogue de
Morphismes, section \ref{morphisms-section-normalization}.

\begin{lemma}
\label{lemma-prepare-normalization}
Soit $S$ un schéma. Soit $X$ un espace algébrique sur $S$.
Les conditions suivantes sont équivalentes :
\begin{enumerate}
\item il existe un morphisme étale surjectif $U \to X$, où $U$
est un schéma tel que tout ouvert quasi-compact de $U$ ait
un nombre fini de composantes irréductibles ;
\item pour tout schéma $U$ et tout morphisme étale
$U \to X$, tout ouvert quasi-compact de $U$ a un nombre fini de
composantes irréductibles ;
\item pour tout espace algébrique quasi-compact $Y$ étale sur $X$,
l'ensemble des points de codimension $0$ de $Y$ (Propriétés des espaces,
définition \ref{spaces-properties-definition-dimension-local-ring})
est fini ;
\item pour tout espace algébrique quasi-compact $Y$ étale sur $X$,
l'espace $|Y|$ a un nombre fini de composantes irréductibles.
\end{enumerate}
Si $X$ est représentable, cela signifie que tout ouvert quasi-compact de $X$
a un nombre fini de composantes irréductibles.
\end{lemma}

\begin{proof}
L'équivalence de (1) et (2), ainsi que la dernière assertion, résultent de
Descente, lemme \ref{descent-lemma-locally-finite-nr-irred-local-fppf}, et de
Propriétés des espaces, lemme \ref{spaces-properties-lemma-type-property}.
Il est clair que (4) implique (1) et (2) en ne considérant que les $Y$
qui sont des schémas. De même, (3) implique (1) et (2), puisque, pour un schéma,
les points de codimension $0$ sont les points génériques de ses
composantes irréductibles ; voir par exemple
Propriétés des espaces, lemme \ref{spaces-properties-lemma-codimension-0-points}.

\medskip\noindent
Réciproquement, supposons (2) et soit $Y \to X$ un morphisme étale d'espaces
algébriques, avec $Y$ quasi-compact. On peut alors choisir un schéma
affine $V$ et un morphisme étale surjectif $V \to Y$
(Propriétés des espaces, lemme
\ref{spaces-properties-lemma-quasi-compact-affine-cover}).
Comme $V$ a un nombre fini de composantes irréductibles d'après (2) et que
$|V| \to |Y|$ est surjectif et continu, on conclut que
$|Y|$ a un nombre fini de composantes irréductibles d'après Topologie, lemme
\ref{topology-lemma-surjective-continuous-irreducible-components}.
La condition (4) est donc satisfaite. De même, d'après
Propriétés des espaces, lemme \ref{spaces-properties-lemma-codimension-0-points},
les images des points génériques des composantes irréductibles de $V$
sont les points de codimension $0$ de $Y$ ; ils sont donc
en nombre fini, ce qui prouve (3).
\end{proof}

\begin{lemma}
\label{lemma-locally-noetherian-can-be-normalized}
Soit $S$ un schéma. Soit $X$ un espace algébrique localement noethérien
sur $S$. Alors $X$ satisfait aux conditions équivalentes du
lemme \ref{lemma-prepare-normalization}.
\end{lemma}

\begin{proof}
Si $U \to X$ est étale et si $U$ est un schéma, alors $U$ est un schéma
localement noethérien ; voir Propriétés des espaces, section
\ref{spaces-properties-section-types-properties}.
L'ensemble des composantes irréductibles d'un schéma localement noethérien est
localement fini (Diviseurs, lemme \ref{divisors-lemma-components-locally-finite}).
On en conclut que $X$ satisfait à la condition (2) du lemme.
\end{proof}

\begin{lemma}
\label{lemma-flat-image-point-codimension-0}
Soit $S$ un schéma. Soit $f : X \to Y$ un morphisme plat d'espaces
algébriques sur $S$. Alors, pour $x \in |X|$, on a : $x$ est de codimension $0$ dans
$X \Rightarrow f(x)$ est de codimension $0$ dans $Y$.
\end{lemma}

\begin{proof}
À l'aide de Propriétés des espaces, lemme
\ref{spaces-properties-lemma-codimension-0-points}, et par localisation étale,
on se ramène au cas d'un morphisme de schémas et de points génériques
de composantes irréductibles. Le résultat découle alors de ce que les générisations
se relèvent le long des morphismes plats de schémas ; voir
Morphismes, lemme \ref{morphisms-lemma-generalizations-lift-flat}.
\end{proof}

\begin{lemma}
\label{lemma-flat-locally-finite-type-points-codimension-0}
Soit $S$ un schéma. Soit $f : X \to Y$ un morphisme d'espaces
algébriques sur $S$. Supposons $f$ plat et localement de type fini, et
supposons que $Y$ satisfasse aux conditions équivalentes du
lemme \ref{lemma-prepare-normalization}.
Alors $X$ satisfait aux conditions équivalentes du
lemme \ref{lemma-prepare-normalization} et, pour $x \in |X|$, on a :
$x$ est de codimension $0$ dans $X \Rightarrow f(x)$ est de codimension $0$ dans $Y$.
\end{lemma}

\begin{proof}
La dernière assertion résulte du
lemme \ref{lemma-flat-image-point-codimension-0}.
Choisissons un morphisme étale surjectif $V \to Y$, où $V$ est un schéma.
Choisissons un morphisme étale surjectif $U \to X \times_Y V$, où $U$
est un schéma. Il suffit de montrer que tout ouvert quasi-compact
de $U$ a un nombre fini de composantes irréductibles.
Nous utiliserons sans autre mention les résultats de
Propriétés des espaces, lemme \ref{spaces-properties-lemma-codimension-0-points}.
D'après ce qui précède,
les points de codimension $0$ de $U$ se trouvent au-dessus de
points de codimension $0$ de $V$, lesquels sont localement en nombre fini par
hypothèse. Il suffit donc de montrer que, pour $v \in V$ de codimension $0$,
les points de codimension $0$ de la fibre schématique $U_v = U \times_V v$
sont localement en nombre fini. Cela résulte de ce que $U_v$ est un schéma
localement de type fini sur $\kappa(v)$, donc localement noethérien ;
et l'on peut appliquer le lemme \ref{lemma-locally-noetherian-can-be-normalized}
par exemple.
\end{proof}

\begin{lemma}
\label{lemma-normalization}
Soit $S$ un schéma. Pour tout espace algébrique $X$ sur $S$ satisfaisant
aux conditions équivalentes du lemme \ref{lemma-prepare-normalization},
il existe un morphisme d'espaces algébriques
$$
\nu_X : X^\nu \longrightarrow X
$$
ayant les propriétés suivantes :
\begin{enumerate}
\item si $X$ satisfait aux conditions équivalentes du
lemme \ref{lemma-prepare-normalization}, alors
$X^\nu$ est normal et $\nu_X$ est entier ;
\item si $X$ est un schéma dont tout ouvert quasi-compact a un nombre fini
de composantes irréductibles, alors $\nu_X : X^\nu \to X$ est la
normalisation de $X$ construite dans
Morphismes, section \ref{morphisms-section-normalization} ;
\item si $f : X \to Y$ est un morphisme d'espaces algébriques sur $S$
dont la source et le but satisfont aux conditions équivalentes du
lemme \ref{lemma-prepare-normalization}, et si tout point de codimension $0$
de $X$ est envoyé par $f$ sur un point de codimension $0$ de $Y$, alors
il existe un unique morphisme $f^\nu : X^\nu \to Y^\nu$ d'espaces
algébriques sur $S$ tel que $\nu_Y \circ f^\nu = f \circ \nu_X$ ;
\item si $f : X \to Y$ est un morphisme étale ou lisse d'espaces
algébriques et si $Y$ satisfait aux conditions équivalentes du
lemme \ref{lemma-prepare-normalization}, alors les hypothèses de (3)
sont satisfaites et le morphisme $f^\nu$ induit un isomorphisme
$X^\nu  \to X \times_Y Y^\nu$.
\end{enumerate}
\end{lemma}

\begin{proof}
Considérons la catégorie $\mathcal{C}$ dont les objets sont les
schémas $U$ sur $S$ tels que tout ouvert quasi-compact de
$U$ ait un nombre fini de composantes irréductibles, et dont les morphismes
sont les morphismes $g : U \to V$ de schémas sur $S$ tels
que tout point générique d'une composante irréductible de $U$
soit envoyé sur le point générique d'une composante irréductible de $V$.
Nous avons déjà montré que :
\begin{enumerate}
\item[(a)] pour $U \in \Ob(\mathcal{C})$, on dispose d'un
morphisme de normalisation $\nu_U : U^\nu \to U$ comme dans
Morphismes, définition \ref{morphisms-definition-normalization} ;
\item[(b)] pour $U \in \Ob(\mathcal{C})$, le morphisme $\nu_U$
est entier et $U^\nu$ est un schéma normal ; voir
Morphismes, lemme \ref{morphisms-lemma-normalization-normal} ;
\item[(c)] pour tout $g : U \to V \in \text{Arrows}(\mathcal{C})$,
il existe un unique morphisme $g^\nu : U^\nu \to V^\nu$
tel que $\nu_V \circ g^\nu = g \circ \nu_U$ ; voir
Morphismes, lemme \ref{morphisms-lemma-normalization-normal},
partie (4), appliqué à la composée $U^\nu \to U \to V$ ;
\item[(d)] si $V \in \Ob(\mathcal{C})$ et si $g : U \to V$ est
étale ou lisse, alors $U \in \Ob(\mathcal{C})$ et
$g \in \text{Arrows}(\mathcal{C})$, et le morphisme $g^\nu$
induit un isomorphisme $U^\nu \to U \times_V V^\nu$ ; voir le
lemme \ref{lemma-flat-locally-finite-type-points-codimension-0}
et Compléments sur les morphismes, lemme
\ref{more-morphisms-lemma-normalization-and-smooth}.
\end{enumerate}
Il s'agit d'étendre cette construction à la catégorie correspondante
des espaces algébriques $X$ sur $S$.

\medskip\noindent
Soit $X$ un espace algébrique sur $S$ satisfaisant
aux conditions équivalentes du lemme \ref{lemma-prepare-normalization}.
Soit $U \to X$ un morphisme étale surjectif, où $U$ est un schéma.
Posons $R = U \times_X U$, avec les projections $s, t : R \to U$, et
$j = (t, s) : R \to U \times_S U$, de sorte que $X = U/R$ ; voir
Espaces, lemme \ref{spaces-lemma-space-presentation}.
Observons que $U$ et $R$ sont des objets de $\mathcal{C}$
en vertu de nos hypothèses sur $X$, et que les morphismes
$s$ et $t$ sont des morphismes étales de schémas sur $S$.
D'après (a), on dispose des morphismes de normalisation
$\nu_U : U^\nu \to U$ et $\nu_R : R^\nu \to R$ ;
d'après (d), on dispose de morphismes $s^\nu : R^\nu \to U^\nu$ et
$t^\nu : R^\nu \to U^\nu$ qui définissent des
isomorphismes $R^\nu \to R \times_{s, U} U^\nu$ et
$R^\nu \to U^\nu \times_{U, t} R$. Il en résulte que
$s^\nu$ et $t^\nu$ sont étales (car ils sont isomorphes
à des changements de base de morphismes étales). Le morphisme induit
$j^\nu = (t^\nu, s^\nu) : R^\nu \to U^\nu \times_S U^\nu$
est un monomorphisme, car il est égal à la composée
\begin{align*}
R^\nu
& \to
(U^\nu \times_{U, t} R) \times_R (R \times_{s, U} U^\nu) \\
& =
U^\nu \times_{U, t} R \times_{s, U} U^\nu \\
& \xrightarrow{j}
U^\nu \times_U (U \times_S U) \times_U U^\nu \\
& =
U^\nu \times_S U^\nu
\end{align*}
Le premier morphisme est le morphisme diagonal de $\nu_R$.
(Cela montre que $R^\nu$ est un sous-schéma de la restriction de
$R$ à $U^\nu$.) Un calcul formel de produits fibrés
utilisant la propriété (d) montre que
$R^\nu \times_{s^\nu, U^\nu, t^\nu} R^\nu$ est la normalisation
de $R \times_{s, U, t} R$. Le morphisme étale
$c : R \times_{s, U, t} R \to R$ se prolonge donc de manière unique en $c^\nu$ d'après (d).
Le morphisme $c^\nu$ est compatible avec la projection
$\text{pr}_{13} : U^\nu \times_S U^\nu \times_S U^\nu \to U^\nu \times_S U^\nu$.
De même, il existe un morphisme $i^\nu : R^\nu \to R^\nu$
compatible avec le morphisme $U^\nu \times_S U^\nu \to U^\nu \times_S U^\nu$
qui échange les facteurs, et un morphisme $e^\nu : U^\nu \to R^\nu$
compatible avec le morphisme diagonal $U^\nu \to U^\nu \times_S U^\nu$.
Il en résulte que $j^\nu : R^\nu \to U^\nu \times_S U^\nu$
est une relation d'équivalence étale.
On peut alors poser, et l'on pose, $X^\nu = U^\nu/R^\nu$
(Espaces, théorème \ref{spaces-theorem-presentation}).
On voit ensuite que $U^\nu = X^\nu \times_X U$ d'après
Groupoïdes, lemme \ref{groupoids-lemma-criterion-fibre-product}.

\medskip\noindent
Voici ce que nous avons montré au paragraphe précédent : pour tout espace
algébrique $X$ sur $S$ satisfaisant aux conditions équivalentes du
lemme \ref{lemma-prepare-normalization}, si l'on choisit un morphisme étale
surjectif $g : U \to X$, où $U$ est un schéma, on obtient
un diagramme cartésien
$$
\xymatrix{
X^\nu \ar[d]_{\nu_X} & U^\nu \ar[l]^{g^\nu} \ar[d]^{\nu_U} \\
X & U \ar[l]_g
}
$$
d'espaces algébriques. Il en résulte immédiatement que $X^\nu$ est un espace
algébrique normal et que $\nu_X$ est un morphisme entier. Cela prouve
la partie (1) du lemme.

\medskip\noindent
Nous montrerons plus bas que le morphisme $\nu_X : X^\nu \to X$
est, à isomorphisme unique près, indépendant du choix de $g$ ;
mais, pour l'instant, si $X$ est un schéma, choisissons $\text{id} : X \to X$,
ce qui établit aussitôt la partie (2) du lemme.

\medskip\noindent
Il reste à prouver les parties (3) et (4). Soient $g : U \to X$,
$\nu_X : X^\nu \to X$ et $g^\nu : U^\nu \to X^\nu$ comme ci-dessus.
Soit $Z$ un schéma normal, et soient $h : Z \to U$ et $a : Z \to X^\nu$
des morphismes sur $S$ tels que $g \circ h = \nu_X \circ a$ et que
toute composante irréductible de $Z$ domine une composante irréductible de $U$
(par $h$). D'après
Morphismes, lemme \ref{morphisms-lemma-normalization-normal}, partie (4),
on obtient un unique morphisme $h^\nu : Z \to U^\nu$ tel
que $h = \nu_U \circ h^\nu$. Voici le diagramme :
$$
\xymatrix{
X^\nu \ar[d]_{\nu_X} &
U^\nu \ar[l]^{g^\nu} \ar[d]^{\nu_U} &
Z \ar[l]^{h^\nu} \ar@/_1em/[ll]_a \ar[dl]^h
\\
X & U \ar[l]_g
}
$$
Observons que $a = g^\nu \circ h^\nu$. En effet, puisque le carré de sommets
$X^\nu$, $X$, $U^\nu$, $U$ est cartésien, cela résulte
immédiatement de l'unicité de $h^\nu$ une fois $h$ donné.
Autrement dit, étant donné $h : Z \to U$ comme ci-dessus (mais non $a$),
il existe un unique morphisme $a : Z \to X^\nu$ tel que
$\nu_X \circ a = g \circ h$.

\medskip\noindent
Soit $f : X \to Y$ comme dans la partie (3) de l'énoncé du lemme.
Supposons choisis des morphismes étales surjectifs $U \to X$ et
$V \to Y$, où $U$ et $V$ sont des schémas, tels que $f$ se relève en un
morphisme $g : U \to V$. Alors $g \in \text{Arrows}(\mathcal{C})$
et l'on obtient un unique morphisme $g^\nu : U^\nu \to V^\nu$ compatible
avec $\nu_U$ et $\nu_V$. Les deux morphismes
$$
R^\nu = U^\nu \times_{X^\nu} U^\nu
\to U^\nu \to V^\nu \to Y^\nu
$$
sont donc égaux d'après les observations du paragraphe précédent (appliquées
avec $Y$ à la place de $X$).
Comme $X^\nu$ est construit comme le quotient de
$U^\nu$ par $R^\nu$, on obtient ainsi un unique
morphisme $f^\nu : X^\nu \to Y^\nu$ comme dans (3).

\medskip\noindent
Pour voir que la construction de $X^\nu$ est indépendante du choix
du morphisme étale surjectif $g : U \to X$, appliquons la construction du
paragraphe précédent à $\text{id} : X \to X$ et à un morphisme
$U' \to U$ entre revêtements étales de $X$. Cela suffit, car,
étant donnés deux revêtements étales de $X$, il en existe un troisième qui
les domine tous deux. On vérifie que le morphisme entre les deux
normalisations construites à l'aide de $U' \to X$ et de $U \to X$
devient un isomorphisme après changement de base à $U'$ et qu'il était donc
déjà un isomorphisme. Nous omettons les détails.

\medskip\noindent
Nous omettons la démonstration de (4), qui est analogue ; indication : utiliser la partie (d) ci-dessus.
\end{proof}

\noindent
Cela conduit à la définition suivante.

\begin{definition}
\label{definition-normalization}
Soit $S$ un schéma. Soit $X$ un espace algébrique sur $S$ satisfaisant aux
conditions équivalentes du lemme \ref{lemma-prepare-normalization}.
On appelle {\it normalisation} de $X$ le morphisme
$$
\nu_X : X^\nu \longrightarrow X
$$
construit dans le lemme \ref{lemma-normalization}.
\end{definition}

\noindent
La définition s'applique aux espaces algébriques localement
noethériens ; voir le lemme \ref{lemma-locally-noetherian-can-be-normalized}.
Habituellement, la normalisation n'est définie que pour les espaces algébriques
réduits. Avec la définition ci-dessus, la normalisation de $X$ est la même
que celle de la réduction $X_{red}$ de $X$.

\begin{lemma}
\label{lemma-normalization-reduced}
Soit $S$ un schéma. Soit $X$ un espace algébrique sur $S$ satisfaisant aux
conditions équivalentes du lemme \ref{lemma-prepare-normalization}.
Le morphisme de normalisation $\nu$ se factorise par la réduction $X_{red}$,
et $X^\nu \to X_{red}$ est la normalisation de $X_{red}$.
\end{lemma}

\begin{proof}
On peut vérifier ceci localement sur $X$ pour la topologie étale et se ramener ainsi au cas
des schémas, qui est
Morphismes, lemme \ref{morphisms-lemma-normalization-reduced}.
Nous omettons quelques détails.
\end{proof}

\begin{lemma}
\label{lemma-normalization-normal}
Soit $S$ un schéma. Soit $X$ un espace algébrique sur $S$ satisfaisant aux
conditions équivalentes du lemme \ref{lemma-prepare-normalization}.
\begin{enumerate}
\item La normalisation $X^\nu$ est normale.
\item Le morphisme $\nu : X^\nu \to X$ est entier et surjectif.
\item L'application $|\nu| : |X^\nu| \to |X|$ induit une bijection entre
les ensembles de points de codimension $0$ (Propriétés des espaces,
définition \ref{spaces-properties-definition-dimension-local-ring}).
\item Soit $f : Z \to X$ un morphisme. Supposons que $Z$ soit un espace algébrique normal
et que, pour $z \in |Z|$, on ait : $z$ est de codimension $0$ dans
$Z \Rightarrow f(z)$ est de codimension $0$ dans $X$. Alors
il existe une unique factorisation $Z \to X^\nu \to X$.
\end{enumerate}
\end{lemma}

\begin{proof}
Les propriétés (1), (2) et (3) résultent des énoncés correspondants
pour les schémas (Morphismes, lemme \ref{morphisms-lemma-normalization-normal}),
combinés au fait qu'un point d'un schéma est un point
générique d'une composante irréductible si et seulement si la dimension
de son anneau local est nulle
(Propriétés, lemme \ref{properties-lemma-generic-point}).

\medskip\noindent
Soit $Z \to X$ un morphisme comme en (4). Soit $U$ un schéma, et soit
$U \to X$ un morphisme étale surjectif. Choisissons un schéma $V$ et
un morphisme étale surjectif $V \to U \times_X Z$. La condition sur les
points de codimension $0$ assure que $V \to U$ envoie les points génériques des
composantes irréductibles de $V$ sur des points génériques de composantes
irréductibles de $U$. On obtient donc une unique factorisation
$V \to U^\nu \to U$ d'après
Morphismes, lemme \ref{morphisms-lemma-normalization-normal}.
L'unicité garantit que les deux applications
$$
V \times_{U \times_X Z} V \to V \to U^\nu
$$
coïncident, car elles donnent l'unique factorisation de l'application
$V \times_{U \times_X Z} V \to V \to U$.
Comme l'espace algébrique $U \times_X Z$ est égal au quotient
$V/V \times_{U \times_X Z} V$
(voir Espaces, section \ref{spaces-section-presentations}),
on obtient un morphisme canonique
$U \times_X Z \to U^\nu$. Voici le diagramme :
$$
\xymatrix{
U \times_X Z \ar[r] \ar[d] & U^\nu \ar[r] \ar[d] & U \ar[d] \\
Z \ar@/_/[rr] \ar@{..>}[r] & X^\nu \ar[r] & X
}
$$
Pour obtenir la flèche pointillée, remarquons que la construction de la
flèche $U \times_X Z \to U^\nu$ est fonctorielle par rapport au morphisme étale $U \to X$
(nous omettons la formulation précise et la démonstration).
Posons donc $R = U \times_X U$, avec les projections $s, t : R \to U$ ;
on obtient alors un morphisme $R \times_X Z \to R^\nu$ compatible avec
$s, t : R \to U$ et $s^\nu, t^\nu : R^\nu \to U^\nu$.
Rappelons que $X^\nu = U^\nu/R^\nu$ ; voir la démonstration du
lemme \ref{lemma-normalization}. Comme $X = U/R$, un argument simple
de théorie des faisceaux
montre que $Z = (U \times_X Z)/(R \times_X Z)$. Ainsi, les morphismes
$U \times_X Z \to U^\nu$ et $R \times_X Z \to R^\nu$ définissent un
morphisme $Z \to X^\nu$, comme voulu.
\end{proof}

\begin{lemma}
\label{lemma-nagata-normalization}
Soit $S$ un schéma. Soit $X$ un espace algébrique de Nagata sur $S$.
La normalisation $\nu : X^\nu \to X$ est un morphisme fini.
\end{lemma}

\begin{proof}
Comme $X$ est de Nagata, il est localement noethérien,
et la définition \ref{definition-normalization} s'applique.
Par la construction de $X^\nu$ dans le lemme \ref{lemma-normalization},
on se ramène immédiatement au cas des schémas, qui est
Morphismes, lemme \ref{morphisms-lemma-nagata-normalization}.
\end{proof}












\section{Morphismes séparés et localement quasi-finis}
\label{section-schemehood}

\noindent
Dans cette section, nous montrons qu'un espace algébrique qui est localement
quasi-fini et séparé sur un schéma est représentable. Il en résulte
qu'un morphisme séparé et localement quasi-fini est
représentable (voir le
lemme \ref{lemma-locally-quasi-finite-separated-representable}).
Mais d'abord...\ un lemme (qui sera rendu caduc par la
proposition \ref{proposition-locally-quasi-finite-separated-over-scheme}).

\begin{lemma}
\label{lemma-neighbourhood-scheme}
Soit $S$ un schéma. Considérons un diagramme commutatif
$$
\xymatrix{
V' \ar[r] \ar[rd] & T' \times_T X \ar[r] \ar[d] & X \ar[d] \\
& T' \ar[r] & T
}
$$
d'espaces algébriques sur $S$. Supposons que
\begin{enumerate}
\item $T' \to T$ soit un morphisme étale de schémas affines,
\item $X \to T$ soit un morphisme séparé et localement quasi-fini,
\item $V'$ soit un sous-espace ouvert de $T' \times_T X$, et
\item $V' \to T'$ soit quasi-affine.
\end{enumerate}
Alors l'image $U$ de $V'$ dans $X$ est un sous-espace ouvert
quasi-compact et représentable de $X$.
\end{lemma}

\begin{proof}
Commençons par quelques observations immédiates.
Le lemme \ref{lemma-quasi-affine-local} montre que $V'$ est représentable.
Il est aussi quasi-compact (comme schéma quasi-affine sur un schéma affine ; voir
Morphismes, lemme \ref{morphisms-lemma-quasi-affine-separated}).
Puisque $T' \times_T X \to X$ est étale
(Propriétés des espaces, lemme \ref{spaces-properties-lemma-base-change-etale}),
l'application $|T' \times_T X| \to |X|$ est ouverte ; voir
Propriétés des espaces, lemme \ref{spaces-properties-lemma-etale-open}.
Soit $U \subset X$ le sous-espace ouvert correspondant à l'image de
$|V'|$ ; voir
Propriétés des espaces, lemme \ref{spaces-properties-lemma-open-subspaces}.
Comme $|V'|$ est quasi-compact, $|U|$ l'est aussi ; par conséquent,
$U$ est un espace algébrique quasi-compact d'après
Propriétés des espaces, lemme \ref{spaces-properties-lemma-quasi-compact-space}.

\medskip\noindent
D'après
Morphismes,
lemme \ref{morphisms-lemma-locally-quasi-finite-qc-source-universally-bounded},
le morphisme $T' \to T$ est universellement borné. On peut donc raisonner par
récurrence sur l'entier $n$ qui borne le degré des fibres de $T' \to T$ ; voir
Morphismes, lemme \ref{morphisms-lemma-etale-universally-bounded}
pour une description de cet entier dans le cas d'un morphisme étale.
Si $n = 1$, alors $T' \to T$ est une immersion ouverte (voir
Morphismes étales, théorème \ref{etale-theorem-etale-radicial-open}),
et le résultat est clair. Supposons $n > 1$.

\medskip\noindent
Considérons le schéma affine $T'' = T' \times_T T'$.
Comme $T' \to T$ est étale, on a une décomposition (en sous-schémas affines
ouverts et fermés) $T'' = \Delta(T') \amalg T^*$. En effet, $\Delta = \Delta_{T'/T}$
est ouverte d'après
Morphismes, lemme \ref{morphisms-lemma-diagonal-unramified-morphism},
et fermée parce que $T' \to T$ est séparé en tant que morphisme de schémas affines.
Par changement de base, les degrés des fibres de la seconde projection
$\text{pr}_1 : T' \times_T T' \to T'$ sont bornés par $n$ ; voir
Morphismes, lemme \ref{morphisms-lemma-base-change-universally-bounded}.
D'autre part, $\text{pr}_1|_{\Delta(T')} : \Delta(T') \to T'$ est
un isomorphisme, et chacune de ses fibres a exactement un point.
En appliquant donc
Morphismes, lemme \ref{morphisms-lemma-etale-universally-bounded},
on conclut que les degrés des fibres de la restriction
$\text{pr}_1|_{T^*} : T^* \to T'$ sont bornés par $n - 1$.
L'hypothèse de récurrence appliquée au diagramme
$$
\xymatrix{
p_0^{-1}(V') \cap X^* \ar[r] \ar[rd] &
X^* \ar[r]_{p_1|_{X^*}} \ar[d] &
X' \ar[d] \\
& T^* \ar[r]^{\text{pr}_1|_{T^*}} & T'
}
$$
montre que $p_1(p_0^{-1}(V') \cap X^*)$
est un schéma quasi-compact. Nous avons posé ici
$X'' = T'' \times_T X$, $X^* = T^* \times_T X$ et $X' = T' \times_T X$,
et $p_0, p_1 : X'' \to X'$ sont les changements de base de $\text{pr}_0, \text{pr}_1$.
La plupart des hypothèses du lemme entraînent,
par changement de base, l'hypothèse correspondante pour le diagramme ci-dessus.
Par exemple, $p_0^{-1}(V') = T'' \times_{T'} V'$
est un schéma quasi-affine sur $T''$, car il s'obtient par changement de base. Nous
omettons quelques vérifications.

\medskip\noindent
D'après
Propriétés des espaces, lemme \ref{spaces-properties-lemma-subscheme},
on en conclut que
$$
p_1(p_0^{-1}(V')) =
V' \cup p_1(p_0^{-1}(V') \cap X^*)
$$
est un schéma quasi-compact. De plus, il est clair que
$p_1(p_0^{-1}(V'))$ est l'image réciproque du
sous-espace ouvert quasi-compact $U \subset X$ considéré dans le
premier paragraphe de la démonstration. Autrement dit, $T' \times_T U$ est un schéma !
Notons que $T' \times_T U$ est quasi-compact,
séparé et localement quasi-fini sur $T'$, puisque
$T' \times_T X \to T'$ est localement quasi-fini et séparé,
étant obtenu par changement de base du morphisme initial $X \to T$ (voir les
lemmes \ref{lemma-base-change-separated} et
\ref{lemma-base-change-quasi-finite}).
Il résulte alors de
Compléments sur les morphismes,
lemme \ref{more-morphisms-lemma-quasi-finite-separated-quasi-affine},
que $T' \times_T U \to T'$ est quasi-affine.

\medskip\noindent
D'après
Descente, lemme \ref{descent-lemma-descent-data-sheaves},
ceci fournit sur $T' \times_T U / T'$ une donnée de descente
relative au recouvrement étale $\{T' \to W\}$, où $W \subset T$
est l'image du morphisme $T' \to T$.
Comme $T' \times_T U$ est quasi-affine sur $T'$, il résulte de
Descente, lemme \ref{descent-lemma-quasi-affine},
que cette donnée est effective ; la dernière partie de
Descente, lemme \ref{descent-lemma-descent-data-sheaves},
implique alors que $U$ est un schéma, comme voulu.
Nous omettons quelques détails mineurs.
\end{proof}

\begin{proposition}
\label{proposition-locally-quasi-finite-separated-over-scheme}
Soit $S$ un schéma.
Soit $f : X \to T$ un morphisme d'espaces algébriques sur $S$.
Supposons que
\begin{enumerate}
\item $T$ soit représentable,
\item $f$ soit localement quasi-fini, et
\item $f$ soit séparé.
\end{enumerate}
Alors $X$ est représentable.
\end{proposition}

\begin{proof}
Soit $T = \bigcup T_i$ un recouvrement ouvert affine du schéma $T$.
Si l'on montre que les sous-espaces ouverts $X_i = f^{-1}(T_i)$ sont
représentables, alors $X$ est représentable ; voir
Propriétés des espaces, lemme \ref{spaces-properties-lemma-subscheme}.
Notons que $X_i = T_i \times_T X$ et que les propriétés d'être localement quasi-fini
et d'être séparé sont toutes deux stables par changement de base ; voir les
lemmes \ref{lemma-base-change-separated} et
\ref{lemma-base-change-quasi-finite}.
On peut donc supposer que $T$ est un schéma affine.

\medskip\noindent
D'après
Propriétés des espaces,
lemme \ref{spaces-properties-lemma-union-of-quasi-compact},
il existe un recouvrement de Zariski $X = \bigcup X_i$
tel que chacun des $X_i$ possède un recouvrement étale surjectif par
un schéma affine. D'après
Propriétés des espaces, lemme \ref{spaces-properties-lemma-subscheme},
il suffit encore de démontrer la proposition pour chaque $X_i$.
On peut donc supposer qu'il existe un schéma affine $U$ et un
morphisme étale surjectif $U \to X$. On se ramène ainsi à la
situation du paragraphe suivant.

\medskip\noindent
Supposons donnés
$$
U \longrightarrow X \longrightarrow T
$$
où $U$ et $T$ sont des schémas affines, $U \to X$ est étale surjectif, et
$X \to T$ est séparé et localement quasi-fini. D'après les
lemmes \ref{lemma-etale-locally-quasi-finite} et
\ref{lemma-composition-quasi-finite}
le morphisme $U \to T$ est localement quasi-fini.
Comme $U$ et $T$ sont affines, il est quasi-fini.
Posons $R = U \times_X U$. Alors $X = U/R$ ; voir
Espaces, lemme \ref{spaces-lemma-space-presentation}.
Puisque $X \to T$ est séparé, le
morphisme $R \to U \times_T U$ est une immersion fermée ; voir le
lemme \ref{lemma-fibre-product-after-map}.
En particulier, $R$ est lui aussi un schéma affine.
Comme $U \to X$ est étale, les morphismes de projection
$t, s : R \to U$ sont eux aussi étales. En particulier, $s$ et $t$ sont
quasi-finis, plats et de présentation finie (voir
Morphismes, lemmes \ref{morphisms-lemma-etale-locally-quasi-finite},
\ref{morphisms-lemma-etale-flat} et
\ref{morphisms-lemma-etale-locally-finite-presentation}).

\medskip\noindent
Soit $(U, R, s, t, c)$ le groupoïde associé à la relation
d'équivalence étale $R$ sur $U$. Soit $u \in U$ un point, et
notons $p \in T$ son image. Nous allons appliquer
Compléments sur les groupoïdes,
lemme \ref{more-groupoids-lemma-quasi-finite-over-base-j-proper},
au groupoïde $(U, R, s, t, c)$ sur le schéma $T$, avec les
points $p$ et $u$ ci-dessus.
D'après le paragraphe précédent, toutes les
hypothèses (1) -- (7) de ce lemme sont satisfaites.
On obtient donc un voisinage étale
$(T', p') \to (T, p)$ et des décompositions en unions disjointes
$$
U_{T'} = U' \amalg W, \quad
R_{T'} = R' \amalg W'
$$
et un point $u' \in U'$ satisfaisant les conclusions
(a), (b), (c), (d), (e), (f), (g) et (h) du lemme précité des
Compléments sur les groupoïdes,
lemme \ref{more-groupoids-lemma-quasi-finite-over-base-j-proper}.
On peut supposer, et l'on suppose, que $T'$ est affine (quitte à rétrécir $T'$).
La conclusion (h) entraîne que $R' = U' \times_{X_{T'}} U'$, les morphismes de
projection s'identifiant aux restrictions de $s'$ et $t'$.
Ainsi, le quintuplet $(U', R', s'|_{R'}, t'|_{R'}, c'|_{R' \times_{t', U', s'} R'})$ de la
conclusion (g) est une relation d'équivalence étale. D'après
Espaces, lemme \ref{spaces-lemma-finding-opens},
on conclut que $U'/R'$ est un sous-espace ouvert de $X_{T'}$. Par la conclusion (d),
les schémas $U'$ et $R'$ sont affines, et les morphismes
$s'|_{R'}, t'|_{R'}$ sont finis étales. La
proposition \ref{groupoids-proposition-finite-flat-equivalence} du chapitre Groupoïdes
s'applique donc et montre que $U'/R'$ est un schéma affine.

\medskip\noindent
On conclut que, pour toute paire de points $(u, p)$ comme ci-dessus, on peut
trouver un voisinage étale $(T', p') \to (T, p)$ tel que
$\kappa(p) = \kappa(p')$, ainsi qu'un point $u' \in U_{T'}$ au-dessus de $u$,
de telle sorte que l'image $x'$ de $u'$ dans $|X_{T'}|$ possède un voisinage ouvert
$V'$ dans $X_{T'}$ qui soit un schéma affine. En appliquant le
lemme \ref{lemma-neighbourhood-scheme},
on obtient un sous-espace ouvert $W \subset X$ qui est un schéma et
qui contient $x$ (l'image de $u$ dans $|X|$).
Comme cela vaut pour tout $x$, on voit que $X$
est un schéma d'après
Propriétés des espaces, lemme \ref{spaces-properties-lemma-subscheme}.
La démonstration est achevée.
\end{proof}



\section{Applications}
\label{section-applications}

\noindent
Une autre démonstration du lemme suivant consiste à le voir comme une conséquence
du théorème principal de Zariski pour les morphismes (non représentables) d'espaces algébriques,
tel qu'il est étudié dans Compléments sur les morphismes d'espaces, section
\ref{spaces-more-morphisms-section-structure-quasi-finite}.
En effet, Compléments sur les morphismes d'espaces, lemme
\ref{spaces-more-morphisms-lemma-quasi-finite-separated-quasi-affine}
entraîne qu'un morphisme quasi-fini et séparé d'espaces algébriques
est quasi-affine, donc représentable.

\begin{lemma}
\label{lemma-locally-quasi-finite-separated-representable}
Soit $S$ un schéma. Soit $f : X \to Y$ un morphisme d'espaces
algébriques sur $S$. Si $f$ est localement quasi-fini et séparé, alors
$f$ est représentable.
\end{lemma}

\begin{proof}
Cela résulte immédiatement de la
proposition \ref{proposition-locally-quasi-finite-separated-over-scheme}
et du fait que les propriétés d'être localement quasi-fini et séparé sont
préservées par tout changement de base ; voir les
lemmes \ref{lemma-base-change-quasi-finite} et
\ref{lemma-base-change-separated}.
\end{proof}

\begin{lemma}
\label{lemma-etale-universally-injective-open}
\begin{slogan}
Les applications étales et universellement injectives sont des immersions ouvertes.
\end{slogan}
Soit $S$ un schéma. Soit $f : X \to Y$ un morphisme étale et universellement
injectif d'espaces algébriques sur $S$. Alors $f$ est une
immersion ouverte.
\end{lemma}

\begin{proof}
Soit $T \to Y$ un morphisme d'un schéma vers $Y$.
Il suffit de montrer que $X \times_Y T \to T$ est une immersion ouverte.
Comme les propriétés d'être étale et universellement injectif sont
stables par changement de base (voir les
lemmes \ref{lemma-base-change-etale} et
\ref{lemma-base-change-universally-injective})
on peut supposer que $Y$ est un schéma. Notons que la
diagonale $\Delta_{X/Y} : X \to X \times_Y X$ est étale, est un monomorphisme et
est surjective d'après le
lemme \ref{lemma-universally-injective}.
On voit donc que $\Delta_{X/Y}$ est un isomorphisme (voir
Espaces, lemme
\ref{spaces-lemma-surjective-flat-locally-finite-presentation}),
et, en particulier, que $X$ est séparé sur $Y$.
Il s'ensuit que $X$ est lui aussi un schéma, d'après la
proposition \ref{proposition-locally-quasi-finite-separated-over-scheme}.
Enfin, $X \to Y$ est une immersion ouverte d'après le théorème fondamental
sur les morphismes étales de schémas ; voir
Morphismes étales, théorème \ref{etale-theorem-etale-radicial-open}.
\end{proof}



\section{Théorème principal de Zariski (cas représentable)}
\label{section-Zariski}

\noindent
Voici la version que l'on peut démontrer en utilisant le fait que la normalisation
commute à la localisation étale. Avant de pouvoir démontrer des versions
plus puissantes (pour les morphismes non représentables), il nous faut
développer davantage d'outils. Voir
Compléments sur les morphismes d'espaces, section
\ref{spaces-more-morphisms-section-structure-quasi-finite}.

\begin{lemma}
\label{lemma-finite-type-separated}
Soit $S$ un schéma. Soit $f : X \to Y$ un morphisme d'espaces
algébriques sur $S$, représentable, de type fini et séparé.
Soit $Y'$ la normalisation de $Y$ dans $X$. On a le diagramme :
$$
\xymatrix{
X \ar[rd]_f \ar[rr]_{f'} & & Y' \ar[ld]^\nu \\
& Y &
}
$$
Alors il existe un sous-espace ouvert $U' \subset Y'$ tel que
\begin{enumerate}
\item $(f')^{-1}(U') \to U'$ soit un isomorphisme, et
\item $(f')^{-1}(U') \subset X$ soit l'ensemble des points où
$f$ est quasi-fini.
\end{enumerate}
\end{lemma}

\begin{proof}
Soit $W \to Y$ un morphisme étale surjectif, où $W$ est un schéma.
Alors $W \times_Y X$ est lui aussi un schéma. D'après le
lemme \ref{lemma-properties-normalization},
l'espace algébrique $W \times_Y Y'$ est représentable et constitue
la normalisation du schéma $W$ dans le schéma $W \times_Y X$. On a le diagramme
$$
\xymatrix{
W \times_Y X \ar[rd]_{(1, f)} \ar[rr]_{(1, f')} & &
W \times_Y Y' \ar[ld]^{(1, \nu)} \\
& W &
}
$$
D'après Compléments sur les morphismes, lemme \ref{more-morphisms-lemma-finite-type-separated},
l'énoncé du lemme vaut sur $W$. Soit $V' \subset W \times_Y Y'$
le sous-schéma ouvert tel que
\begin{enumerate}
\item $(1, f')^{-1}(V') \to V'$ soit un isomorphisme, et
\item $(1, f')^{-1}(V') \subset W \times_Y X$ soit l'ensemble des points où
$(1, f)$ est quasi-fini.
\end{enumerate}
D'après le lemme \ref{lemma-locally-finite-type-quasi-finite-part},
il existe un plus grand ouvert $U \subset X$ constitué de points où $f$
est quasi-fini, et $W \times_Y U = (1, f')^{-1}(V')$.
Le morphisme $f'|_U : U \to Y'$ est une immersion ouverte
d'après le lemme \ref{lemma-closed-immersion-local},
car son changement de base à $W$ est l'isomorphisme $(1, f')^{-1}(V') \to V'$
suivi de l'immersion ouverte $V' \to W \times_Y Y'$.
Poser $U' = \Im(U \to Y')$ achève la démonstration
(nous omettons la vérification de $(f')^{-1}(U') = U$).
\end{proof}

\noindent
Dans le lemme suivant, on peut supprimer l'hypothèse de
représentabilité, puisque nous avons montré qu'un morphisme
localement quasi-fini et séparé est représentable.

\begin{lemma}
\label{lemma-quasi-finite-separated-quasi-affine}
Soit $S$ un schéma.
Soit $f : X \to Y$ un morphisme d'espaces algébriques sur $S$.
Supposons $f$ quasi-fini et séparé.
Soit $Y'$ la normalisation de $Y$ dans $X$.
On a le diagramme :
$$
\xymatrix{
X \ar[rd]_f \ar[rr]_{f'} & & Y' \ar[ld]^\nu \\
& Y &
}
$$
Alors $f'$ est une immersion ouverte quasi-compacte et $\nu$ est entier.
En particulier, $f$ est quasi-affine.
\end{lemma}

\begin{proof}
D'après le lemme \ref{lemma-locally-quasi-finite-separated-representable},
le morphisme $f$ est représentable. On peut donc appliquer le
lemme \ref{lemma-finite-type-separated}. Il existe ainsi un sous-espace
ouvert $U' \subset Y'$ tel que
$(f')^{-1}(U') = X$ (!) et que $X \to U'$ soit un isomorphisme ! Autrement
dit, $f'$ est une immersion ouverte. Notons que $f'$ est quasi-compact, car
$f$ est quasi-compact et $\nu : Y' \to Y$ est séparé
(lemme \ref{lemma-quasi-compact-permanence}).
Ainsi, pour tout schéma affine $Z$ et tout morphisme $Z \to Y$, le
produit fibré $Z \times_Y X$ est un sous-schéma ouvert quasi-compact
du schéma affine $Z \times_Y Y'$. Par définition, $f$ est donc
quasi-affine.
\end{proof}








\section{Homéomorphismes universels}
\label{section-universal-homeomorphisms}

\noindent
La classe des homéomorphismes universels de schémas est stable par
composition et par tout changement de base, et elle est locale pour la topologie fppf sur le but. Voir
Morphismes, lemmes
\ref{morphisms-lemma-composition-universal-homeomorphism} et
\ref{morphisms-lemma-base-change-universal-homeomorphism}, ainsi que
Descente, lemme \ref{descent-lemma-descending-property-universal-homeomorphism}.
Ainsi, en appliquant à cette notion la discussion de la
section \ref{section-representable},
on sait ce que signifie, pour un morphisme représentable
d'espaces algébriques, être un homéomorphisme universel.

\begin{lemma}
\label{lemma-universal-homeomorphism-representable}
Soit $S$ un schéma.
Soit $f : X \to Y$ un morphisme représentable d'espaces algébriques sur $S$.
Alors $f$ est un homéomorphisme universel
(au sens de la section \ref{section-representable}) si et seulement si,
pour tout morphisme d'espaces algébriques $Z \to Y$, l'application obtenue par
changement de base $Z \times_Y X \to Z$ induit un homéomorphisme
$|Z \times_Y X| \to |Z|$.
\end{lemma}

\begin{proof}
Si, pour tout morphisme d'espaces algébriques $Z \to Y$, l'application obtenue par
changement de base $Z \times_Y X \to Z$ induit un homéomorphisme
$|Z \times_Y X| \to |Z|$, il en va de même lorsque $Z$ est un schéma,
ce qui entraîne formellement que $f$ est un homéomorphisme universel au
sens de la
section \ref{section-representable}.
Réciproquement, si $f$ est un homéomorphisme universel au sens de la
section \ref{section-representable},
alors $X \to Y$ est entier, universellement injectif et surjectif
(d'après Espaces, lemme
\ref{spaces-lemma-representable-transformations-property-implication}
et
Morphismes, lemme \ref{morphisms-lemma-universal-homeomorphism}).
Ainsi, $f$ est universellement fermé ; voir le
lemme \ref{lemma-integral-universally-closed} ;
il est aussi universellement injectif et (universellement) surjectif, c'est-à-dire que
$f$ est un homéomorphisme universel.
\end{proof}

\begin{definition}
\label{definition-universal-homeomorphism}
Soit $S$ un schéma.
Un morphisme $f : X \to Y$ d'espaces algébriques sur $S$
est appelé un {\it homéomorphisme universel}
si et seulement si, pour tout morphisme d'espaces algébriques $Z \to Y$,
le changement de base $Z \times_Y X \to Z$ induit un homéomorphisme
$|Z \times_Y X| \to |Z|$.
\end{definition}

\noindent
Cette définition est compatible avec la définition antérieure pour les
morphismes représentables d'espaces algébriques, d'après le
lemme \ref{lemma-universal-homeomorphism-representable}.
Pour les morphismes d'espaces algébriques, les homéomorphismes
universels ne sont pas toujours entiers.

\begin{example}
\label{example-universal-homeomorphism}
Poursuivons la
remarque \ref{remark-universally-injective-not-separated}.
Considérons l'espace algébrique
$X = \mathbf{A}^1_k/\{x \sim -x \mid x \not = 0\}$.
Il existe des morphismes
$$
\mathbf{A}^1_k \longrightarrow X \longrightarrow \mathbf{A}^1_k
$$
tels que la première flèche soit étale surjective, la seconde
universellement injective, et que leur composée soit l'application $x \mapsto x^2$.
La composée est donc universellement fermée. Il s'ensuit que
l'application $X \to \mathbf{A}^1_k$ est un homéomorphisme universel, mais que
$X \to \mathbf{A}^1_k$ n'est pas séparé.
\end{example}

\noindent
Soit $S$ un schéma.
Soit $f : X \to Y$ un homéomorphisme universel d'espaces algébriques
sur $S$. Alors $f$ est universellement fermé, donc quasi-compact ; voir le
lemme \ref{lemma-universally-closed-quasi-compact}.
Cependant, $f$ n'est pas nécessairement séparé (voir l'exemple ci-dessus), ni même
quasi-séparé : on peut prendre pour exemple l'espace affine de dimension infinie
$\mathbf{A}^\infty = \Spec(k[x_1, x_2, \ldots])$ modulo
la relation d'équivalence donnée par le changement de signe d'un nombre fini de
coordonnées non nulles (nous omettons les détails).

\medskip\noindent
Énonçons d'abord les lemmes usuels.

\begin{lemma}
\label{lemma-base-change-universal-homeomorphism}
Le changement de base d'un homéomorphisme universel d'espaces algébriques
par un morphisme quelconque d'espaces algébriques est un homéomorphisme universel.
\end{lemma}

\begin{proof}
Cela résulte immédiatement de la définition.
\end{proof}

\begin{lemma}
\label{lemma-composition-universal-homeomorphism}
La composée de deux homéomorphismes universels
d'espaces algébriques est un homéomorphisme universel.
\end{lemma}

\begin{proof}
Démonstration omise.
\end{proof}

\begin{lemma}
\label{lemma-reduction-universal-homeomorphism}
Soit $S$ un schéma. Soit $X$ un espace algébrique sur $S$.
L'immersion fermée canonique $X_{red} \to X$ (voir
Propriétés des espaces, définition
\ref{spaces-properties-definition-reduced-induced-space})
est un homéomorphisme universel.
\end{lemma}

\begin{proof}
Démonstration omise.
\end{proof}

\noindent
Nous plaçons ici le résultat suivant, faute pour le moment d'un meilleur
emplacement.

\begin{lemma}
\label{lemma-integral-universally-injective-push-pull}
Soit $S$ un schéma. Soit $f : Y \to X$ un morphisme universellement injectif
et entier d'espaces algébriques sur $S$.
\begin{enumerate}
\item Le foncteur
$$
f_{small, *} : \Sh(Y_\etale) \longrightarrow \Sh(X_\etale)
$$
est pleinement fidèle, et son image essentielle est formée des faisceaux d'ensembles
$\mathcal{F}$ sur $X_\etale$ dont la restriction à $|X| \setminus f(|Y|)$
est isomorphe à $*$ ;
\item le foncteur
$$
f_{small, *} : \textit{Ab}(Y_\etale) \longrightarrow \textit{Ab}(X_\etale)
$$
est pleinement fidèle, et son image essentielle est formée des faisceaux abéliens sur
$X_\etale$ dont le support est contenu dans $f(|Y|)$.
\end{enumerate}
Dans les deux cas, $f_{small}^{-1}$ est un inverse à gauche du foncteur $f_{small, *}$.
\end{lemma}

\begin{proof}
Comme $f$ est entier, il est universellement fermé
(lemme \ref{lemma-integral-universally-closed}).
En particulier, $f(|Y|)$ est une partie fermée de $|X|$,
et les énoncés ont un sens.
La suite de la démonstration est identique à celle du
lemme \ref{lemma-closed-immersion-push-pull},
si ce n'est que l'on utilise
Cohomologie étale, proposition
\ref{etale-cohomology-proposition-integral-universally-injective-pushforward}
au lieu de
Cohomologie étale, proposition
\ref{etale-cohomology-proposition-closed-immersion-pushforward}.
\end{proof}























\begin{multicols}{2}[\section{Autres chapitres}]
\noindent
Préliminaires
\begin{enumerate}
\item \hyperref[introduction-section-phantom]{Introduction}
\item \hyperref[conventions-section-phantom]{Conventions}
\item \hyperref[sets-section-phantom]{Théorie des ensembles}
\item \hyperref[categories-section-phantom]{Catégories}
\item \hyperref[topology-section-phantom]{Topologie}
\item \hyperref[sheaves-section-phantom]{Faisceaux sur les espaces}
\item \hyperref[sites-section-phantom]{Sites et faisceaux}
\item \hyperref[stacks-section-phantom]{Champs}
\item \hyperref[fields-section-phantom]{Corps}
\item \hyperref[algebra-section-phantom]{Algèbre commutative}
\item \hyperref[brauer-section-phantom]{Groupes de Brauer}
\item \hyperref[homology-section-phantom]{Algèbre homologique}
\item \hyperref[derived-section-phantom]{Catégories dérivées}
\item \hyperref[simplicial-section-phantom]{Méthodes simpliciales}
\item \hyperref[more-algebra-section-phantom]{Compléments d'algèbre}
\item \hyperref[smoothing-section-phantom]{Lissage des morphismes d'anneaux}
\item \hyperref[modules-section-phantom]{Faisceaux de modules}
\item \hyperref[sites-modules-section-phantom]{Modules sur les sites}
\item \hyperref[injectives-section-phantom]{Objets injectifs}
\item \hyperref[cohomology-section-phantom]{Cohomologie des faisceaux}
\item \hyperref[sites-cohomology-section-phantom]{Cohomologie sur les sites}
\item \hyperref[dga-section-phantom]{Algèbre différentielle graduée}
\item \hyperref[dpa-section-phantom]{Algèbre à puissances divisées}
\item \hyperref[sdga-section-phantom]{Faisceaux différentiels gradués}
\item \hyperref[hypercovering-section-phantom]{Hyperrecouvrements}
\end{enumerate}
Schémas
\begin{enumerate}
\setcounter{enumi}{25}
\item \hyperref[schemes-section-phantom]{Schémas}
\item \hyperref[constructions-section-phantom]{Constructions de schémas}
\item \hyperref[properties-section-phantom]{Propriétés des schémas}
\item \hyperref[morphisms-section-phantom]{Morphismes de schémas}
\item \hyperref[coherent-section-phantom]{Cohomologie des schémas}
\item \hyperref[divisors-section-phantom]{Diviseurs}
\item \hyperref[limits-section-phantom]{Limites de schémas}
\item \hyperref[varieties-section-phantom]{Variétés}
\item \hyperref[topologies-section-phantom]{Topologies sur les schémas}
\item \hyperref[descent-section-phantom]{Descente}
\item \hyperref[perfect-section-phantom]{Catégories dérivées de schémas}
\item \hyperref[more-morphisms-section-phantom]{Compléments sur les morphismes}
\item \hyperref[flat-section-phantom]{Compléments sur la platitude}
\item \hyperref[groupoids-section-phantom]{Schémas en groupoïdes}
\item \hyperref[more-groupoids-section-phantom]{Compléments sur les schémas en groupoïdes}
\item \hyperref[etale-section-phantom]{Morphismes étales de schémas}
\end{enumerate}
Thèmes de théorie des schémas
\begin{enumerate}
\setcounter{enumi}{41}
\item \hyperref[chow-section-phantom]{Homologie de Chow}
\item \hyperref[intersection-section-phantom]{Théorie de l'intersection}
\item \hyperref[pic-section-phantom]{Schémas de Picard des courbes}
\item \hyperref[weil-section-phantom]{Théories cohomologiques de Weil}
\item \hyperref[adequate-section-phantom]{Modules adéquats}
\item \hyperref[dualizing-section-phantom]{Complexes dualisants}
\item \hyperref[duality-section-phantom]{Dualité pour les schémas}
\item \hyperref[discriminant-section-phantom]{Discriminants et différentes}
\item \hyperref[derham-section-phantom]{Cohomologie de de Rham}
\item \hyperref[local-cohomology-section-phantom]{Cohomologie locale}
\item \hyperref[algebraization-section-phantom]{Géométrie algébrique et formelle}
\item \hyperref[curves-section-phantom]{Courbes algébriques}
\item \hyperref[resolve-section-phantom]{Résolution des surfaces}
\item \hyperref[models-section-phantom]{Réduction semi-stable}
\item \hyperref[functors-section-phantom]{Foncteurs et morphismes}
\item \hyperref[equiv-section-phantom]{Catégories dérivées de variétés}
\item \hyperref[pione-section-phantom]{Groupes fondamentaux de schémas}
\item \hyperref[etale-cohomology-section-phantom]{Cohomologie étale}
\item \hyperref[crystalline-section-phantom]{Cohomologie cristalline}
\item \hyperref[proetale-section-phantom]{Cohomologie pro-étale}
\item \hyperref[relative-cycles-section-phantom]{Cycles relatifs}
\item \hyperref[more-etale-section-phantom]{Compléments de cohomologie étale}
\item \hyperref[trace-section-phantom]{La formule des traces}
\end{enumerate}
Espaces algébriques
\begin{enumerate}
\setcounter{enumi}{64}
\item \hyperref[spaces-section-phantom]{Espaces algébriques}
\item \hyperref[spaces-properties-section-phantom]{Propriétés des espaces algébriques}
\item \hyperref[spaces-morphisms-section-phantom]{Morphismes d'espaces algébriques}
\item \hyperref[decent-spaces-section-phantom]{Espaces algébriques décents}
\item \hyperref[spaces-cohomology-section-phantom]{Cohomologie des espaces algébriques}
\item \hyperref[spaces-limits-section-phantom]{Limites d'espaces algébriques}
\item \hyperref[spaces-divisors-section-phantom]{Diviseurs sur les espaces algébriques}
\item \hyperref[spaces-over-fields-section-phantom]{Espaces algébriques sur des corps}
\item \hyperref[spaces-topologies-section-phantom]{Topologies sur les espaces algébriques}
\item \hyperref[spaces-descent-section-phantom]{Descente et espaces algébriques}
\item \hyperref[spaces-perfect-section-phantom]{Catégories dérivées d'espaces}
\item \hyperref[spaces-more-morphisms-section-phantom]{Compléments sur les morphismes d'espaces}
\item \hyperref[spaces-flat-section-phantom]{Platitude sur les espaces algébriques}
\item \hyperref[spaces-groupoids-section-phantom]{Groupoïdes dans les espaces algébriques}
\item \hyperref[spaces-more-groupoids-section-phantom]{Compléments sur les groupoïdes dans les espaces}
\item \hyperref[bootstrap-section-phantom]{Amorçage}
\item \hyperref[spaces-pushouts-section-phantom]{Sommes amalgamées d'espaces algébriques}
\end{enumerate}
Thèmes de géométrie
\begin{enumerate}
\setcounter{enumi}{81}
\item \hyperref[spaces-chow-section-phantom]{Groupes de Chow des espaces}
\item \hyperref[groupoids-quotients-section-phantom]{Quotients de groupoïdes}
\item \hyperref[spaces-more-cohomology-section-phantom]{Compléments sur la cohomologie des espaces}
\item \hyperref[spaces-simplicial-section-phantom]{Espaces simpliciaux}
\item \hyperref[spaces-duality-section-phantom]{Dualité pour les espaces}
\item \hyperref[formal-spaces-section-phantom]{Espaces algébriques formels}
\item \hyperref[restricted-section-phantom]{Algébrisation des espaces formels}
\item \hyperref[spaces-resolve-section-phantom]{Résolution des surfaces revisitée}
\end{enumerate}
Théorie des déformations
\begin{enumerate}
\setcounter{enumi}{89}
\item \hyperref[formal-defos-section-phantom]{Théorie formelle des déformations}
\item \hyperref[defos-section-phantom]{Théorie des déformations}
\item \hyperref[cotangent-section-phantom]{Le complexe cotangent}
\item \hyperref[examples-defos-section-phantom]{Problèmes de déformation}
\end{enumerate}
Champs algébriques
\begin{enumerate}
\setcounter{enumi}{93}
\item \hyperref[algebraic-section-phantom]{Champs algébriques}
\item \hyperref[examples-stacks-section-phantom]{Exemples de champs}
\item \hyperref[stacks-sheaves-section-phantom]{Faisceaux sur les champs algébriques}
\item \hyperref[criteria-section-phantom]{Critères de représentabilité}
\item \hyperref[artin-section-phantom]{Axiomes d'Artin}
\item \hyperref[quot-section-phantom]{Espaces de Quot et de Hilbert}
\item \hyperref[stacks-properties-section-phantom]{Propriétés des champs algébriques}
\item \hyperref[stacks-morphisms-section-phantom]{Morphismes de champs algébriques}
\item \hyperref[stacks-limits-section-phantom]{Limites de champs algébriques}
\item \hyperref[stacks-cohomology-section-phantom]{Cohomologie des champs algébriques}
\item \hyperref[stacks-perfect-section-phantom]{Catégories dérivées de champs}
\item \hyperref[stacks-introduction-section-phantom]{Introduction aux champs algébriques}
\item \hyperref[stacks-more-morphisms-section-phantom]{Compléments sur les morphismes de champs}
\item \hyperref[stacks-geometry-section-phantom]{La géométrie des champs}
\end{enumerate}
Thèmes de théorie des espaces de modules
\begin{enumerate}
\setcounter{enumi}{107}
\item \hyperref[moduli-section-phantom]{Champs de modules}
\item \hyperref[moduli-curves-section-phantom]{Espaces de modules de courbes}
\end{enumerate}
Divers
\begin{enumerate}
\setcounter{enumi}{109}
\item \hyperref[examples-section-phantom]{Exemples}
\item \hyperref[exercises-section-phantom]{Exercices}
\item \hyperref[guide-section-phantom]{Guide bibliographique}
\item \hyperref[desirables-section-phantom]{Desiderata}
\item \hyperref[coding-section-phantom]{Style de programmation}
\item \hyperref[obsolete-section-phantom]{Obsolète}
\item \hyperref[fdl-section-phantom]{Licence de documentation libre GNU}
\item \hyperref[index-section-phantom]{Index généré automatiquement}
\end{enumerate}
\end{multicols}


\bibliography{my}
\bibliographystyle{amsalpha}

\end{document}
